% 本文件由 LaTeX 排版 Agent 根据有效 translation.json 自动生成。
% author=Augustine of Hippo; work_id=1409; title=答多纳徒派佩提利安

\chapter{答多纳徒派佩提利安}

% structure: p0001 source=h1 target=omitted reason=page-title-already-rendered-by-parent

\Person{圣奥斯定}三卷书\newline{}答\Person{佩提利安}的书信，\newline{}\OriginalTerm{多纳徒派}{Donatist}、\Place{基尔塔}主教。\par

\Emph{写于约公元400年，有人说398年，但\Person{奥斯定}将其定于\newline{}论\WorkTitle{圣洗圣事}之后：\WorkTitle{更正篇}第二卷第二十五章。从\newline{}相同来源，我们收集以下关于此论著起源的要点：\newline{}在\Person{奥斯定}完成其论\WorkTitle{天主圣三}的著作和其\newline{}逐字注释\BibleBook{}{创世纪}之前，一封回复信件（\newline{}\Person{佩提利安}寄给其追随者的），只有一小部分\newline{}落入\Person{奥斯定}手中，需要立即准备。\newline{}这构成了第一卷。随后全部文件被\newline{}获取，他着手准备第二卷，约公元401年；\newline{}但即使在获得\Person{佩提利安}全部论著之前，\newline{}后者已获得\Person{奥斯定}的第一卷，并随后发表一封\newline{}辱骂\Person{奥斯定}的书信。对此的答复是第\newline{}三卷，约公元402年。\Person{佩提利安}原本是律师。反对者\newline{}指控他是被迫成为\OriginalTerm{多纳徒派}{Donatist}，并\newline{}擅自使用护慰者称号，并竭力阻止所有\newline{}接触其著作。}\par

第一卷\newline{}第二卷\newline{}第三卷\par

\SourceRecord{Translated by J.R. King and revised by Chester D. Hartranft.；Nicene and Post-Nicene Fathers, First Series；Vol. 4.；Edited by Philip Schaff.；Buffalo, NY: Christian Literature Publishing Co.,；1887.；<http://www.newadvent.org/fathers/1409.htm>.}

% structure: p0001 source=h1 target=section reason=page-title

\section{\WorkTitle{答复多纳徒派彼得利安（第一卷）}}

\EditorialNote{约写于公元400年，亦有说398年．但奥斯定在《论圣洗》之后的《订正录》卷二第二十五章中将其定于该年之后．从同书中我们得知本论著的缘起如下：在奥斯定尚未完成其《论天主圣三》及逐字注释《创世纪》时，答复一封彼得利安写给其追随者之信函（然仅一小部分落入奥斯定之手）的工作迫在眉睫．这便是第一卷．其后，奥斯定获得了全文，并于约401年着手准备第二卷；但即使在获得彼得利安全文之前，彼得利安已得到了奥斯定的第一卷书，并随后散发了一封辱骂奥斯定的书信．对这封书信的答复即是第三卷，约写于402年．彼得利安原是一位律师．反对者指控他是被迫成为多纳徒派，且僭称\Quote{护慰者}之衔，并企图阻止他们接触其著作．}\par

\Emph{以书信形式写就，致天主教徒；文中审查并驳斥了彼得利安写给他追随者之信函的首部分。}\par

\Person{奥斯定}，致属我们照管之亲爱的弟兄们，在主内问安：\par

% structure: p0005 source=h2 target=subsection reason=relative-heading-rank; base=h2

\subsection{第一章}

1. 你们知道，我们常希望将多纳徒派异端者亵渎神明的错误公开揭露并驳斥，与其说是靠我们自己的论证，不如说是靠他们自己的言论；为此，我们甚至致信他们的一些领袖——并非为与他们共融，因为那早已因他们脱离教会而被他们自己玷污；也非斥责，而是以和解的语气，旨在使他们与我们讨论那导致他们与普世之神圣共融分裂的问题后，因思虑真理而愿意改正，不以更愚蠢的顽固捍卫其前人的刚愎乖戾，而是重归于公教之根，结出爱德之果。但正如经上所载：\Quote{我与那憎恶和平的人更加和平。}，他们便拒绝了我们的信函，正如他们憎恶和平之名本身，而信函正是为了和平之益而写。然而，如今我在\Place{君士坦丁纳}的教会中，有同僚\Person{福图纳图斯}，他的主教\Person{阿布森提乌斯}在场，弟兄们向我提及一封书信，据他们说，那分裂派的一位主教曾致信其司铎们，如书信标题所示。我读后，惊见在他最初几句话中，便砍断了他整个派别所有共融要求之根基，以至于我不愿相信这是那人的书信——若传言属实，他在他们中间尤以学识和口才著称。但一些在我阅读时在场的人，熟悉其文笔的精致与修饰，逐渐说服我，那无疑是他写的。然而我想，无论作者是谁，都需加以驳斥，免得作者在缺乏经验者中间，自以为写了什么反对公教会的重话。\par

2. 那么，他在信中首先提出的论断是：\Quote{我们责备他们重复施行圣洗，而我们自己却以沾染罪污的沐浴玷污自己的灵魂。}但全面复述他所有侮辱性言辞有何益处？因为，既然证明论点是一回事，以驳斥方式处理辱骂之词是另一回事，让我们转向他如何试图证明我们不拥有圣洗，因此他们不需重复已存在的东西，而是赋予迄今所欠缺的东西。因他说：\Quote{我们所寻求的是施予者的良心来洁净领受者的良心。}但假设施予者的良心是隐藏的，或许被罪恶玷污，它如何能洁净领受者的良心？若如他所说：\Quote{我们所寻求的是施予者的良心来洁净领受者的良心。}若他说，领受者不因施予者良心中隐藏多少罪恶而受影响，也许无知有如此效力，即一个人不会因施洗者的良心之罪而被玷污，只要他不知道。那么，姑且承认邻人的罪过良心不能玷污一个人，只要他不知道；但难道因此就能洁净他自身的罪过吗？\par

% structure: p0008 source=h2 target=subsection reason=relative-heading-rank; base=h2

\subsection{第二章}

3. 那么，当施予者的良心被玷污而领受者不知情时，领受圣洗者将如何被洁净？尤其当他接着说：\Quote{因为从无信者领受信德的人，所领受的不是信德，而是罪过。}一个无信者站在我们面前准备施洗，而要领洗者不知其无信：你以为他将领受什么？信德，还是罪过？若你回答：信德，那么你就得承认，人有可能从无信者领受信德，而非罪过；而前一句话：\Quote{从无信者领受信德的人，所领受的不是信德，而是罪过}便为虚假。因为我们发现，人即使从无信者也能领受信德，只要他不知施予者的无信。因为他没有说：从公开且臭名昭著的无信者领受信德；而是说：\Quote{从无信者领受信德的人，所领受的不是信德，而是罪过}；这若某人由隐藏其无信者施洗，则显然是假的。但若他说：即使施洗者的无信被隐藏，领受者也不会从他领受信德，而是罪过；那么，让他们重洗那些众所周知曾由那些在自己身体内长期隐藏有罪生活、最终被揭露、定罪、判决的人所施洗的人吧。\par

% structure: p0010 source=h2 target=subsection reason=relative-heading-rank; base=h2

\subsection{第三章}

因为，只要他们未被发现，他们便不能将信德赋予任何他们施洗的人，而只有罪过——若果真：\Quote{凡从无信者领受信德的人，所领受的不是信德，而是罪过}。因此，让他们由善人施洗，好使他们能领受信德而非罪过。\par

4. 然而，他们又如何能确实知道那些将要施予他们信德的人是好是坏呢？如果我们所看重的是施予者的良心，而这是受施者所看不见的？因此，按照他们的判断，灵魂的救恩变得不确定了，因为他们违背了圣经，圣经说：\BibleQuote{信赖上主胜过信赖人，}以及\BibleQuote{信赖人的，受诅咒，}\BibleRef{耶 17:6}他们却将受洗者的希望从他们的主天主身上移开，劝他们将希望寄托在人身上；其结果就是，他们的救恩不仅变得不确定，而且实际上是虚无无效的。因为\BibleQuote{救恩属于上主，}而\BibleQuote{人的援助是虚无的。}所以，凡信赖人的，即便他明知所信赖的人是正直无辜的，仍是受诅咒的。为此，宗徒\Person{保禄}责备那些说自己是属保禄的人说：\BibleQuote{保禄为你们被钉死吗？你们是受洗归于保禄的名下吗？}\BibleRef{格前 1:13}\par

% structure: p0013 source=h2 target=subsection reason=relative-heading-rank; base=h2

\subsection{第四章}

5. 因此，如果他们错了，并本会丧亡，除非被纠正——那些愿意归属于保禄的人——那么，我们必须认为那些愿意归属于多纳图的人有什么希望呢？因为他们竭尽全力证明，受洗者的根源、根本和元首无非是给他施洗的那个人。结果是，既然施洗者是什么样的人常常是不确定的，那么受洗者的希望便有着不确定的根源、不确定的根本、不确定的元首，本身便完全不确定。而且，既然施予者的良心可能处于受诅咒且被玷污的状态而受施者不知情，其结果便是，受洗者有着受诅咒的根源、受诅咒的根本、受诅咒的元首，他的希望可能会证明是虚无且无根据的。因为\Person{培提利安}在他的书信中明确声称：\Quote{一切都是由根源和根本构成的；如果它没有东西作为元首，它就是虚无。}而且，既然他想要将受洗者的根源、根本和元首理解为施洗的那个人，那么不幸的受施者不知道他的施洗者实际上有多坏，又能得到什么好处呢？因为他不知道他自己有一个坏元首，或者实际上根本没有元首。然而，一个人，无论他知道与否，如果他有一个非常坏的元首或根本没有元首，他还能有什么希望呢？当他的施洗者要么是坏元首要么根本不是元首时，我们能说他的无知本身就构成了一个元首吗？但凡这样想的人，无疑是没有元首的。\par

% structure: p0015 source=h2 target=subsection reason=relative-heading-rank; base=h2

\subsection{第五章}

6. 因此，我们问，既然他说：\Quote{从无信德者获得信德的人，获得的不是信德，而是罪债，}并随即进一步声明：\Quote{一切都是由根源和根本构成的；如果它没有东西作为元首，它就是虚无，}——我们问，我说，在施洗者的无信德未被察觉的情况下：如果这样，那么他所施洗的人获得信德，而非罪债；如果这样，施洗者不是他的根源、根本和元首，那么他从谁获得信德呢？他从哪里生发呢？他作为枝条所扎根的根部在哪里？他的出发点元首在哪里？难道可能是这样：当受洗者不知道他施洗者的无信德时，那时是基督赐予信德，那时是基督是根源、根本和元首？哀哉人的鲁莽和自负！为什么你们不承认，永远是基督赐予信德，以藉此使人成为基督徒？为什么你们不承认，基督永远是基督徒的根源，基督徒总是扎根于基督，基督是基督徒的元首？那么，我们可否坚持认为，即使是当属神的恩宠藉由一位圣洁有信德的施行人之手分施给那些信者时，成义的仍然不是施行人本人，而是那一位，关于衪曾说：\BibleQuote{祂使不虔敬的人成义？}\BibleRef{罗 4:5}但除非我们承认这点，要么宗徒\Person{保禄}是他所栽种之人的元首和根源，要么\Person{阿波罗}是他所浇灌之人的根本，而不是那在信仰中赐予他们信德的一位；然而，同一位保禄却说：\BibleQuote{我栽种了，阿波罗浇灌了，但使之生长的是天主；所以，栽种的不算什么，浇灌的也不算什么，只有使之生长的天主才算什么。}\BibleRef{格前 3:6}宗徒本人也不是他们的根本，而是那一位说：\BibleQuote{我是葡萄树，你们是枝条。}\BibleRef{若 15:5}再者，他岂能是他们的元首呢？因他说：\BibleQuote{我们许多人，在基督内是一个身体，}\BibleRef{罗 12:5}并在多处明确宣称基督自己是整个身体的元首？\par

% structure: p0017 source=h2 target=subsection reason=relative-heading-rank; base=h2

\subsection{第六章}

7. 因此，无论一个人从有信德还是无信德的施行人领受圣洗圣事，他整个的希望都在基督内，以免受那谴责：\Quote{信赖人的，受诅咒。}不然，如果每个人重生时，在属神的恩宠中与他受洗时所藉的那人同类，并且如果施洗者明显是好人，那么就是他亲自赐予信德，他亲自是那重生之人的根源、根本和元首；而如果施洗者无信德且未被察觉，那么受洗者从基督获得信德，他从基督获得根源，他扎根于基督，他以基督为元首而夸耀——在这种情况下，所有受洗者都该希望他们有无信德的施洗者，并对他们的无信德一无所知：因为不论他们的施洗者可能有多好，基督无疑更美好无比；而如果施洗者的无信德未被察觉，那么基督将成为受洗者的元首。\par

% structure: p0019 source=h2 target=subsection reason=relative-heading-rank; base=h2

\subsection{第七章}

8. 但若持此见实属至极疯狂（因为常是基督使不虔敬者成义，以其不虔敬转变为基督徒；常是从基督接受信德，基督常是重生之人的根源和教会之首），那么，那些被轻率读者按其音声而非探究其内义来珍视的言辞，又有何分量？因为那不以耳闻为足、而思忖措辞之义的人，当他听到\Quote{我们所看重的是施予者的良心，好能洁净领受者的良心}时，将回答：人的良心于我常属未知，但我确知基督的仁慈；当他听到\Quote{从无信德者领受信德者，领受的不是信德，而是罪责}时，将回答：基督不是无信者，我从祂领受的不是罪责，而是信德；当他听到\Quote{一切皆由根源与本根组成；若没有首，便一无所有}时，将回答：我的根源是基督，我的本根是基督，我的首是基督。当他听到\Quote{除非由好种子重生，否则没有事物能善得第二次出生}时，他将回答：我重生所凭的种子是天主圣言，即便传讲者自己不遵行其所教导，我也被警告要专注聆听；按照使我在此安全的主的话语：\BibleQuote{凡他们对你们所说的，你们要行要守；但不要照他们的行为去做，因为他们只说不做。}\BibleRef{玛 23:3} 当他听到\Quote{这是何等乖谬，那因自己罪愆而有罪者，竟能使他人无罪！}时，他将回答：除祂以外，无人能使我脱罪；祂为我们的罪而死，并为使我们成义而复活。因为我所信的，不是那施行洗礼之人的手，而是那使不虔敬者成义的那一位，好使我的信德被算为我的正义。\par

% structure: p0021 source=h2 target=subsection reason=relative-heading-rank; base=h2

\subsection{第八章}

9. 当他听到\Quote{凡好树都结好果子，而坏树则结坏果子：人岂能从荆棘上摘葡萄呢？}和\BibleQuote{善人从他心里所存的善就发出善来，恶人从他心里所存的恶就发出恶来；}\BibleRef{玛 12:35}时，他将回答：因此好果子就是，我应当为一棵好树，即一个好人，我应当显出好果子，即善工。但这并非由栽种者或灌溉者赐给我，而是由使之生长的天主所赐。因为若好树是好施洗者，以致其好果子就是他施洗的人，那么任何由坏人施洗的人，即使他的邪恶不显露，也无力成为好人，因为他是从坏树生出的。因为好树是一回事；品质被隐藏却仍然是坏树，是另一回事。或者，若树虽是坏的，却隐藏其坏处，那么凡由它施洗的人，并非生于它，而是生于基督；那么，那些由隐藏其恶性的坏人施洗的人，比那些由明显的好人施洗的人，要以更成全的圣德而成义。\par

% structure: p0023 source=h2 target=subsection reason=relative-heading-rank; base=h2

\subsection{第九章}

10. 再者，当他听到\Quote{与死人接触而洗涤的，他的洗涤毫无益处}时，他将回答：\BibleQuote{基督既从死者中复活，就不再死；死不再统治祂了。}\BibleRef{罗 6:9} 论到祂曾说过：\BibleQuote{祂就是那以圣神施洗的人。}\BibleRef{若 1:33} 但那些在偶像庙宇中受洗的人，是被死人施洗。因为连他们自己也不认为他们所期望的圣化是从他们的司祭领受，而是从他们的神领受；既然这些神曾是人，且已死亡，既不在世上也不在天上的安息中，他们真是被死人施洗。同样的回答也适用，如果神圣圣经的这些话能以其他方式得以查考、有益地讨论和理解。因为若我在此处将施洗之人理解为罪人，就会产生同样的荒谬：凡由不虔敬之人施洗的人，即使他的不虔敬未被发现，也如同被死人施洗一样，白白地被洗涤。因为经上不是说\Quote{被明显死去的人施洗}，而是绝对地\Quote{被死人施洗}。若他们认为他们所知道的罪人是死人，而凡在他们共融中的人即使最巧妙地隐藏罪恶生活，却仍被视为活人，那么首先，他们以可咒骂的骄傲为自己要求比归于天主的更多：当罪人向他们显露时，他被称为死人；而当罪人被天主所知时，却被视为活人。其次，若那罪人因被人所知而要被称为死人，那么对于\Person{奥普塔图斯}，他们虽久知其邪恶却不敢定罪，他们将如何回答？为什么那些由他施洗的人不被说成是由死人施洗？难道是因为伯爵是他的信德而活着？——这是他们某些早期同僚的一句优雅而巧妙的说法，他们自己常骄傲地引用，却不明白当骄傲的\Person{哥肋雅}死亡时，割下他头颅的正是他自己的剑。\BibleRef{撒上 17:51}\par

% structure: p0025 source=h2 target=subsection reason=relative-heading-rank; base=h2

\subsection{第十章}

11. 最后，如果他们乐意将\Quote{死人}之名既不给予那罪孽隐藏的恶人，也不给予那罪孽显露但尚未被他们定罪的人，而只给予那罪孽显露且已被定罪的人，以致凡由他受洗的人本身就是由死人受洗，他的洗涤对他毫无益处，那么，对于他们自己的同党已\Quote{藉一次全体会议无可指摘的声音}与\Person{马克西米安}及那些为他祝圣的人一同定罪的人——我暂且说的是\Person{穆斯蒂的费利西安}和\Person{阿苏拉的普雷特克斯塔图斯}，他们被列为\Person{马克西米安}的十二位祝圣者之一，如同在他们的祭台（\Person{普里米安}所在的祭台）对面另立祭台——我们又该说什么呢？他们当然被他们算在死人之中。对此，他们那场著名会议的庄严法令提供了明确见证，该法令当初在他们中间宣读以获通过时，曾引起毫无保留的欢呼，但现在若我们偶然在他们耳边宣读，它却会静默地接受；其实他们当初本应对其雄辩慢些欢喜，以免日后因它信誉尽失而悲伤。因为其中他们以下列言辞谈及那些被排除出他们共融的\Person{马克西米安}追随者：\Quote{既然某些人的船破肢体被真理的波浪抛向尖石，如同埃及人的习俗，海岸上覆满了垂死者的尸体；其刑罚在死亡本身中加剧，因为当他们的生命被复仇之水夺去后，他们连埋葬也得不到。}确实，他们以如此粗鄙的言辞侮辱那些因分裂而脱离他们团体的人，称他们为无埋葬的死人；但他们理应希望他们获得埋葬，哪怕仅仅是为了不让他们看见\Person{奥普塔图斯·吉尔多尼亚努斯}率领军队如一波越过同类的汹涌波浪，将\Person{费利西安}和\Person{普雷特克斯塔图斯}再次吸回他们的范围，从岸上无数无埋葬的尸体中。\par

% structure: p0027 source=h2 target=subsection reason=relative-heading-rank; base=h2

\subsection{第十一章}

12. 对此我要问：当他们回到他们的海中时，他们是否复活了？还是他们仍在那里是死人？如果仍是死尸，那么洗池绝不能有益于任何由这些死人受洗的人。但如果他们已复活，那么对于那些他们在外面尚未有生命时所施洗的人，洗池又怎能有益呢？如果\Quote{凡由死人受洗的人，他的洗礼对他有什么益处}这句话要按他们所想的方式理解？因为\Person{普雷特克斯塔图斯}和\Person{费利西安}尚在\Person{马克西米安}共融时所施洗的人，如今连同施洗者——即\Person{费利西安}和\Person{普雷特克斯塔图斯}本人——一起被保留在他们中间，分享他们的共融，而未重新受洗。借此事实，若非他们珍惜自己顽固执拗的开端，而不考虑他们属灵救恩的确定结局，他们必定会警醒，并应当恢复理智的健全，以便在天主教平安中重新呼吸；只要他们放下骄傲的膨胀，克服顽固的疯狂，留心并看出这是何等可怕的亵渎：咒诅我们从圣经中得知在初期建立的外地教会的洗礼，却接受他们亲口定罪的\Person{马克西米安}追随者的洗礼。\par

% structure: p0029 source=h2 target=subsection reason=relative-heading-rank; base=h2

\subsection{第十二章}

13. 然而，上述教会的子女——我们的弟兄们——当时既不知道，至今仍不知道许多年前在非洲所行的事；因此他们无论如何不能被多纳特派对非洲人所提出的指控所玷污，甚至不知道这些指控是否属实。但多纳特派公开分离并分裂出去，尽管他们甚至据说参与了\Person{普里米安}的祝圣，却谴责了\Person{普里米安}，祝圣了另一位主教与\Person{普里米安}对立，在\Person{普里米安}共融之外施洗，在\Person{普里米安}之后重洗，然后带着那些被他们自己在外面施洗、从未在内部被任何人重洗的门徒回到\Person{普里米安}那里。如果与\Person{马克西米安}派的这种联合并未玷污多纳特派，那么仅仅关于非洲人的传闻又怎能玷污外地人呢？如果那些互相谴责的嘴唇在平安礼中相遇而不冒犯，为什么那些被他们定罪的、位于极其遥远的大海相隔的教会中的人，不被作为忠信的天主教徒以亲吻问候，反而被作为不敬神的外邦人怒气冲冲地驱逐呢？而如果他们接纳\Person{马克西米安}派是为了他们自己的合一而讲和，那么我们绝不怪责他们，只是他们用自己的决定割断了自己的喉咙：因为为了保持他们分裂中的合一，他们重新聚拢了从自己分离出去的部分，却藐视将自己的分裂本身重新与教会的真正合一联合。\par

% structure: p0031 source=h2 target=subsection reason=relative-heading-rank; base=h2

\subsection{第十三章}

14. 如果为了多纳特派团体的合一，没有人重新为那些在邪恶分裂中领受圣洗的人施洗，而那些犯有如此严重罪行——他们在自己的公议会中将此罪比作被活活吞入地下的古老分裂始作俑者\BibleRef{户 16:31}——的人，在分离之后不受惩罚，或在定罪后又恢复原位；那么，为什么为了基督的合一——这合一遍及整个居住世界，正如预言所说，它要从这海到那海，从大河直到地极\BibleRef{咏 2:8}，而这预言似乎正在实际证实中——他们却不承认那作为继承产业的真正普世法则，这法则从我们共有的经卷中发出：\Quote{我将使外邦人成为你的产业，使地极成为你的基业。}为了多纳特的合一，他们不必重新召集他们已分散的人，而是被警告要听从圣经的呼声：他们为什么不明白，他们遭遇这样的对待是出于天主的仁慈？因为他们诬告天主教会，仿佛不愿沾染其过度圣洁，从而被奥普塔图斯·吉尔多尼亚努斯的最高权威迫使重新接纳并联合自身真正的大罪，这些罪被他们自己的全体公议会以真实的声音——如他们所说——所定罪。让他们最终明白，他们是如何被自己党派的真实罪行所充满，而他们却捏造虚构的罪行来控告弟兄们；即使那些指控是真的，他们也该最终明白，为了和平的缘故，为了基督的和平，应当忍受多少才能回归到并不谴责未被发现之罪行的教会，而为了多纳特的和平，他们却准备赦免已被定罪的人。\par

% structure: p0033 source=h2 target=subsection reason=relative-heading-rank; base=h2

\subsection{第十四章。}

15. 因此，弟兄们，让我们满足于这一点：他们应在对待马克西米安追随者这件事上被劝诫和纠正。我们不翻查古代档案，我们不揭示古老图书馆的内容，我们不向远方公布我们的证据；但我们引入祖先的所有证据作为我们之间的仲裁者，我们传播那响彻全世的见证。\par

% structure: p0035 source=h2 target=subsection reason=relative-heading-rank; base=h2

\subsection{第十五章。}

16. 请看穆斯蒂和阿苏拉两城：许多人仍活在今世和本省，他们已分离，也有许多人他们已与之分离；许多人立起了祭台，也有许多人反对那祭台；许多人定了罪，也有许多人被定罪；他们接纳了，也被接纳了；他们在外面领受了圣洗，却没有在里面重新领受圣洗。如果这一切在合一的缘故上玷污人，就让被玷污的人闭口；如果这一切在合一的缘故上并不玷污人，就让他们接受纠正，终止争论。\par

% structure: p0037 source=h2 target=subsection reason=relative-heading-rank; base=h2

\subsection{第十六章。}

17. 至于他信中所跟随之言，连作者自己也几乎不能不嘲笑它们：当他无知而虚假地使用他所引用的圣经经文\Quote{被死人洗濯的人，他的洗濯有什么益处？}时，他力图向我们说明\Quote{一个活着的\OriginalTerm{背信者}{traditor}在何种程度上可被视为死人。}然后他接着说：\Quote{那人不配在真正的圣洗中重生，他就是死人；同样，那人虽在真正的圣洗中重生，却与一个背信者联合，他也是死人。}因此，如果马克西米安的追随者不是死人，为什么多纳特派在他们的全体公议会中说\Quote{海岸上布满了他们垂死的身体}？但如果他们是死人，他们所施与的圣洗中生命从何而来？再者，如果马克西米安不是死人，为什么由他领受圣洗的人要重新领受圣洗？但如果他是死人，为什么与他一齐在穆斯蒂被祝圣的费利西安——他祝圣了马克西米安，并可能已与某个非洲的背信者同事死于海外——不是死人？或者，如果他自己也是死人，那么，在你的团体中，那些由那位死人而在外面领受圣洗、却从未在里面重新领受圣洗的人，如何与他同在生命之中？\par

% structure: p0039 source=h2 target=subsection reason=relative-heading-rank; base=h2

\subsection{第十七章。}

18. 然后他进一步说：\Quote{两人都没有圣洗的生命，既包括从未拥有圣洗的人，也包括曾拥有但失去圣洗的人。}因此，费利西安——马克西米安或普雷特克斯塔图的追随者——在外面施洗的人从未拥有圣洗；而这些人在拥有之时失去了圣洗，因此，当这些人连同他们的追随者被接纳时，谁给了他们所施洗的人以前所没有的？谁又恢复了他们自己所失去的圣洗？他们带走了圣洗的形式，却因邪恶的分裂失去了圣洗的真正卓越。为什么你们在你们未曾听过的天主教徒身上否定这形式本身——它在任何时代和任何地方都是神圣的——却在你们已惩罚的马克西米安追随者身上愿意承认它？\par

19. 但是，无论他似乎关于叛徒犹大说了什么指控的话，我看不出这与我们有何相干，因为我们并未被他们证明背叛了我们的信任；而且，即使在我们之前已死于我们共融中的任何人身上证明了这样的背叛，那背叛也绝不会玷污我们，因为我们不认同它，并且不喜悦它。因为，如果他们自己不被自己所定罪并后来宽恕的过犯所玷污，那么，我们一旦听到就立即不认同的事，又怎能玷污我们呢？因此，无论他对背信者的控诉多么严厉，让他知道，我也以同样的措辞谴责他们。但我仍做一个区分：他指控我这边一个早已死去的人，此人未经我所做的任何调查便被定罪。而我指出一个紧靠他身边、曾被他自己定罪或至少因亵渎的分裂而分离、而他又以完好无损的尊荣接纳的人。\par

% structure: p0042 source=h2 target=subsection reason=relative-heading-rank; base=h2

\subsection{第十八章。}

20. 他说：\Quote{你这最邪恶的\Emph{traditor}，竟以迫害者和杀人者的面目出现，来对付我们这些守法的人。}如果\Person{马克西米安}的追随者在他们与你们分离时遵守了法律，那么当你们与遍布普世的教会分离时，我们也可以承认你们是法律的遵守者。但如果你提出迫害的问题，我立即回答：如果你们曾遭受不公正的对待，这与那些虽然不赞同如此行事的人、却为了合一中的和平而忍受他们、并以值得称颂的方式行事的人无关。因此，你们没有什么可以拿来反对主的麦子，他们忍受他们中间的糠秕直到最后的扬净；若不是你们在扬场者到来之前，因诱惑之风而飞走，显明自己比糠秕更轻，你们绝不会与主麦子分离。但不妨只举这个例子，就是主塞在他们牙齿间的，为要封住这些人的口，如果他们显为聪明，为要纠正他们；但如果他们固执于愚妄，则为要使他们困惑：如果受迫害的人比施加迫害的人更正义，那么那些\Person{马克西米安}的追随者就更正义，因为他们的教堂被彻底拆毁，他们受到\Person{奥普塔图斯}军事随从的严重虐待，其时总督的命令——命令将他们全部驱逐出教堂——明显是由\Person{普里米安}的追随者促成的。因此，如果当皇帝憎恨他们的共融时，他们竟敢于采取如此暴力的措施来迫害\Person{马克西米安}的追随者，那么若他们得以与君王共融而能随心所欲，他们又会做什么呢？如果他们做了我提到的这些事情是为了纠正恶人，为什么他们惊讶于公教皇帝以更大的权力下令来对付和纠正那些企图重施圣洗整个基督徒世界的人，而后者却没有理由与他们不同？因为他们自己见证，即使在有真正指控的情况下，为了和平也有必要容忍恶人，因为他们接纳了那些自己曾谴责的人，承认在他们中间授予的荣誉，以及在分裂中施行的圣洗。让他们最终思考，他们在世间基督徒权力手中应得何种待遇，这些权力是整个世界基督徒合一的敌人。因此，如果纠正令人痛苦，仍让他们不要羞于感到羞愧；免得他们开始读自己所写的东西时，被笑声压倒，因为他们在别人身上找不到自己所希望的东西，并且在别人身上指责的，却在自己身上认不出来。\par

% structure: p0044 source=h2 target=subsection reason=relative-heading-rank; base=h2

\subsection{第十九章。}

21. 那么，他在信中引用我主对犹太人说的一句话，是什么意思：\Quote{所以，看哪，我派遣先知、智者、经师到你们这里来；你们要杀害并钉死他们中的一些，鞭打其中的另一些。}\BibleRef{玛 23:34} 因为如果他们将自身理解为智者、经师和先知，而我们就好像先知和智者的迫害者，那么他们为什么不愿意与我们交谈，既然他们是被派遣到我们这里来的？事实上，如果我们现在正在回应的那封信的作者被我们强迫承认那是他所写的，并在上面签名验证，我很怀疑他是否愿意这样做，他们如此害怕我们拥有他们的任何文字。因为当我们急于以某种方式获取同一封信的后半部分时，因为提供给我们的人无法描述整封信，我们问及的人没有一个人愿意将它给我们，一旦他们知道我们正在对我们所拥有的部分进行答复。因此，当他们读到主对先知说：\Quote{你要大声呼喊，不要制止，用我的笔写下他们的罪过。}\BibleRef{依 58:1} 这些作为先知被派遣到我们这里来的人，在这点上毫无恐惧，反而采取一切预防措施，使他们的呼喊不被我们听到；如果他们所说关于我们的话是真的，他们当然不会害怕。但他们的恐惧并非没有根据，正如\BibleBook{}{圣咏}所记载：\BibleQuote{说谎之人的口必被堵住。} 因为如果他们没有接受我们的圣洗的理由是我们是一窝毒蛇——用他信中的话来说——那么他们为什么接受\Person{马克西米安}派的圣洗呢？他们的会议对此用以下言辞描述：\Quote{因为中毒的子宫的包裹早已隐藏了毒蛇种子的有害后代，而所孕育的邪恶之湿滑凝结物以缓慢的热量流向蛇的肢体。} 因此，难道不也是因为他们自己而在同一会议中说：\Quote{蝮蛇的毒液在他们的嘴唇下，他们的口充满咒骂和苦毒，他们的脚快速流人血，毁灭和不幸在他们的路上，平安的道路他们未曾认识。} 然而，他们现在仍保持这些人自身的荣誉不减，并接纳这些人在团体之外所付洗的人进入他们的团体。\par

% structure: p0046 source=h2 target=subsection reason=relative-heading-rank; base=h2

\subsection{第二十章。}

22. 因此，所有这些关于毒蛇的种类、唇下有毒蛇的毒汁，以及他们针对那些不认识和平道路之人所说的一切，如果他们说真话，实际上更严格地适用于他们自己，因为他们为了\Person{多纳图}的和平而接受了这些人的\OriginalTerm{圣洗}{baptism}，并且在议会法令的措辞中使用了上述表达；但是他们却拒绝接受普世各地基督的教会——和平本身正是由此来到非洲——的圣洗，从而亵渎地伤害了基督的和平。那么，谁才是穿着羊皮来到你们中间、内里却是残暴的豺狼的假先知呢？\BibleRef{玛 7:15}——是那些在天主教会中未能察觉恶人、并以无罪之心与他们共融的人？还是为了合一的和平而忍受那些在簸扬者到来之前无法从主的打谷场上分离出去的人？还是那些在分裂中做出他们在天主教会中所谴责之事，并在自己的分离中接纳那些他们声称在教会合一中避之不及的人——即使这些事在合一中已被众人知晓、被他们自己的声音定罪，并且甚至不一定存在？\par

% structure: p0048 source=h2 target=subsection reason=relative-heading-rank; base=h2

\subsection{第二十一章。}

23. 最后，正如他自己也引用的，经上说过：\BibleQuote{你们凭他们的果实就可以认出他们来。}\BibleRef{玛 7:16}让我们因此检视他们的果实。你们指控我们的前辈交出圣书。我们更能以这同样的指控，合理地反对他们的控告者本身。为免查考过远，在同一个\Place{君士坦提纳}城中，你们的前辈在西尔瓦努斯分裂之初就祝圣他为主教。他当时还是副执事，被明确无误地作为\OriginalTerm{叛教者}{traditor}登录在该城的档案中。如果你们一方提交不利于我们前辈的文件，我们只要求公平的条件：要么我们相信两者都真，要么两者都假。若两者都真，你们无疑犯了分裂之罪，因为你们假装在普世教会的共融中避开过错，而这些过错在你们自己小教派的一隅却比比皆是。但若两者都假，你们也无疑犯了分裂之罪，因为你们因交出圣书的虚假指控，用与教会分离的滔天大罪玷污了自己。但如果我们有可指控之事而你们没有，或者我们的指控属实而你们的虚假，那你们的口究竟是如何被封住的，便无需讨论了。\par

% structure: p0050 source=h2 target=subsection reason=relative-heading-rank; base=h2

\subsection{第二十二章。}

24. 倘若神圣而真实的基督的教会要说服并胜过你们，即使我们没有任何支持我们事业的文件，或只有虚假的文件，而你们却拥有某些交出圣书的真实证据，那又会如何？那时你们还能做什么，除非你们愿意表示对和平的热爱，否则就应闭嘴？因为无论如何，在那种情况下，你们可能提出的证据，我都能最轻松、最真实地说，你们有义务向那已传播开来、建立在如此众多民族中的圆满大公教会证明这些，目的是你们应留在其中，而你们所定罪的人应被逐出。如果你们曾试图这样做，当然你们未能证实；你们被击败或激怒，以亵渎的重罪与那些不能因证据不足而定罪的无辜者分离。但如果你们甚至没有试图这样做，那么你们便以最可咒诅、最不自然的盲目，从基督的麦子中割断了你们自己——这麦子生长在他的整个田地，即整个世界，直到终结——因为你们在非洲因几棵莠子而跌倒。\par

% structure: p0052 source=h2 target=subsection reason=relative-heading-rank; base=h2

\subsection{第二十三章。}

25. 总之，据说在迫害时期，某些人将圣约投诸火焰。现在，无论圣约从何处被发掘，让它的教训被宣读。确实，在立约者应许的开头，就对亚巴郎说了这话：\BibleQuote{地上万国都要因你的后裔得福，}\BibleRef{创 22:18}并被宗徒正确地解释为：\BibleQuote{你的后裔，他说，就是指基督。}\BibleRef{迦 3:16}任何人的背叛都没有使天主的应许失效。与地上万国共融，然后你们可以夸耀你们保存了圣约，未被火焰毁灭。但如果你们不这样做，那么哪一方更应被相信是坚持烧毁圣约的？除了那一方，即当圣约被揭示时，它不愿同意其教导的那一方。因为，一个现在用舌头迫害圣约的人——据说圣约曾被他们用火焰迫害——更确定无疑地，无需任何亵渎的鲁莽，就可以被认定为加入了\OriginalTerm{叛教者们}{traditors}的行列。你们指控我们迫害：主的真麦子回答你们：\Quote{或者它是公正地做的，或者它是我们中间的糠秕做的。}你们对此有何话说？你们反对说我们没有圣洗：同样的主的真麦子回答你们：圣事的标记即使在教会内也对某些人无益，如对\Person{西满术士}施洗时没有益处，更何况对那些在教会之外的人。然而，当他们离开时，那圣洗仍留在他们里面，这由这一点证明：当他们回归时，圣洗并不重新赐予他们。因此，除非有最无耻的厚颜，你们绝不能向那麦子喊叫，或称他们是披着羊皮、内里却是残暴的豺狼的假先知；因为他们要么在天主教会的合一中不认识恶人，要么为了合一的缘故忍受他们所认识的人。\par

% structure: p0054 source=h2 target=subsection reason=relative-heading-rank; base=h2

\subsection{第二十四章。}

26. 但让我们来审视你们的果实。且不提在城市中，特别是在他人地产上的专制权威；且不提巡回团的疯狂，以及那些自愿跳崖者的身体所受的亵渎和亵圣的崇拜；且不提酗酒的狂欢，以及整个非洲在奥普塔图斯·吉尔多尼亚努斯一人的残暴下呻吟了十年：这一切我都略过不提，因为你们中间有些人会喊叫说，这些事他们从来就不喜欢，现在也不喜欢。但他们说，他们为了和平的原因而容忍了这些事，因为他们无法制止它们；他们的判断本身就定了自己的罪：因为如果他们真的如此热爱和平，就绝不会撕裂合一的纽带。还有什么疯狂比这更大呢？在和平之中抛弃和平，而在纷争之中却竭力要维持和平？因此，为了那些假装看不见这同一多纳徒派别之恶事的人们——这些事所有人都看见并谴责，他们甚至对奥普塔图斯本人也说：\Quote{他做了什么？}\Quote{谁控告了他？}\Quote{谁定了他的罪？}\Quote{我什么都不知道，}\Quote{我什么都没看见，}\Quote{我什么都没听见，}——为了这些人，我说，他们假装对众所周知的恶事一无所知，马克西米安派兴起了，通过他们，这些人的眼睛被打开，嘴巴被堵住：因为他们公开分裂自己；他们公开竖起祭台对抗祭台；他们在一个会议中被公开称为亵渎者和毒蛇，是乐意流人血的，与达堂、阿彼兰和科辣黑相比，并以严厉的憎恶言辞被谴责；而他们又同样公开地，连同他们所施洗的人一起，以不减的尊荣被接纳回来。这就是这些人的果实，他们为多纳徒的和平做这一切，为自己披上羊皮，却拒绝基督在世界上的和平，好成为圈内凶暴的豺狼。\par

% structure: p0056 source=h2 target=subsection reason=relative-heading-rank; base=h2

\subsection{第二十五章。}

27. 我认为我没有留下任何多纳徒信函中的陈述未作答复，至少就我们所拥有的那部分而言。如果他们也能拿出另一部分就好了，万一其中有无法反驳的内容呢。但至于我们靠着天主的助佑对这些所作的答复，我劝告你们的基督徒爱德，不仅要把它们传达给那些寻求它们的人，也要强加给那些不渴望它们的人。让他们随意回答任何问题；如果他们不愿回复我们，那么至少让他们写信给自己派别的人，只要不禁止把内容给我们看。因为他们若这样做，就最公开地显示出他们的果实，由此被证明是披着羊皮的凶暴的豺狼——他们暗中为我们的羊设置陷阱，却公开回避给牧人任何答复。我们只指控他们分裂之罪，他们全都深陷其中——而非他们中间某些人的过犯，这些人声称那些过犯连他们自己也不喜欢。另一方面，如果他们自己不指控我们别人的罪，他们就没有什么可指控我们的，因此他们完全无法为自己辩护以脱免分裂之罪；因为他们是以邪恶的分裂离开了主的打谷场，离开了在全世界生长的无辜麦粒的陪伴，而他们提出的指控，要么是虚假的、自己捏造的，要么即便是真的，也仅涉及糠秕而已。\par

% structure: p0058 source=h2 target=subsection reason=relative-heading-rank; base=h2

\subsection{第二十六章。}

28. 但你们可能期望我继续反驳他关于摩尼的论述。现在，关于这一点，唯一冒犯我的事是，他以完全不充分的严厉措辞谴责了一个极其恶毒和有害的异端——即摩尼派的异端——如果这些措辞真能算作谴责的话，尽管天主教会以最有力的真理证据击垮了他的防御。因为基督的遗产，建立在万民之中，对于那些被排斥在遗产之外的异端来说是安全的；但是，正如主所说：\BibleQuote{撒殚怎能驱逐撒殚？}\BibleRef{谷 3:23}，那么，多纳徒派的错误又怎能推翻摩尼派的错误呢？\par

% structure: p0060 source=h2 target=subsection reason=relative-heading-rank; base=h2

\subsection{第二十七章。}

29. 因此，我亲爱的弟兄们，尽管那错误以多种方式被揭露和战胜，且不敢以任何理性的样子反对真理，只敢以无耻的固执顽抗；然而，为了不加重你们记忆的负担，愿你们记住马克西米安派的一个行动，用这一事实来面对他们，塞住他们的口，以此摧毁他们的诽谤，如同用三叉戟摧毁三头怪物。他们指控我们背叛圣书；他们指控我们迫害；他们指控我们假圣洗：用关于马克西米安派的同一事实来回答他们所有的指控。因为他们认为证明他们的前辈焚烧圣书卷的证据已经丢失；但至少这一点他们无法隐藏：他们已接纳了那些沾染裂教亵渎的人，且未减损其尊荣。他们也认为那些最暴力的迫害是隐藏的，当他们有能力时，他们便对任何反对者施加这种迫害；但是，精神的迫害超越身体的迫害，他们却接受了那些自己曾身体迫害过的人——马克西米安派，且未减损其尊荣，而他们自己曾论及这些人说：\Quote{他们的脚急於流人血；}而这一点至少他们无法隐藏。\par

% structure: p0062 source=h2 target=subsection reason=relative-heading-rank; base=h2

\subsection{第二十八章。}

最后，他们认为关于圣洗的问题是隐藏的，他们以此欺骗可怜的灵魂。但是，当他们说凡是在一个教会共融之外受洗的人都没有领受圣洗时，他们却接纳了马克西米安派及其在分裂中、在多纳徒共融之外所施洗的人，且未减损其尊荣，而这一点至少他们无法隐藏。\par

30. 他们说：\Quote{但为了和平的缘故，这些事并不带来玷污；在极其严厉的惩罚中怀有仁慈是好的，好使折断的枝子可以重新被接上。}因此，整个问题就此解决了——他们失败，而我们不可能失败；因为如果和平之名被用来作为哪怕最微弱的辩护借口，以证明在分裂中忍受恶人是正当的，那么毫无疑问，在普世合一中，真正的和平本身遭到破坏，就是可憎的罪，没有任何辩解可言。\par

% structure: p0065 source=h2 target=subsection reason=relative-heading-rank; base=h2

\subsection{第二十九章}

31. 弟兄们，我希望你们把这些作为你们行动和宣讲的基础，以不懈的温良：爱世人，同时消灭错误；毫无骄傲地接受真理；毫不残忍地为真理而战。为那些你们驳斥并说服其错误的人祈祷。因为先知为这样的人向天主祈求仁慈说：\BibleQuote{上主，求祢使他们的脸上充满羞耻，好叫他们寻求祢的名。}事实上，主已经这样做了，使\Person{马克西米安}追随者的脸在全人类面前充满羞耻：现在只待他们学会为自己的灵魂健康而羞愧。如此他们便能寻求主的名，他们曾因极度自毁而背离主名，却将自己的名高举以代替基督的名。愿你们在基督内生活并坚持，繁衍增多，满溢天主的爱和彼此之间的爱，以及对所有人的爱，亲爱的弟兄们。\par

\SourceRecord{Translated by J.R. King and revised by Chester D. Hartranft.；Nicene and Post-Nicene Fathers, First Series；Vol. 4.；Edited by Philip Schaff.；Buffalo, NY: Christian Literature Publishing Co.,；1887.；<http://www.newadvent.org/fathers/14091.htm>.}

% structure: p0001 source=h1 target=section reason=page-title

\section{\WorkTitle{答多纳徒派培提利安（卷二）}}

在此，\Person{奥古斯丁}答复\Person{培提利安}书信中的各项陈述，仿佛与论敌当面辩论。\par

% structure: p0003 source=h2 target=subsection reason=relative-heading-rank; base=h2

\subsection{第一章}

1. 凡能读到或听到我们答复的人，都会记得我们曾对\Person{培提利安}书信的第一部分（那是我们所能找到的全部）作出了充分而圆满的回答。但后来，弟兄们找到了该书信的全本并抄录下来，寄给我们，希望我们整体答复；这项任务是我们无法推卸的——并非他在其中提出了什么新的论点，我们已经以多种方式和在不同时期多次答复过了；但为了理解较慢的弟兄们，他们读到一个问题时，不能总是参考同一主题之前说过的一切，我将顺从那些力劝我务必逐点答复之人的意见，而且仿佛我们以对话形式面对面地讨论一样。我将把他在书信中的话记在他的名下，把我的答复记在我的名下，仿佛这一切都是由记录员在我们辩论时记录下来的。这样，就没有人能抱怨说我遗漏了什么，或者由于不能区分讨论双方而无法理解；同时，那些不愿在我们面前讨论问题的多纳徒派信徒（正如他们在本派内流传的信件所表明的），也能像与我们面对面讨论一样，听到真理逐点答复他们。\par

\Person{培提利安}说：\Quote{一位主教，致我们亲爱的弟兄、同工司铎和执事们——这些同我们一同在整个教区内担任福音职务的执事——愿恩宠与平安，由天主我们的父及主耶稣基督赐予你们。}\par

\Person{奥古斯丁}答复：我认出这宗徒的问候。你看见你是什么样的人，竟使用它；但你也看见你是从哪里学到了你所说的话。因为\Person{保禄}就是用这些话问候罗马人，用同样的话问候格林多人、迦拉达人、厄弗所人、哥罗森人、斐理伯人、得撒洛尼人。那么，不愿与那些教会共享平安的救恩，是何等疯狂？你正是从他们的书信中学会了这问候的形式。\par

% structure: p0007 source=h2 target=subsection reason=relative-heading-rank; base=h2

\subsection{第二章}

\Person{培提利安}说：\Quote{那些以罪恶的圣洗盆、以圣洗之名玷污了自己灵魂的人，指责我们施两次圣洗——其实，任何污秽都比他们的污秽更洁净，因为他们因清洁的扭曲，竟因洗濯而变得更肮脏。}\par

\Person{奥古斯丁}答复：我们的圣洗并未使我们更肮脏，你们的圣洗也未使我们更洁净。但当任何人因父、及子、及圣神之名施予圣洗之水时，那水既不是我们的，也不是你们的，而是属于那一位的，他曾对若翰说：\BibleQuote{你看见圣神降下，停在谁身上，谁就是那要以圣神施洗的。}\BibleRef{若 1:33}\par

% structure: p0010 source=h2 target=subsection reason=relative-heading-rank; base=h2

\subsection{第三章}

\Person{培提利安}说：\Quote{因为我们所看重的是施予者的良心，以洁净领受者的良心。}\par

\Person{奥古斯丁}答复：因此，我们无需为基督的良心担忧。但如果你断言任何人是施予者，无论他是谁，领受者的洁净便无确定性，因为施予者的良心无确定性。\par

% structure: p0013 source=h2 target=subsection reason=relative-heading-rank; base=h2

\subsection{第四章}

\Person{培提利安}说：\Quote{因为从无信者接受信德的人，领受的不是信德，而是罪责。}\par

\Person{奥古斯丁}答复：基督不是无信者，从他那里，有信德的人领受的不是罪责，而是信德。因为那人信靠那使不虔敬者成义的天主，好使他的信德被算为义。\BibleRef{罗 4:5}\par

% structure: p0016 source=h2 target=subsection reason=relative-heading-rank; base=h2

\subsection{第五章}

\Person{培提利安}说：\Quote{因为万物皆由根源和根基构成；若没有什么为头，就什么都不是；任何事物也不能完好地领受第二次出生，除非它由好种子重生。}\par

\Person{奥古斯丁}答复：你为什么要把自己放在基督的位置上，而不愿将自己置于祂之下？祂是那重生之人的根源、根基和头；在祂身上，我们无需担忧任何一个人——无论他是谁——可能会是虚伪的、品行败坏的人，而我们发现自己竟出自一个败坏的根源、生长于一个败坏的根基、连接于一个败坏的头。因为哪个人能对人放心呢？经上记载：\BibleQuote{凡信赖世人的人，是可咒骂的。}\BibleRef{耶 17:5}但使我们重生的种子是天主的话，即福音。因此宗徒说：\BibleQuote{我在基督耶稣内藉福音生了你们。}\BibleRef{格前 4:15}然而，他竟允许那些不纯洁地传福音的人也传福音，并为他们的宣讲而喜乐\BibleRef{斐 1:17}；因为他们虽然不纯洁地传福音，只寻求自己的利益，而非耶稣基督的事\BibleRef{斐 2:21}，但他们所传的福音是纯洁的。主也曾论及类似的人说：\BibleQuote{凡他们对你们所说的，你们要行要守；但不要照他们的行为去做，因为他们只说不做。}\BibleRef{玛 23:3}因此，如果纯粹的事本身被纯洁地宣讲，那么宣讲者本人，就他与圣言的关系而言，也在生育信徒的事上有份；但如果他本人未重生，而他所宣讲的却是纯洁的，那么这位信徒并非由职务的荒芜而生，乃是由圣言的丰硕而生。\par

% structure: p0019 source=h2 target=subsection reason=relative-heading-rank; base=h2

\subsection{第六章}

\Person{培提利安}说：\Quote{如果情况如此，弟兄们，这是何等悖逆：有罪的人因自己的罪而有罪，却能叫别人无罪脱罪？主耶稣基督说：\BibleQuote{凡好树结好果子，坏树结坏果子。人岂能从荆棘上摘葡萄呢？}又说：\BibleQuote{善人从心里所存的善，就发出善来；恶人从心里所存的恶，就发出恶来。}\BibleRef{玛 12:35}}\par

13. 奥古斯丁回答：没有人，即使他自己并没有因自己的罪而有罪，能使他邻居无罪，因为他不是天主。否则，如果我们期望从施行圣洗者的无罪产生受圣洗者的无罪，那么每个人将因他找到的施行圣洗者越无罪而自己越无罪；而如果施行圣洗者不那么无罪，他自己也将不那么无罪。如果施行圣洗者碰巧对另一个人怀有怨恨，这也要算在受圣洗者头上。那么，这个可怜的人为什么要急忙去受圣洗呢——是为了让他自己的罪被赦免，还是为了让别人的罪记在他身上？他就像一艘商船，卸下一个负担，又装上另一个？但通过好树和它的好果实，以及坏树和它的坏果实，我们通常理解为人及其行为，正如你在所引的其他话中显示的：\Quote{善人从心里所存的善就发出善来；恶人从心里所存的恶就发出恶来。}但一个人宣讲天主圣言，或施行天主的圣事时，如果他是一个坏人，他并不是从自己的宝库里宣讲或施行；他将被算作那些关于他们所说的话的人：\Quote{凡他们所吩咐你们的，你们都要谨守遵行；但不要效法他们的行为：}因为他们吩咐你们遵守的是天主的事，而他们的行为是他们自己的。因为如果像你所说，即施行圣洗者的果实在于受圣洗者本身，那么你就在非洲宣布了一个大祸，如果奥普塔图斯所施行圣洗的每一个人都产生了一个小奥普塔图斯的话。\par

% structure: p0022 source=h2 target=subsection reason=relative-heading-rank; base=h2

\subsection{第七章}

14. 彼得利安说：\Quote{再者，\InnerQuote{那由死人施行圣洗的，他的洗涤对他毫无益处。}他并不是说施行圣洗者是尸体，没有生命的身体，准备埋葬的人的遗骸，而是指缺乏天主圣神的人，他被比作尸体，正如主在另一处对门徒所宣告的，根据福音的见证。因为祂的门徒说：\BibleQuote{主，请准许我先去埋葬我的父亲。但耶稣对他说：\InnerQuote{跟随我，让死人去埋葬他们的死人。}}\BibleRef{玛 8:21}那门徒的父亲没有受圣洗。他指出他作为一个外邦人属于外邦人的行列；除非祂这样说不信的人：死人不能埋葬死人。因此，他不是因某种死亡而死去，而是甚至在活着时就被击打致死。因为如此生活以致注定永死的人，是在生命中受死亡折磨。所以，由死人施行圣洗，就是领受死亡而非生命。因此我们必须考虑并宣布，\OriginalTerm{叛徒}{traditor}在活着时应在多大程度上被视为死人。他是死人，他没有配得以真正的圣洗重生；同样，他也是死人，他虽以真正的圣洗重生，却与\OriginalTerm{叛徒}{traditor}有牵连。两者都缺少圣洗的生命——既从未拥有过它的人，又曾拥有却失去了它的人。因为主耶稣基督说：\BibleQuote{那人的景况就比先前更不好了。}\BibleRef{玛 12:45}}\par

15. 奥古斯丁回答：要更仔细地查考你为证明自己论点所引用的圣经经文是在什么意义上说的，以及如何理解。因为一切不义的人在奥秘意义上常被称为死人，这是足够清楚的；但基督，真圣洗属于祂，你却说因人的过犯，圣洗是假的，祂是活着的，坐在圣父的右边，祂不会再因任何肉身的软弱而死：\BibleQuote{死再不能辖制祂。}\BibleRef{罗 6:9}那些受祂的圣洗的人，不是由死人施行圣洗。如果某些牧者，是欺诈的工人，寻求自己的事，而非耶稣基督的事，不纯正地宣讲福音，出于嫉妒和纷争宣讲基督，因他们的不义而被称为死人，那么，活天主的圣事甚至在死人中也不会死去。因为在撒玛黎雅受圣洗的\Person{西满}是死的，他想要用钱购买天主的恩赐；但他所领受的圣洗仍活在他身上，为他的惩罚。\par

16. 但你所做的论断是多么错误，你说\Quote{两者都缺少圣洗的生命——既从未拥有过它的人，又曾拥有却失去了它的人}，你可以从这一点看出：对于那些在受圣洗后背教，然后通过补赎回来的人，圣洗并不重新赋予他们，如果它丢失了，就会重新赋予。事实上，按照你的解释，你们的死人怎样施行圣洗呢？我们难道不该把醉酒者也算在死人中吗（且不说其他，只提众人皆知且日常可见的事）？因为宗徒关于寡妇说：\BibleQuote{但那好宴乐的寡妇，虽活着也是死的。}\BibleRef{弟前 5:6}其次，在你们那个议会中，你们谴责了\Person{马克西米安}及其顾问或牧者，难道你忘记了你曾多么雄辩地说：\Quote{即使像埃及人那样，岸边布满了垂死者的尸体，那更重的惩罚落在死亡本身，因为他们的生命被复仇之水夺去后，他们甚至没有找到埋葬？}然而你们自己可以看到，他们中的一位，\Person{费利西安}，是否已经复活；但他仍与你们身体内那些他在外面施行圣洗的人同在共融中。因此，正如披上活基督的圣洗的人是由活着的一位施行圣洗，那么，由死人施行圣洗的人乃是裹着死去的萨图尔或任何类似的人的圣洗；为了暂时以我们所能的简要说明，你所引用的经文可以在什么意义上理解，而我们中间没有人吹毛求疵。因为，按照你们接受的意义，你们并不努力解释它们，而只是试图将我们与你们自己纠缠在一起。\par

% structure: p0026 source=h2 target=subsection reason=relative-heading-rank; base=h2

\subsection{第八章}

17. \Person{Petilianus}说：\Quote{我们必须思量并宣告，那背信的\OriginalTerm{叛徒}{traditor}虽生犹死，其程度如何。\Person{犹达斯}背叛\Person{基督}时还是宗徒；但他已经死了，灵性上已失去了宗徒的职分，后来注定要自缢而死，正如经上所载：\InnerQuote{他说：\InnerQuote{我犯了罪，因为我出卖了无辜的血}；他就离去，上吊自尽了。}\BibleRef{玛 27:4}叛徒被绳索绞死：他把绳索留给了其他像他一样的人，主\Person{基督}曾向圣父大声呼喊：\InnerQuote{父啊，祢赐给我的人，我保全了他们，其中没有一个丧亡，除了那丧亡之子，使经书得以应验。}\BibleRef{若 17:12}因为\Person{达味}早已对那将要向不信者出卖\Person{基督}的人下了判决：\InnerQuote{让别人取他的职位。愿他的儿女成为孤儿，他的妻子成为寡妇。}看，先知们的神力何等伟大，竟能将一切未来的事视如当下，使一个日后才出生的叛徒在多个世纪前就被定罪。最后，为使那判决得以完成，圣\Person{玛弟亚}领受了那位失落的宗徒的主教职。愿无人如此迟钝，无人如此无信，以致争论这点：\Person{玛弟亚}为自己赢得的是胜利，而非冤屈，因为他从主\Person{基督}的胜利中夺取了叛徒的战利品。那么，此后你们为何以更坏叛徒的继承者自居，声称主教的职位呢？\Person{犹达斯}在肉身中将\Person{基督}出卖给不信者；而你们却在精神上疯狂地将神圣福音出卖给亵渎的火焰。\Person{犹达斯}将立法者出卖给不信者；而你们，仿佛是出卖了他留下的一切，将天主的法律交给人类毁灭。然而，倘若你们爱慕法律，像年轻的玛加伯一样，你们本应为了天主的法律而欣然受死（如果人因主而死而获得不朽，这还能称为死亡的话）；因为关于那些兄弟，我们得知其中一位用这些充满信德的话回答那亵渎的暴君：\InnerQuote{你如狂怒般将我们带出此生；但世界的君王（祂永远为王，祂的国度无终）必将使我们这些为祂的法律而死的人复活，得享永生。}如果你们用火烧毁一个死者的遗嘱，难道不会被当作伪造遗嘱者而受罚吗？那么，你们这些烧毁我们天主和审判者最神圣法律的人，又会怎样呢？\Person{犹达斯}甚至在死时还悔恨他的行为；而你们不仅不悔改，反而成为我们这些持守法律者的迫害者和屠夫，你们是最邪恶的\OriginalTerm{叛徒}{traditors}。}\par

18. \Person{奥斯定}回答说：请看你们诽谤的言语与我们所坚持的真理之间有多么大的差别。请稍听片刻。看你们如何夸大了交出圣书的罪过，像某个诡辩的指控编造者一样，用最可憎的比喻将我们与叛徒\Person{犹达斯}相提并论。但当我以最简短的方式回答你们这一点时——我没有做你们所断言的事；我没有交出圣书；你们的指控是虚假的；你们永远无法证明它——那么所有那些浓烈的言辞烟雾岂不立即消散？或者你们或许试图证明你们所说的为真？那么，你们应该先做这件事；然后你们可以用你们所选的任何大量辱骂来攻击我们，如同对付已定罪的人一样。这是一个荒谬之处：请看第二个。\par

19. 你们自己在谈到预言\Person{犹达斯}被定罪时，用了这些话：\Quote{看，先知们的神力何等伟大，竟能将一切未来的事视如当下，使一个日后才出生的叛徒在多个世纪前就被定罪}；然而你们却没有看到，在那同样可靠的预言、确定不移的真理中，既然预言了其中一个门徒日后要出卖\Person{基督}，也预言了全世界日后要信仰\Person{基督}。你们在预言中留意出卖\Person{基督}的人，却在同一处忽略了\Person{基督}为之被出卖的世界。谁出卖了\Person{基督}？\Person{犹达斯}。他把祂出卖给谁？给犹太人。犹太人对祂做了什么？\Quote{他们刺穿了我的手和我的脚，}圣咏作者说，\Quote{我数尽我的骨骼，他们注视着我。他们分开我的衣袍，为我的衣服拈阄。}那么，那以如此代价买来的事物何等重大，我愿你们在圣咏稍后几节中读到：\Quote{大地四极都要记起，归向上主；万邦万族都要在祢面前敬拜。因为王国属于上主，祂治理万邦。}可是，谁能足以引证所有其他无数预言经文，这些经文都见证那注定要信仰的世界？然而，你们引用一个预言，因为你们在其中看到了出卖\Person{基督}的人；你们却没有在其中看见\Person{基督}因被出卖而买得的产业。这是第二个荒谬之处：请看第三个。\par

20. 在您的谩骂中的许多其他表达中，您说过：\Quote{若你用火焚烧死人的遗嘱，难道你不会因伪造遗嘱而受罚吗？那么，你们焚烧了我们天主和审判者最神圣的法律，又会怎样呢？} 在这些话中，您没有注意到那些确实本该打动您的事，即我们如何可能焚烧了遗嘱，却仍坚守在那遗嘱中所描述的继承中；但令人惊奇的是，您保存了遗嘱却失去了继承。那遗嘱中岂不是写着：\Quote{你向我请求，我必将外邦人赐你为产业，将地极赐你为基业}？请参与这继承，那么您就可以随意指控我关于遗嘱的事。这是何等的疯狂：您不敢把遗嘱投入火中，却反对立遗嘱者的话！另一方面，我们手中虽然持有教会和国家的记录，从中读到那些对立于塞西利安而祝圣对立主教的人，实际上是圣书的背叛者，但我们并不因此侮辱您，或用谩骂追赶您，或哀悼您手中圣页的灰烬，或将玛加伯的烈火折磨与您恐惧的亵渎相比较，说：\Quote{你应当将自己的肢体交于火焰，而不是交托天主的言语}。因为我们不愿如此荒谬，因他人的行为而起空泛的喧闹反对您，那些行为您要么一无所知，要么予以否认。但我们看到您与普世教会的共融分离（这是最严重的罪，且对全人类显而易见，是你们共同的罪），如果我想要夸大，时间会比言语更先耗尽。如果您想就此辩护，那只能通过指控全世界，但若指控成立，您只是提供进一步指控您自己的材料；若不能成立，它们对您毫无辩护之力。那么，您为何因圣书背叛之事向我自高，那件事若我们遵守不相互指控他人之罪的协议，就与您我无关；若协议不成立，它影响您甚于我？然而，即使不违反协议，我认为可以完全公正地说：凡不与全世界一同将自己交付于基督的人，应被视为与出卖基督的人同伙。\Quote{那么，}宗徒说，\Quote{你们就是亚巴郎的后裔，按照恩许成为继承人。}\BibleRef{迦 3:29} 他又说：\Quote{天主的继承人，与基督同为继承人。}\BibleRef{罗 8:17} 同一宗徒表明，亚巴郎的后裔属于万民，源于赐给亚巴郎的恩许：\Quote{地上万民都要因你的后裔蒙受祝福。}\BibleRef{创 22:18} 因此，我认为我只是要求我们片刻思考天主的遗嘱，它早已开启，并认为凡不找见自己与所背叛的那位同为继承人的人，就是背叛者自己的继承人；凡否认基督赎买了全世界的人，就属于那出卖基督的人。因为祂复活后显现给门徒，并把自己肢体给怀疑的人抚摸时，对他们说：\Quote{经上曾这样记载：基督必须受苦，第三天从死者中复活；并且必须从耶路撒冷开始，因祂的名向万邦宣讲悔改，以得罪之赦。}\BibleRef{路 24:46} 看，你们使自己远离了怎样的继承！看，你们反对的是怎样的继承人！难道一个在地上行走的人会放过基督，而一个坐在天上的人却反对祂的话？你们还不明白吗？你们指控我们的一切，都是在指控祂的话。基督徒世界曾被应许并相信：应许应验了，却被否认。我恳求你们，想想你们应因这种不敬受何惩罚。然而，若我不知道你们受了什么——若我没有看见，没有造成——那么你们今天，没有受我逼迫暴力的，请向我说明你们分离的理由。但你们很可能一再说那些话，它们除非你们证明，否则不影响任何人；即使你们证明，也与我无关。\par

% structure: p0031 source=h2 target=subsection reason=relative-heading-rank; base=h2

\subsection{第九章}

培提利安说：\Quote{因此，被这些过犯所困，你不能成为真正的主教。}\par

奥古斯丁答：什么过犯？您展示了什么？您证明了什么？如果您证明了某人的指控，那与亚巴郎的后裔有何关系？万民都因那后裔蒙受祝福。\par

% structure: p0034 source=h2 target=subsection reason=relative-heading-rank; base=h2

\subsection{第十章}

培提利安说：\Quote{宗徒曾迫害过任何人吗？或者基督曾出卖过任何人吗？}\par

24. 奥古斯丁回答说：我确实可以说，撒殚本身比所有恶人都更坏；然而宗徒把一个人交给撒殚，为的是败坏他的肉体，使他的精神在主耶稣的日子得以得救。\BibleRef{格前 5:5} 同样，他也把其他人交出去，他说：\Quote{我已把他们交给撒殚，使他们学会不再亵渎。}\BibleRef{弟前 1:20} 主基督也用鞭子把不敬虔的商人赶出圣殿；与此相关，我们也看到圣经的见证说：\Quote{我为祢的殿所耗的忧虑将我吞没了。}\BibleRef{若 2:15} 因此，我们确实发现宗徒把人交给定罪，而基督是迫害者。我本可以说这一切，并使你陷入不小的激动和不安，以至于你不得不追问的不是那些受苦者的抱怨，而是造成痛苦之人的意图。但你不要为此烦恼；我不说这个。我确实说，如果有什么不应该做的事做在了你身上，那与亚巴郎的后裔无关，这后裔存在于万民中，也许是由主庄稼中的糠秕做的，尽管这样，这糠秕仍存在于万民中。因此，你当为你自己的分离做解释。但首先，你要想一想，你们中间有什么样的人，你不想被他们责备；然后看看你行事多么不公正，当你把别人的行为丢在我们脸上时，即使你证明了你的主张，也是这样。因此，将会发现你的分离没有任何理由。\par

% structure: p0037 source=h2 target=subsection reason=relative-heading-rank; base=h2

\subsection{第11章。}

25. \Person{培提里安}说：\Quote{然而有人会说：\InnerQuote{我们不是\OriginalTerm{背教者}{traditor}的儿子。}任何人若是仿效那个人的行为，他就是那人的儿子。因为那些在父母的性情和行为上肖似父母的人，无疑就是他们的儿子，不仅因为他们是他们的骨肉，而且因为他们的品性和行为。}\par

26. \Person{奥古斯丁}回答说：刚才你说的没有一句与我们相悖，现在你甚至开始说一些有利于我们的话。因为你的这个主张约束你如下：如果你今天不能指证我们——即你与之争论的人——是\OriginalTerm{背教者}{traditor}和杀人犯，以及你指控我们的任何其他罪过，那么你对过去的人可能证明的任何此类指控，都不能伤害我们。因为我们不能成为行为与我们不相肖的人的儿子。你且看你自己承诺了什么。如果你侥幸指证某个甚至我们同时代、与我们生活在一起的人犯了此类罪过，那也绝不损害在亚巴郎的后裔中蒙福的万民，你与他们分离就被证明犯有亵渎之罪。因此，除非你（这是完全不可能的）认识世上存在的所有人，不仅熟悉他们所有的品性和行为，而且证明他们像你描述的那样坏，你没有理由以不知什么样的出身来指责整个世界的圣徒们，你证明他们与那些出身相似。即使你能证明那些品行不端的人与品行端正的人共同领受圣事，那也一点帮不了你。首先，因为你们自己应该看看与你们一同庆祝那些圣事的人——你们给他们圣事、从他们领受圣事、也不愿他们被当作责备扔回你们脸上的人。此外，如果所有效法\Person{犹大}（即宗徒中的魔鬼）行为的人都是\Person{犹大}的儿子，那么为什么我们不把那些让这样的人分沾——不是分沾他们自己的行为，而是分沾主的圣事——的人称为宗徒的儿子呢？因为宗徒们也与那个叛徒一同领受了主的晚餐。在这方面，你们与他们非常不同，你们把那些为保存合一而努力的人所做的事，扔回去责备他们，而你们所行的却是撕裂合一。\par

% structure: p0040 source=h2 target=subsection reason=relative-heading-rank; base=h2

\subsection{第12章。}

27. \Person{培提里安}说：\Quote{主耶稣对犹太人论及自己说：\InnerQuote{我若不做我父的工，你们就不要信我。}}\BibleRef{若 10:37}\par

28. \Person{奥古斯丁}回答说：我前面已经回答过了，这是真的，并且对我们有利，对你们不利。\par

% structure: p0043 source=h2 target=subsection reason=relative-heading-rank; base=h2

\subsection{第13章。}

29. \Person{培提里安}说：\Quote{他屡次这样责备说谎者和撒谎者：\InnerQuote{你们是魔鬼的儿女，因为他从起初就是毁谤者，没有站在真理中。}}\par

30. \Person{奥古斯丁}回答说：我们不习惯说\Quote{他是毁谤者}，而是说\Quote{他是杀人者}\BibleRef{若 8:44}。但我们问魔鬼如何从起初就是杀人者；我们发现他杀死了第一个人，不是用刀，也不是用任何身体的暴力，而是引诱他犯罪，从而将他逐出乐园的幸福。那么，乐园现在由教会代表。所以，那些把人从教会中拉走而杀死他们的人，就是魔鬼的儿女。但正如我们藉天主的言语知道乐园的位置，现在也藉基督的言语知道教会在哪里：他说：\Quote{从耶路撒冷开始，直到万民。}因此，凡把一个人从那个整体中分离出来，放到任何一个部分的人，就被证明是魔鬼的儿女和杀人者。但是，再看你自己使用\Quote{他是毁谤者，没有站在真理中}这句话时的应用。因为你因别人的罪过而控告全世界，尽管你更有能力控告那些人而不是定他们的罪；而且你没有站在基督的真理中。因为他说教会是\Quote{从耶路撒冷开始，直到万民}，而你说教会是在\Person{多纳徒}的党派中。\par

% structure: p0046 source=h2 target=subsection reason=relative-heading-rank; base=h2

\subsection{第14章。}

31. Petilianus said: \Quote{在第三处，祂也同样用这个名称称呼迫害者的疯狂：\BibleQuote{毒蛇的种类！你们怎能逃避地狱的处罚？为此，我派遣先知、贤士和经师到你们这里来；有的你们要杀死、钉死，有的你们要在你们的会堂里鞭打，从这城追逼到那城，叫世上所流的一切义人的血，都归到你们身上，从义人\Person{亚伯}的血起，直到你们在\Place{圣所}与\Place{祭台}之间所杀的\Person{贝勒基雅}的儿子\Person{匝加利亚}的血。}} \BibleRef{玛 23:33}\par

32. Augustine answered: 如果我说这是指像你们这样品性的人，你们会回答说：\Quote{证明给我们看。}那么，你们证明了什么呢？或者，如果你们认为只要说出口就算证明，那就没有必要重复同样的话。你们从我们口里听到同样的判决，就把它加在你们自己身上吧。你们难道没有看到，如果这算证明，我也证明了？但我仍希望你们知道证明的真正含义。事实上，我甚至不需要从外部寻找证据来证明你们是毒蛇。因为请确信，正是这点在你们身上显明毒蛇的本性：你们口中没有真理的基础，只有诽谤的毒液，正如经上所记：\BibleQuote{蝮蛇的毒液在他们的嘴唇下。}因为这一点可能被任何人随意用于别人，就好像被问：\Quote{在谁的嘴唇下？}他立刻加上：\BibleQuote{他们的口充满咒骂和苦毒。}所以，当你们对散布在全世界的、你们一无所知的人说这样的话，许多人甚至从未听说过\Person{塞希利安}或\Person{多纳图}的名字，而你们没有听到他们在沉默中回答：\Quote{我们说的与你们无关；我们从未见过，从未做过；我们完全不明白你们在说什么}——既然你们除了说这些你们根本无法证明的话之外，别无他求，你们怎能不承认你们的口充满咒骂和苦毒呢？因此，你们看，除非你们能证明全世界所有民族中的所有基督徒都是\Emph{\OriginalTerm{背信者}{traditors}}、凶手，而不是基督徒，否则你们能否证明自己不是毒蛇？不，实际上，即使你们能够知道并向我们展示世上每个人的生活和行为，但在你们能做到之前，鉴于你们如此轻率行事，你们的口就是毒蛇的口，你们的口充满咒骂和苦毒。现在向我们证明，如果你们能，我们杀了、钉死了、或在会堂里鞭打了哪一位先知、贤士或经师。看看你们花了多少力气却根本无法证明\Person{多纳图}和\Person{马库路斯}是先知、贤士或经师，因为他们实际上根本不是。但即使你们能证明这一点，你们又能在多大程度上证明他们是被我们杀害的呢？因为我们自己甚至都不认识他们。而你们用有毒的口诽谤的全世界就更不认识了？或者，你们从哪里能证明我们有像杀害他们的人那样的精神呢？而事实上你们甚至无法证明他们是被任何人杀害的。仔细考虑所有这些点，看你们能否证明其中任何一点，无论是关于全世界，还是为了全世界的满足——而你们对全世界坚持不懈的诽谤中，你们恰恰证明了那些你们虚假散布给世界的指控，在你们身上是真实的。\par

33. 此外，即使我们想要证明你们是杀害先知的人，收集所有关于杀戮的证据——你们愤怒的Circumcelliones领袖和一群被酒和疯狂点燃的人从分裂开始直到现在不仅犯了而且继续犯下的——将是一项太长的任务。以最近的情况为例。让神圣的言语被提出，这些言语通常在双方手中。让我们认为那些与先知的话语相矛盾的人是杀害先知的人。还有什么比这更有学问的定义呢？什么能更快地得到证明？你们用剑刺穿先知的尸体不如用舌头试图摧毁先知的话语残忍。先知说：\BibleQuote{普世万民都要想起上主，都归向他。} 看这如何正在实现，如何正在完成。但你们不仅闭耳不信所说的话，反而疯狂地伸出舌头来反对已经正在实现的事。亚巴郎听到了应许：\BibleQuote{地上万民都要因你的后裔蒙受祝福，}\BibleRef{创 22:18} 并且\BibleQuote{他相信了上主，上主就以此算为他的正义。}\BibleRef{罗 4:3} 你们看见事实已经完成，却大声反对；你们不愿这算为你们的不义，正如它理所当然会算为不义一样，即使你们的不信不是发生在预言应验之后，而仅仅是在预言发出之时。不仅如此，你们不仅不愿这算为你们的不义，甚至你们因这种不虔诚而遭受的惩罚，你们却愿它算为你们的正义。或者，如果你们的行为不是对先知的迫害，因为你们的工具不是刀剑而是舌头，那么圣经上何以会说：\BibleQuote{人子啊，他们的牙齿是枪和箭，他们的舌头是快刀？} 但我哪来的时间从所有先知那里收集所有关于教会分散在全世界的见证？——所有你们试图通过反对来摧毁和使之无效的见证。但你们被捉住了，因为\BibleQuote{他们的声音传遍了全地，他们的言语传到了地极。} 然而，我将从主的口中提出这一句话，他是见证人的见证。他说：\BibleQuote{凡在梅瑟法律、先知和圣咏上关于我所记载的，都必须应验。} 让我们从他本人那里听见这些是什么：\BibleQuote{于是他开了他们的理智，叫他们理解圣经，又对他们说：经上这样记载：基督必须受苦，第三天从死人中复活，并且要因他的名从耶路撒冷开始，向万邦宣讲悔改和赦罪。}\BibleRef{路 24:44} 看，在梅瑟法律、先知和圣咏上关于主所写的是什么。看主自己如何启示他自己和教会，使自己显现，发出关于教会的应许。但对你们来说，你们这样抗拒这些明显的证据，既然不能摧毁它们，便试图歪曲它们。如果你们遇到先知们的尸体，你们会做什么？你们对先知们的言论如此疯狂，甚至不听从主，当他在实现、显明并解释先知时？因为你们岂不是尽力试图杀害主自己吗？因为你们连他也不屈服。\par

% structure: p0050 source=h2 target=subsection reason=relative-heading-rank; base=h2

\subsection{第十五章。}

\Quote{达味也曾以下列的话论你们为迫害者：\InnerQuote{他们的咽喉是敞开的坟墓；用他们的舌头行欺骗；他们的嘴唇下有毒蛇的毒汁；他们的口充满咒骂与苦楚；他们的脚飞快地流血；在他们的路上是毁灭与不幸，和平之道他们不知道；在他们眼前没有天主的敬畏。所有作恶的人，难道没有知识，他们吞食我的百姓如同吃饼？}}\par

35. \Person{奥斯定}回答说：他们的咽喉是敞开的坟墓，他们藉谎言呼出死亡。因为\BibleQuote{说谎的口杀害灵魂}\BibleRef{智 1:11}。但是，如果基督所说的——祂的教会要从耶路撒冷开始遍及万民——再真实不过，那么你所说的——教会在多纳图派中——就再虚假不过了。而那些欺骗人的舌头，正是那些明知自己的行为，不仅自称义人，还自称使人成义之人的舌头——这只有那位\BibleQuote{使不虔敬的人成义}\BibleRef{罗 4:5}的那一位才如此说，因为祂\BibleQuote{是正义的，也是那使人成义者}\BibleRef{罗 3:26}。至于毒蛇的毒汁和充满咒骂与苦毒的嘴，我们已经说得够多了。然而你们自己说过，马克西米安的支持者有快脚步流人血，这是你们全体会议的判决所见证的，该判决在行省总督府和国家档案中多次被引用。但是，据我们所知，他们从未在身体上杀害过任何人。显然，你们明白，灵魂的血是通过分裂的剑在精神谋杀中被流出的，而你们在马克西米安身上谴责了这一点。那么，看看你们的脚是否不流人血，当你们将人们从全世界的合一中割裂时——如果你们对马克西米安的支持者这样说是因为他们割断了多纳图派的一些人，那么你们是否也有同样的罪？难道我们不知道和平的道路吗？我们努力在和平的纽带中保持圣神的合一。而你们却拥有那样的知识吗？你们拒绝基督复活后与门徒所讲的话，那番话以\Quote{愿你们平安}的问候开始，那样充满和平，以至于你们被证明对祂说的无非是：\Quote{你所说万民的合一乃是假的；我们所说万民的冒犯才是真的。}谁会在敬畏天主的情况下说出这样的话呢？因此，看吧，若你们天天说这样的话，岂不是在试图摧毁散居世界各地的天主子民，如同吞吃面包一样吞吃他们吗？\par

% structure: p0053 source=h2 target=subsection reason=relative-heading-rank; base=h2

\subsection{第十六章}

36. \Person{培提利安}说：\Quote{主基督也警告我们说：\BibleQuote{你们要提防假先知，他们来到你们这里，外面披着羊皮，里面却是残暴的狼；你们可凭他们的果实分辨他们。}\BibleRef{玛 7:15}}\par

37. \Person{奥斯定}回答说：如果我问你，你凭哪些果实认出我们是残暴的狼，你一定会回答，用其他人的罪来控告我们，而且这些罪从未在那些被指控的人身上得到证实。但如果你问我，我们凭哪些果实认出你们更是残暴的狼，我就要用分裂的罪来控告你们，你们会否认，但我马上会证明；因为事实上，你们不与世上所有的民族共融，也不与宗徒们辛劳建立的教会共融。对此你会说：\Quote{我不与\OriginalTerm{叛教者}{traditors}和杀人者共融。}亚巴郎的后裔回答你说：\Quote{这些都是你提出的指控，要么不真实，要么与我无关。}但我暂时把这些放在一边；你同时请指出教会。现在，那声音将在我耳中响起，即主所指示要在假先知中避免的，他们炫耀各自的派别，竭力使人们脱离天主教会，说：\Quote{看，基督在这里，或在那里。}但是，你认为基督的真羊群是如此缺乏理智——他们被告知\BibleQuote{不要相信}\BibleRef{玛 24:23}——以至于他们会听从狼说\Quote{看，基督在这里}，却不听从牧者说\Quote{从耶路撒冷开始，遍及万民}吗？\par

% structure: p0056 source=h2 target=subsection reason=relative-heading-rank; base=h2

\subsection{第十七章}

38. \Person{培提利安}说：\Quote{就是这样，就是这样，你这邪恶的迫害者，不论你披上怎样的义袍隐藏自己，不论你以和平之名用亲吻发动战争，不论你以合一之名企图陷害人类——你，至今仍在欺骗和欺诈，你是魔鬼的真正儿子，用你的本性显明你的出身。}\par

39. \Person{奥斯定}回答说：请回想，这些话是我们反对你们的；并且，为使你们知道这些话对我们中谁更为恰当，请记起我之前所说的。\par

% structure: p0059 source=h2 target=subsection reason=relative-heading-rank; base=h2

\subsection{第十八章}

40. 培提里安说：\Quote{这本来也不足为奇，你竟无凭无据擅自窃取主教的称号。这正是魔鬼的习惯，他宁愿选择一种欺骗方式，借此篡夺圣洁语词的含义，正如宗徒向我们宣告说：\BibleQuote{这并不希奇，因为他［撒殚］说：\InnerQuote{因为撒殚自己也装作光明的天使，所以他的仆役装作正义的仆役，也不算什么大事。}}\BibleRef{格后 11:14}所以，你若谎称自己是主教，也不足为奇。因为那些堕落的天使，那些贪爱世间少女、因肉身的败坏而败坏了自身的天使，虽因失掉了神圣的德性而不再是天使，却仍保留着天使的名号，常自以为天使；然而，他们既从侍奉天主的职务中被解除，便由其品格的相似转入魔鬼的军队，正如伟大的天主所言：\InnerQuote{我的神不会常与人相争，因为人也是血肉。}\BibleRef{创 6:3}对这些有罪者和你，主基督必要说：\InnerQuote{你们这些被咒骂的，离开我，到那为魔鬼和他的天使预备的永火里去吧！}\BibleRef{玛 25:41}若没有邪恶的天使，魔鬼就没有天使；关于这些天使，宗徒说，在复活的审判中，他们将被圣人们定罪：\InnerQuote{你们不知道吗？我们要审判天使。}\BibleRef{格前 6:3}若他们是真正的天使，人就没有权柄审判天主的天使。同样，那六十位宗徒，当十二位单独与主基督同在时，因背弃信仰而离去，至今仍被不幸的人们视为宗徒，以致摩尼及其余诸派从他们出发，在许多魔鬼异端中缠累了许多灵魂；他们毁了这些灵魂，为要将其网罗其中。的确，堕落的摩尼，即便他曾堕落，也不该算在那六十人之内，除非我们能在十二宗徒中找到他的名字，或他是由基督的声音被立，当玛弟亚被选填补叛徒犹达斯的空缺时，或是另一位第十三位的保禄，他自称是宗徒中最后的一位，特意说明任何晚于他的人不得视为宗徒。因为他的话是：\InnerQuote{我是宗徒中最后的一位，我原不配称为宗徒，因为我曾迫害过天主的教会。}\BibleRef{格前 15:9}你们不可在此自夸：他原是犹太人，做了这事。你们作为外邦人，也可以毁灭我们。因为你们未经许可便发动战争，我们却不能反过来攻击你们。因为你们渴望在我们杀害你们之后仍能活着；但我们的胜利，或是逃脱，或是被杀。}\par

41. 奥斯定回答说：你看你如何引用了圣经的见证，或如何理解了它，这根本无关于当前争论的要点。因为你所提出的一切，不过是为了证明存在假主教，正如存在假天使和假宗徒一样。我们也很清楚，有假天使、假宗徒、假主教，而且正如真宗徒所说，还有假弟兄\BibleRef{格后 11:26}。但是，既然双方都可能彼此提出这样的指控，就需要某种程度的证据，而非空洞的言辞。但若你愿看见哪一方更真实地犯有虚妄之罪，请回想我们先前所说，你在那里会看到这一点的阐述，以免我们重复同样的话而使读者厌烦。然而，普世教会又怎会在乎你关于它的糠秕所说的任何话呢？这糠秕搀杂在普世教会之中；或者你在摩尼及其它魔鬼异端上所讲的，又怎能影响教会呢？因为若麦子不受任何关于仍同它搀杂的糠秕的言论所影响，那么，那些早已且公开与基督的肢体分离的怪物，又怎能影响散布于普世世界的基督肢体呢？\par

% structure: p0062 source=h2 target=subsection reason=relative-heading-rank; base=h2

\subsection{第十九章}

42. 培提里安说：\Quote{主耶稣基督命令我们说：\InnerQuote{有人在这城里迫害你们，你们就逃到另一城里去；若在那城里迫害你们，就再逃到第三城。因为我实实在在告诉你们：你们还没有走完以色列的城邑，人子就来了。}\BibleRef{玛 10:23}若祂在犹太人和外邦人的情形下给我们这个警告，那么你这自称为基督徒的人，就不该效法外邦人的可怕行径。难道你事奉天主的方式，就是要让我们死在你的手下吗？你错了，你错了，你若不幸持有这种信仰的话。因为天主不会有屠夫作祂的司铎。}\par

43. \Person{奥古斯丁}回答说：从一地逃往另一地以躲避迫害的面目，这并不是作为诫命或许可而加给异端者或分裂者，正如你们这样的人；而是加给福音的宣讲者，你们所反抗的正是他们。这一点我们可以轻易证明如下：你们现在在你们自己的城里，并没有人迫害你们。因此你们必须出来，对你们的分离作出交代。因为不能主张，正如肉身的软弱在迫害的暴力前屈服时可以原谅，同样真理也应当向虚假屈服。此外，如果你们遭受迫害，为什么不从你们所在的城里退出，以便履行你们从福音中引用的训示呢？但如果你们没有遭受迫害，为什么不愿意回答我们呢？或者如果事实是你们害怕一旦回答了就会遭受迫害，那么你们如何效法那些宣讲者的榜样呢？他们对那些人说过：\Quote{看，我派遣你们好像羊进入狼群中。}对他们还进一步说过：\Quote{你们不要害怕那些杀害肉身而不能杀害灵魂的。}你们又如何逃避违反宗徒伯多禄训诫的指控呢？他说：\Quote{你们要时常准备好，对每一个问你们心中所怀希望的理由的人，作出答辩。}\BibleRef{伯前 3:15}最后，为什么你们一有机会就热衷于以最暴力的骚乱骚扰天主教会，这已由无数简单事实的实例所证明？但你们说你们必须保卫你们的地方，并且你们用棍棒、屠杀以及任何其他手段抵抗。既然如此，你们为什么没有听从主的声音呢？祂说：\Quote{我却对你们说，不要抵抗恶人。}\BibleRef{玛 5:39}或者，允许在某些情况下，用暴力抵抗暴力的人是合理的，并且不违反我们从主领受的诫命\Quote{我却对你们说，不要抵抗恶人}，那么为什么一个虔敬的人不能藉合法正当设立的权威，将不虔敬的人，或义人将不义的人，从非法篡夺或冒犯天主而占据的席位上驱逐出去呢？因为你们不会说，假先知在\Person{厄里亚}手中遭受迫害，与\Person{厄里亚}遭受最邪恶的国王的迫害是同一回事。\EditorialNote{列王纪上第十八章}也不会因为主被祂的迫害者鞭打，就认为那些被祂本人用鞭子赶出圣殿的人可以与祂的受苦相提并论。因此，我们应当承认，没有其他需要解决的问题，除了你们是否在从普世教会的共融中分离出来时是虔敬还是不虔敬。因为如果发现你们的行为是不虔敬的，那么你们就不会惊讶于天主的仆人不会缺乏，藉他们你们可能被鞭打，因为你们所受的迫害不是来自我们，而是如经上所记，出自他们自己的可憎之事。\BibleRef{智 12:23}\par

% structure: p0065 source=h2 target=subsection reason=relative-heading-rank; base=h2

\subsection{第二十章。}

44. \Person{佩蒂利安}说：\Quote{\InnerQuote{扫禄，扫禄，你为什么迫害我？你向刺棒踢去是难的。}\BibleRef{宗 9:4}他那时被称为扫禄，为的是在圣洗中后来领受他真实的名字。但你们如此频繁地在祂司铎的人身上迫害基督，却不算难，尽管主自己呼喊：\InnerQuote{不可触摸我的受傅者。}清点所有圣人的死亡，你们就多次谋杀了基督，祂活在每个人内。最后，如果你们没有犯亵圣罪，那么一个圣人不能是杀人者。}\par

45. \Person{奥古斯丁}回答说：你们为自己辩护，以免受那些与\Person{马克西米安}追随者一起从你们中分离出去的人在你们一方手下所受的迫害的指控，这样你们就会找到我们的辩护。因为如果你们说你们没有犯下这样的行为，我们只需向你们宣读代行执政官省和国家的记录。如果你们说你们迫害他们是正确的，为什么你们不愿意自己遭受同样的事呢？如果你们说：\Quote{但我们没有引起分裂}，那就让这受到调查，在决定是否如此之前，没有人可以对迫害者提出指控。如果你们说甚至分裂者也不应该遭受迫害，我问，是否他们也不应该被赶出那些他们设陷阱迷惑弱者的圣堂，即使是由合法设立的权威所做？如果你们说这事也不该做，那就先把圣堂归还给\Person{马克西米安}的追随者，然后再与我们讨论。如果你们说那是正确的，那么看看那些反抗合法设立权威的人应该遭受什么，因为他们反抗天主的命令。因此宗徒明确地说：\Quote{因为他不是白白佩剑，他是天主的仆人，是向作恶的人施行愤怒的报复者。}但即使经过精密的调查后发现，在公开审判后也不应该对分裂者施加任何刑罚，或将他们从因背叛和欺骗而占据的位置上赶走；如果你们说你们对\Person{马克西米安}的追随者中有一部分人遭受这样的对待感到烦恼——为什么不从主的整个田地，即从全世界，主的麦子更大声地呼喊说：\Quote{我们自己完全不受我们中间莠子和糠秕所行之事的影响，因为那是违背我们的意愿的？}如果你们承认，澄清你们责任的因素在于你们一方所做的一切恶事都是违背你们意愿的，那么你们为什么还要分离呢？如果你们不从不义者中分离出来，是因为各人承担自己的负担，那么为什么你们从那些你们认为（或声称认为）不义的人中分离出来呢？是为了让你们都平等地分担分裂的负担吗？\par

46. 当我们问你，你们中哪一位你能证明曾被我们杀害时，我的确记不得有任何皇帝颁布的法律，规定应将你们处死。你们最常引用来归咎于我们的那些死者——\Person{马尔库鲁斯}和\Person{多纳图斯}——提出了一个重大问题：他们是自己跳下悬崖的（如同你们的教导毫不犹豫地以日常事例来鼓励的那样），还是由某个权威的真正命令扔下去的？因为\OriginalTerm{巡行派}{Circumcelliones}的首领按他们的习惯自尽的说法难以置信，那么罗马当局竟然能判他们一种不合常规的刑罚就更难以置信了！因此，在考虑这件你们认为特别可憎的事时，假设你们所说属实，其中哪一点关系到主的麦子呢？让飞到外面的糠秕去控告尚且留在里面的糠秕吧，因为不可能在最后审判簸扬之前把它们全部分开。但若你们所言为假，那么糠秕仿佛被一阵轻微的分裂之风刮走时，甚至用虚假的控告攻击主的麦子，又有什么奇怪呢？因此，在考虑所有这类可憎的指控时，基督的麦子——它被命令在整个田地即整个世界与稗子一同生长——用自由无畏的声音回答你们：你们若不能证明你们所说的话，那就不涉及任何人；你们若能证明，仍不涉及我。其结果是，任何人因稗子或糠秕的过错而离开了麦子的合一，就无法为谋杀之罪辩护，这罪仅因分裂和分党的过错而犯了，正如圣经所说：\BibleQuote{凡恨自己弟兄的，就是杀人者。}\BibleRef{若一 3:15}\par

% structure: p0069 source=h2 target=subsection reason=relative-heading-rank; base=h2

\subsection{第二十一章}

47. \Person{Petilianus}说：\Quote{因此，正如我们所说，主基督喊道：\InnerQuote{扫禄，扫禄，你为什么迫害我？你向刺棒踢去是难的。他说：主，你是谁？主说：我是纳匝肋的基督，你所迫害的。他战战兢兢地惊讶道：主，你要我做什么？主对他说：起来，进城去，必有人告诉你你该做的事。}接着立即说：\InnerQuote{扫禄从地上起来；当他睁开眼时，什么也看不见。}你看，由于疯狂而来的盲目怎样惩罚性地遮蔽了迫害者眼中的光明，除非借着洗礼，绝不能再被驱除！因此，我们看他进城后做了什么。经上说：\InnerQuote{阿纳尼雅进了扫禄的家，按手在他身上说：扫禄弟兄，在你来的路上向你显现的主耶稣打发我来，叫你能看见，并充满圣神。立刻有鳞片从他眼中掉下；他立刻看见了，起来，受了洗。}\BibleRef{宗 9:4} 既然保禄借着洗礼脱离了迫害的过犯，重新获得了无罪的眼睛，为什么你这迫害者和\OriginalTerm{背教者}{traditor}，被虚假的洗礼弄瞎了，却不接受你所迫害者的洗礼呢？}\par

48. 奥古斯丁答：你不能证明我——你希望重新为我施洗的人——是迫害者或\OriginalTerm{背教者}{traditor}。即使你能向任何人证明这一指控，但迫害者和背教者不应重新受洗，如果他已领受过基督的洗礼。因为保禄必须受洗的原因是他从未接受过那种洗礼。因此你选择插入关于保禄的段落与你和我们争论的案件没有丝毫相似之处。但若你没有插入此段，你就没有地方放你那幼稚的演说：\Quote{看，疯狂惩罚而来的盲目，除非借着洗礼，绝不能再被驱除！}因为有人能用更大的力量向你呼喊：\Quote{看，疯狂惩罚而来的盲目，它在西满而不在保禄身上找到相似，即使你领受了洗礼也未从你身上被驱除！}因为如果迫害者应由他们所迫害的人施洗，那么就让\Person{普里米安}由\Person{马克西米安}的追随者施洗，因为普里米安曾极其热切地迫害他们。\par

% structure: p0072 source=h2 target=subsection reason=relative-heading-rank; base=h2

\subsection{第二十二章}

49. \Person{Petilianus}说：\Quote{有人可能主张，基督曾对宗徒们说——正如你们经常引用对抗我们的——：\InnerQuote{洗过澡的，只需洗脚，就全身洁净了。}现在如果你们详尽讨论这些话，你们必受随后经文的约束。因为以下是祂的原话：\InnerQuote{洗过澡的，只需洗脚，就全身洁净了；你们是洁净的，但不都是。}祂这话是因为犹达斯，即那要出卖祂的人；因此祂说，你们不都是洁净的。\BibleRef{若 13:10} 因此，凡犯了背叛之罪的，与你们一样，已丧失了洗礼。再者，基督的背叛者自己被定罪后，主就这样更充分地证实了祂对十一宗徒所说的话：\InnerQuote{现在你们因我给你们讲的话，已是洁净的了。你们住在我内，我也住在你们内。}\BibleRef{若 15:3} 并且祂又对同这十一人说：\InnerQuote{我把平安留给你们，我将我的平安赐给你们。}\BibleRef{若 14:27} 既然这些话是对十一宗徒说的——当时背叛者已被定罪，如我们所见的——你们同样，作为\OriginalTerm{背教者}{traditors}，便同样丧失了平安和洗礼。}\par

50. \Person{奥斯定}答说：因此，如果每一个\OriginalTerm{叛教者}{traditor}都已经丧失了他的圣洗，那么，每一个被你们施洗后又成为\OriginalTerm{叛教者}{traditor}的人，都应重新受洗。如果你们不这样做，你们自己就充分证明了以下说法的虚假：\Quote{因此，凡犯有叛国罪的人，都像你们一样丧失了他的圣洗。}因为如果他已丧失，就让他回来，重新领受；但如果他回来却不领受，显然他并没有丧失。再者，如果对宗徒们所说\BibleQuote{现在你们洁净了}和\BibleQuote{我将我的平安赐给你们}是因为叛徒已经离开房间，那么，祂在出去以前分给众人的那如此伟大圣事的晚餐，岂不是洁净并能赐予平安的吗？如果你们闭目不顾真理，胆敢这样说，我们除了更大声呼喊还能做什么？看，盲目如何降临作为对那些人的惩罚，他们愿如宗徒所说的，\BibleQuote{教师们却不懂自己所说的，也不明自己所主张的。}\BibleRef{弟前 1:7}然而，若非盲目阻挡了他们的固执，你们理解并看到主在犹达斯面前并没有说\BibleQuote{你们还没有洁净}，而是说\BibleQuote{现在你们洁净了}，并不是一件难事。然而，祂加上了\BibleQuote{但不是全部}，因为那里有一个不洁净的人；但如果他的存在污染了别人，就不会对他们宣告\BibleQuote{现在你们洁净了}，而是如我先前所说，你们还没有洁净。但是，犹达斯出去后，祂对他们说\BibleQuote{现在你们洁净了}，并没有加上\BibleQuote{但不是全部}这句话，因为那个在他们中间被说\BibleQuote{你们已经洁净，但不是全部}的人已经离去，因为那里有一个不洁净的人。因此，主借此话宣告，在领受同一圣事的同一群人中，某些成员的不洁不能伤害洁净的人。当然，如果你们认为我们当中有像犹达斯这样的人，你们可以把\BibleQuote{你们是洁净的，但不是全部}这句话应用在我们身上。但这并不是你们所说的；你们却说，因为有些人不洁净，所以我们全都不洁净。主在犹达斯面前并没有对门徒说这话，因此，谁这样说，就没有从善师那里学到祂所说的是什么。\par

% structure: p0075 source=h2 target=subsection reason=relative-heading-rank; base=h2

\subsection{第二十三章}

51. \Person{白提利安}说：\Quote{但是，如果你们说我们施洗两次，实际上倒是你们这样做，你们杀害受过洗的人；我们这样说，不是因为你们给他们施洗，而是因为你们使他们每一个人通过被杀的行动，在自己血中受洗。因为水或圣神的洗礼，当殉道者的血被榨出时，就仿佛加倍了。同样，我们的救主自己，在最初受若翰的洗礼后，声称祂必须再次受洗，这次不是用水或圣神，而是用血的洗礼，即受苦的十字架，正如经上所载：\BibleQuote{有两个门徒，即载伯德的儿子，来到祂跟前说：\InnerQuote{主，当祢进入祢的国时，请使我们一个坐在祢右边，一个坐在祢左边。}耶稣对他们说：\InnerQuote{你们求的是艰难的事：你们能饮我将饮的杯吗？你们能受我将受的洗吗？}他们回答说：\InnerQuote{我们能。}耶稣便对他们说：\InnerQuote{你们固然要饮我的杯，也要受我所受的洗。}}\BibleRef{谷 10:35}等等。如果这是两种洗礼，那么你们因恶意而赞扬我们，我们必须承认。因为当你们杀害我们的身体时，我们就举行第二次洗礼；那就是我们以自己的洗礼和血受洗，如同基督一样。惭愧，惭愧，你们这些迫害者。你们使殉道者像基督一样，在真正的水洗之后，又洒上了血的洗礼。}\par

52. \Person{奥斯定}答说：首先，我们立即回答，我们没有杀害你们，而是你们自己以真正的死亡杀害了自己，当你将自己从活生生的合一之根上砍下来时。其次，如果所有被杀的人都在自己血中受洗，那么所有强盗、所有不义、不虔、可咒骂的人，他们被法律判死，都应被视为殉道者，因为他们是在自己血中受洗。但是，如果只有那些为义而被杀的人是在自己血中受洗，因为天国是他们的，\BibleRef{玛 5:10}你们已经看到，首先要问的是你们为何受苦，然后才问你们受什么苦。那么，为什么你们在尚未证明自己行为的义德之前就鼓起腮帮子？为什么在你们的品格被认可之前，你们的舌头就先发出声音？如果你们造成了分裂，你们是不虔敬的；如果你们不虔敬，当你们因不虔敬而受罚时，你们就作为亵圣者而死；如果你们作为亵圣者而死，你们如何在自己的血中受洗？或者你们说：\Quote{我没有造成分裂？}那就让我们查究此事。为什么在证实你们的案件之前就大声呼喊？\par

53. 或者你辩称：\Quote{即使我犯了亵渎之罪，你也不该杀死我？} 这是一个问题，关乎我行为之严重性——你从未如实证明过；另一个问题关乎你们所谓的血洗——那是你们自夸的根源。因为我从未杀过你，你也未能证明你是谁杀的。即使你能证明，也与我毫无关系，无论杀你的是谁，无论他是靠着主合法赐予的权力公正行事，还是像主的庄稼中的糠秕一样，出于邪恶的欲望犯下谋杀之罪；正如你与那个人毫无关系一样，那人近年来以不堪忍受的暴政，甚至带了一队士兵——不是因为他怕谁，而是要让所有人都怕他——压迫寡妇，毁掉学生，出卖他人的产业，废除他人的婚姻，策划变卖无辜者的财产，与哀哭的原主分赃。我似乎是在凭空捏造这些事，但我不提名字，大家也清楚我说的是谁。如果这一切千真万确，那么正如你与此无关一样，我们也与你所说的一切无关，即使你所说的是真的。但如果你的那位同僚确实是个公义无辜的人，却被谣言诽谤，那么我们同样也要学会不要相信那些关于无辜者的传闻，仿佛他们交出过圣经或杀害了同伴似的。此外，我指的是一个与你们同住的人，你们惯常举行盛大集会庆祝他的生日，在圣事中与他行平安之吻，将圣体放在他手中，又伸手从他手中领受；当整个非洲都在叹息时，他的耳朵是聋的，你们却不敢以直言冒犯他；因为有人间接地恭维他，说他在伯爵那里有一位神作同伴，你们中有人就被捧上了天。但你们指责我们，说的是我们从未与之同住、从未见过其面的人的行为，那些人活着时我们或是孩童，或许尚未出生。那么，你们如此不公、如此悖谬，是什么意思？你们想把那些我们素不相识之人的重担强加于我们，却不愿承担你们朋友的重担？圣经大声疾呼：\BibleQuote{你见了盗贼，就与他同伙}。如果你看见的那个人并没有玷污你，你为什么用我未曾看见的人来指责我？或者你辩称：\Quote{我没有与他同谋，因为他的行为令我不悦。} 但无论如何，你曾与他一同走向天主的祭台。好，如果你要为自己辩护，请区分你的两个立场：你说，像两位长老合谋玷污苏撒纳的贞洁那样，共同同意犯罪是一回事；而与盗贼一同领受主的圣事，就像宗徒们与犹大一同领受那第一次晚餐，是另一回事。我完全赞成你的辩护。但你为何不考虑，在你的辩护过程中，你何等轻易地赦免了普天下万国万邦——基督的产业遍布其中？因为如果你看见一个盗贼，并与你所见的盗贼共享圣事，却并未分担他的罪，那么，即使你们的指控和断言为真，即使这些远方的万国与那些非洲的\Emph{叛教者}和迫害者共享圣事，他们又怎能与后者的罪有任何共通之处呢？或者你辩称：\Quote{我在他身上看见的是主教，不是盗贼。} 随你怎么说。我也允许这种辩护；如此一来，整个世界就被洗清了你们对它的指控。因为如果你们可以无视一个你们所认识之人的品性，为什么全世界不能忽略那些他们从未认识的人呢？除非多纳徒派可以忽略他们不愿知道的事，而地上的万国却不能忽略他们不可能知道的事？\par

54. 或者你会说：盗窃是一回事，移交圣书或迫害是另一回事？我承认有区别，现在也不值得指出区别何在。但请听辩论的要点。如果他不能使你成为盗贼，因为他的盗窃行为在你眼中是不悦的，那么谁能对那些憎恶背叛或杀人的人，使他们成为\OriginalTerm{移交者}{traditor}或杀人犯呢？首先，那么，承认你参与了\Person{奥普塔图斯}的一切恶行，你认识他，却还以我所不认识的人身上的任何恶行来责备我。不要对我说：\Quote{我的指控是严重的，你的指控是微不足道的。}你必须先承认它们，无论在你身上是多么微不足道，不是在我承认我的指控之前，而是在我允许你对我谈论这些严重之事之前。你认识的\Person{奥普塔图斯}，他作为你的同僚，使你成为盗贼，还是没有？请回答二者之一。如果你说他没有，我问为什么没有——因为他自己不是盗贼？还是因为你不认识他？还是因为你反对他的行为？如果你说：\Quote{因为他自己不是盗贼}，那么我们就更不应该相信你用来责备我们的那些人是你所断言的那种人。因为如果我们不能相信关于\Person{奥普塔图斯}的事，无论是基督徒、异教徒、犹太人，甚至我们的人和你们的人都断言的事，那么我们怎能相信你对任何人所作的断言呢？但如果你说：\Quote{因为你不认识他}，普天下的万民回答你：\Quote{我们更不认识你用来责备我们的那些人中的一切事。}但如果你说：\Quote{因为你反对他的行为}，他们以同样的声音回答你：\Quote{虽然你从未证实你所说的话的真实性，但我们却以不赞同的眼光来看待这类行为。}但如果你说：\Quote{看，我认识的\Person{奥普塔图斯}，他作为我的同僚使我成为盗贼，因为他做那些事时我习惯与他同到祭台前；但我并不太在意，因为过错是微不足道的；而你们的人却使我成为\OriginalTerm{移交者}{traditor}和杀人犯}，——我回答说，我不承认我也因别人的罪成为\OriginalTerm{移交者}{traditor}和杀人犯，只因为你承认你因另一个人的罪成为盗贼；因为必须记住，你被证明是盗贼，不是藉着我们的判断，而是藉着你自己的承认。因为按我们所说，各人必担当自己的担子，宗徒是我们的见证。\BibleRef{迦 6:5} 但你却自愿将\Person{奥普塔图斯}的担子扛在自己肩上，不是因为你犯了盗窃罪或同意它，而是因为你声明了自己的信念，即别人所作的也适用于你。因为宗徒论及食物时说：\BibleQuote{我凭主耶稣知道并确信，没有一样事物本身是不洁的；但若有人以为某物是不洁的，那么对他就是不清的。}\BibleRef{罗 14:14}依同一原则，可以说别人的罪不能牵连那些反对它们的人；但若有人以为它们影响他，那么他就被它们影响。因此，你无法使我们成为\OriginalTerm{移交者}{traditor}或杀人犯，即使你能证明那些与我们共享圣事的人有这类行为；但盗窃之罪却粘在你身上，即使你反对\Person{奥普塔图斯}所做的一切，这不是由于我们的控诉，而是出于你自己的决定。而且为免你认为这是微不足道的过错，请读宗徒所说：\BibleQuote{盗贼不得继承天主的国。}\BibleRef{格前 6:10}但那些不得继承天主国的人肯定不会在祂的右边，属于那些被说\BibleQuote{我父所祝福的，来承受自创世以来为你们预备的国度}的人。\BibleRef{玛 25:34}他们若不在那里，往何处去呢？岂非在左边吗？因此，属于那些被说\BibleQuote{可咒骂的，离开我，到那为魔鬼及其使者预备的永火里去吧}的人。因此，你沉溺于你的安全感，认为这是使你与天国隔绝，并将你投入永火的微不足道的缺点，是徒然的。你最好将自己交托于真正的羞愧，说：\Quote{我们各人必担当自己的担子，在最后的日子，簸扬扇会将糠秕与麦子分开。}\par

55. 但显然你害怕立刻有人对你说：\Quote{那么，当你们试图把别人的担子放在一些人背上时，竟敢在最后一天的簸扬之前，与分散普天下的主的麦子分离？} 因此，你们反对你们一方之行为的人，一边提防被指控你们所制造的裂教，一边却把自己卷入你们没有犯的别人的罪中；而精明的\Person{培提利安}害怕我能够说\Quote{难道我不像他所认为的\Person{凯齐利亚努斯}那样吗？}他就被迫承认他自己就像他所认识的\Person{奥普塔图斯}那样。要么你们不正如\Place{阿非利加}的众声所宣告他的那样吗？那么，我们也不像你们所责备我们的那些人，他们要么是被你们的错误所怀疑，要么是被你们的疯狂所诽谤，要么是被真理所证明的那样。更不用说普天下万国中主的麦子有这样的品性，因为它们从未听说过你们所提之人的名字。因此，你们没有理由在如此分离之罪和如此裂教之亵渎中丧亡。然而，若你们因天主的审判为这大不敬而受苦，你们却说你们甚至在自己的血中受洗；如此，你们不满足于对分裂毫无悔意，反而必须在惩罚中夸耀。\par

% structure: p0081 source=h2 target=subsection reason=relative-heading-rank; base=h2

\subsection{第二十四章}

56. \Person{培提利安}说：\Quote{但你会回答说，你持守同一宣告：\BibleQuote{一次被洗的人，除洗脚外，不需要再洗。}\BibleRef{若 13:10} 现在，这个\InnerQuote{一次}是有权威的一次，是真理所证实的一次。}\par

57. \Person{奥古斯丁}回答说：因父、及子、及圣神之名而施行的圣洗\BibleRef{玛 28:19}以基督为权威，而非任何人，不论他是谁；基督就是真理，而非任何人。\par

% structure: p0084 source=h2 target=subsection reason=relative-heading-rank; base=h2

\subsection{第二十五章}

58. 培提利安说：\Quote{因为当你们犯罪时，你们施行虚假的仪式，而我并没有两次举行圣洗，因为你们从未举行过一次。}\par

59. 奥斯定答：首先，你并未证明我们有罪。如果一个有罪的人施洗是虚假的圣洗，那么你们派中由那些——我不说公开，甚至是暗中——有罪的人所施洗的人，没有一个拥有真正的圣洗。因为如果施洗者给予的是属于天主的恩赐，而他在天主面前已经有罪，那么一个有罪的人怎能给予属于天主的恩赐呢？如果他有罪就不能给予真正的圣洗？但实际上，你等到他在你面前也有罪的时候才说，仿佛他所要赋予的恩赐是属于你的。\par

% structure: p0087 source=h2 target=subsection reason=relative-heading-rank; base=h2

\subsection{第廿六章}

60. 培提利安说：\Quote{因为若你将虚假与真实混合，虚假常常模仿真理，踏着它的足迹。正如图画描绘出自然人的真实，用色彩描绘出真理的虚假相似。同样，镜子的光辉捕捉面容，仿佛映照出看镜者的眼睛。这样，它向每个人呈现他自己的面容，以至于看镜者的特征彼此相遇；明镜的虚假如此有效，以至于眼睛看见自己，仿佛在别人身上认出自己。甚至当一个影子站在它面前时，它使反射加倍，通过虚假大大分裂了它的统一。那么，我们是否应当因为虚假的再现而认为某物是真实的？但画一个人与生一个人是两回事。难道有人会给一个想要后嗣的人描绘虚构的儿女？或者有人会在图画的虚假中寻找真正的后嗣？确实，迷恋图画而放弃真实是疯狂的表现。}\par

61. 奥斯定答：那么，你难道不羞于将基督的圣洗称为谎言，即使它存在于最虚假的人中？但愿没有人认为，主的麦子——它被命令在整个田地，即整个世界，与莠子一同生长直到收获，即世界末日——会因为你的恶言而败坏。不，甚至在莠子本身（它们被命令不要收集，而要容忍直到末日）和糠秕（它只会在最后一天的扬场中与麦子分离，\BibleRef{玛 3:12}）中，难道有人敢说任何奉父、子、圣神之名施与和领受的圣洗是虚假的吗？你是否会说那些因诱奸妇女的证词（这类例子在任何地方都不少见）而被免职的你的同僚或同伴司铎，在他们的罪行被证明之前，是虚假的或是真实的？你必定回答：虚假。那么，为什么他们能够拥有并给予真正的圣洗呢？为什么他们的虚假作为人没有败坏他们中天主的真理？不是最真实地写着：\BibleQuote{因为训诫的圣神必远避欺诈？}\BibleRef{智 1:5} 既然圣神远避了他们，圣洗的真理又怎能在他们内呢？除非是因为圣神所远避的是人的虚假，而不是圣事的真理。再者，如果连欺诈的人也拥有真正的圣洗，那么那些在真理中拥有它的人是如何拥有的呢？由此你应当观察，倒是你的谈话涂上了孩童般的颜料；因此，忽略生活的圣言而取悦于这种颜料的人，正是爱图画而取代了实体。\par

% structure: p0090 source=h2 target=subsection reason=relative-heading-rank; base=h2

\subsection{第廿七章}

62. 培提利安说：\Quote{有人会用宗徒保禄的话来反对我们：\InnerQuote{\BibleQuote{一个主，一个信德，一个圣洗}}\BibleRef{弗 4:5} 我们宣称只有一个；因为那些说有两个的人无疑是疯了。}\par

63. 奥斯定答：你的这些话是反对你自己的论证；但在你的疯狂中，你并未察觉。因为那些说有两个圣洗的人，正是那些声称义人与不义之人有不同的圣洗的人；然而圣洗既不属这一方也不属那一方，而是在两者中都是同一个，因为它是基督的，尽管他们自己不是一体；然而那同一个圣洗，义人得救，不义之人自取灭亡。\par

% structure: p0093 source=h2 target=subsection reason=relative-heading-rank; base=h2

\subsection{第廿八章}

64. 培提利安说：\Quote{然而，若允许我用比喻，可以肯定，太阳在疯狂的人看来是双倍的，虽然这只是因为一片深蓝色的云常常阻挡它，而云变色表面被光芒照射，太阳的光线从上面反射，仿佛发出它自己的光线。同样，在圣洗的信德中，追求反射是一回事，认识真理是另一回事。}\par

65. 奥斯定答：请问，你在说什么？当一片深蓝色的云反射照射它的太阳光线时，是否只有疯狂的人，而不是所有观看者，才看到两个太阳？但既然它如此呈现在疯狂的人面前，那只是他们独见。然而，若我可直言不讳，我愿你小心，免得说这些话、这样谈论本身就是疯狂的表征。我想，你本意是说，义人拥有圣洗的真理，不义之人只有它的反射。如果真是这样，我敢说，那反射存在于我们派中的那个人身上，对他而言，不是天主，而是一位伯爵成了他的天主；但真理要么在你身上，要么在说出针对奥普塔图斯的那句隽语的人身上，他说：\Quote{在那位伯爵身上，他拥有一个作为同伴的神。} 请区分那些由这两人中任何一位施洗的人：一方面赞成真正的圣洗，另一方面排除反射，引入真理。\par

% structure: p0096 source=h2 target=subsection reason=relative-heading-rank; base=h2

\subsection{第廿九章}

66. \Person{Petilianus}说：\Quote{但是要快速通过这些次要问题：一个不是法庭官员的人能说他制定法律吗？或者当他以私人身份扰乱公共权利时，他所制定的东西能被看作法律吗？难道不是他不仅使自己有罪，而且被视为伪造者，他所写的东西被视为伪造吗？}\par

67. \Person{Augustine}回答说：\Quote{如果你的私人——你视他为伪造者——向任何人陈述皇帝的法律，那人把它与持有真正法律的人的法律比较后，发现完全相同，他不是会放下对他从何处获得法律的关心，而只考虑他得到了什么吗？因为伪造者给出的是假的当他用自己的虚假给出时；但当任何一个人给出真的东西时，即使他是伪造者，尽管给出者不真实，但礼物仍然是真实的。}\par

% structure: p0099 source=h2 target=subsection reason=relative-heading-rank; base=h2

\subsection{第30章}

68. \Person{Petilianus}说：\Quote{如果有人偶然忆起一位司铎的咏唱，他就能被认为是司铎吗？因为他用亵渎的口诵读了司铎的曲调？}\par

69. \Person{Augustine}回答说：在这问题上，你说得好像我们目前是在探询什么是真司铎，而非什么是真圣洗。因为人要成为真司铎，他不仅需要有圣事，还需要有义，正如经上所载：\Quote{愿你的司铎披上正义。}但若一个人仅凭圣事就成为司铎，如同最真的大司祭的迫害者大司祭\Person{Caiaphas}，那么即使他自己不真实，他所给予的仍是真实的，如果他给予的不是他自己的，而是天主的；正如关于\Person{Caiaphas}本人所说的：\Quote{这话不是由他自己说的，但他是那年的大司祭，他预言了。\BibleRef{若 11:51}}然而，用你自己所用的同样比喻：若你甚至从任何亵渎者那里听到司铎的祈祷，措辞符合福音奥迹，你能对他说：\Quote{你的祈祷不真实}吗？尽管他本人可能不仅不是真司铎，甚至根本不是司铎？既然\Person{保禄}宗徒曾说某个我不知道的克里特先知的见证是真的，尽管他不被视为天主的先知，因为他说：\Quote{他们中有一位，即他们自己的先知，说：\InnerQuote{克里特人常说谎，是恶兽，是懒惰的肚腹。}这个见证是真实的。\BibleRef{铎 1:12}}因此，如果宗徒本人也因那外来族的某无名先知的见证是真实的而加以证实，那么当我们发现在任何人身上有属于基督的、且是真实的东西时，即使那人可能是诡诈和乖僻的，我们为何不区分那人身上的过失和他并非出于自己而是出于天主的真理？我们为何不说：\Quote{这圣事是真实的}，正如\Person{保禄}说的：\Quote{这见证是真实的}？难道我们因此就说那人本身也是诚实的，因为我们说这圣事是真实的？如同我要问：宗徒是否因为在他身上发现真理而证实那先知列于上主的先知中？同样，同一位宗徒在雅典时，在假神祭坛中发现一座祭坛，上面刻着：\Quote{献给未识之神。}他利用这个见证来在基督内建造他们，以至于在他的讲道中引用了这铭文，并加上：\Quote{你们所不认识而敬拜的，我现在向你们宣告。}难道因为他发现那祭坛在偶像中间，或由亵渎之手所立，就因此谴责或拒绝其中他所发现的真理吗？或者因为其上发现的真理，就因此劝说他们也该跟随外教人的亵渎习俗？显然，他两者都没有做；而是当时，他按自己所判断的合适方式，希望将主自己（他们不认识，但他认识）介绍给他们认识时，他在其他事情中说：\Quote{他离我们每人都不远，因为我们生活、行动、存在都在他内，就如你们自己的某些诗人也说过。}能说这里因为他发现在亵渎中有真理的见证，他要么因见证而认可了他们的邪恶，要么因邪恶而谴责了见证？但你们必定总是错的，只要你们因人的过失而轻慢天主的圣事，或者为了天主的圣事（我们不愿在你们身上轻慢它们）而以为我们承担你们甚至分裂的亵渎。\par

% structure: p0102 source=h2 target=subsection reason=relative-heading-rank; base=h2

\subsection{第31章}

70. \Person{Petilianus}说：\Quote{因为没有权柄不是出于天主的，\BibleRef{罗 13:1}} 任何有权柄的人也不例外；正如主耶稣基督回答般雀·比拉多说：\Quote{若不是从上头赐给你的，你毫无权柄办我。\BibleRef{若 19:11}} 又如若望的话：\Quote{人不能领受什么，除非是天上赐给他的。\BibleRef{若 3:27}} 因此，告诉我们吧，你这\Emph{traditor}，你何时领受了模仿奥迹的权柄。\par

71. \Person{Augustine}回答说：不如你告诉我们，整个传播基督产业的普世万民，以及宗徒建立教会的众多民族，何时丧失了施洗的权柄？你永远无法告诉我们——不仅因为你诽谤了他们，并没有证明他们是\Emph{traditor}，而且因为，即使你证明了这一点，任何恶人的罪责，无论他们是可疑的、欺骗的，还是被容忍如莠子或糠秕，都不能推翻那些应许，以致地上万民不再因亚巴郎的后裔蒙福；在这些应许中，当你不愿与地上万民共融于合一时，你就剥夺了他们的份。\par

% structure: p0105 source=h2 target=subsection reason=relative-heading-rank; base=h2

\subsection{第32章}

72. 培提利阿努斯说：\Quote{因为虽然只有一个圣洗，却按三级祝圣。若翰没有呼号天主圣三之名而施水，如他自己所宣称：\InnerQuote{我固然以水洗你们，为叫你们悔改；但在我以后来的那一位，比我更强，我连给祂提鞋也不配；祂要以圣神和火洗你们。}\BibleRef{玛 3:11} 基督赐予圣神，如经上所载：\InnerQuote{祂向他们嘘了一口气，说：你们领受圣神罢！}\BibleRef{若 20:22} 护慰者亲自以烈火般哗啦作响的火焰降到宗徒身上。哦，真神性，看似燃烧却不焚毁！如经上所载：\InnerQuote{忽然从天上来了一阵响声，好像暴风刮来，充满了他们所在的整座房屋。有散开的火舌，如同火焰一般，分别停留在每人头上。众人都充满了圣神，照圣神赐给他们的话，说起外方话来。}\BibleRef{宗 2:2} 但你这迫害者，连悔改的水也没有，因为你所掌握的权柄不是来自被杀的若翰，而是来自杀人的黑落德。因此，你，\OriginalTerm{叛徒}{traditor}，没有基督的圣神；因为基督并没有出卖别人至死，而是祂自己被出卖。为你，阴间的火灵是充满生命的——那火舌翻腾饥渴，将能永远焚烧你的肢体却不烧尽，如经上论地狱中罪人的惩罚所载：\InnerQuote{他们的火也不熄灭。}\BibleRef{依 66:24}}\par

73. 奥古斯丁答道：你是诽谤者，不是真诚的辩论者。你何时才停止作出那种论断？你若不能证明，它们就不适用于任何人；即使你能证明，也肯定不适用于整个世界在圣人们中如同在天主麦子中的合一。假如我们也乐意以诽谤还诽谤，我们或许也能发出雄辩的诽谤。我们也可以使用\Quote{哗啦作响的火焰}这种说法；但对我来说，用法不当的说法无论如何没有雄辩之音。我们也可以说\Quote{火舌翻腾饥渴}；但我们并不希望我们著作中的火舌，当任何有理智的人读到时，被认为因缺乏重量之汁而饥渴，或者读者自己，在他们当中找不到有益情感的食粮，而遭受极度空虚的饥渴之苦。看，我宣告，你们的\OriginalTerm{色康多派}{Circumcelliones}正在燃烧，不是哗啦作响的，而是猛烈的火焰。如果你回答：这与我们何干？那么当你以任意某人来指责我们时，你为何不也听听我们的回答：我们也一无所知？如果你回答：你并未证实这事，为何整个世人不也回答你说：你也未证实？因此，我们乐意达成一致：你不得以你认为属于我们的恶人之罪来指责我们，我们也停止类似的指控。这样，你会看到，通过这个公正的、确认并批准了的协议，你没有任何指控可以针对亚巴郎的苗裔，即普世万民中的亚巴郎苗裔。但我毫不困难地找到一项严重指控加于你：为何你们不虔地脱离了亚巴郎的苗裔，即普世万民中的亚巴郎苗裔？针对这个指控，你们肯定无法为自己辩护。因为我们各自不承担他人的罪；但这一点——你们不与普世万民共融，这些万民在亚巴郎的苗裔中蒙受祝福——是一项极其严重的罪行，不是你们中的某些人，而是你们全体都犯下的。\par

74. 然而，你们知道，如同你们通过引文证明的那样，圣神如此降临，以致当时被充满的人说起各种方言：那记号和奇迹有何意义？为何现在圣神赐予的方式是，没有领受者说起各种方言？除非因为那奇迹当时预表了地上万民要相信，并因此福音要在每一种语言中找到。正如在诗篇中早已预言的那样：\BibleQuote{没有语言，没有方言，听不到他们的声音。}这话是针对那些注定在领受圣神后说各种方言的人。但因为这段经文本身意指福音此后要在万民和万语中找到，基督的身体要在全世界用每一种语言发出声音，所以诗篇继续写道：\BibleQuote{它们的声音传遍普世，它们的言语达到地极。}因此，真教会对任何人都不是隐藏的。因此，主自己在福音中说：\BibleQuote{建在山上的城是不能隐藏的。}\BibleRef{玛 5:14} 因此，达味在同一篇诗篇中继续写道：\BibleQuote{祂在太阳中安置了祂的帐幕，}即在白天的光天化日之下；正如我们在《列王纪》中所读的：\BibleQuote{你暗中行了这事，我却要在以色列众人面前，在太阳下做这事。}\BibleRef{撒下 12:12} 祂自己是\BibleQuote{如新郎出洞房，又如勇士欢然奔跑。祂的起头从天上极处起}：这里你看到主降生成人。\BibleQuote{祂的循环直到天的尽头}：这里你看到祂的复活与升天。\BibleQuote{没有一物能隐藏祂的热气}：这里你看到圣神的降临，祂以火舌派遣圣神，为显明爱德的热火，那爱德是那些不与教会在和平的纽带中持守圣神合一的人所绝不能拥有的；而教会存在于各种语言中。\par

75. 其次，关于你所说确实只有一个圣洗\BibleRef{弗 4:5}，但它却由三个阶段祝圣，并且你将这三种形式如此分配给三个人：你把水归于若翰，圣神归于主耶稣基督，第三，火归于从上降临的护慰者——请想想你陷入何等的错误。因为你之所以持有这种意见，仅仅是由于若翰的话：\BibleQuote{我固然以水洗你们，但那在我以后来的，比我更强，他以圣神和火洗你们。}\BibleRef{玛 3:11} 而你却不考虑，这三样并非分别归于三个人——水归于若翰，圣神归于基督，火归于护慰者——而应归于两个人：其中之一归于若翰，另外两个归于我们的主。因为经上并没有说：\Quote{我固然以水洗你们，但那在我以后来的，比我更强，我连提他的鞋也不配；他以圣神洗你们，而那在我以后来的护慰者，他以火洗你们。}而是说：\BibleQuote{我固然以水洗你们，但那在我以后来的，他以圣神和火洗你们。}他把一样归于自己，两样归于那以后来的。因此，你看你在数量上受了何等的欺骗。请再听。你说有一个圣洗由三个阶段祝圣——水、圣神和火；并且你给这三个阶段分别指派了三个人：若翰归于水，基督归于圣神，护慰者归于火。因此，如果若翰的水属于同一个被赞扬为唯一的圣洗，那么宗徒保禄命令那些他发现曾被若翰施洗的人再次受洗就不对了。因为他们已经有了水，如你所说，属于同一个圣洗；所以他们只需要接受圣神和火，因为这些在若翰的圣洗中缺乏，以使他们的圣洗得以完成，按你所说，由三个阶段祝圣。但由于他们由一位宗徒的权柄命令受洗，这就充分表明若翰施洗所用的水与基督的圣洗无关，而是属于另一种适应时代需要的分施。\par

76. 最后，当你想要证明圣神是由基督赐予，并从福音中引证说，耶稣从死者中复活后向祂的门徒脸上吹了一口气，说：\BibleQuote{你们领受圣神吧！}\BibleRef{若 20:22} 并且当你想要证明那最后的与圣洗有关的火，就是在圣神降临时所显示的舌形火焰中发现的，你怎么会想到说：\Quote{护慰者亲自降临在宗徒们身上，如同燃烧着噼啪火焰的火，}就好像有一个圣神，是祂通过向门徒脸上吹气而赐予的，另有另一个圣神，是在祂升天后降临在宗徒们身上的？那么我们是否该认为有两个圣神？谁能疯狂到如此地步，竟敢断言这事？因此，基督自己赐予了同一个圣神，无论是通过向门徒脸上吹气，还是在五旬节日从天上派遣祂降下，并无疑地明示了祂神圣的圣事。因此，并不是基督赐予圣神，而护慰者赐予火，为应验\BibleQuote{以圣神和火}这句话；而是同一位基督在这两种情况下都赐予了圣神，当祂还在世上时通过吹气，当祂升天后通过火焰的舌头显明。因为为使你知道若翰的话\BibleQuote{他要以圣神洗你们}并不是在祂向门徒脸上吹气时就应验了，以致他们需要等到护慰者降临时再受洗，不再是受圣神的洗，而是受火的洗，我要你回想圣经最明确的言论，并看主升天时亲自对他们说的话：\BibleQuote{若翰固然以水洗你们，但你们在不多几天以后，在五旬节，要受圣神的洗，就是你们所要领受的。}\BibleRef{宗 1:5} 还有什么比这证言更清楚的？但按照你的解释，祂应该说：若翰固然以水洗你们，但当我在你们脸上吹气时，你们已经受了圣神的洗；接下来按顺序你们要受火的洗，你们在不多几天后就要领受——这样，由三个阶段完成的那个唯一圣洗就得以圆满。结果证明，你仍然不明白\BibleQuote{他要以圣神和火洗你们}这句话的含义，而且你够鲁莽，愿意教导自己所不懂的事。\par

% structure: p0111 source=h2 target=subsection reason=relative-heading-rank; base=h2

\subsection{第三十三章}

\Quote{但是，为彻底查考因天主圣三之名而施行的圣洗，主基督曾对祂的宗徒们说：\BibleQuote{你们去，给万民授洗，因父、及子、及圣神之名；教训他们遵守我所命令你们的一切。}\BibleRef{玛 28:19} \Emph{\OriginalTerm{叛徒}{traditor}}，你教训谁呢？就是你定罪的那个人吗？\Emph{\OriginalTerm{叛徒}{traditor}}，你教训谁呢？就是你杀害的那个人吗？再问一次，你教训谁呢？就是你使其成为凶手的那个人吗？那么，你怎么能因天主圣三之名授洗呢？你不可称天主为你的父。因为主基督曾说：\BibleQuote{缔造和平的人是有福的，因为他们要称为天主的子女。}\BibleRef{玛 5:9} 你们这些没有灵魂平安的人，不能以天主为父。或者，你怎么能因圣子之名授洗呢？你出卖了圣子本身，你没有在任何苦难或十字架上效法天主子。或者，你怎么能因圣神之名授洗呢？因为圣神只降临在那些没有犯叛逆罪的宗徒们身上。因此，既然天主不是你的父，你也没有真正藉着圣洗的水重生。你们中没有一个人是完美重生的。你们这些不敬虔的人，既无父也无母。既然你们是这样的，难道我不应该给你们授洗吗？即使你们像犹太人那样洗自己一千次，他们仿佛是在洗肉体。}\par

奥古斯丁答道：你原说彻底查考因天主圣三之名而施行的圣洗，也使我们聚精会神地听；但你似乎遵循着对你最容易的方式，多么快就回到了你惯常的辱骂！你对此确实流利。因为你给自己设定无论怎样的牺牲品，随你喜好用任意苦涩来谩骂他们；在你这种话语的广阔自由中，如果有人只是用了\Quote{证明它}这个小词，你就会被逼入极大的困境。因为这是亚巴郎的子孙对你说的；既然万民都因他而蒙福，他们不太在意被你咒骂。但是，既然你在讨论圣洗，你认为圣洗在义人身上时才是真的，在不义之人身上时就是假的，那么请看，我也按照你的规则来查考因天主圣三之名而施行的圣洗，我可以很充分地说，在我看来，那与一位伯爵一同以天主为伴的人，不能以天主为父；他也不信任何人是他的基督，除非是为了他的缘故而忍受苦难的那位；他没有圣神，因为那位如此不同地用火舌焚烧了可怜的阿非利加的人。那么，他们怎么能有圣洗，或者他们怎么能因父、及子、及圣神之名施行圣洗呢？无疑你现在应当明白，圣洗可以存在于一个不义的人身上，也可以由一个不义的人施行，而且这种圣洗不是不义的，而是正义和真实的——不是因为属于不义的人，而是因为属于天主。在此，我没有对你进行任何诽谤，像你不断对全世界所做的那样，以某种借口诽谤；更不可容忍的是，你甚至对你诽谤所依据的那些要点毫无证明。但我不知道这如何可能被容忍，因为你不仅对圣人们进行关于不义之人的诽谤，而且甚至对神圣的圣洗本身提出指控，这圣洗无论在任何人身上，无论他如何不义，都必须是神圣的，你从恶人的罪所带出的感染进行比较，从而说圣洗具有其拥有者、施行者或领受者的品性。此外，如果一个人具有与他一同接近神圣奥迹者的品性，如果圣事本身具有它们所在之人的品性，那么圣人们也许可以满足于从中得到安慰，想到他们只是像神圣的圣洗本身一样从你口中听到虚假的指控。但你最好看看你如何被自己的口定罪，如果你们中清醒的人因你们中醉酒者的感染而算为醉酒，你们中有怜悯的人因抢劫者的感染而变成抢劫者，以及无论什么恶在你们中间存在于恶人身上，也不得不被那些非恶人分沾；如果圣洗本身在你们所有不洁之人中是不洁的，并且根据不洁本身不同的性质而有不同的种类，因为如果它不得不具有拥有者或施行者的品性，那么它就是如此。这些假设无疑是假的，因此当你们提出这些反对我们而不回顾你们自己时，它们丝毫不能伤害我们。但它们伤害你们，因为当你们虚假地提出它们时，它们不落到我们身上；但由于你们想象它们是真实的，它们就反噬你们自己。\par

% structure: p0114 source=h2 target=subsection reason=relative-heading-rank; base=h2

\subsection{第34章}

\Quote{因为如果宗徒们被允许给那些被若翰用悔改的洗礼所洗的人施洗，那么同样不也应该允许我给像你们这样犯亵渎罪的人施洗吗？}\par

那么，你上面说的若翰的洗礼和基督的圣洗不是两个，而是一个圣洗，分三个阶段祝圣，其中若翰给水，基督给圣神，护慰者给火，这话到哪里去了？那么，为什么宗徒们对那些已经接受若翰所给的水的人重新给水呢？这水属于那个分三个阶段祝圣的同一个圣洗？无疑你必须明白每个人都有必要理解他所讨论的内容。\par

% structure: p0117 source=h2 target=subsection reason=relative-heading-rank; base=h2

\subsection{第35章}

\Quote{确实，圣神不可能通过司铎覆手被植入任何人的心中，除非良心洁净的水先行使他重生。}\par

八十二。\Person{奥古斯丁}答：在你的这几句话中，包含了两个错误；其中一个固然对我们的讨论无大影响，但却有助于证明你们缺乏技巧。因为圣神降临在了一百二十人身上，没有任何人的覆手，又降临在\Person{科尔乃略}百夫长和与他在一起的人身上，甚至在他们领受圣洗之前。但是，你这些话中的第二个错误完全推翻了你的整个论点。因为你宣称，纯洁良心的水必须先产生新生，然后圣神才能随之而来。因此，要么所有因父、及子、及圣神之名祝圣的水都是纯洁良心的水，不是由于施行者或领受者的功绩，而是由于设立此圣洗者的无玷功劳；要么，如果只有施行者和领受者双方的纯洁良心才能产生纯洁良心的水，那么，对于那些由良心尚被未发现的罪污染的人所施洗的人，你怎么办？特别是在这些领洗者中，如果有人承认他在领洗时良心不好，因为他可能想利用那个机会去犯某种罪？因此，当向你表明，无论是施行圣洗的人还是领受圣洗的人都没有纯洁良心时，你是否会判定他应当重新受洗？你肯定既不会这样说也不会这样做。因此，圣洗的纯洁与施洗者或领洗者良心的纯洁或不纯洁完全无关。那么，你敢说那欺骗者、强盗、压迫孤儿寡妇者、拆散婚姻者、叛卖者、贩卖者、分割他人遗产的人是纯洁良心的人吗？或者，你还敢说那些人是纯洁良心的人吗？在如此时代，难以想象缺少这样的人，他们与我描述的人结交以便受洗，不是为了基督，也不是为了永生，而是为了获得世俗的友谊和满足世俗的欲望。再者，如果你不敢说这些人是纯洁良心的人，那么如果你发现他们中有任何人受过洗，就给他们尚未领受的纯洁良心的水吧；如果你不愿这样做，那就不要再把你所不明白的事强加给我们，免得你被迫回答我们关于你所深知的事。\par

% structure: p0120 source=h2 target=subsection reason=relative-heading-rank; base=h2

\subsection{第三十六章。}

八十三。\Person{佩提利安}说：\Quote{那圣神当然不能降临在你们身上，因为你们甚至没有被悔改的圣洗所洗；但\Emph{\OriginalTerm{叛教者}{traditor}}的水，最确实需要为之悔改，却只带来污染。}\par

八十四。\Person{奥古斯丁}答：事实上，你们不仅不能证明我们是叛教者，你们的祖先也没有证明我们的祖先犯有这项罪；即使证明了，按照你们早先的说法，他们也不会是我们的祖先，因为我们没有效法他们的行为：然而我们也不应因他们而被切断与合一团体的联系，以及与\Person{亚巴郎}后裔的联系，\BibleQuote{地上万民都要因他而蒙祝福}。\BibleRef{创 22:18} 然而，如果基督的水是一回事，\OriginalTerm{叛教者}{traditor}的水是另一回事，因为基督不是\OriginalTerm{叛教者}{traditor}，那么为什么基督的水不能是一回事，强盗的水是另一回事呢？因为基督当然不是强盗。因此，你们在你们的强盗之后再次施洗，而我将在\OriginalTerm{叛教者}{traditor}之后再次施洗，这\OriginalTerm{叛教者}{traditor}既不是我的也不是你们的；或者，如果必须相信所提出的文件，那既是我的也是你们的；或者，如果我们相信普世教会的共融而非\Person{多纳图}派，那他就不是我的，而是你们的。但是，按照更美好、更健全的判断，因为符合宗徒的话，\BibleQuote{我们每人要背负自己的重担}；\BibleRef{迦 6:5} 那强盗不是你们的，如果你们自己不是强盗；也没有任何一个\OriginalTerm{叛教者}{traditor}属于我们或你们中的任何人，如果他自己不是\OriginalTerm{叛教者}{traditor}。然而，我们这些公教徒，遵循那判断的精神，不离开教会的合一；但你们是异端者，你们因对某些人提出的控告（无论真假），不愿与\Person{亚巴郎}的后裔保持基督徒的爱德。\par

% structure: p0123 source=h2 target=subsection reason=relative-heading-rank; base=h2

\subsection{第三十七章。}

八十五。\Person{佩提利安}说：\Quote{但为证明此事的真理来自宗徒，我们由他们的行动得知，正如所记载的：\InnerQuote{当\Person{阿颇罗}在\Place{格林多}的时候，\Person{保禄}走遍了上边一带地方，来到\Place{厄弗所}；遇见几个门徒，问他们说：你们信从的时候，曾领受圣神没有？他们回答说：我们连圣神的存在都没有听说过。\Person{保禄}说：那么，你们受的是什么洗呢？他们说：是\Person{若翰}的洗。\Person{保禄}说：\Person{若翰}所施的是悔改的洗，他告诉人民要相信在他以后来的那一位，就是\Person{耶稣}。他们听了这话，就因主\Person{耶稣}之名受了洗。\Person{保禄}给他们覆手，圣神便降临在他们身上，他们就讲各种语言，也说预言。他们一共约有十二人。} \BibleRef{宗 19:1} 因此，如果他们受洗是为了领受圣神，那么你们如果希望领受圣神，为什么不在你们的谎话之后，设法获得真正的更新呢？如果我们在此事上催促得不对，你们又为何寻找我们呢？或者，如果这是一种冒犯，首先谴责\Person{保禄}吧；那\Person{保禄}确实洗去了已经存在的东西，而我们却在你们身上给予尚未存在的圣洗。因为你们，如我们常说的，并不是以真正的圣洗施洗；而是以虚假圣洗的空名，给人们带来恶名。}\par

86. \Person{奥斯定}回答说：\Quote{我们不控告\Person{保禄}，他之所以给人授基督的圣洗，是因为他们没有基督的圣洗，而只有若翰的圣洗，按照他们自己的回答；因为当他们被问：\InnerQuote{你们受的是什么洗？}他们回答：\InnerQuote{若翰的圣洗。}这同基督的圣洗毫无关系，既不是它的一部分，也不是迈向它的一步。否则，要么当时基督圣洗的水被更新了第二次，要么如果基督的圣洗当时是由两种水完成的，那么现在所施的圣洗就不那么完美，因为它没有用若翰所施的水来施予。但持有这两种看法中的任何一种都是邪恶和亵渎的。因此，\Person{保禄}给那些没有基督圣洗、只有若翰圣洗的人施予了基督的圣洗。}\par

87. 但是，为何若翰的圣洗——现在已不必要——在当时却是必要的，我在别处已经解释过了；这个问题与我们目前争论的焦点无关，除了它表明若翰的圣洗是一回事，基督的圣洗是另一回事——正如宗徒所说的，我们的祖先在云中和海中受的洗是另一种洗，他们在\Person{梅瑟}的带领下穿过红海时受的洗。\BibleRef{格前 10:1}因为法律和先知直到洗者若翰，都有预示未来事物的圣事；但我们时代的圣事却证明那先前的圣事所预言要来的已经来了。因此，若翰是基督的预告者，在时间上比所有在他之前的人都更接近基督。因为所有古时的义人和先知都渴望看到那藉着圣神的启示他们预知将要发生的事的实现——因此主自己说：\BibleQuote{有许多先知和义人愿意看你们所看见的，却没有看见；愿意听你们所听见的，却没有听见。}\BibleRef{玛 13:17}——所以论到若翰，说他比先知还大，凡妇人所生的，没有一个比他更大；\BibleRef{玛 11:9}因为在他之前的义人只获准预告基督的来临，但若翰被赐予既在基督不在时预告他，又目睹他的亲临，以致发现向他显明了别人所渴望的。因此，他圣洗的圣事仍然与预告基督的来临有关，尽管是非常即将来临的事，因为直到他的时代，仍然有对我们主第一次来临的预告，而对此我们现在只有宣告，不再有预言。但是，主教导谦卑之道，屈尊使用了他在这里所找到的、与预告他的来临相关的圣事，不是为了帮助他洁净的行动，而是作为我们虔诚的榜样，以此向我们展示，我们应当以何等尊敬接受那些见证他已来临的圣事，当他并未不屑于使用那些预示他即将来临的圣事时。因此，若翰虽然在时间上最接近基督，与他年龄相差不到一年，但当他施洗时，他走在尚未到来的基督的道路前面；为此论到他说：\BibleQuote{看，我派遣我的使者在你面前，他要在你前面预备你的道路。}他自己也宣讲说：\BibleQuote{有一位比我更有能力的在我以后来。}\BibleRef{谷 1:7}同样地，第八日的割损礼，赐给祖先，预示了我们的成义，藉着我们主的复活除去肉体的情欲，这复活发生在第七日即安息日之后，第八日即主日，亦即他埋葬后的第三日；然而婴孩基督也接受了这具有先知意义的肉身割损礼。正如犹太人通过宰杀羔羊所庆祝的逾越节，预表了我们主的受难和他离开这个世界回到父那里，然而主自己却与门徒一同庆祝了这个逾越节，当门徒提醒他时，说：\BibleQuote{你愿意我们在哪里为你预备吃逾越节晚餐？}\BibleRef{玛 26:17}同样，他自己也接受了若翰的圣洗，这构成了对他来临最晚近预告的一部分。但是，犹太人的肉身割损礼是一回事，我们在人受洗后第八日所举行的礼仪是另一回事；犹太人仍然通过宰杀羔羊所庆祝的逾越节是一回事，我们在主的体和血中所领受的逾越节是另一回事——如此，若翰的圣洗是一回事，基督的圣洗是另一回事。因为前一系列礼仪预表了后者注定要到来；后一系列则宣告前者已应验。即使基督接受了前者，它们对我们并不必要，因为我们接受了在它们中预告的主自己。但是，当我们主的来临还是最近的事时，任何接受了前者的人还必须受后者的浸透；但已经受后者浸透的人，没有必要再被迫回到前礼仪。\par

88. 因此，不要试图从若翰的圣洗中制造混乱，它的来源和意图要么如我在此所述的那样；或者如果能够给出任何其他更好的解释，这一点仍然是清楚的：若翰的圣洗与基督的圣洗是两种截然不同且分开的事物，前者明确被称为若翰的圣洗，正如你所引用的那些人的回答和主亲自说的话所清楚表明的那样，祂说：\BibleQuote{若翰的圣洗是从哪里来的？是从天上来的，还是从人来的？}\BibleRef{玛 21:25}。但后者从未被称为\Person{塞西利安}、\Person{多纳图}、\Person{奥斯定}或\Person{培提利安}的圣洗，而是基督的圣洗。因为如果你认为我们是厚颜无耻的，因为我们不允许任何人在我们这里领受圣洗之后再受洗，尽管我们看到那些领受了若翰圣洗的人（此人当然比我们大得不可比拟）又受了洗，那么你会坚持\Person{若翰}和\Person{奥普塔塔}是同等尊贵的吗？这看起来很荒谬。然而，我想你并不认为他们平等，反而认为\Person{奥普塔塔}更伟大。因为宗徒在若翰的圣洗之后为人施洗：你不敢在\Person{奥普塔塔}的圣洗之后为任何人施洗。是因为\Person{奥普塔塔}与你们合一吗？我不知道这是怎样保持的一种理论，如果伯爵的朋友（他在伯爵身边有一个神作为同伴）被称为在合一中，而新郎的朋友却被排除在外。但如果\Person{若翰}卓越地处于合一中，且比我们所有人和你们所有人都更卓越和伟大，而宗徒\Person{保禄}却在他之后为人施洗，那么你为什么不在\Person{奥普塔塔}之后施洗呢？除非是你的盲目使你陷入困境，以至于你竟然说\Person{奥普塔塔}有赐予圣神的能力，而\Person{若翰}却没有！如果你不说这话（因为怕你的疯狂甚至被疯子嘲笑），那么当有人问你为什么人在领受了\Person{若翰}的圣洗之后还需要受洗，而没有人需要在领受\Person{奥普塔塔}的圣洗之后受洗时，你能做出什么回答呢？除非是因为前者受的是\Person{若翰}的圣洗，而每当一个人受的是基督的圣洗时，无论是由\Person{保禄}还是由\Person{奥普塔塔}施洗，他的圣洗的本质都没有区别，尽管\Person{保禄}和\Person{奥普塔塔}之间有天壤之别？因此，\BibleQuote{你们这些悖逆的人，回头归向正心罢！}\BibleRef{依 46:8}不要试图用人的品性和行为来衡量天主的圣事。因为圣事因着其所归属的那一位而是圣的；当被相称地领受时，它们带来赏报；当不相称地领受时，带来审判。虽然相称或不相称地领受天主圣事的人不是同一个，但领受的圣事本身却是相同的；因此它本身不会变得更好或更坏，而只是在两种情况下对领受者或产生生命或产生死亡。至于你所说的，\Quote{在保禄为那些领受了若翰圣洗的人施洗时，他洗去了已经存在的东西}，你肯定是一点也没有考虑你在说什么。因为如果若翰的圣洗需要洗去，那么毫无疑问它里面有一些污秽。那我为什么还要催逼你呢？回想或阅读，看看\Person{若翰}从哪里领受了它，这样你就会明白你说了那个亵渎的话是针对谁；当你发现这一点时，你的心一定会被击打，如果你不给你的舌头套上缰绳的话。\par

89. 接下来，关于你自以为对我说的如此尖刻的话：\Quote{如果我们在这件事上做错了，你为什么寻找我们？}难道你仍然不能想起，只有那些丧亡的人才会被寻找吗？或者，看不见这一点正是你们毁灭的一部分？因为羊也可以同样荒谬地对牧人说：\Quote{如果我离开羊群做错了，你为什么寻找我？}不明白它被寻找恰恰是因为它认为不需要被寻找。但是，除了那一位通过祂的圣经、通过大公的与和解的声音、通过暂时苦难的鞭打向你们施行怜悯的天主，还有谁在寻找你们呢？因此，我们寻找你们，为要找到你们；因为我们爱你们，愿你们获得生命，其强度与我们恨你们的错误相同，好使它被消灭，这错误要毁灭你们，但它本身并不牵涉在你们的毁灭中。但愿我们能够这样寻找你们，甚至找到，并能为你们中的每一个人高兴地说：\BibleQuote{他死而复生，失而复得！}\BibleRef{路 15:32}\par

% structure: p0129 source=h2 target=subsection reason=relative-heading-rank; base=h2

\subsection{第三十八章}

\Person{培提利安}说：\Quote{如果你声称你持有大公教会，那么\OriginalTerm{大公的}{catholic}一词仅仅是希腊文\Quote{全部}或\Quote{整体}的同义词。但显然，你不在整体之中，因为你已经滑向部分。}\par

91. 奥古斯丁答：我确实对希腊语只有极其浅薄的了解，几乎谈不上知识，但我并不无耻地说，我知道\OriginalTerm{全体}{\GreekText{ὅλον}}的意思不是\Quote{一}，而是\Quote{全体}；而\OriginalTerm{按照全体}{\GreekText{καθ}' \GreekText{ὅλον}}的意思是\Quote{按照全体}：由此，天主教会得名，正如主所说：\Quote{不是你们可以知道父凭自己的权柄所定的时候、日期；但当圣神降临于你们时，你们将得到能力，并要在耶路撒冷、犹太全地、撒玛黎雅，直到地极作我的见证。}\BibleRef{宗 1:7} 这里你们有了\Quote{天主教}之名的来源。但你们如此执意闭眼冲向那小石长成的大山——按达内尔的预言——充满全地，\BibleRef{达 2:35} 以至于你们竟告诉我们，我们陷入了部分，而不在全地之中，那些共融遍布全地的人。但正如假设你们说我是\Person{Petilianus}，我无法找到反驳你们的方法，除非嘲笑你们是在开玩笑，或哀叹你们疯了；同样，在此我看出我别无选择；既然我不相信你们是在开玩笑，你们就明白剩下的是什么。\par

% structure: p0132 source=h2 target=subsection reason=relative-heading-rank; base=h2

\subsection{第39章。}

92. \Person{Petilianus}说：\Quote{但黑暗与光明没有共融，苦与蜜甜没有共融；生命与死亡、无辜与有罪、水与血没有共融；酒渣与油没有共融，尽管作为沉淀物彼此相关，但凡可弃绝的都会流走。那正是罪恶的阴沟；按若望的话：\InnerQuote{他们从我们中间出去，却并非属于我们；因为如果他们属于我们，就必会与我们同在。}\BibleRef{若一 2:19} 他们的污秽中没有金子：一切宝贵的都已炼净。因为经上记着：\InnerQuote{金子在炉中试炼，义人也同样在困苦的折磨中受试炼。}残忍不属于温柔，宗教不属于亵圣；\Person{Macarius}一派绝不能属于我们，因为他玷污了我们礼仪的表象。因为敌人的队伍——填满敌人之名——不是它所反对的力量的一部分；但如果它真是称为一部分，它将在撒罗满的判断中找到合适的座右铭：\InnerQuote{愿他们的份从地上剪除。}\BibleRef{箴 2:22}}\par

93. 奥古斯丁答：不证实任何事就发出这些辱骂，岂不是纯粹的疯狂？你们看全世界的莠子，却不看麦子，尽管两者都被命令在整个田地中一起生长。你们看恶者所撒的种子，在收割时将被分离，\BibleRef{玛 13:24} 却不看亚巴郎的种子，万邦将因他而得福。\BibleRef{创 22:18} 如同你们已是炼净的团块、纯蜜、精炼的油、纯金，或者更确切地说，是粉刷的白墙的表象。因为，且不提你们的其他过失，醉酒者算是清醒者的一部分吗？贪婪者被算在智慧者的一份中吗？如果温良的人占有光明这词，那么\OriginalTerm{周游者}{Circumcelliones}的疯狂该被置于何处，除非在黑暗中？那么，由这些人所施的洗礼在你们中为何有效，而基督的同一个洗礼，不论由世上的任何人施行，在你们中为何无效？事实上，你们看出自己与普世教会分离，仅此一点，即你们并不都是醉酒者，也不都是贪婪者，也不都是暴力者，但你们都是异端者，并因此都是不虔敬和亵圣的。\par

94. 至于你说普世欢欣于基督徒共融的那一派是\Person{玛卡里}一派，任何头脑尚存些许理智的人，谁能说出这样的话？但因我们说你们是\Person{多纳徒}一派，你们便寻找一个人，好说我们是属他的一派；处此极大困境中，你们提及某个籍籍无名之辈的名字，此人若在非洲为人所知，则在地球任何其他角落必然无人知晓。因此，请听遍布大地四极的\Person{亚巴郎}所有后裔向你的回答：你们声称我们所属的那个\Person{玛卡里}，我们全然不知。你们能否反过来声称你们不知\Person{多纳徒}？但即使我们说你们是\Person{奥普塔徒}一派，你们中有谁能说他不认识\Person{奥普塔徒}，除非他不认识他本人，如同或许他也不认识\Person{多纳徒}？然而你们承认你们以\Person{多纳徒}之名为荣，你们也以\Person{奥普塔徒}之名为乐吗？那么，当你们全都同样被\Person{奥普塔徒}玷污时，\Person{多纳徒}之名对你们有何益处？当你们被西尔孔塞利奥人的酗酒所玷污时，\Person{多纳徒}的清醒对你们有何益处？按照你们的看法，当你们被\Person{奥普塔徒}的贪婪所玷污时，\Person{多纳徒}的无辜对你们有何益处？因为你们的错误在于，你们认为一个人的不义在感染邻人方面，比一个人的义在洁净周围的人方面更有能力。因此，若两人共享天主的圣事，一人是义人，另一人是不义之人，但前者并不效法后者的不义，后者也不效法前者的义，你们却说结果不是两人都成为义，而是两人都成为不义；连那两人共同领受的圣物也变得不洁，失去其原有的圣洁。不义何时为自己找到了这样的辩护者，以致通过他们的疯狂，不义竟被看作得胜？那么，在如此错谬的乖戾中，你们又如何因\Person{多纳徒}之名而自夸，实则它并非显明\Person{培提利安}值得成为\Person{多纳徒}那样，而是迫使\Person{多纳徒}成为\Person{奥普塔徒}那样？但愿\Place{以色列}家说：\BibleQuote{天主永是我的福分}；愿\Person{亚巴郎}的后裔在万民中说：\BibleQuote{上主是我产业的福分}。因为他们知道如何透过福音述说荣耀之天主的光荣。因为你们，藉着在你们内的圣事，如同迫害主的\Person{盖法}一样，在不知不觉中预言。\BibleRef{若 11:51}因为希腊文用\OriginalTerm{玛卡里}{\GreekText{Μαχάριος}}表达的，在我们的语言中就是\Quote{有福者}；这样，我们当然是属于\Person{玛卡里}、即这有福者的一派。因为谁比基督更蒙福？我们是基督一派，地极都因他而被呼召，都归向他，万国万民都在他面前朝拜。因此，这\Person{玛卡里}（即这有福者）的一派，并不畏惧你们从\Person{撒罗满}的话中扭曲而来的最后诅咒，即它必从地上灭亡。因为\Person{撒罗满}论不敬虔者所说的话，你们却试图应用于基督的产业，并以无法言说的不虔敬竭力证明这事已经成就；因为他论不敬虔者时说：\BibleQuote{他们的福分必从地上消失。}\BibleRef{箴 2:22}但是，当你们引用圣经的话\BibleQuote{我要将外邦人赐你为产业}和\BibleQuote{地极都要记念上主，并归向他}，说其中所包含的应许已从地上消失时，你们是在试图将关于不敬虔者命运的预言转过来反对基督的产业；然而，只要基督的产业持续并增长，你们在说这样的话时就在消亡。因为你们并非在一切情况下都藉着天主的圣事预言，因为在此你们只是出于自己的疯狂发出恶愿。但是，真先知的预言比假先知的说坏话更有能力。\par

% structure: p0136 source=h2 target=subsection reason=relative-heading-rank; base=h2

\subsection{第四十章。}

95. \Person{培提利安}说：\Quote{\Person{保禄}宗徒也命令我们：\InnerQuote{你们不要与不信者同负一轭：因为义与不义有何相通？光明与黑暗有何相通？基督与贝里雅耳有何和谐？信者与不信者有何相干？}\BibleRef{格后 6:14}}\par

96. \Person{奥古斯丁}回答：我认出宗徒的话，但我完全看不出它们如何能帮助你们。因为我们当中哪一位说，义与不义之间有任何相通，即使义人和不义之人，如\Person{犹达斯}和\Person{伯多禄}的例子，共同领受圣事？因为从同一个神圣之物，\Person{犹达斯}领受了审判，\Person{伯多禄}领受了救恩，正如你和\Person{奥普塔徒}一起领受圣事，你若与他不同，就不会因此有分于他的抢劫。难道抢劫不是不义吗？谁会疯狂到如此断言？那么，当你们一同走近同一祭台时，你的义与他人的不义有何相通？\par

% structure: p0139 source=h2 target=subsection reason=relative-heading-rank; base=h2

\subsection{第四十一章。}

97. \Person{培提利安}说：\Quote{而且，他以下列话语教导我们不应有分裂：\InnerQuote{我的意思是，你们每人各说：我是属\Person{保禄}的，我是属\Person{阿颇罗}的，我是属\Person{刻法}的，我是属基督的。基督被分裂了吗？\Person{保禄}为你们被钉了吗？或者你们受洗归于\Person{保禄}的名下吗？}\BibleRef{格前 1:12}}\par

98. \Person{奥古斯丁}回答：所有读这话的人，请记住是\Person{培提利安}引用了宗徒的这些话。因为谁能相信他会引用如此有利于我们反对他的话呢？\par

% structure: p0142 source=h2 target=subsection reason=relative-heading-rank; base=h2

\subsection{第四十二章。}

99. \Person{培提利安}说：\Quote{如果\Person{保禄}对这些无知的人和义人说了这些话，那么我对你们这些不义的人说：基督被分裂了吗，你们竟从教会分离？}\par

100. \Person{奥斯定}答道：我担心有人会认为，在本作品中，抄写员犯了错误，写下了标题\Emph{佩提利安说}，而他原本应该写\Emph{奥斯定答}。但我明白你的意图：你似乎要抢先占据立场，以免我们引用那些话作为不利于你的证词。然而你究竟做了什么？不过是使它们被重复引用了两次。因此，如果你如此喜欢听那些对你不利的话，以至于必须重复它们，那么，我恳请你，把这些话当作出于我、佩提利安所说的来听：\BibleQuote{基督被分裂了吗，你们竟要从教会中分开？}\par

% structure: p0145 source=h2 target=subsection reason=relative-heading-rank; base=h2

\subsection{第四十三章}

101. \Person{佩提利安}说：\Quote{难道是叛徒犹达斯为你自缢，或是他将他的品性注入了你，使你效法他的行为，夺取教会的财宝，并将我们这些基督的继承人卖给世俗的权势以换取金钱？}\par

102. \Person{奥斯定}答道：犹达斯没有为我们而死，但基督为我们死了，普世教会向祂说：\BibleQuote{这样我便有了回答那辱骂我者的理由，因为我信赖祢的话}。因此，当我听到主的话说：\BibleQuote{你们将要在耶路撒冷、全犹太和撒玛黎雅，直到地极，为我作证人}，\BibleRef{宗 1:8} 并且通过祂先知的声音说：\BibleQuote{他们的声音传遍全地，他们的言语达于地极}，任何邪恶的肉身掺杂都不能扰乱我，如果我知道向祂说：\BibleQuote{求祢为祢仆人作保，使我得福，不要让骄傲的人欺压我}。因此，当我有了确实的应许，我就不理会虚妄的诽谤。但如果你们抱怨那些你们曾经拥有、如今不再拥有的属于教会的事务或地方，那么犹太人也可能说自己是义人，并指责我们为非义，因为基督徒如今占据了他们从前不敬虔地统治的地方。那么，如果按照主类似的旨意，如今主教徒拥有从前异端者所拥有的事物，又有何不妥？因为所有这类人——即所有不敬虔和不义的人——主的话都有效力：\BibleQuote{天主的国必从你们手中夺去，赐给那能结果实的民族}，\BibleRef{玛 21:43} 或者岂非徒然记载：\BibleQuote{义人要吃不虔之人的劳苦}？\par

% structure: p0148 source=h2 target=subsection reason=relative-heading-rank; base=h2

\subsection{第四十四章}

103. \Person{佩提利安}说：\Quote{因为我们受洗时，正如经上所载，披戴了被出卖的基督；\BibleRef{迦 3:27} 而你们受感染时，披戴了叛徒犹达斯。}\par

104. \Person{奥斯定}答道：我也可以说，你们受感染时披戴了叛徒、强盗、压迫者、分离夫妻的\Person{奥普塔图斯}；但远非如此，免得以恶言相报的欲望诱使我说谎：因为你们并没有披戴\Person{奥普塔图斯}，我们也没有披戴犹达斯。因此，如果每一个来到我们这里的人回答我们的提问说，他是在\Person{奥普塔图斯}的名下受洗的，他必须在基督的名下受洗；如果你们为任何从我们这里来的人施洗，而他们说他们是在叛徒犹达斯的名下受洗的，那么我们对你们所做的便无可指责。但如果他们是在基督的名下受洗的，难道你们看不出，你们犯了一个多大的错误，竟以为天主的圣事能因人类罪孽的可变性而改变，或能因任何人的生命污秽而玷污？\par

% structure: p0151 source=h2 target=subsection reason=relative-heading-rank; base=h2

\subsection{第四十五章}

105. \Person{佩提利安}说：\Quote{但如果这些是派别，派别的名号并不对我们构成偏见。因为有两条路，一条是窄路，我们行走其上；另一条是为不虔敬的人，他们将在其中灭亡。然而，尽管名称相似，实际上却有很大的差别，免得因名称的共融而玷污了义路。}\par

106. \Person{奥斯定}答道：你害怕你的数目与普世中的众多相比，因此为了使你派稀少的数目赢得赞誉，你试图引入自己行走窄路的比较。但愿你坚持的不是其赞誉，而是那窄路本身！你确实会看到，在万民的教会中也有同样稀少的人；但义人相对于不义之人的众多而被说成是少数，正如相对于糠秕，在最丰盛的庄稼中可以说有少量的麦粒，而这些麦粒本身当聚成堆时，就装满谷仓。因为马克西米安的追随者自己就会在数目稀少上超越你，如果你认为义在于此，也在于他们失去所拥有地方所遭受的迫害。\par

% structure: p0154 source=h2 target=subsection reason=relative-heading-rank; base=h2

\subsection{第四十六章}

\Quote{在第一篇圣咏中，达味将有福的人与不虔敬的人分开，他并没有使他们分成党派，而是将所有不虔敬的人排除于圣洁之外。\InnerQuote{不随从不虔敬的人的计谋、不站罪人的道路的人，是有福的。}让那迷失了正义之道以致丧亡的人，再回到这道路上来吧。\InnerQuote{不坐于轻慢者的座位上。}当他发出这警告时，可悲的人哪，你们为何坐于那座位上？\InnerQuote{他喜爱上主的法律，昼夜默思祂的法律；他像一棵栽在溪水旁的树，按时结果，叶子不枯干；凡他所做的，都顺利。不虔敬的人却不是这样，他们像被风吹散的糠秕。}他蒙蔽了他们的眼睛，使他们看不见。\InnerQuote{因此，不虔敬的人在审判中站立不住，罪人在义人的集会中也站立不住。因为上主认识义人的道路，但不虔敬的人的道路必致丧亡。}}\par

奥古斯丁回答说：\Quote{圣经中谁不会区分这两类人呢？但你们却以糠秕的过犯来诬蔑麦粒；而你们自己不过是糠秕，却夸耀自己是唯一的麦粒。但是真正的先知宣告说，这两类人从起初就混合在一起，遍及整个世界，即遍及主的整个麦田，直到审判日要进行的簸扬。但我劝你们去读那第一篇圣咏的希腊文译本，这样你们就不敢以全体世界是\Person{玛卡里乌斯}派来责备世界了；因为你们也许会明白，在亚巴郎的后裔中蒙受祝福的万民中，所有圣人中间有一个\Person{玛卡里乌斯}派。因为我们语言中所谓\Quote{有福的人}，在希腊文是\OriginalTerm{有福的人}{\GreekText{Μαχάριος} \GreekText{ἀνήρ}}。但如果那冒犯你们的\Person{玛卡里乌斯}是个坏人，他既不属于这个区分，也不对它有所损害。但如果他是个好人，让他证明自己的工作，好使他在自己身上夸耀，而不在他人身上。\BibleRef{迦 6:4}}\par

% structure: p0157 source=h2 target=subsection reason=relative-heading-rank; base=h2

\subsection{第四十七章}

\Quote{但同一位圣咏作者歌颂了我们的圣洗圣事。\InnerQuote{上主是我的牧者，我必一无所缺。祂使我卧在青草地上，领我到幽静的溪水旁。祂使我的灵魂苏醒，为祂的名引导我走正义的路。纵使我走过死荫的幽谷——他指的是即使迫害者杀害我——我也不怕凶险，因为祢与我同在；祢的牧棒和牧杖安慰了我。}正是借着这受傅的油，它击败了哥肋雅。\InnerQuote{在我敌人面前，祢为我摆设了筵席；祢用油傅了我的头，我的杯爵满溢。愿恩宠与慈爱终生伴随我，我且要住在上主的殿中，直到永远。}}\par

110. 奥斯定答曰：这圣咏论及那些正确领受圣洗，并善用如此神圣之物的人。因为那些话甚至不涉及西满·玛古斯，他也领受了同一神圣的圣洗；但正因为他没有神圣地使用它，他并未玷污它，也未表明在此情况下应重复圣洗。但既然你提到了哥肋雅，请听那论及哥肋雅本人的圣咏，看看他是如何在\Quote{新歌}中被描绘的；因为那里说：\BibleQuote{我要向祢，天主，唱新歌：在十弦琴和瑟上，我要向祢歌咏。}并要看看，谁拒绝与普世教会共融，是否属于这歌。因为别处说：\BibleQuote{你们要向主唱新歌；普世大地，都要向主歌唱。}因此，普世大地——你并不与之合一——唱这新歌。而这也是普世大地的话：\BibleQuote{主是我的牧者，我实在一无所缺}等等。这些不是莠子的话，虽然它们在同一庄稼中被容忍直到收获。这些不是糠秕的话，而是麦子的话，虽然它们被同一雨水滋养，在同一打谷场上同时被敲打，直到末日在簸箕中彼此分离。然而，这两者确实拥有同一圣洗，尽管它们本身不同。但如果你们的党派也是天主的教会，你们肯定会承认这圣咏不适用于狂怒的\Quote{川西雍克派}群体。或者，如果他们自己也被引导行过正义的道路，那么当你们被责备与他们为伍时，为何否认他们是你们的同伙？尽管在大多数情况下，你们用\Quote{川西雍克派}的棍棒——而非主的杖和棍——来安慰你们派系的稀少，并自以为即使对抗罗马法律也是安全的——而与之冲突无疑是行过死荫的幽谷？但主与他同在的人，不怕邪恶。然而，你们必不敢说，这歌中所唱的话甚至属于那些狂怒的人，而你们不仅承认，还炫耀他们拥有圣洗。因此，这些话不被任何未被圣水滋养的人所用，正如所有天主的义人那样；也不被那些因使用它而走向毁灭的人所用，正如那术士在受斐理伯圣洗后那样：然而，两种人中的水本身是相同的，具有同样的神圣性。这些话只被那些将属于右边的人所用；但羊和山羊都在同一牧人下的同一牧场吃草，直到他们被分开，接受应得的报酬。这些话只被那些像伯多禄一样，从主的筵席获得生命而非审判的人所用，如犹达斯那样；然而，同一晚餐对两者是相同的，但对他们并非同等有益，因为他们不是一体。这些话只被那些藉圣油傅油，也在精神上蒙福的人所用，如达味；而非仅在身体上被祝圣，如撒乌耳：然而，两者既接受了同一外在标记，圣事并无不同，只是个人功过有别。这些话只被那些以归正的心领受主杯以致永生的人所用；而非那些吃喝自己审判的人，如宗徒所说：\BibleRef{格前 11:29}然而，虽然他们不是一体，他们所领受的杯却是一体，其能力在殉道者身上使他们获得天上的赏报，而非在\Quote{川西雍克派}身上，使他们以死亡标记悬崖。因此，请记住，恶人的品性绝不损害圣事的德能，以至于其神圣性被摧毁或甚至削弱；而是他们伤害不义之人本身，使他们拥有这些圣事作为他们受罚的见证，而非健康的帮助。因为无疑，你们本应考虑这圣咏的实际结尾词，并理解，由于那些在受圣洗后背弃信仰的人，所有领受圣洗的人不能说：\BibleQuote{我要永远住在主的殿中。}然而，无论他们坚守信仰，还是已经堕落，虽然他们本身不是一体，他们的圣洗却是一体，虽然他们本身不都是圣的，但两者中的圣洗都是圣的；因为即使背教者若回归，也不会重新受圣洗，仿佛失去了圣事，而是接受补赎，因为他们对仍存于他们内的圣事行了侮辱。\par

% structure: p0160 source=h2 target=subsection reason=relative-heading-rank; base=h2

\subsection{第四十八章。}

111. 培提利安说：\Quote{但为叫你们不自称圣洁，首先，我声明：谁若没有过无罪的生活，就不拥有圣洁。}\par

112. 奥斯定答曰：请给我们指出你被立为审判者的法庭，使普世世界都要在你面前受审，以及你用什么样的眼睛审查并讨论了——我不说良心——而是所有人的行为，以致你说普世世界失去了无罪。那位被提到第三层天上的人说：\BibleQuote{连我自己也不审判自己。}\BibleRef{格前 4:3}而你竟敢对普世世界——基督的产业遍布其中——作出判决？其次，如果在你看来，你所说的\Quote{谁若没有过无罪的生活，就不拥有圣洁}已足够确定，我要问你：如果撒乌耳没有圣事的圣洁，那么达味在他身上敬畏的是什么？但如果他有无罪，为何他迫害无罪之人？正是由于他傅油的圣洁，达味在他在世时尊敬他，在他死后为他复仇；并且因为他割下了他衣袍的一块，就心怀恐惧，战栗不已。这里你看到，撒乌耳没有无罪，但他却有圣洁——不是圣洁生活的个人圣洁（因为没有无罪，没有人能拥有它），而是天主圣事的圣洁，这圣洁甚至在恶人身上也是圣的。\par

% structure: p0163 source=h2 target=subsection reason=relative-heading-rank; base=h2

\subsection{第四十九章。}

113. 培提利安说：\Quote{因为，就算你们这些无信的人认识法律，我这么说绝无损害法律之意，魔鬼也认识法律。因为在义人约伯的事上，魔鬼曾回答主天主，好像自己是义人似的，正如经上所载：\BibleQuote{主对撒殚说：你曾用心察看我的仆人约伯没有？世上没有人像他那样完美、正直、敬畏天主，远离邪恶；他仍然坚守自己的忠诚，虽然你激动我攻击他，无故地毁灭他？}\BibleRef{约 2:3}撒殚回答主，以皮换皮，人为了自己的性命，情愿舍弃一切。看，他说话时引用了法律条文，即使他在对抗法律。第二次，他企图用这样的话来试探主基督，正如经上所记：\BibleQuote{魔鬼把耶稣带到圣城，叫祂站在圣殿顶上，对祂说：\InnerQuote{你若是天主子，就跳下去，因为经上记着：\InnerQuote{祂会为你吩咐祂的天使，他们要用手托着你，免得你的脚碰在石头上。}} 耶稣对他说：\InnerQuote{经上又记着：\InnerQuote{不可试探主，你的天主。}}}\BibleRef{玛 4:5}你们认识法律，我说，如同魔鬼一样，它在图谋中被打败，在行为上蒙羞。}\par

114. 奥斯定回答：我本来可以问你，当魔鬼对义人约伯口出诽谤时，它所用的那些话是记在哪部法律中，如果我要证明的是你自己不熟悉你所声称魔鬼认识的法律的话；但这不是我们之间的争论点，故略过不提。但你想方设法证明魔鬼精通法律，好像我们认为一切认识法律的人都是义人似的。因此，我不明白你引用魔鬼的例子对你有什么帮助——除非也许它能提醒我们你是如何效法魔鬼本身的。因为正如魔鬼用法律的话语来反对法律的颁布者一样，你也用法律的话语来控告你不认识的人，以抗拒那同一法律中所包含的天主的许诺。然后，我倒想知道，你的那些宣信者从悬崖跳下殉道，究竟是光荣谁？是光荣那位曾在类似建议前推开魔鬼的基督，还是光荣那向基督建议这种行为的魔鬼本身？自杀者所采用的两种特别卑劣而常见的死法是上吊和跳崖。你确实在这封信的前面部分说过：\Quote{叛徒吊死了：他把这种死亡留给了所有像他一样的人。}这话与我们毫不相干，因为我们拒绝尊崇任何自缢的人为殉道者。我们更有理由对你说：所有叛徒的主子——魔鬼，曾想说服基督从高处跳下，但被击退了！那么，对于那些被它成功说服的人，我们该说什么呢？的确，除了说他们是基督的仇敌、魔鬼的朋友、诱惑者的门徒、叛徒的同门吗？因为两者都从同一个师傅学会了自杀——犹达斯通过上吊，其他人通过跳崖。\par

% structure: p0166 source=h2 target=subsection reason=relative-heading-rank; base=h2

\subsection{第五十章}

115. 培提利安说：\Quote{但是，为了逐一驳斥你们的论点，如果你们以司铎自称，那么主天主曾通过祂先知的口说过：\InnerQuote{上主的报应临于假司铎。}}\par

116. 奥斯定回答：与其寻求你可以真实地说什么，不如寻求你说中伤话的来源；与其教你什么，不如教你用怎样的辱骂来唾弃我们。\par

% structure: p0169 source=h2 target=subsection reason=relative-heading-rank; base=h2

\subsection{第五十一章}

117. 培提利安说：\Quote{如果你们这些可怜的人为自己要求一个座位，正如我们之前所说的，你们无疑拥有那先知兼圣咏作者达味所称之为\InnerQuote{亵慢人的座位}的那一个。因为那座位恰当地留给了你们，因为圣人不能坐在其中。}\par

118. 奥斯定回答：你再次没有看到这不是什么辩论，而是空洞的辱骂。因为这是我刚才说的：你们口出法律的话语，却不注意对象是谁；正如魔鬼口出法律的话语，却未察觉它在对谁说。它想推翻我们的头，祂即将升到高处；但你却想把那同一位头的分布在全世界的身子缩减成一小部分。当然，你自己刚才说过，我们认识法律，口出法律术语，但在行为上蒙羞。你这样说了，却没有任何证据；但即使你证明某些人是这样，你也不能断言其他人都是这样。然而，如果全世界的人都像你凭空指控的那样，那么罗马教会的座位做了什么得罪了你？伯多禄曾坐在那里，今日由阿纳斯塔修斯占据；或者耶路撒冷教会的座位做了什么？雅各伯曾坐在那里，今日由若望占据，我们在公教合一中与这些座位联系，而你却以疯狂怒气与之分离。你为什么称宗座为亵慢人的座位？如果是因为你相信那些只用法律话语而不实行的人，那么你发现我们的主耶稣基督因法利塞人——关于他们祂说：\BibleQuote{他们只说不做}——而对他们的座位有任何不敬吗？祂不是称赞梅瑟的座位，维护座位的尊荣，同时谴责坐在其上的人吗？因为祂说：\BibleQuote{经师和法利塞人坐在梅瑟的座位上：凡他们对你们所说的，你们要行要守；但不要照他们的行为去做，因为他们只说不做。}\BibleRef{玛 23:2}如果你想到这些，你就不会因你所诽谤的人而对宗座不敬，而你在宗座上毫无份子。但像你这样的行为，除了不知该说什么、又无力禁止恶言之外，还有什么呢？\par

% structure: p0172 source=h2 target=subsection reason=relative-heading-rank; base=h2

\subsection{第五十二章}

\Person{Petilianus} 说：\Quote{如果你以为你能献祭，天主自己这样论及你们这些最败坏的罪人：\BibleQuote{恶人，祂说，献上牛犊，好似砍断狗的颈项；奉献素祭，好似献上猪血。}\BibleRef{依 66:3} 你们在此认出你们的祭品吧，你们已经流了人的血。祂又说：\BibleQuote{他们的祭物必如丧者的饼，凡吃这饼的必被玷污。}\BibleRef{欧 9:4}}\par

\Person{奥古斯丁} 回答：我们说，在每一个人身上，所献的祭品具有那前来献祭或前来领受者的品性；那些吃这样的人的祭品的人，是在前来领受时分享了献祭者的品性。因此，如果恶人向天主献祭，而善人从他手中领受，这祭品对每个人的品性就如他本人所显示的那样，因为我们也读到经上记载：\BibleQuote{对洁净的人，一切都是洁净的。}\BibleRef{铎 1:15} 根据这个真实而大公的评判，如果你们不赞成\Person{奥普塔图斯}的行为，你们也不会被他的祭品所玷污。因为他的饼确实是丧者的饼，因为整个阿非利加都在他的不义之下哀悼。但你们全体党派的分裂所涉及的邪恶，使这丧者的饼成为你们共同的饼。因为根据你们教务会议的判决，\Person{穆斯提的斐理西安}是一个流人血者。因为你们在谴责他们时说：\Quote{他们的脚飞快流人血。} 因此，请看看你们视为司铎的人献的是什么样的祭品，而你们自己已判定他犯有亵渎之罪。如果你们认为这丝毫不损害你们，那么我问你们，你们诽谤的空洞性又如何能损害全世界呢？\par

% structure: p0175 source=h2 target=subsection reason=relative-heading-rank; base=h2

\subsection{第五十三章。}

\Person{Petilianus} 说：\Quote{如果你向天主祈祷或发出恳求，这对你完全没有任何益处。因为你的血污良心使你的微弱祈祷无效；因为上主天主更看重良心的纯洁，而非祈求的话语，正如主基督所说：\BibleQuote{不是凡对我说\InnerQuote{主啊，主啊}的人，都能进入天国；而是那遵行我在天之父旨意的人。}\BibleRef{玛 7:21} 天主的旨意无疑是良善的，因此我们在圣祷中如此祈求：\BibleQuote{愿你的旨意奉行在人间，如同在天上。}\BibleRef{玛 6:10} 这样，既然祂的旨意是良善的，它就能赐予我们一切良善。因此，你们不做天主的旨意，因为你们每天都做邪恶的事。}\par

\Person{奥古斯丁} 回答：如果我们这边也把你们加于我们的一切反加于你们，那么任何听见的人，如果他自己头脑清醒，难道不会认为我们是疯狂的争吵者，而非基督信仰的辩论者吗？因此，我们不辱骂相报。因为\BibleQuote{主的仆人不应争辩；而应温良对待众人，乐意学习，用温柔的心教导那些反对的人。}\BibleRef{弟后 2:24} 因此，如果我们因你们中间每天行恶的人而责备你们，我们就有不适当的争辩之罪，用别人的罪来控告一个人。但如果我们提醒你们，既然你们不愿这些事被用来反对你们，所以你们也应停止用别人的罪来反对我们，那么我们就是在温柔地教导你们，唯一的希望是有一天你们会回到更好的心意。\par

% structure: p0178 source=h2 target=subsection reason=relative-heading-rank; base=h2

\subsection{第五十四章。}

\Person{Petilianus} 说：\Quote{但如果碰巧——虽然是否如此我不能说——你们驱逐魔鬼，这对你们也毫无益处；因为魔鬼本身既不屈服于你们的信德，也不屈服于你们的功德，而是因主耶稣基督的名被驱逐。}\par

\Person{奥古斯丁} 回答：感谢天主，你终于承认，即使由罪人呼号，基督之名的呼求也可能为别人的救恩带来益处。因此，你们可以明白，当基督之名被呼求时，一个人的罪不会妨碍另一个人的救恩。但关于我们如何呼求基督之名，我们不需要你们的判断，而是需要被我们呼求的基督自己的判断；因为只有祂能知道我们在何种精神中被呼求。然而，从祂自己的话中，我们确信，祂被所有在\Person{亚巴郎}后裔中蒙受祝福的万民呼求而得救恩。\par

% structure: p0181 source=h2 target=subsection reason=relative-heading-rank; base=h2

\subsection{第五十五章。}

\Person{Petilianus} 说：\Quote{即使你们做非常德善的行为，并显奇迹，但因你们的邪恶，主不认识你们；正如主自己所说：\BibleQuote{到那一天，有许多人要向我说：\InnerQuote{主啊，主啊，我们不是因你的名字说过预言吗？不是因你的名字驱逐过魔鬼吗？不是因你的名字行过许多奇迹吗？}那时我必要向他们声明说：\InnerQuote{我从来不认识你们；你们这些作恶的人，离开我吧！}}\BibleRef{玛 7:22}}\par

126. 奥古斯丁答：我们承认主的话。因此宗徒也说：\BibleQuote{我若有全备的信德，甚至能移山，但我若没有爱德，我什么也不算。}\BibleRef{格前 13:2} 因此我们应当探究谁拥有爱德：你会发现，没有别人，只有那些爱慕合一的人。至于驱魔和行奇迹，既然有许多不属于天国的人也做这些事，而许多属于天国的人却不做，所以无论我们一方还是你们一方，若有人偶然有此能力，都不可夸耀，因为主认为连那些能真正行这些事以利人得救的宗徒，也不应当以此夸耀，当祂对他们说：\BibleQuote{你们不要因为魔鬼屈服于你们而喜欢，你们应当喜欢的，乃是因为你们的名字，已经登记在天上了。}\BibleRef{路 10:20} 因此，你们从福音书中所提出的一切，我也可以对你们重复，如果我看见你们行了奇迹异能的伟大作为；同样，如果你们看见我作了类似的事，你们也可以对我重复。所以，我们不要彼此说那些同样可以站在对方立场上说的话；抛开一切诡辩，既然我们探究的是何处能寻得基督的教会，就让我们听从基督自己的话，祂以自己的血赎回了教会：\BibleQuote{你们要在耶路撒冷、犹太全地、撒玛黎雅，直到地极，为我作见证。}\BibleRef{宗 1:8} 那么，你们看，谁若不愿与这个传遍天下的教会共融，就是拒绝与谁共融——如果你们至少听见了这是谁的话。因为还有什么比领受主的圣事，却拒绝与主的话共融，更证明疯狂呢？这样的人很可能要说：\BibleQuote{因祢的名我们吃过喝过}，并且听到这些话：\BibleQuote{我从来不认识你们。}\BibleRef{玛 7:22} 因为他们领受了祂的体血，在福音中却不承认祂散布天下的肢体，因此他们在审判中也不被算在肢体之内。\par

% structure: p0184 source=h2 target=subsection reason=relative-heading-rank; base=h2

\subsection{第五十六章}

佩提利安说：\Quote{但是，即使如你们自己所设想的那样，你们纯洁地遵守主的法律，然而让我们以法律的形式考虑至圣的法律本身。宗徒保禄说：\BibleQuote{法律原是好的，只要人合法地使用它。}\BibleRef{弟前 1:8} 那么法律说什么？\BibleQuote{不可杀人。}杀人犯加音只做了一次，而你们在杀害弟兄的事上却常常做。}\par

奥古斯丁答：我们不愿像你们一样：因为并不缺少可以说出的话，如同你们也说出这些；也不缺少已知的事，因为你们不知道这些；也不缺少生活中所表现的行为，因为你们没有表现出这些。\par

% structure: p0187 source=h2 target=subsection reason=relative-heading-rank; base=h2

\subsection{第五十七章}

佩提利安说：\Quote{经上记载：\BibleQuote{不可奸淫。}你们每个人，即使身体上保持贞洁，心灵上却是奸淫者，因为他玷污了自己的圣洁。}\par

奥古斯丁答：这些话也可以真实地用来指责我们和你们中的某些人；但如果他们的行为受到我们和你们同样的谴责，那么他们既不属于我们，也不属于你们。但你们却希望，你们针对某些人所说的话，即使没有在他们的具体个案中证明，也应当被当作已经成立——不是针对一些从亚巴郎的后裔中坠落的人，而是针对天下所有在亚巴郎的后裔中获得祝福的万民。\par

% structure: p0190 source=h2 target=subsection reason=relative-heading-rank; base=h2

\subsection{第五十八章}

佩提利安说：\Quote{经上记载：\BibleQuote{不可作假见证陷害你的近人。}当你们向这世界的君王假称我们持有你们的意见时，你们岂不是在制造谎言吗？}\par

132. 奥斯定回答说：如果你们所持的并不是我们的意见，那么你们从马克西米安派那里接受的也不是你们的意见。但如果那些意见之所以是你们的，是因为他们犯了亵渎圣事的裂教之罪，不与多纳特派共融，那你们就要注意自己所处的立场，以及你们拒绝与谁的产业共融，并思考你们将如何向你们的主基督，而不是向今世的君王，作出答复。论到祂，经上说：\BibleQuote{祂的统治权从这海到那海，从大河直到地极。}这里指的是哪条河呢，难道不是祂受洗、鸽子降在祂身上的那条河吗？那正是爱德与合一的伟大标记。但你们拒绝与这合一共融，却至今仍占据着合一的场所；你们利用总督的法令将你们的裂教者从多纳特派的场所驱逐出去，借此使我们不得今世君王的欢心。这些话并非凭空在虚空中飞驰：那些人还活着，各城邑为此作证，总督和各城镇的档案也引以为证。因此，让那毁谤的声音终于沉默吧，它本想将今世的君王，就是那些通过他们你们这些碎片也不肯宽恕从你们那里分裂出去的碎片的总督们，拿来反对整个大地。当我们说你们持有我们的意见时，我们并不显得是在作假见证，除非你们能证明我们不在基督的教会内，而这正是你们不断声称却永远无法证明的；事实上，当你们说这话时，你们不是在控告我们作假见证，而是在控告主自己。因为我们是在教会内，这是祂亲自的见证所预言的，祂向自己的见证人作证说：\BibleQuote{你们在耶路撒冷、犹太全地、撒玛黎雅，甚至直到地极，要为我作证人。}但你们不仅因为抗拒这真理而显出自己是假见证人，而且在你们与马克西米安派裂教争讼的案件中也显出如此。因为如果你们是依照基督的法律行事，那么某些信奉基督的皇帝依此制定法令岂不更为合理？如果连外教总督在判案时都能遵循祂的旨意，那又当如何？但即使你们认为地上的帝国的法律也应被召来协助你们，我们并不为此责怪你们。这正是保禄所做的，他在敌人面前见证自己是罗马公民。\BibleRef{宗 22:25}但我想问，是哪条地上的法律规定将马克西米安派从他们的场所驱逐出去？你们根本找不到任何这样的法律。事实上，你们选择根据反对异端的法律来驱逐他们，而你们自己也属于那些法律所反对的人之列。你们仿佛以更强的力量胜过了弱者。因此，他们完全无力，声称自己是无辜的，就像狮子爪下的狼一样。然而，你们若没有假见证的帮助，你们一定不能利用那些反对你们自己的法律来反对别人。因为如果那些法律基于真理，那么你们就该从你们所占据的位置下来；但如果基于虚假，你们为何用它们将他人从教会驱逐？但若它们既基于真理，而你们又只能在假见证的帮助下才能用来驱逐别人，那又当如何？为使法官服从它们的权威，他们愿意将异端者从教会驱逐出去，而他们本应先驱逐你们；但你们宣称自己是公教徒，以便逃避你们用来压迫别人的法律的严苛。你们在自己中间看自己如何，由你们决定；无论如何，根据那些法律，你们不是公教徒。那么，你们为何要么通过假见证使它们变成虚假（如果它们是真实的），要么利用它们（如果它们是虚假的）来压迫别人？\par

% structure: p0193 source=h2 target=subsection reason=relative-heading-rank; base=h2

\subsection{第五十九章}

133. 培提利安说：\Quote{经上说：\BibleQuote{不可贪图你邻人的任何东西。}\BibleRef{出 20:13}}\par

134. 奥斯定回答说：合一所拥有的一切事物，不属于别人，只属于我们这些留在合一中的人，并非根据人的毁谤，而是根据基督的话，万邦万民都因祂而蒙受祝福。我们也不因那些不义的人而远离麦子的团体（在审判扬场之前，我们不能将他们从主的麦子中分别出来）；如果有任何事物是你们这些被砍下来的已经开始的拥有，我们并不因为主已将从你们那里夺去的赐给我们，而贪图邻人的财物，因为这一切已因万物所属之主的权威而成为我们的；它们理当属于我们，因为你们惯于将它们用于分裂，而我们却用于促进合一。否则，你们一派甚至可以责斥天主的首批子民贪图邻人的财物，因为他们在天主的大能面前被驱逐出去，因为他们妄用了土地；而犹太人自己（根据主的话，天国从他们那里夺去，赐给了那结果实的民族）\BibleRef{玛 21:43}也可以控告那个民族贪图邻人的财物，因为如今在基督的教会占据的地方，曾是基督的迫害者统治的地方。此外，当有人对你们说：\InnerQuote{你们是在贪图别人的财物，因为你们将马克西米安派从圣堂中驱逐出去}，你们却找不到任何可回答的话。\par

% structure: p0196 source=h2 target=subsection reason=relative-heading-rank; base=h2

\subsection{第六十章}

135. 培提利安说：\Quote{那么，你们凭着什么法律证明自己是基督徒，既然你们所行与法律相反？}\par

136. 奥斯定回答说：你们渴望争辩，而非讨论。\par

% structure: p0199 source=h2 target=subsection reason=relative-heading-rank; base=h2

\subsection{第六十一章}

137. 帕提里亚努斯说：\Quote{但主基督说：\InnerQuote{谁若实行并教导这些诫命，他在天国里将称为大的。}但祂这样谴责你们这些可怜的人：\InnerQuote{谁若废除这些诫命中的一条，他在天国里将称为最小的。}}\par

138. 奥古斯丁答：当你碰巧引用《圣经》的证词，而其内容与真实状况不同，且与我们所争论的问题无关时，我并不太在意；但当它涉及当前的问题，而又未正确引用时，如果你不介意的话，我认为我有权利提醒你那段经文究竟是怎么说的。因为你知道，你所引用的经文并非如你所引，而是这样：\Quote{谁若废除这些最小诫命中的一条，并这样教导人，他在天国里将称为最小的；但谁若实行并教导它们，他在天国里将称为大的。}紧接着祂又说：\Quote{我告诉你们：你们的义德若不超过经师和法利塞人的义德，你们决不能进入天国。}\BibleRef{玛 5:19}因为在别处祂曾指出并证明法利塞人\InnerQuote{他们说而不做}。因此，祂在这里所指的也正是这些人，当祂说：\Quote{谁若废除这些诫命中的一条，并这样教导人}——也就是说，用言语教导他在行为上所违反的；祂说，我们的义德必须超过这些人的义德，即我们必须遵守诫命并这样教导人。然而，即便对于那些法利塞人——你们拿他们来与我们相比，不是出于谨慎而是出于恶意——我们主并非命令要遗弃梅瑟的座位，这座位无疑意指祂自己的一个预表；因为祂确实说过，坐在梅瑟座位上的那些人总是说而不做，但祂警告百姓要实行他们所说的，而不要效法他们所行的\BibleRef{玛 23:2}，以免这座位连同其全部神圣性被遗弃，并且羊群的合一因牧者的不忠而分裂。\par

% structure: p0202 source=h2 target=subsection reason=relative-heading-rank; base=h2

\subsection{第六十二章}

139. 帕提里亚努斯说：\Quote{又经上写着：\InnerQuote{凡人所犯的一切罪都在身体以外；但那在圣神内犯罪的，无论是在今世或是在来世，都不会被赦免。}}\par

140. 奥古斯丁答：这也不是如你所引的那样写的，并且你看它把你引向何处。宗徒在致格林多人书中说：\Quote{凡人所犯的一切罪都在身体以外；但犯邪淫的，却是得罪自己的身体。}\BibleRef{格前 6:18}但这是一回事，而主在福音中所说的是另一回事：\Quote{一切罪过和亵渎都可得赦免；但说话干犯圣神的，无论是在今世或是在来世，都不会被赦免。}\BibleRef{玛 12:31}但你却以宗徒书信中的语句开头，而以福音中的语句结束，我想你这样做并非有意欺骗，而是由于错误；因为这两段经文都与当前的问题无关。你说这话的原因和意义，我完全无法理解，除非是因为，你之前说过凡违反任何一条诫命的人都为主所定罪，而如今你考虑到你们党派中有多少人违反的不仅一条而是许多条诫命；为了避免在这方面受到反驳，你试图通过引入一种罪的区分来超越困难，即违反一条容易获得宽恕的诫命是一回事，而干犯圣神、在今生和来世都不得赦免的罪是另一回事。因此，出于对罪的感染的恐惧，你不愿对此事保持沉默；而又出于对你能力不足以应对的深刻问题的恐惧，你希望顺带粗略地提及它，以至于如此慌乱，就像那些匆忙着手一项任务而陷入混乱的人，常会把衣服或鞋子穿歪一样，你也不屑于去查看那段经文属于什么，或者它出现的上下文和含义是什么。至于那种在今生和来世都不得赦免的罪的性质，你如此无知，以至于虽然你相信我们正处于这种罪中，你们却通过你们的圣洗向我们承诺赦免它。然而，如果这罪的性质是它在今生和来世都不得赦免，那怎么可能呢？\par

% structure: p0205 source=h2 target=subsection reason=relative-heading-rank; base=h2

\subsection{第六十三章}

141. 帕提里亚努斯说：\Quote{然而你们在哪些方面遵守了天主的诫命？主基督说：\InnerQuote{神贫的人是有福的，因为天国是他们的。}但你们在迫害中因恶毒而散播疯狂的财富。}\par

142. 奥古斯丁答：这话你还是对你们自己的\OriginalTerm{circumcelliones}{Circumcelliones}说吧。\par

% structure: p0208 source=h2 target=subsection reason=relative-heading-rank; base=h2

\subsection{第六十四章}

143. 帕提里亚努斯说：\Quote{\InnerQuote{温良的人是有福的，因为他们要承受大地。}你们既不温良，便同时丧失了天堂和大地。}\par

144. 奥古斯丁答：你可以一次又一次地听到主说：\Quote{你们将为我作证，在耶路撒冷、全犹太和撒玛黎雅，并直到地极。}\BibleRef{宗 1:8}那么，那些为了避免与地上万民共融而藐视坐在天上者的话的人，他们怎么就没有丧失天堂和大地呢？至于证明你们的温良，不是你们的话，而是\OriginalTerm{circumcelliones}{Circumcelliones}的棍棒需要被检查。你们会说：这与我们何干？就好像我们这样说是为了从你们那里引出那个回答似的。因为你们的裂教是你们的罪责，其理由在于你们不承认自己应承担他人的罪责，而你们之所以与我们分离，唯一的原因就是你们指控我们承担了别人的罪。\par

% structure: p0211 source=h2 target=subsection reason=relative-heading-rank; base=h2

\subsection{第六十五章}

145. 帕提里亚努斯说：\Quote{\InnerQuote{哀恸的人是有福的，因为他们要受安慰。}你们，我们的刽子手，是他人哀恸的原因：你们自己并不哀恸。}\par

146. \Person{奥斯定}回答说：请稍加思量，从你的武装人员口中发出的\Quote{赞美归于天主}的呼声，是何等强烈地使多少人哀恸。你又说，\Quote{这与我们何干？}那么我也要再次用你的话回敬你：\Quote{这与我们何干？}这与地上的万民何干？这与那些从日出之地到日落之处赞美主名的人何干？这与那歌唱新歌的全地何干？这与\Person{亚巴郎}的后裔何干，地上万民都要因他而得福？\BibleRef{创 22:18}因此，你们分裂的亵渎之罪要归在你们身上，正是因为你们的同伴的恶行不归在你们身上；也正因为如此，那些你们因其而将自己与世人分离之人的行为，即使你们证明你们的指控属实，也不能使世人有罪。\par

% structure: p0214 source=h2 target=subsection reason=relative-heading-rank; base=h2

\subsection{第六十六章}

147. \Person{培提利安}说：\Quote{\BibleQuote{饥渴慕义的人是有福的，因为他们要得饱饫。} 在你们看来，你们渴求我们的血似乎就是义德。}\par

148. \Person{奥斯定}回答说：人啊，除了说你是诽谤者，我还能对你说什么呢？\Person{基督}的合一确实正渴望并渴求你们众人；我希望它能将你们吞没，因为那样你们就不再是异端了。\par

% structure: p0217 source=h2 target=subsection reason=relative-heading-rank; base=h2

\subsection{第六十七章}

149. \Person{培提利安}说：\Quote{\BibleQuote{怜悯人的人是有福的，因为他们要受怜悯。} 但你们既然惩罚义人，我怎能称你们为怜悯人呢？我岂不更应当称你们为一个极不义的团体，只要你们还在玷污灵魂？}\par

150. \Person{奥斯定}回答说：你们两点都没有证明——既没有证明你们自己是义人，也没有证明我们惩罚了不义的人；然而，正如虚假的谄媚通常是残忍的，正义的纠正永远是仁慈的。因为你们所不明白的那句话：\BibleQuote{愿义人击打我，这算是仁慈；愿他责备我。}他这话是对仁慈纠正的严厉而言的；诗人紧接着又说到毁灭性谄媚的柔和：\BibleQuote{但罪人的油不可折断我的头。}因此，请你们想一想，你们被召往何处，又从何处被唤回。因为你们怎么知道，你们认为残忍的人对你们怀有怎样的感情呢？无论他怀有何种感情，每个人都要背负自己的重担，无论在我们这边还是在你们那边。但我希望你们放弃你们所有人所背负的分裂的重担，好使你们在合一中背负你们美好的重担；我愿你们仁慈地纠正所有背负邪恶重担的人，如果你们有能力的话；如果这超出了你们的能力，我愿你们在平安中容忍他们。\par

% structure: p0220 source=h2 target=subsection reason=relative-heading-rank; base=h2

\subsection{第六十八章}

151. \Person{培提利安}说：\Quote{\BibleQuote{心里洁净的人是有福的，因为他们要看见天主。} 你们这些心怀不洁恶毒、瞎眼的人，什么时候才能看见天主呢？}\par

152. \Person{奥斯定}回答说：你为何这样说？难道我们用那些人所传述的隐晦不明之事来斥责万民，却不选择去理解\Person{天主}在古时对地上万民所说的明明白白的话吗？这真是心中极大的瞎眼；如果你们在自己身上认不出这一点，那就是更大的瞎眼。\par

% structure: p0223 source=h2 target=subsection reason=relative-heading-rank; base=h2

\subsection{第六十九章}

153. \Person{培提利安}说：\Quote{\BibleQuote{缔造和平的人是有福的，因为他们要称为天主的子女。}\BibleRef{玛 5:3} 你们以你们的邪恶假装和平，却以战争寻求合一。}\par

154. \Person{奥斯定}回答说：我们不是以邪恶假装和平，而是从福音宣讲和平；如果你们与福音和好，你们也会与我们和好。复活的主在显现给门徒时，不仅让他们用眼注视祂，也让他们用手触摸祂，祂以这话开始对他们讲论：\BibleQuote{愿你们平安。}随后祂又向他们揭示了这平安当如何持守。因为\BibleQuote{祂开启了他们的明悟，叫他们理解圣经，又对他们说：经上这样记载，基督必须受苦，第三天从死者中复活；并且要因祂的名，从耶路撒冷开始，向万民宣讲悔改和罪过的赦免。}如果你们与这些话保持平安，你们就不会与我们相争。因为如果我们寻求合一是通过战争，那么我们的战争不可能用更光荣的言辞来赞美，因为经上记载：\BibleQuote{你应爱你的近人如你自己。}\BibleRef{玛 22:39}又记载：\BibleQuote{从来没有人恨过自己的肉身。}\BibleRef{弗 5:29}然而，\BibleQuote{肉身与圣神相争，圣神与肉身相争。}\BibleRef{迦 5:17}但若从来没有人恨过自己的肉身，而人却与自己的肉身相争，这就是你们所见的通过战争寻求合一，使身体因受纠正而服从。但圣神对肉身所行的，与它作战，不是出于仇恨而是出于爱，这正是属神的人对属血肉的人所行的，好使他们对他们所行的正如对自己所行的一样，因为他们爱近人如同真正的近人。但属神的人所进行的战争是爱中的纠正：他们的剑是天主的话。宗徒的号角以极大的力量吹响，唤醒了这样的战争：\BibleQuote{宣讲圣言，无论顺境逆境，务要坚持；以各种忍耐和教导，斥责、警告、劝勉。}\BibleRef{弟后 4:2}可见，我们不是用剑，而是用圣言行事。但你们回答的不是真话，反而诬告我们。你们不纠正自己的过错，却把别人的过错归咎于我们。\Person{基督}为地上的万民作真实的见证；你们却反对\Person{基督}，对地上的万民作假见证。如果我们相信你们而不信\Person{基督}，你们就会称我们为缔造和平的人；因为我们相信\Person{基督}而不信你们，我们就被称为以邪恶假装和平的人。而你们在说和做这样的事的同时，还厚颜无耻地引用这话：\BibleQuote{缔造和平的人是有福的，因为他们要称为天主的子女。}\par

% structure: p0226 source=h2 target=subsection reason=relative-heading-rank; base=h2

\subsection{第七十章}

\Person{Petilianus} 说：\Quote{虽然宗徒保禄说：\BibleQuote{所以我这为主被囚的劝你们，弟兄们，你们行事为人应与所蒙的召叫相称，凡事要谦逊、温和、忍耐，在爱德中彼此担待，尽力以和平的联系保持心神的合一。}}\BibleRef{弗 4:1}\par

\Person{奥斯定} 回答说：如果你们不仅说出这些话，而且肯听从它们，你们就会为了和平而容忍已知的恶行，而不是为了争吵而制造新的恶行，哪怕仅仅是因为你们后来为了多纳徒的和平，学会了容忍奥普塔徒最昭彰的恶行。这是何等疯狂？他们容忍已知的，以免碎片进一步分裂；他们诽谤无人知晓的，以免他们留在未分裂的整体中。\par

% structure: p0229 source=h2 target=subsection reason=relative-heading-rank; base=h2

\subsection{第七十一章}

\Person{Petilianus} 说：\Quote{先知对你们说：\BibleQuote{和平，和平；哪有和平？}}\BibleRef{耶 8:11}\par

\Person{奥斯定} 回答说：是你们对我们说这话，不是先知。因此我们回答你们：若你们问和平在哪里，睁开你们的眼睛，看看论及谁曾说过：\Quote{他使战争在全地止息。} 若你们问和平在哪里，睁开你们的眼睛看看那建在山上的城，是不能隐藏的；睁开你们的眼睛看看那山本身，让\Person{达尼尔}指给你们看，它从一颗小石头长出，充满了全地。\BibleRef{达 2:35} 但当先知对你们说：\Quote{和平，和平；哪有和平？} 时，你们将展示什么？你们将展示多纳徒派，不为无数认识基督的万民所知？那城确实不能隐藏；但它从哪里来，岂不是因为不是建在山上？\BibleQuote{因为基督是我们的和平，他使双方合而为一，}\BibleRef{弗 2:14}——不是多纳徒，他使一分为二。\par

% structure: p0232 source=h2 target=subsection reason=relative-heading-rank; base=h2

\subsection{第七十二章}

\Person{Petilianus} 说：\Quote{\BibleQuote{为义而受迫害的人是有福的，因为天国是他们的。}\BibleRef{玛 5:10} 你们不是有福的；但你们使殉道者成为有福的，他们的灵魂充满了诸天，大地因他们的纪念而繁茂。因此你们自己不尊敬他们，却为我们提供了尊敬的对象。}\par

\Person{奥斯定} 回答说：显然，如果没有说\Quote{为义而受迫害的人是有福的}，而是说什么\Quote{从悬崖跳下去的人是有福的}，那么诸天就会充满你们的殉道者。诚然我们看到许多鲜花从他们的身体在大地上绽放；但正如俗语所说，那花是灰烬。\par

% structure: p0235 source=h2 target=subsection reason=relative-heading-rank; base=h2

\subsection{第七十三章}

\Person{Petilianus} 说：\Quote{既然你们不是有福的，反而歪曲天主的诫命，主基督以他神圣的法令谴责你们：\BibleQuote{祸哉，你们经师和法利塞人，假善人！因为你们给人封闭了天国：你们自己不进去，也不让那要进去的人进去。祸哉，你们经师和法利塞人，假善人！因为你们走遍海洋和陆地，为使一个人成为归依者；当他成了，你们却使他成为比你们更坏的地狱之子。祸哉，你们经师和法利塞人，假善人！因为你们捐献薄荷、茴香和莳萝的十分之一，却忽略了法律上更重要的事：公义、怜悯和信德；这些你们该做，那些也不可忽略。盲目的向导！你们滤出蚊蚋，却吞下了骆驼。祸哉，你们经师和法利塞人，假善人！因为你们好像粉饰的坟墓，外面看来美丽，里面却充满死人的骨骸和一切污秽。同样，你们外面在人前显出正义，里面却充满伪善和不法。}}\par

\Person{奥斯定} 回答说：请告诉我，你们所说的哪一句话不能同样被任何一个诽谤和恶意的舌头对你们说呢？但从我先前所说的话，任何人若愿意就可以发现，这些话可以针对你们说，不是空洞的辱骂，而是有真实见证的支持。然而，既然有机会，我们不可放过。毫无疑问，对古代天主的子民来说，割损相当于圣洗。因此我问，假如那些你们所引述的经文斥责的法利塞人，曾使某人成为归依者，这人若模仿他们，就如经文所说，会变成比他们更坏的地狱之子；假设他后来悔改，渴望效法\Person{西默盎}、\Person{匝加利亚}或\Person{纳塔乃耳}，那么他们是否需要为他重新行割损？若此假设是荒谬的，那么为什么，尽管你们空洞地把我们比作这样的人，你们却仍在我们之后施洗？但如果你们真是这样的人，那么我们不从你们之后施洗，岂不是更好、更符合真理？正如我所提及的那些人不应在最坏的法利塞人之后行割损一样。此外，当这样的人坐在\OriginalTerm{梅瑟的座位}{the seat of Moses}上（主曾为这座位保留了应有的尊荣），你们为何因着一些人（无论你们有正当理由还是不正当理由把他们比作法利塞人）而亵渎\OriginalTerm{宗徒之座}{apostolic chair}呢？\par

% structure: p0238 source=h2 target=subsection reason=relative-heading-rank; base=h2

\subsection{第七十四章}

\Person{Petilianus} 说：\Quote{但这些事并不吓倒我们基督徒；因为我们早已从主基督那里得到关于你们将行恶事的警告。他说：\Quote{\BibleQuote{看，我派遣你们像羊进入狼群中。}\BibleRef{玛 10:16}} 你们充满狼的疯狂，你们布设或准备布设陷阱攻击教会，正是以狼的方式，它们对羊栏张着大口，甚至带着毁灭的贪婪，从它们的下颚呼出充满怒气的喘息，下颚浸透了鲜血。}\par

一六四．奥古斯丁答道：我很乐意对你说同样的看法，但不是用你所用的言辞：它们太不当，甚至可说疯狂。但所需的是，你应表明我们是狼而你们是羊，不是靠最空洞的毁谤，而是靠某些明确的证据。因为当我也说：\Quote{我们是羊，你们是狼}时，你认为因你用浮夸的言辞表达这一想法就会有差别吗？但请听我证明我所断言的事。因为主在福音中说道，你知道得很清楚，不管你愿意与否：\BibleQuote{我的羊听我的声音，也跟随我。}\BibleRef{若 10:27}主在不同的事上说许多话；但假设，例如，有人怀疑同一主是否肉身复活，而要引用他的话：\BibleQuote{你们摸摸我，看看；因为鬼魂没有肉身和骨骼，如同你们看见我所具有的}——即使在此之后，他仍不愿相信他的身体从死者中复活，那么这人必定不能算作主的羊，因为他不听他的声音。所以现在也是这样，当我们之间的问题是：教会在哪里？我们引用同一段福音中随后的话，其中，在他复活后，他将自己的身体甚至给那些怀疑的人触摸，藉此他显示了教会未来的广阔范围，说：\BibleQuote{经上这样记载：基督必须受苦，第三天从死者中复活；并且要因他的名，从耶路撒冷开始，向万民宣讲悔改和罪赦。}然而你们不愿与万民共融——在这些民族中这些话已经应验——你们如何算是这位牧人的羊呢？你们不仅听到他的话而不服从，甚至还要反对它们？因此我们由此向你们表明你们不是羊。但再听我们从何处向你们表明你们相反却是狼。因为当用他自己的话指明教会应在何处时，也自然清楚我们须在哪里寻找基督的羊栈。因此，每当任何羊群从这个羊栈（由主明确的宣告清楚指示给我们的羊栈）分离出来——我不说是由于虚假的指控，而是由于不可否认地、带着极大不确定性对同侪的指控，并因此杀害那些被他们撕碎、从合一生活和基督之爱中疏离的羊——难道不明显他们是凶残的狼吗？但或许有人说，这些人自己赞美并宣讲主基督。他们因此就是他所说的：\BibleQuote{他们来到你们面前，外面披着羊皮，内里却是凶残的狼。凭着他们的果实你们可认出他们来。}\BibleRef{玛 7:15}羊皮见于对基督的赞美；狼性的果实则见于他们毁谤的牙齿。\par

% structure: p0241 source=h2 target=subsection reason=relative-heading-rank; base=h2

\subsection{第七十五章。}

一六五．Petilianus说：\Quote{不幸的\OriginalTerm{叛教者}{traditors}！的确，圣经必须这样应验。但在你们身上，我为此悲痛，因为你们显示自己配得上履行邪恶的角色。}\par

一六六．奥古斯丁答道：我或许更要说，不幸的\OriginalTerm{叛教者}{traditors}！如果我有意，或者不如说如果正义促使我向你们所有人提出你们中间某些人的行为。但就涉及到你们全体而言，不幸的异端者们，我这边要引用你余下的话；因为经上记载：\BibleQuote{在你们中间难免有异端，好叫那些经得起考验的人，在你们中间显明出来。}\BibleRef{格前 11:19}因此\Quote{圣经必须这样应验。但在你们身上，我为此悲痛，因为你们显示自己配得上履行邪恶的角色。}\par

% structure: p0244 source=h2 target=subsection reason=relative-heading-rank; base=h2

\subsection{第七十六章。}

一六七．Petilianus说：\Quote{但对我们而言，主基督，与你们致命的命令相反，命令了单纯的忍耐与无害。因为祂说什么？\BibleQuote{一条新诫命我给你们，你们该彼此相爱；如同我爱了你们，你们也该照样彼此相爱。}\BibleRef{若 13:34}又说：\BibleQuote{如果你们彼此相爱，众人就可认出你们是我的门徒。}}\par

一六八．奥古斯丁答道：如果你们不把这些与你们的品格如此迥异的言辞转移到你们谈话的表面，你们如何能用羊皮遮盖自己呢？\par

% structure: p0247 source=h2 target=subsection reason=relative-heading-rank; base=h2

\subsection{第七十七章。}

一六九．Petilianus说：\Quote{宗徒保禄，在受到万民可怕的迫害时，还受了假弟兄更痛苦的折磨，正如他亲自证实的，常常受苦：\BibleQuote{在危险中，有外邦人的危险，有本国人的危险，有假弟兄的危险。}\BibleRef{格后 11:26}他又说：\BibleQuote{你们效法我，如同我效法基督一样。}\BibleRef{格前 11:1}因此，当像你们这样的假弟兄攻击我们时，我们效法我们的师傅保禄在危险中的忍耐。}\par

一七〇．奥古斯丁答道：当然，你所说的那些人是假弟兄，宗徒在另一处地方抱怨他们，当时他赞扬弟茂德天性的真诚：\Quote{我没有一个像他那样诚心关切你们的事的人，}他说，\Quote{因为众人都谋求自己的事，不求耶稣基督的事。}\BibleRef{斐 2:20}无疑他是在说他写这封信时与他在一起的人；因为不可能各地所有的基督徒都谋求自己的事，而不求耶稣基督的事。因此，正如我所说，他是在说这哀叹时与他在一起的那些人。因为他在另一处地方说：\Quote{外有战斗，内有恐惧}\BibleRef{格后 7:5}，除了因那些在内部而更加惧怕的人，他还能指谁呢？因此，如果你们要效法保禄，你们就会容忍内部的假弟兄，而不是诽谤外部的无辜者。\par

% structure: p0250 source=h2 target=subsection reason=relative-heading-rank; base=h2

\subsection{第七十八章。}

171. \Person{培提利安}说：\Quote{你们中间那缺乏爱德的信德算什么？因为\Person{保禄}自己说：\InnerQuote{\BibleQuote{我若能说人间的语言，和能说天使的语言，但我若没有爱德，我就成了个发声的锣，或发响的钹。我若有先知之恩，又明白一切奥秘和各种知识；我若有全备的信德，甚至能移山，但我若没有爱德，我什么也不算。我若把我所有的财产全施舍了，又舍身投火被焚，但我若没有爱德，为我毫无益处。}}}\par

172. \Person{奥古斯丁}回答说：这正是我刚才所说的，你们渴望披上羊皮，好让羊群在尚未意识到你们的接近之前，就感受到你们的咬伤。难道不是你们所陶醉的这爱德赞歌，常常在最清晰的光照下显明你们的诽谤吗？你们能设法使这些武器不再属于我们，因为你们先企图将它们据为己有吗？此外，这些武器是有生命的：无论从何处发出，它们都认得出应该摧毁谁。如果它们是从我们手中发出的，它们会钉住你们；如果你们瞄准它们，它们会反弹到你们身上。因为在这段颂扬爱德优越性的宗徒话语中，我们惯常向你们指明，人若拥有信德或圣事却没有爱德，这对他多么无益，好使你们来到公教会合一中时，能明白加给你们的恩赐是什么，以及你们至少在某些方面所缺乏的是多么大的事物；因为基督徒的爱德除非在教会的合一中是无法持守的：从而你们可以看见，没有爱德你们什么也不是，即使你们拥有圣洗和信德，并且借着后者甚至能移山。但如果你们也持这种观点，那么让我们既不否认也不拒绝你们所知道的、属于天主的圣事，也不否认信德本身，而是让我们持守爱德，没有它，即使有圣事和信德，我们也什么都不是。我们持守爱德，如果我们紧抓合一；我们紧抓合一，如果我们不是以自己的言语在党派中制造虚假的合一，而是借着基督的话语在合一的整体中认出它。\par

% structure: p0253 source=h2 target=subsection reason=relative-heading-rank; base=h2

\subsection{第七十九章。}

173. \Person{培提利安}说：\Quote{再者，\BibleQuote{爱德是含忍的，慈祥的；爱德不嫉妒，不夸张，不自大，不作无礼的事，不求己益。}但你们却求他人的东西。\BibleQuote{不动怒，不图谋恶事；不以不义为乐，却与真理同乐；凡事包容，凡事忍耐。爱德永不止息。}\BibleRef{格前13:1}简而言之，意思就是说：爱德不迫害人，不煽动皇帝夺走他人的性命；不抢夺他人的财物；不会在掠夺之后继续杀人。}\par

174. \Person{奥古斯丁}回答说：我必须对你们重复同样的话多少次？如果你们不证实这些指控，它们不针对世上的任何人；如果你们证实了它们，它们也与我们无关；正如那些每天由疯子的狂暴行为、醉汉的奢侈、自杀者的盲目、强盗的暴虐所做的事情与你们无关一样。因为谁看不出我说的是真话呢？但如果你们有爱德，它会以真理为乐。因为披着羊皮的遮掩说得多么巧妙：\Quote{爱德凡事包容，凡事忍耐！}但当你们接受考验时，狼的牙齿就无法隐藏了。因为当你们顺从圣经的话，\BibleQuote{以爱德互相担待，竭力保持圣神内的合一，以和平的联系}，\BibleRef{弗4:2}爱德会迫使你们，即使你们知道教会内有什么恶事，我不说要同意它们，但至少如果你们不能阻止，就容忍它们，以免因那些在末日要被簸扬出去的恶人，你们现在因脱离善人的团体而割断和平的联系；你们抗拒她的影响，被轻浮的风吹出去，指责麦子是糠秕，并声称你们所虚构的恶事借着传染之力甚至在义人身上也成立。当主说过：\BibleQuote{田地是世界，收获是世界的终结，}尽管祂论及麦子和莠子时说：\BibleQuote{让两者一起生长直到收获，}你们却企图用你们的言语使人相信麦子已在田地的主要部分消失了，只残存在你们的小角落里——希望基督被证明是说谎者，而你们是真理之人。你们实在是违背自己的良心说话；因为任何稍加真实看福音的人都不敢从心里说，在那么多民族中——其中可以听到阿们的回应，并且几乎同声歌唱阿肋路亚——没有基督徒存在。然而，为了不让\Person{多纳徒}派显得有错误，他们不与世界上的各个民族共融，如果有一位能看见全世界的天使从天而降，宣告在你们的共融之外，任何地方都找不到善良无罪的人，那么毫无疑问，你们会因人类的罪恶而欢喜，并在未获得确证之前就夸口说了真话。那么，你们那里岂能有不以不义为乐的爱德呢？但不要受骗。在整个田地，即整个世界，直到世界的终结，都会找到主的麦子生长。基督说了这话：基督是真理。愿爱德在你们内，并以真理为乐。即使从天上来的天使传给你们与祂的福音不同的福音，也当受诅咒。\BibleQuote{即使从天上来的天使传给你们与祂的福音不同的福音，也当受诅咒。}\BibleRef{迦1:8}\par

% structure: p0256 source=h2 target=subsection reason=relative-heading-rank; base=h2

\subsection{第八十章。}

175. \Person{培提利安}说：\Quote{最后，迫害的理由是什么？我问你们，你们这些可怜的人，是否你们认为你们的罪建立在任何法律的权威之上？}\par

176. \Person{奥斯定}回答说：那犯罪的人，不是凭着法律的权威犯罪，而是违反法律的权威犯罪。但既然你问逼迫的正当理由是什么，我也反过来问你，在圣咏中是谁的声音说：\Quote{凡暗中诽谤邻舍的，我要将他铲除。}因此，你要追究逼迫的理由或尺度，而不要因笼统地指责那些逼迫不义者的人而显露出你的极端无知。\par

% structure: p0259 source=h2 target=subsection reason=relative-heading-rank; base=h2

\subsection{第八十一章}

177. \Person{佩提里安}说：\Quote{但我要反过来回答你，\Person{耶稣基督}从未逼迫任何人。当宗徒指责某些人，并暗示祂应诉诸逼迫时（祂自己来是为了藉着邀请人归向祂来建立信德，而非通过强迫），那些宗徒说：\InnerQuote{许多人奉祢的名按手，却不与我们同在。}但\Person{耶稣}说：\InnerQuote{由他们去吧；如果他们不反对你们，就是支持你们的。}}\par

178. \Person{奥斯定}回答说：你说你将以更丰富的内容从你的仓库里拿出圣经中没有记载的东西，这话说得不错。因为如果你想从圣经中拿出证据，你甚至能拿出圣经中找不到的东西吗？但凭你自己的意愿，你可以随意增加你的谎言。因为你所引用的那句话写在哪里？什么时候有人向我们的主提过那种建议，或者我们的主回答过那种话？\Quote{许多人奉祢的名按手，却不与我们同在。}这话从来没有一个门徒对天主之子说过；因此，祂也不可能回答：\Quote{由他们去吧；如果他们不反对你们，就是支持你们的。}但在福音中我们确实读到过一些类似的话——有人向主提起过一个人，那人奉祂的名驱魔，却不跟随祂和祂的门徒；主在那情况下说道：\BibleQuote{不要禁止！因为谁不反对你们，就是支持你们的。}\BibleRef{路 9:49}但这与指出那些主被认为宽恕了的人毫无关系。如果你被一种表面相似的说法所欺骗，这不是谎言，而只是人性的软弱。但如果你想让那些不熟悉圣经的人陷入谎言的迷雾中，那么愿你扎心，蒙羞并改正。然而，关于这个向主提起的人，我们要强调一点。因为正如当时在门徒的共融之外，基督的圣德仍然非常有效，同样，现在在教会的共融之外，圣事的圣德也是有效的。因为圣洗圣事除了因父、及圣子、及圣神之名以外，不能祝圣。但谁会疯狂到宣称，圣子的名字在教会的共融之外仍然有效，而圣父和圣神的名字却不可能？或者它在医治人时有效，而在祝圣圣洗时却无效？但显而易见，在教会的共融、至圣的合一纽带、以及最卓越的爱德恩赐之外，无论是驱魔的人，还是受洗的人，都不能获得永生；正如那些因圣事共融而看似在内、却因品德的败坏而被视为在外的人，也不能获得永生一样。但是，基督曾经用身体上的责罚来逼迫那些祂用鞭子赶出圣殿的人，我们已经在上面说过了。\par

% structure: p0262 source=h2 target=subsection reason=relative-heading-rank; base=h2

\subsection{第八十二章}

179. \Person{佩提里安}说：\Quote{但圣宗徒说了这话：\InnerQuote{无论如何，总当传扬基督。}}\par

180. 奥斯定回答说：你是在反对你自己；然而，既然你站在真理一边说话，如果你爱真理，就让你所说的话算作有利于你吧。因为我问你，宗徒保禄这话是对谁说的？我们若愿意，不妨稍微追溯一下前文。他说：\BibleQuote{有的传基督，是出于嫉妒和纷争；有的却是出于善意，出于爱德，知道我是为辩护福音而设立的。但那等传基督，是出于争执，并不诚心，想要给我的锁链加添苦楚。那又怎样呢？无论如何，或是假意，或是诚心，基督终究被传开了；我为此欢喜，而且还要欢喜。}\BibleRef{斐 1:15}我们看见，他们所传的本身是神圣、纯洁和真实的，但他们的方式却不纯洁，出于嫉妒和争执，没有爱德，没有纯洁。确实，不久之前，你似乎还在引用宗徒的见证，来向我们大加赞扬爱德，说没有爱德，任何事物都无济于事；然而你看，那些人没有爱德，但他们却宣讲基督，而宗徒在这里说，他为此欢喜。因为他欢喜的不是他们里面的恶，而是耶稣基督之名里面的善。在他身上，的确有那\BibleQuote{不以不义为乐，却与真理同乐}的爱德。\BibleRef{格前 13:6}此外，他们里面的嫉妒是来自魔鬼的邪恶，因为魔鬼正是藉此既杀害又推翻了。那么，这些被宗徒如此谴责的恶人，他们里面有那么多好事使他欢喜，他们是在内，还是在外？你来选择吧。如果他们在内，保禄认识他们，但他们并没有玷污他。同样，在普世教会的合一中，那些你指控的人——无论这些指控是真的，或是你自己捏造的谎言——也不会玷污你。既然如此，你为何要分离？为何要在分裂的亵渎之罪中毁灭自己？但如果他们在外面，那么你看，即使在外面的人——他们当然不能获得永生，因为他们没有爱德，也不住在合一中——仍然有基督之名的神圣，以致宗徒因这名的内在神圣而喜悦地确认他们的宣讲，尽管他拒绝他们本人。因此，当外面的人来到我们这里时，我们不损害这圣名本身，是正确的；但我们改正个人，同时尊敬这名。所以你要留心，看看你对待那些被你似乎谴责其行为的人是多么邪恶：你轻看在他们身上神圣的基督之名的圣事。而你，如你的话所示，认为宗徒所说的那些人是在教会界限之外。因此，当你害怕天主教徒的迫害时（你提到这迫害是为了让我们招人憎恨），你在异端者中确认了基督之名，而你在重行圣洗时却藐视这圣名。\par

% structure: p0265 source=h2 target=subsection reason=relative-heading-rank; base=h2

\subsection{第八十三章}

181. 培提利安说：\Quote{如果当真没有一些人，信仰的一切权能都与之对立，那你凭什么逼迫人，强迫他们玷污自己呢？}\par

182. 奥斯定回答说：我们并不逼迫你们，除非是真理逼迫谬误；即使有人用其他方式逼迫了你们，那也与我們无关，正如你方若有人行同样的事，也与你们无关；我们也不强迫你们玷污自己，而是劝你们获得痊愈。\par

% structure: p0268 source=h2 target=subsection reason=relative-heading-rank; base=h2

\subsection{第八十四章}

183. 培提利安说：\Quote{但如果某种法律曾授权强迫人行善，那么你们这些不幸的人本应被我们强迫接受最纯洁的信仰。但绝不可，绝不可让我们良心强迫任何人接受我们的信仰。}\par

184. \Person{奥斯定}答复说：的确，没有人应被强迫违背其意愿接受信德；但是，借着天主的严厉，或者更确切地说，借着天主的仁慈，背信弃义常受到患难鞭笞的惩罚。难道因为最好的道德是由自由意志选择的，所以最坏的道德就不受法律完整的惩罚吗？但是，除非当已学得的生活准则被藐视时，否则不可能有纪律来惩罚邪恶的生活方式。因此，如果制定了任何反对你们的法律，你们并不是因此被迫行善，而只是被阻止作恶。因为没有人能行善，除非他经过深思熟虑的选择，并且爱那属于自由意志的事物；但对惩罚的恐惧，即使不能分享良心的喜乐，至少阻止邪恶的欲望逃出思想的界限。然而，是谁制定了这些对你们不利的法律，以压制你们的狂妄呢？岂不正是使徒所说的那些人，他们\BibleQuote{不空空佩剑；因为他们是天主的仆役，是复仇者，执行忿怒以惩罚作恶的人？}\BibleRef{罗 13:4}因此，全部问题在于，你们是否在作恶，因为全世界都控告你们犯有如此重大分裂的亵渎之罪。然而，你们忽视这个问题的讨论，反倒谈论无关的事情；你们像强盗一样活着，却夸耀自己像殉道者一样死去。而且，由于惧怕法律本身，或可能招致的憎恶，或是因为你们无力抵抗，我不说那么多的人，而是说那么多的公教民族，你们甚至夸耀你们的温和，说你们不强迫任何人加入你们的党派。按照你们的说法，当鹰因为飞走而未能抓走家禽时，它可能称自己为鸽子。因为你们何时有过权力却不使用呢？因此你们表明，只要你们能够，你们会做得更多。当\Person{尤利安}嫉妒基督的和平，将属于统一的教会交还给你们时，谁能数清你们所犯下的所有屠杀？当时连魔鬼也与你们一同高兴，因为他们的庙宇重新开放了。在与\Person{菲尔穆斯}及其党派的战争中，让\Place{凯撒毛里塔尼亚}自己告诉我们，摩尔人\Person{罗加图斯}在你们手中遭受了什么。在\Person{吉尔多}时期，因为你们的一位同工是他的密友，让\Person{马克西米安}派的人作证他们所受的苦。因为如果有人可以呼唤现在与你们同在的\Person{费利西安}本人，让他发誓说\Person{奥普塔图斯}没有强迫他违背意愿回到你们的共融中，他不敢开口，尤其是如果\Place{穆斯蒂}的人民能看到他的脸，他们是所有事情的见证者。但是，如同我所说的，让他们作证他们在那些如此对待\Person{罗加图斯}的人手中所遭受的。公教会自身，尽管得到统治陆地与海洋的公教王侯的援助，却遭到\Person{奥普塔图斯}率领的武装敌对军队的野蛮攻击。正是这一点，首先使得有必要在代理官\Person{塞拉努斯}面前敦促执行一项对你们的法令，即处以十磅黄金的罚金，而你们直到今天从未支付过，然而你们却控告我们残忍。但是，你们在哪里能找到比用罚款来惩罚你们如此重大的罪行更温和的处理方式呢？或者谁能数算你们在你们所占据的地方，凭自己的主权和喜乐，各尽其所能，没有任何法官或其他掌权者的友谊，所行的一切事迹呢？在我们这边，城镇的居民中，有谁没有从前辈那里听说过这类事情，或者亲身经历过呢？难道不是这样吗？在我所在的\Place{希波}，有不少人记得，你们的领袖\Person{福斯提努斯}在他掌权的时候，由于当地天主教徒人数稀少，下令任何人都不得为他们烤面包，以至于一个面包师，我们一位执事的佃户，竟将地主的生面团扔掉，虽然他没有被任何法律判处流放，但他不仅在一个罗马国家中，而且甚至在自己的家乡，不仅在自己的家乡，而且在自己的家中，剥夺了他所有的生活必需品？为什么，甚至最近，正如我自己至今仍悲伤地回忆，你们中的一位，来自\Place{卡拉马}的\Person{克里斯皮努斯}，购买了一处地产，而且只是佃户，竟大胆毫不犹豫地将不少于八十个灵魂浸入第二次圣洗的水中，他们在恐惧的唯一影响下发出可怜的呻吟；而且这是在属于公教皇帝的农场上，根据他们的法律，你们甚至被禁止进入任何罗马城市？但是，除了你们这些行为之外，还有什么使得必须通过你们所抱怨的那些法律呢？的确，这些法律远远不足以与你们的罪行相称；但是，尽管如此，你们应该感谢自己使它们存在。的确，若不是我们在城镇中把你们当作人质，你们这些连人民的注视和所有可敬之人的谴责都不愿忍受的人（出于羞愧，若不是出于恐惧），我们岂不早就被你们那些在你们指挥下疯狂作战的\OriginalTerm{围栅者}{Circumcelliones}从各地赶出乡村了吗？因此，不要说：\Quote{远离吧，远离我们的良心，强迫任何人接受我们的信德。}因为你们能做的时候就做；当你们不做时，是因为你们做不到，要么出于对法律的恐惧，要么由于随之而来的憎恶，要么因为抵抗的人数众多。\par

% structure: p0271 source=h2 target=subsection reason=relative-heading-rank; base=h2

\subsection{第八十五章}

185. 培提里安说：\Quote{因为主基督说：\InnerQuote{除非那派遣我的圣父吸引他，没有人能到我这里来。}\BibleRef{若 6:44}但我们为什么不允许各人随从他的自由意志呢？因为主天主自己赐予人自由意志，却将正义之路指示给他们，免得有人因无知而丧亡。因为祂说：\InnerQuote{我在你面前放了好坏。我把火和水放在你面前；你选择哪一个。}从这选择中，你们这些可怜的人为自己选择了不是水，而是火。\InnerQuote{然而，}祂说，\InnerQuote{选择好，好使你存活。}\BibleRef{德 15:16}你们这些不愿选择好的人，凭自己的判决声明你们不愿存活。}\par

186. 奥古斯定答覆：若我向你提出一个问题：圣父如何吸引人归向圣子，同时又让他们留在自由行动中？你或许会难以回答。因为若祂任凭他们自行选择所喜欢的，祂又如何吸引他们呢？然而这两点都是真实的；只是这真理少有智慧能洞悉。因此，正如圣父在赐予人自由意志后，仍能吸引人归向圣子；同样，那些通过法律惩罚而来的警告，也并非剥夺自由意志。因为当人遭受任何严厉和不悦之事时，他便受警告去思考为何受苦：若他发现是为正义而受苦，他可能选择那存在于为正义受苦行为本身中的善；但若他看到自己所受的痛苦是不义的，他便可能因考虑到自己徒然受苦受折磨，而改变心意趋向更好，同时摆脱那徒然的烦恼和不义本身——这邪恶在作恶上更为有害和危险。所以，当君王制定反对你们的法令时，你们应当视之为警告，思考为何如此对待你们。因为若为正义之故，那么他们真是你们的迫害者；而你们是有福的，因义受迫害者将承受天国。但若因你们分裂的邪恶，那么他们不过是你们的纠正者；而你们，如同其他犯各种罪行并受法律惩罚的人一样，无疑在此世和来世都是不幸的。因此，无人剥夺你们的自由意志。但我恳劝你们认真思考：你们宁愿选择什么？——是在和平中受纠正而存活，还是坚持恶意，以殉道的虚名承受真正的惩罚？但我对你们说话，好像你们正遭受与罪行相称的苦楚，而你们实则犯下如此大罪却逍遥法外。你们如此疯狂，以至于\Quote{\Emph{赞美天主}}的呼喊比战争号角更令人恐惧；你们如此充满诽谤，甚至当你们无故跳下悬崖时，也归咎于我们的迫害。\par

187. 他又像最慈爱的导师般说：\Quote{你们这些不愿选择好的人，凭自己的判决声明你们不愿存活。}据此，若我们相信你们的指控，我们便活在仁慈中；但因我们相信天主的许诺，我们凭自己的判决声明我们不愿存活。在我看来，你们似乎记得宗徒们被禁止宣讲基督时如何回答犹太人。因此我们也要说：你们应回答我们，我们该听天主的，还是听人的？\BibleRef{宗 5:29}\Emph{\OriginalTerm{叛教者}{Traditors}}、焚香者、迫害者：这些都是人反对人的言语。基督只存留在\Person{多纳图}的爱中：这是人的言语，以基督之名高抬人，而使基督本身的荣耀受损。因为经上记载：\Quote{民众众多，是君王的光荣；若民众稀少，是王侯的衰败。}\BibleRef{箴 14:28}所以，这些都是人的言语。但福音中的话：\Quote{基督必须受苦，第三天从死人中复活，并且因祂的名从耶路撒冷开始，向万邦宣讲悔改和罪过的赦免。}\BibleRef{路 24:46}这些是基督的话，显示祂从圣父所得的荣耀，显现在祂国度的广袤中。当我们两者都听见时，我们宁愿选择教会的共融，将基督的话置于人的话之上。我问：谁能否认我们选择了善？除非有人要说基督教导了恶。\par

% structure: p0275 source=h2 target=subsection reason=relative-heading-rank; base=h2

\subsection{第八十六章。}

188. 培提里安说：\Quote{那么天主甚至命令对分裂者进行屠杀？若祂要发出这样的命令，你们理应被某些蛮族和斯基泰人杀死，而非被基督徒杀死。}\par

189. 奥古斯定答覆：请你们的\OriginalTerm{乱党}{Circumcelliones}保持安静，我恳求你们不要用蛮族来恐吓我们。至于我们或你们谁是分裂者，既不要问你们，也不要问我，而要问基督，让祂指示何处能找到祂的教会。那么，读福音，你在那里找到答案：\Quote{在耶路撒冷、全犹太、撒玛黎雅，直到地极。}\BibleRef{宗 1:8}因此，若有人不在教会内，不必再对他提问，要么纠正或归化他，要么揭露他，他就不该抱怨。\par

% structure: p0278 source=h2 target=subsection reason=relative-heading-rank; base=h2

\subsection{第八十七章。}

190. 培提里安说：\Quote{因为主天主从未欢喜于人的血，甚至祂愿意杀害兄弟的加音继续活在他的杀人生活中。}\par

191. 奥古斯丁回答说：如果天主不愿意将死亡加于杀害兄弟者身上，而宁愿让他在凶手的生命中继续存在，那么请看，这是否就是原因，因为君王的心在天主手中，祂亲自制定了许多法律来纠正和责备你们，然而没有一条君王的法律命令将你们处死，也许正是出于这个目的：凡在亵渎疯狂中执拗固执己见的人，要承受杀害兄弟者\Person{加音}的惩罚，即凶手的生命。因为我们读到，上主的仆人\Person{梅瑟}曾慈悲地杀了许多人；因为他这样为他们的邪恶亵渎向主代祷说：\InnerQuote{上主，如果你愿意赦免他们的罪……；若不然，求你将我从你所写的书中抹去吧。} \BibleRef{出 32:28}他的不可言喻的爱德和慈悲清楚显明。那么，难道他忽然转为残忍，当\Person{梅瑟}从山上下来时，命令杀死了数千人？因此，请考虑，这是否是来自天主的更愤怒的标记：虽然制定了那么多针对你们的法律，却没有一位皇帝下令处死你们。或者你们认为自己不能与那杀害兄弟者相比？请听主藉先知说的话：\BibleQuote{从日出之地到日落之地，我的名在外邦人中要大显扬；在每一处，要向我的名奉献香品和洁净的祭品，因为我的名在外邦人中要大显扬，万军的上主说。} \BibleRef{拉 1:11}对于这位兄弟的祭献，你们显露出恶意的眼光，超越了天主对它的尊重；如果你们曾听说过\Quote{从日出之地到日落之地，上主的名是应当赞美的}，这就是那生活的祭献，经上说：\Quote{你们应以感谢祭献给天主}，那么你们的面容就会像那凶手的面容一样沉落。但因你们不能杀死全世界，你们单凭仇恨就陷入同样的罪中，依照\Person{若望}的话：\BibleQuote{凡恨自己兄弟的，就是凶手。} \BibleRef{若一 3:15}我却希望，任何一个无辜的兄弟宁可落入你们的\OriginalTerm{寻器派}{Circumcelliones}手中，被他们的武器杀死，也不愿遭受你们舌头的毒害和重洗。\par

% structure: p0281 source=h2 target=subsection reason=relative-heading-rank; base=h2

\subsection{第八十八章}

192. \Person{Petilianus}说：\Quote{我们因此劝告你们，如果你们愿意听从的话，即使你们不愿意接受，我们仍警告你们：主基督为基督徒制定的，并不是任何形式的杀害，而只是死亡的形式。因为如果祂喜爱这样喜爱战斗的人，祂就不会为我们而被杀。}\par

193. 奥古斯丁回答说：但愿你们的殉道者会遵行祂所规定的形式！他们就不会跳下悬崖，这是祂拒绝魔鬼命令时所做的事。\BibleRef{玛 4:6}但当你们用假见证迫害我们现在已死去的祖先时，你们从哪里得到这形式？你们试图用我们不认识的人的罪行来玷污我们，同时却不愿意你们自己一方的恶行被算在你们头上，你们从哪里得到这形式？但如果我们过于自以为是，当我们看到你们亲自对主说假见证时，我们埋怨自己，因为祂亲自应许并彰显祂的教会将扩展到万民中，而你们却反对。因此，你们甚至没有从犹太迫害者那里得到这形式，因为他们迫害祂在地上的身体；你们却迫害祂在天上坐席的福音。这福音更能忍耐愤怒君王的火焰，却无法忍耐你们的舌头；因为当火焰燃烧时，合一仍然存在，而这对你们的话语是不可能做到的。那些希望天主的话在火焰中消灭的人，不相信它被阅读时会受轻视。因此，如果你们让他们用舌头反对福音，他们就不会用火焰烧福音。在早期的迫害中，基督的福音被一些人愤怒地寻找，被另一些人恐惧地出卖；被一些人愤怒地焚烧，被另一些人热爱地隐藏；它受到攻击，但没有找到人反对它的真理。更受诅咒的迫害部分为你们保留，当对异教徒的迫害耗尽时。那些迫害基督名字的人相信了基督；而现在那些因基督名字受尊敬的人却被发现反对祂的真理。\par

% structure: p0284 source=h2 target=subsection reason=relative-heading-rank; base=h2

\subsection{第八十九章}

194. \Person{Petilianus}说：\Quote{这里你们有最充分的证据，证明基督徒不可参与毁灭他人。但这原则的首创是在\Person{伯多禄}的事例中，正如所记载的：\InnerQuote{西满伯多禄有一把剑，就拔出来，砍了大司祭的仆人，削掉了他的右耳。耶稣遂对\Person{伯多禄}说：把你的剑收入鞘内，因为凡持剑的，必死于剑下。}}\par

195. 奥古斯丁答：那你们为什么不用这样的话来制止\Quote{\OriginalTerm{圆场党}{Circumcelliones}}的武器呢？倘若你们说：\Quote{凡持棍棒的，必死于棍棒之下}，你们是否认为这就超出了福音的话？如果我们的祖先未能制止你们所抱怨将\Person{马尔库卢斯}推下悬崖的人，那你们就不要吝惜宽恕；因为福音上也并没有写着：\Quote{凡惯于把人推下悬崖的，也要从那里被推下去}。但愿你们的控诉若非虚假便是陈旧，你们那些朋友的棍棒也能停止！尽管如此，你们或许还会因我们在言语上——即便不是凭借法律的力量——剥夺了你们军团的武装而感到不悦，因为我们说他们只以棍棒发狂。因为那是他们邪恶的古老方式，但现在他们已经走得太远了。因为在他们酗酒狂欢、随意集会、游荡街头、嬉笑饮酒、整夜与无夫的女子共处之中，他们不仅学会了挥舞棍棒，还学会了挥剑和投石。可是，我为什么不向他们说（天主知道我以何种心情说，他们以何种心情听！）：\Quote{疯子们，\Person{伯多禄}的剑，虽然是在尚未摆脱血肉之污的动机下拔出的，但终究是为了保卫基督的身体以对抗迫害者的身体；而你们的武器却是用来反对基督的事业。但以他为元首的身体，即他的教会，却遍布万民。他亲自说了这话，并升了天，犹太人的暴怒无法跟随他，而你们的暴怒却攻击他升天后托付给我们照顾的肢体。为了保护这些肢体，但凡在天主教会内、信德尚小、与\Person{伯多禄}当年为主之名拔剑时动机相同的人，都向你们发怒，抵抗你们。但你们的迫害与他们的抵抗之间有很大的区别。你们如同大司祭的仆役；为服事你们的首领，你们武装起来反对天主教会，即反对基督的身体。而他们却如当年的\Person{伯多禄}，甚至用肉体的力量为基督的身体——教会——争战。但如果他们被命令安静下来，如同当年的\Person{伯多禄}一样，那么你们更应受到警告，放下异端的疯狂，加入那些肢体——他们正是为此争战——的合一。然而，被这样的人伤害后，你们也恨我们；仿佛失去了右耳，你们听不见坐在圣父右边的基督的声音。但我向谁说话，或如何向他们说话呢？因为在他们身上我找不到说话的时机。即便清早，他们已酒气熏天，或已醉酒终日，或仍宿醉未醒。不仅如此，他们还发出威胁，而且不仅是他们，连他们自己的主教也针对他们发出威胁，随时准备否认他们的所作所为与自己有关。愿主赐给我们一首登阶歌，使我们能说：\Quote{我与憎恨和平的人在一起时，我是和平的。我正要与他们说话，他们却无端地与我作战}。因为基督的身体，即遭到异端分子在世界各地攻击（这里一些人，那里另一些人，一切异端者，无论他们在哪里）的身体，如此说道。}\par

% structure: p0287 source=h2 target=subsection reason=relative-heading-rank; base=h2

\subsection{第九十章}

196. \Person{培提里安}说：\Quote{因此我说，他规定我们应当为信仰而受死，每人当为教会的共融而这样做。因为基督信仰藉其殉道者的死亡而前进。倘若信众惧怕死亡，就没有人能以完美的信德生活。因为主基督说：\BibleQuote{一粒麦子若不落在地里死了，仍旧是一粒；若是死了，就结出许多子粒来。}\BibleRef{若 12:24}}\par

197. 奥古斯丁答：我倒想知道，你们中间是谁第一个跳下悬崖的。因为那粒麦子真是多产，由此生出了如此众多类似的自杀者。告诉我，当你们引用主的话说一粒麦子应当死了结出许多子粒时，你们为何嫉妒那真正的果实——它确实已遍及全世界——并拿你们所听说或编造的一切莠子或糠秕的罪名来攻击它呢？\par

% structure: p0290 source=h2 target=subsection reason=relative-heading-rank; base=h2

\subsection{第九十一章}

198. \Person{培提里安}说：\Quote{但你们撒的是荆棘和莠子，而不是麦种，所以你们应与它们一同在最后的审判中被焚烧。我们并不诅咒；但每一荆棘般的心，都因天主所宣布的判决而受此惩罚。}\par

199. 奥古斯丁答：的确，当你们提到莠子时，你们也该想到麦子；因为两者都被命令在田里一同生长直到收割。但你们凶残地只盯着莠子，并违背基督的明确宣告，坚持说唯有莠子长满了全地，只有\Place{非洲}除外。\par

% structure: p0293 source=h2 target=subsection reason=relative-heading-rank; base=h2

\subsection{第九十二章}

200. \Person{培提里安}说：\Quote{主基督的话在哪里：\BibleQuote{凡打你右脸的，连左脸也转过来由他打}？\BibleRef{玛 5:39} 当他受人唾面时，他所显示的忍耐在哪里？他正是用至圣的唾液开了盲人的眼睛。宗徒\Person{保禄}的话在哪里：\InnerQuote{若有人打你的脸}？同一宗徒的另一句话在哪里：\InnerQuote{论鞭打是过量的，在监牢里是更频繁的，在死亡中是屡次的}？他所提及的是他所受的苦，而非他所行的事。对于基督信仰来说，这些事由犹太人来做已经足够：可悲的人哪，你们为什么又加添这些事呢？}\par

201. \Person{奥斯定}回答说：难道真是这样，当人打你一边脸颊时，你就把另一边也转过来给他打吗？你们狂暴的团伙以可怕的邪恶到处游荡，横行整个非洲，并未为你们带来这样的名声。我倒情愿人们与你约定，要你们依照旧法律，只求\Quote{以眼还眼，以牙还牙}\BibleRef{申 19:21}，而不是在你们听到不顺耳的话时便举起棍棒。\par

% structure: p0296 source=h2 target=subsection reason=relative-heading-rank; base=h2

\subsection{第九十三章}

202. \Person{佩蒂利安}说：\Quote{但你们与这世界的君王有何相干？基督教在她身上除了嫉妒从未找出任何东西。为简短地告诉你们我所言的真实：一位君王迫害了玛加伯弟兄们。\EditorialNote{达尼尔第三章}一位君王也将那三位青年定罪于圣化之火，他不知道自己在做什么，因为他自己与天主对抗。\EditorialNote{达尼尔第六章}一位君王图谋婴儿救主的性命。\BibleRef{玛 2:16}一位君王将达尼尔投入，如他所想，使野兽吞噬。\EditorialNote{达尼尔第六章}主基督自己也被一位君王最邪恶的审判官处死。\BibleRef{玛 27:26}因此，宗徒呼喊：\InnerQuote{我们在成全的人中讲智慧，但不是这世代的智慧，也不是这世代将要消灭的执政者的智慧；我们讲的，乃是天主奥秘的智慧，这智慧是隐藏的，天主在万世之前预定使我们得荣耀的；这世代执政者中没有一个知道这智慧，因为他们若知道，就不至于把荣耀之主钉在十字架上了。\BibleRef{格前 2:6}}但假设这是关于以前的异教君王说的。然而你们，这一时代的统治者，因为你们渴望成为基督徒，却不允许人们成为基督徒，反而当人诚心相信时，你们用虚假的污秽和迷雾将他们全然拖入你们的邪恶，使他们用为抵御国家敌人而准备的武器来攻击基督徒，并认为，在你们的煽动下，他们是在做基督的工作，如果他们杀死你们所仇恨的我们，正像主基督所说：\InnerQuote{时候将到，凡杀你们的，就以为是事奉天主。\BibleRef{若 16:2}}因此，虚假的教师们，无论这世界的君王是愿意成为异教徒（愿天主禁止），还是基督徒，对你们无关紧要，只要你们不停止武装他们反对基督的家族。但你们难道不知道，或者难道没有读过，煽动谋杀者的罪大于实行者的罪？依则贝尔煽动她的丈夫君王杀害一位穷困的义人，然而丈夫和妻子都同样灭亡于同等的惩罚。\EditorialNote{列王纪上第二十一章}而且你们煽动君王的方式，与女人狡猾的说服常常煽动君王犯罪的方式并无不同。因为黑落德的妻子通过她的女儿求得并获得了恩惠，使若翰的头被盛在盘子里端到桌上。\BibleRef{玛 14:8}同样，犹太人迫使本丢·比拉多钉死主耶稣，比拉多祈求他的血归在他们和他们的子孙身上。\BibleRef{玛 27:24}所以，你们也同样因你们的罪以我们的血自加沉重。因为，尽管是审判官的手击打，这并不表明你们的诽谤不是更有罪过。因为先知达味以基督的身份说话：\InnerQuote{外邦人为什么嚣张？万民为什么谋算虚妄的事？世上的君王一齐起来，掌权者一同商议，要敌抗上主，敌抗祂的受傅者，说：我们要挣开他们的捆绑，脱去他们的绳索。那坐在天上的发笑，上主必讥笑他们。那时祂要在怒中对他们说话，在烈怒中惊吓他们。说：我已立我的君王在熙雍我的圣山上。我要传述上主的旨意：祂对我说：你是我的儿子，我今日生了你。你向我求，我就将外邦人赐你为基业，将地极赐你为产业。你必用铁杖管辖他们，将他们如同陶器摔碎。}他也在以下训诫中警告君王们，使他们不要像无知无识的人一样，寻求逼迫基督徒，免得他们自己毁灭——我愿我们能教训他们，因为他们对此无知；或者至少你们会指给他们看，如果你们希望他们活着，你们无疑会这样做；或者，如果两种做法都不允许，你们的恶意会允许他们自己阅读。达味的第一篇诗篇无疑会说服他们，他们应当作基督徒生活并统治；但与此同时，你们欺骗他们，因为他们信任你们。因为你们向他们代表邪恶之事，向他们隐藏良善之事。那么，让他们终于读到这些，他们早就应该读了。因为他说什么？\InnerQuote{所以，你们君王应当明智；你们地上的掌权者，应受教训！以敬畏事奉上主，以战兢欢喜快乐。要接受管教，免得上主发怒，你们从正道上灭亡；因为祂的怒气快要燃起。凡投靠祂的都是有福的。}我告诉你们，你们用你们的劝说煽动皇帝，正如我们上面所表明的，犹太人煽动比拉多一样，虽然他在众人面前洗手，大声说：\InnerQuote{流这义人的血，罪不在我，\BibleRef{玛 27:24}}——好像一个人能因自己犯的罪而免于罪责一样。但是，且不说古代的例子，从你们自己人的例子中观察，有多少你们的皇帝和审判官因逼迫我们而灭亡。尼禄是第一个逼迫基督徒的，多米提安几乎以与尼禄同样的方式灭亡，还有图拉真、盖塔、德西乌斯、瓦勒良、戴克里先；马克西米安也灭亡了，在他的命令下，使人向他们的神烧香，焚烧神圣的书卷，玛策利诺首先，但之后迦太基的门苏里乌斯和塞西利亚努斯，从亵渎的火焰中逃脱了死亡，像一些余烬或灰烬一样存活。因为烧香的罪责涉及你们所有同意门苏里乌斯的人。玛卡里乌斯灭亡了，乌尔萨西乌斯灭亡了，你们所有的伯爵都同样因天主的报应灭亡。因为乌尔萨西乌斯在与野蛮人的战斗中被杀，之后飞鸟用它们的利爪，狗用它们贪咬的牙齿将他肢解。他不也是因你们的建议而成为凶手吗？他像阿哈布王（我们已表明他被一个女人劝说）一样，杀害了一位穷困的义人？\EditorialNote{列王纪上第二十一章}所以你们也不停止杀害我们这些正义而贫穷的人（贫穷，指世上的财富；因为在天主的恩宠中，我们中没有一人是贫穷的）。因为即使你们不亲手杀人，你们也用你们屠夫般的舌头不停地杀人。因为经上记着：\InnerQuote{死亡和生命在舌头的权下。\BibleRef{箴 18:21}}因此，所有被杀害的人，你们这些煽动者都杀死了。而且，屠夫的手除了因你们舌头的煽动而发热外，不会发烫；那胸中可怕的热情被你们的话点燃，去取别人的血——那血将公正地报复在流血者身上。}\par

203. \Person{奥古斯丁}回答说：如果我要充分且尽己所能地回答这段被你如此夸大和长篇大论的文字，你在其中以恶意的言辞论及我们与此世君王的关系，我甚为担心你也会指责我存心激起君王对你的愤怒。然而，当你以自己的方式被这场反对所有天主教徒的猛烈抨击所裹挟时，你确实没有放过我。然而，我将尽力，若能的话，表明正是你因你所说的话而犯了此过，而非我将要作答的我。首先，看看你是如何自相矛盾；因为你确实在你所引用的段落前面加了这样的话：\Quote{你们与此世的君王有何相干？基督信仰从未在他们身上找到任何东西，唯有忌恨而已。}这些话确实断绝了我们与此世君王的一切往来。稍后你又说：\Quote{他也以下列训诫警告君王本身，他们不应像无知无识的人那样，试图迫害基督徒，免得他们自己反被消灭——我愿我们能将这些训诫教导他们，因为他们对此一无所知；或者，至少你们能向他们指明，如果你们希望他们活着，你们无疑会这样做。}那么，你希望我们以何种方式成为君王的导师？事实上，我们团体中凡与基督徒君王有友谊的人，若善用此友谊，并不犯罪；但若有人因此自傲，他们所犯的罪也远比你轻。因为你们这样指责我们——你们与一位外教君王有何相干？更糟的是，与背教者、基督之名的敌人\Person{尤利安}有何相干？当你们乞求将大殿归还给你们，仿佛它们是你们自己的时，你们将这样的称赞归于他：\Quote{在他身上，唯独正义找到了栖身之所。}——这些话（我相信你们懂拉丁文）将\Person{尤利安}的偶像崇拜和背教都称为正义。我手中持有你们祖先所呈递的请愿书；体现他们请求的备忘录；以及他们陈述的编年史。请观看并留意。对基督的敌人、背教者、基督徒的对头、魔鬼的仆役，你们的那位朋友、代表、\Person{蓬提乌斯}如此哀求：\Quote{你们去，对我们说：你们与此世的君王有何相干？}这样，你们就像聋子一样向聋哑的万民宣读你们和他们都不愿听的话；\BibleQuote{你看见兄弟眼中的木屑，却不理会自己眼中的大梁。}\BibleRef{玛 7:3}\par

204. 你说：\Quote{你们与这世界的君王有何相干？在他们身上，基督信仰除了嫉妒外，从未找到任何东西。}说了这话，你便尽力列举义人曾发现哪些君王是他们的敌人，却不考虑有更多可以列举出来证明是他们的朋友。圣祖亚巴郎曾受到一位君王最友善的对待，并获赠友谊的凭证，这位君王曾被上天警告不可玷污他的妻子。\EditorialNote{创世纪第20章}他的儿子依撒格同样遇到一位对他极为友善的君王。\BibleRef{创 26:11}雅各伯在埃及受到一位君王的盛情接待，甚至祝福了他。\EditorialNote{创世纪第47章}我还能说什么关于他的儿子若瑟呢？他在监狱的磨难中，其贞洁如金子被火锻炼，被法老提拔至高位，甚至以法老的性命起誓，\BibleRef{创 42:15}——并非因虚荣自负，而是不忘法老对他的恩惠。一位君王的女儿收养了梅瑟。\BibleRef{出 2:10}达味因以色列君王的不义所迫，投奔了异族的君王。\EditorialNote{撒慕尔纪上第27章}厄里亚在最邪恶的君王的车前奔跑——并非王命所迫，而是出于自己的忠诚。\BibleRef{列上 18:44}厄里叟认为最好主动向那曾接待他的妇人提出，她可能希望经由他代求从君王那里得到任何东西。\BibleRef{列下 4:13}但我现在要谈到天主之民被掳掠的真实时期，在那时，说得温和些，你产生了一种奇怪的健忘。因为你想要证明基督信仰在君王身上除了嫉妒外从未找到任何东西，你便提到那三位青年和达尼尔，他们在迫害的君王手下受苦，而你却不能从那些并非隔代而是完全相同的段落中获得教训，即：在那不受伤害的火焰奇迹之后，君王本人的表现，无论是在赞美和显扬天主的圣名上，还是在他本人尊重那三位青年上，或是从君王对达尼尔的器重以及授予他的礼物（他毫不勉强地接受）上，当他给予君王权力应得的尊重时，正如从他的话语中充分显示的，他毫不迟疑地运用天主赐予他的恩赐来解释君王的梦境。因此，当嫉妒这位圣先知的那些人以亵渎的疯狂诽谤他，迫使君王极不情愿地将他投入狮子坑中时，君王虽然悲伤地那样做了，但他仍坚信藉着天主的帮助和保护，他会安然无恙。于是，当达尼尔因奇迹般压制了狮子之怒而毫发无损时，君王首先向他发出友好的声音，语气焦虑，达尼尔从坑中回应祝福：\Quote{大王，万岁！}\EditorialNote{达尼尔书第3-6章}怎么会这样：当你的论证正围绕同一主题，当你自己引用天主仆人们的例子，在他们的身上这些事发生了，你或者没有看见，或者不愿意看见，或者看见并知道却保持沉默，以一种我不知道你将如何辩护的方式，对君王对圣人们所怀有的友谊的那些事例？但若不是你，作为最卑劣事业辩护者，被建造谎言的欲望所阻碍，因而要么不愿意要么无知地偏离真理之光，那么毫无疑问，你可以毫不困难地记起一些好君王和一些坏君王，一些对圣人们友善的和一些不友善的。我们不禁感到惊奇，你们的克氏团就是这样将自己从悬崖抛下。请问，谁曾追赶你们？哪个玛加略，哪个士兵曾追逐你们？当然，我们中没人将你们推入这谎言的深渊。那么，为何你这样闭着眼睛猛冲，以致当你说\Quote{你们与这世界的君王有何相干？}时，你没有加上\Quote{在他们身上，基督信仰常常找到嫉妒}，却胆敢说\Quote{在他们身上，基督信仰除了嫉妒外从未找到任何东西}？难道你真的既没有自己思考，也没有考虑那些读你著作的人会思考，有多少君王的实例与你的观点相悖？他难道不知道自己在说什么？\par

205. 或者你以为，因我所提到的人属于古代，所以他们对你不构成任何驳斥？因为你并没有说：\Quote{义德在她身上从未找到任何东西，除了嫉妒}，而是说：\Quote{基督信仰在她身上从未找到任何东西，除了嫉妒}，——这意味着，或许应当理解为，从他们开始拥有基督徒之名时，就对义人产生了嫉妒？那么，你甚至更加轻率地想要用这些古代的例子来证明你曾如此轻率地断言的东西，又有什么意义呢？因为\Person{玛加伯}、\Person{三青年}和\Person{达尼尔}所行所受的，岂不都是在基督降生于世之前吗？再者，正如我刚才所问，为什么你们向基督信仰的确定无疑的敌人\Person{尤利安}呈递请愿书？为什么你们要从他那里收回大殿？为什么你们宣称只有义德在他那里获得立足之地？如果基督信仰的敌人听到这样的话，那么从他那里听到这话的又是什么人呢？但应当注意，\Person{君士坦丁}肯定不是基督徒之名的敌人，反而因这名而显耀，他谨记自己在基督内所持守的望德，并为基督的合一作出极公正的裁决，却未被你们承认，甚至当你们自己向他上诉时也是如此。这两位都是基督徒时期的皇帝，但并非两人都是基督徒。但如果两人都是基督信仰的敌人，你们为什么向其中一位上诉？又为什么向另一位呈递请愿书？因为当你们的祖先呈递请愿书时，\Person{君士坦丁}已在\Place{罗马}和\Place{阿尔勒}作出了主教裁决；然而，你们将其中第一人在他面前控告，又从另一人向他上诉。但如果事实如此，其中一人信了基督，另一人背弃了基督，为什么推进合一利益的基督徒被鄙视，而支持诡诈的背教者却被赞扬？\Person{君士坦丁}下令从你们手中收回大殿，\Person{尤利安}则下令归还。你们想知道哪个行动有助于基督徒的和平吗？一个是由信基督的人所作的，另一个是由离弃基督的人所作的。哦，你们多么希望你们能说：\Quote{向\Person{尤利安}如此呈递请愿书，确实做得不妥，但这与我们何干？}但如果你们这么说，天主教会同样会以这些话得胜，其圣人遍布天下，他们远不像你们那样在意你们对那些你们可能任意对待之人的评说。但你们无力说出\Quote{向\Person{尤利安}如此呈递请愿书，确实做得不妥}。你们的喉咙被封闭；你们的舌头被身边的权威所阻止。那是\Person{庞提乌斯}所为。\Person{庞提乌斯}呈递了请愿书；\Person{庞提乌斯}宣称那背教者最是正义；\Person{庞提乌斯}阐明只有义德在那背教者那里获得立足之地。\Person{庞提乌斯}以这些话向他请愿，我们有\Person{尤利安}本人明确的证据，他毫无掩饰地指名道姓。你们的陈情仍然存在。这不是不确定的传闻，而是证明事实的公开文件。难道因为那背教者对你们的祈求作出一些让步，损害了基督的合一，你们就认为所说的话——只有义德在他那里获得立足之地——是真实的？但因为基督徒皇帝违背你们的意愿作出裁决，因这在他们看来最有助于基督的合一，所以他们就被称为基督信仰的敌人？愿所有异端者都显出这样的愚昧；并愿他们重获智慧，好让他们不再成为异端者。\par

206. 你会说，主所宣告的什么时候应验呢？他说：\BibleQuote{时候将到，凡杀你们的，就以为是事奉天主}\BibleRef{若 16:2}。但无论如何，这话不能用于外邦人，他们迫害基督徒不是为了天主的缘故，而是为了他们的偶像。你们没有看到，如果这话是对那些以基督徒之名为喜乐的皇帝说的，那么他们的首要命令必定是你们应被处死；而他们从未下过这样的命令。但是你们同党的人，以敌对的态度对抗法律，将应得的惩罚招致自身；他们自己的自愿死亡，只要他们以为能给我们带来憎恶，他们就绝不认为对自己有害。但如果他们以为基督的那句话是指那些尊崇基督之名的君王，就让他们问问，当亚略派的\Person{瓦伦斯}作皇帝时，天主教会在东方所受的苦难。在那里，我确实可以找到我认为足以应验主的话\BibleQuote{时候将到，凡杀你们的，就以为是事奉天主}\BibleRef{若 16:2}的情况，免得异端者将天主教皇帝针对他们错误所颁布的禁令当作自己的特殊荣耀。但我们记得，那句话是在我主升天后应验的，圣经是众所周知的见证。犹太人以为当他们处死宗徒时，是在事奉天主。在那些以为自己在事奉天主的人当中，甚至还有我们的\Person{扫禄}——虽然那时他还不是我们的。所以在他那些已过去且应被遗忘的自信理由中，他列举如下：\BibleQuote{他是希伯来人中的希伯来人；就法律说，是法利塞人；就热忱说，是迫害教会的}\BibleRef{斐 3:5}。这里有一个人，当他迫害教会时，以为自己在事奉天主，而后来他自己也遭受了同样的迫害。因为有四十个犹太人起誓要杀他\BibleRef{宗 23:12}，他使这事被千夫长知道了，于是在武装士兵的保护下逃脱了他们的埋伏。但那时还没有人对他说：你和千夫长及君王的武器有什么关系（不是和君王，而是和千夫长及君王的武器）？也没有人对他说：你竟敢从士兵手中寻求保护，而你的主正是被他们拖去受难的吗？那时还没有像你们这样的疯狂先例；但已经准备好了足以反驳他们的榜样。\par

207. 此外，你竟以何等可怕的力量胆敢说出并宣告以下的话：\Quote{至于古代的榜样，暂且不提；从你们自己一方的事例可以看到，有多少你们的皇帝和法官因迫害我们而丧亡。}当我读到你的信中的这段话时，极其迫切地等待着看你将要说什么，列举谁；可是，看哪！好像略过它们似的，你开始向我引述尼禄、图密善、图拉真、盖塔、德西乌斯、瓦勒良、戴克里先、马克西米安。我承认还有更多；但你完全忘记了你在和谁辩论。这些人岂不都是外教人，为了他们的偶像而普遍迫害基督徒的名号吗？那么，你们要警惕；因为你所提到的人并非与我们共融的。他们迫害的是合一的整体本身，而按照你的想法，或按照基督的教导，是你或我们从这个整体中分离出去。但你原本打算表明，我们的皇帝和法官因迫害你们而丧亡。或者，你自己并不要求我们计算这些人，因为你在提到他们时略过了，说\Quote{暂且不提尼禄}；你带着这个保留，打算涵盖所有其余的人吗？那么，引用他们又有什么用，如果它们与事情无关？但这与我何干？我现在和你一起略过它们。接下来，请把你答应我们的更大的数目列举出来，除非也许找不到他们，因为你曾说他们已经丧亡。\par

208. 因为现在你继续提到那些主教，你惯常控告他们交出圣书；关于他们，我们这方惯常回答说：要么你的证据不足，那么这与任何人无关；要么你成功了，那么与我们仍然无关。因为他们已经担负了自己的重担，无论好坏；而我们确实相信那是好的。但无论性质如何，那总是他们自己的；正如你们中的恶人担负了自己的重担，你们不担他们的，他们也不担你们的。但你们所有人的共同且最恶的重担就是分裂。这一点我们已经多次说过。因此，请向我们展示，不是主教的名字，而是我们的皇帝和法官的名字，他们因迫害你们而丧亡。因为这才是你提议的，这才是你承诺的，这才是你使我们极其迫切地期待的。\Quote{听啊，}他说，\Quote{玛卡里丧亡了，乌尔撒奇丧亡了，你们所有的伯爵也照样因天主的惩罚而丧亡。}你只提到两个名字，而他们两个都不是皇帝。我问，谁会对此满意呢？难道你自己不深深不满吗？你承诺将提到我们一方大量因迫害你们而丧亡的皇帝和法官；然后，对皇帝只字不提，你提到了两个可能是法官或伯爵的人。至于你补充的\Quote{你们所有的伯爵也照样因天主的惩罚而丧亡}，这与事情无关。因为照此原则，你早就该结束你的论证，根本不需要提任何人的名字。那么你为什么没有提到我们的皇帝，即与我们共融的皇帝呢？难道你害怕被控告犯叛国罪吗？标志圆环派的勇气在哪里？此外，你如此多地引入那些你上面提到的人，用意何在？他们更有权对你说：你为什么把我们找出来？因为他们对你的理由毫无帮助，而你却一一提名。那么，你是什么样的人，竟害怕提名那些据你说已经丧亡的人？至少，你可以多提到一些法官或伯爵，你似乎对他们并不害怕。然而你却止于玛卡里和乌尔撒奇。你所提到的这两个人，就是你所说的那个巨大数目吗？你是在想我们小时候学过的功课吗？因为如果你问我二这个数是单数还是复数，我能回答什么呢，除了它是复数？但即使如此，我仍然不是没有回答的方法。我从你的名单中除去玛卡里；因为你肯定没有告诉我们他是如何丧亡的。或者你是否主张，任何迫害你的人，除非他在这地上不朽，否则当他死时，都被认为是因为你而死？如果君士坦丁没有享受如此长时间的统治和如此长久的繁荣，他可是第一个颁布许多法令反对你们谬误的人，那又怎样？如果尤利安，他把大殿归还给你们，没有如此迅速地被夺去生命，那又怎样？那样的话，你们什么时候才会停止说这种胡话呢？因为你们现在还不愿住口。然而我们也不说尤利安死得这么快是因为他把大殿归还给你们。因为在这一点上，我们同样可以和你们说得同样冗长，但我们不愿同样愚蠢。那么，正如我开始说的，从这两个人中，我们将除去玛卡里。因为当你提到两个名字，玛卡里和乌尔撒奇时，你重复了乌尔撒奇的名字，目的是向我们展示他如何罪有应得；你说：\Quote{因为乌尔撒奇在与蛮族的战斗中被杀，之后猛禽用它们凶残的利爪，和贪婪的狗牙撕咬他，将他肢解。}由此十分清楚，既然你惯于因玛卡里而激起对我们更大的仇恨，甚至你称我们为玛卡里派而非乌尔撒奇派，那么如果你能对他的死说出任何类似的事情，你一定会说得最多。因此，在这两个之中，当你使用复数时，如果除去玛卡里，剩下乌尔撒奇一人，一个单数的专有名词。那么，你那个威胁性的、惊人的、许诺有众多支持你论点的人的承诺，又在哪里实现呢？\par

209. 我认为，到了现在，所有对词语含义稍有了解的人都应该明白，这是多么可笑：当你说了\Quote{马卡里乌斯死了，乌尔撒奇乌斯死了，你们所有那些伯爵也都同样因天主的报复而死了}，仿佛人们要求你证明事实——然而事实上，无论是听众还是读者，都没有再要求你任何东西——你立刻串起了一大段论证，试图证明我们所有的伯爵都同样因天主的报复而死。你说：\Quote{因为乌尔撒奇乌斯在与野蛮人的战斗中阵亡，之后猛禽用它们凶残的利爪，和狗贪婪的牙齿撕咬他，将他肢解。}同样地，任何其他同样不晓其所言含义的人，都可能断言你们所有的主教都因天主的报复而死在狱中；当被问及如何证明这一事实时，他可能立刻补充说：因为奥普塔图斯被指控属于吉尔多的同党，以类似的方式被处死。我们不得不听、思考并驳斥这样轻率的指控；但我们为软弱者担忧，唯恐因他们智力更迟钝，会迅速落入你们的圈套。至于你所说的乌尔撒奇乌斯，如果他是度圣善生活的，并确实如你所断言的那样死去，他将从天主的应许中得到安慰，天主说：\BibleQuote{我必追讨你们生命的血债；向一切野兽我必追讨。}\BibleRef{创 9:5}\par

210. 至于你向我们提出的诽谤性的指控，说我们激起世上君王的愤怒来反对你们，因为我们没有教导他们圣经的教训，反而进献我们自己的战争欲望，我不认为你完全聋到听不见圣书本身的雄辩，以至于你反而不担心他们会熟悉这雄辩。但无论你愿意与否，他们进入教会；即使我们缄口不言，他们也注意听诵读者；且不说别的，他们特别留意听你所引用那篇圣咏。因为你说我们不教导他们，也尽可能不让他们熟悉圣经的话：\BibleQuote{所以，你们君王，应自觉受教；你们世上裁判官，应受教。以敬畏事奉主，战战兢兢地向祂欢呼。接受教导，免得主发怒……}相信这首圣咏被咏唱，他们也听到了。但是，无论如何，他们听到了同一圣咏中前面所写的话，你，除非我弄错了，只不愿跳过它，以免被人认为你害怕。因此他们也听到了这话：\BibleQuote{主对我说：你是我的儿子；我今天生了你。你向我请求，我便将外邦人赐给你作为产业，将地极赐给你作为属你。}听到这些，他们不能不惊奇，竟有人反对基督的这个产业，企图将其缩小到大地的一角；在他们惊奇之际，也许由于听到下面的话：\BibleQuote{以敬畏事奉主}，他们便询问，在作为君王的能力范围内，他们如何能事奉主。因为所有人都应该事奉天主——一方面，因他们共同的人性条件；另一方面，因他们不同的恩赐，其中这个人在地上有一个职务，另一个人有另一个职务。因为没有一个人作为私人，能够命令将偶像从地上除去，而这早在很久以前就预言必将实现。因此，当我们考虑到人类的社会状况时，我们发现君王，正因为他们是君王，有一种他们可以用来事奉主的方式，这是任何没有王权的人所不能做到的。\par

211. 因此，当他们思考你所引述的内容时，他们也听到了你自己所引述的关于三个青年的事迹，并且是在极其庄严的场合下听到的。因为同一段圣经在教会中演唱最多的时候，正是节庆本身的氛围使那些在一年中其余时间更怠惰的人也激发起额外的热忱。那么，你认为基督徒皇帝们听到三个青年因为不愿意同意敬拜国王雕像的邪恶而被投入燃烧的火窑时，会作何感想？除非你认为他们考虑的是：圣徒的虔敬自由既不能被国王的权力所征服，也不能被任何严厉的惩罚所压制，并且他们庆幸自己不属于那些惩罚藐视偶像者如同犯渎圣罪一样的国王之列。但是，此外，当他们听到随后发生的同一国王被神奇景象吓坏——不仅是三个青年，而且火焰本身也在服务天主——他也开始敬畏地事奉天主，虔诚地喜乐，并接受教训时，他们难道不明白之所以如此公开地记载和陈述这些，是为了给天主的仆人们树立榜样，以免他们顺从国王而犯渎圣之罪，同时也给国王们自己树立榜样，使他们因信仰天主而显示虔诚吗？因此，他们基于你本人插入你著作中的同一圣咏的劝诫，愿意变得明智，接受教训，以敬畏事奉天主，以虔诚向他喜乐，并接受教训，那么他们后来听到那国王所说的话时会多么关注！因为他宣布，他要为他所统治的所有人民颁布一条法令：凡说话亵渎\Person{沙得辣客}、\Person{默沙客}、\Person{阿贝得乃哥}的天主的人，必被处死，他们的房屋也必被彻底摧毁。如果他们知道颁布这条法令是为了禁止亵渎那位减弱火焰力量并解放三个青年的天主，那么他们当然会进一步考虑，他们应当在自己的王国中颁布什么样的法令，以使那位赐予罪赦、给予全地自由的天主，不至于在他们王国中的信徒中被藐视。\par

212. 因此，当基督徒皇帝们为维护公教合一而颁布任何反对你们的法令时，要小心：不要嘴上指责他们似乎不懂圣经，心中却悲伤他们如此熟悉圣经的教导。因为谁能容忍你们在同一个达尼尔事件中提出的亵渎可憎的谬论——一方面挑剔国王，因为他被投入狮子坑；另一方面却拒绝赞扬国王，因为他被提升至崇高尊荣？因为即使当他被投入狮子坑时，国王本人更倾向于相信他会安全，而不是会被毁灭，并且因担忧他而拒绝进食。然后你们竟敢对基督徒说：\Quote{你们与世上的国王有何相干？}因为达尼尔曾遭受国王的迫害，但你们却不回顾同一个达尼尔忠实地为国王解梦、称国王为主、接受国王的礼物和尊荣吗？同样，你们竟敢在上述三个青年的事件中，煽动对国王的仇视火焰，因为他们拒绝崇拜雕像而被投入火焰，而同时你们却闭口不言，一点不提他们后来被国王如此赞扬和尊崇？就算国王是迫害者，当他将达尼尔投入狮子坑时；但当他平安地接待他出来，在喜乐和庆贺中把达尼尔的敌人扔进同一群狮子被撕裂吞吃时，他那时是什么——是迫害者，还是不是？\EditorialNote{达尼尔书第二至六章}我要求你们回答我。因为如果他是，为什么达尼尔本人不反抗他，他本可以借着与国王的深厚友谊轻易做到，而你们却命令我们制止国王迫害人？但如果他不是迫害者，因为他以迅捷的正义惩罚了对圣人的暴行，那么请问：如果先知的肢体因为遭受危险而需要这样的报复，那么对基督圣事的侮辱，国王应该施加什么样的报复呢？我再次承认，国王确实是迫害者，当他把三个青年投入火窑因为他们拒绝崇拜他的雕像；但我要问，当他颁布法令说凡亵渎唯一真天主的都要被消灭，全家被废时，他仍是迫害者吗？因为如果他是迫害者，你们为什么要对迫害者的话回答\Emph{阿们}？但如果他不是迫害者，你们为什么称那些阻止你们亵渎疯狂的人为迫害者？因为如果他们强迫你们崇拜偶像，那么他们就像不虔敬的国王，而你们像三个青年；但如果他们阻止你们与基督作战，那么你们若试图这样做，就是不虔敬的。但他们如果以可怕的威胁禁止这事，他们可能是什么，我不敢说。如果你们不愿称他们为虔敬的皇帝，就给他们另找一个名称吧。\par

213. 如果是我提出达尼尔和三个青年的这些例子，你们也许会抗拒，并声明这些例子不应从那些时代引用来类比我们的时代；但感谢天主，你们自己提出了它们，确实是为了证明你们想要确立的观点，但你们看到它们的力量反而更有利于你们最不愿证明的。也许你们会说，这不是出于你们的欺骗，而是人性的易错性。但愿如此！那么就改正吧。你们不会失去名誉；不，明确地标记更高超的心智的是：通过坦诚的认错来熄灭敌意的火焰，胜过仅凭理解力的敏锐来逃避谎言的迷雾。\par

% structure: p0309 source=h2 target=subsection reason=relative-heading-rank; base=h2

\subsection{第九十四章}

214. 培提里安说：\Quote{天主的法律在哪里？你们的基督信仰在哪里？如果你们不仅杀人害命，而且还下令做这样的事？}\par

215. 奥斯定答：对此，请看基督的共同继承人如何在全世界发言。我们既不杀人害命，也不下令做这样的事；而你们比那些做这样事的人更加疯狂，因为你们把这些想法植入人心，反对永生的希望。\par

% structure: p0312 source=h2 target=subsection reason=relative-heading-rank; base=h2

\subsection{第九十五章}

216. 培提里安说：\Quote{如果你们希望我们做你们的朋友，为什么违背我们的意愿强拉我们？但如果你们希望我们做你们的敌人，为什么杀死你们的敌人？}\par

217. 奥斯定答：我们既不违背你们的意愿强拉你们，也不杀死我们的敌人；但我们在与你们的交往中所做的一切，虽然可能违反你们的意愿，却是出于对你们的爱，好使你们自愿改正，过改过自新的生活。因为没有人是违背自己的意愿而活的；然而，一个男孩为了学会自由意志的这一课，却是违背自己的意愿挨打，而且往往是那个最亲他的人打的。这正是君王们想要对你们说的话，如果他们打你们的话，因为他们的权力是由天主所设定的。但即使他们不打你们，你们也大声喊叫。\par

% structure: p0315 source=h2 target=subsection reason=relative-heading-rank; base=h2

\subsection{第九十六章}

218. 培提里安说：\Quote{但有什么理由，或是怎样的空虚矛盾，使你们如此急切地渴望与我们共融，而同时却用异端这个虚假的头衔称呼我们？}\par

219. 奥斯定答：如果我们如此急切地渴望与异端共融，我们就不会焦虑你们能从异端的错误中悔改；但当我们与你们交涉的目的正是要你们不再做异端时，我们又怎能急切地渴望与异端共融呢？事实上，正是分裂和分离使你们成为异端；而和平与合一使人成为天主教徒。因此，当你们从你们的异端归向我们时，你们就不再是我们所憎恨的，而开始成为我们所爱的。\par

% structure: p0318 source=h2 target=subsection reason=relative-heading-rank; base=h2

\subsection{第九十七章}

220. 培提里安说：\Quote{简言之，请在两者中选择一个：如果无辜在你们一边，为什么用刀剑迫害我们？或者，如果你们称我们有罪，那么你们自己既然无辜，为什么寻求我们的团契？}\par

221. 奥斯定答：最巧妙的二难推理，或者说，最愚蠢的废话！当对方不可能同时采纳两者时，才向他提供两个选择；不是吗？如果你向我提供两个命题的选择，让我说我们是无辜的还是有罪的；或者，再就另一对命题，即关于你们的，我无法避免选择其中一个。但事实上，你向我提供这两个选择：我们是无辜的，还是你们是有罪的，并希望我说出我选择哪个来回答。但我拒绝选择；因为我两者都主张：我们是无辜的，你们是有罪的。我说，对于你们加给我们的虚假和诽谤的指控，我们中凡是属于天主教会的人，都可以凭平安的良心说，我们没有交出圣书、没有参与拜偶像、没有杀人、也没有犯你们指控我们的其他任何罪过；而任何可能犯过此类过错的人（尽管你们从未在任何案件中证明过），都因此关上了天国的大门，不是针对我们，而是针对他们自己；\BibleQuote{因为各人要背负自己的重担。}\BibleRef{迦 6:5}这里你有了第一个回答。我进一步说，你们都是有罪的和被诅咒的——不是你们中某些人因他人所犯的罪（这些罪在你们中间由某些人施行，并由另一些人谴责），而是你们全体因分裂之罪；从这最严重的亵渎中，你们中没有人可以说自己自由，只要他拒绝与普世万民的合一共融，除非他被迫说基督在论及从耶路撒冷开始、传遍万民的教会时说了谎。\BibleRef{路 24:47}于是你有了我的第二个回答。看，我给了你两个回答，而你原本希望我们被迫选择其中一个。无论如何，你应当注意到我们可以同时作出两个断言；当然，如果这是你所希望的，你应当请求我们选择其中一个，而当你看到我们有权选择两者时。\par

222. 但\Quote{如果无辜在你们一边，为什么用刀剑迫害我们？}请回头看看你们的队伍，他们现在不仅像古代他们祖先那样装备了棍棒，而且还增添了斧头、长矛和刀剑，然后自己判断我们中谁更适合这个问题：\Quote{为什么用刀剑迫害我们？}\Quote{或者，如果你们称我们有罪，}你们说，\Quote{那么你们自己既然无辜，为什么寻求我们的团契？}这里我回答得非常简洁。为什么有罪的你们被无辜者寻求？是为了使你们不再有罪，而开始成为无辜。这样，我在关于我们的两个选项中选择了两者，并回答了关于你们的两个选项；只是你们反过来选择一个。你们是无辜还是有罪？这里你们不能选择同时作出两个断言，但可以两者都选，如果你们乐意的话。因为至少你们不能在同一情境中同时是无辜和有罪的。因此，如果你们是无辜的，不要惊讶于你们被邀请与弟兄们和平相处；但如果你们是有罪的，不要惊讶于你们被君王们寻求惩罚。但既然你们为自己选定了这两个选项中的一个，而另一个则是由我们加给你们的——你们自认无辜，而我们指证你们生活不虔敬——那么请再听一次我将就每个方面说的话。如果你们是无辜的，为什么反对基督的见证？但如果你们是有罪的，为什么不去投奔他的仁慈？因为他的见证，是一方面，关乎世界的合一；他的仁慈，另一方面，在于兄弟之爱。\par

% structure: p0322 source=h2 target=subsection reason=relative-heading-rank; base=h2

\subsection{第九十八章}

223. \Person{佩提利安}说：\Quote{最后，正如我们此前多次说过的，你们竟敢如此狂妄，竟敢谈论君王，而\Person{达味}却说：\InnerQuote{信赖主比信赖人更好；信赖主比信赖君王更好？}}\par

224. \Person{奥古斯丁}回答说：我们不信赖人，而是尽力劝人信赖主；我们也不信赖君王，而是尽力劝君王信赖主。尽管我们可能寻求君王的帮助以促进教会的利益，但我们并不信赖他们。因为宗徒本人也不信赖那位千夫长，像圣咏作者谈论信赖君王那样；宗徒从千夫长那里获得指派一队武装士兵护送他；他也不信赖那些武装士兵，靠他们的保护逃脱恶人的圈套，像圣咏作者谈论信赖人那样。\BibleRef{宗 23:12}但我们也不责备你们自己，因为你们向皇帝请求将大殿归还给你们，好像你们信赖了\Person{朱利安}君王一样；但我们责备你们，因为你们对基督的见证绝望了，而你们将那些大殿与基督的合一分离了。因为你们是按照基督的仇敌的命令接受了它们，为的是在它们里面藐视基督的命令，而你们却在\Person{朱利安}所颁布的命令中找到了力量和真理，他说：\Quote{此外，根据\Person{罗加提亚努斯}、\Person{庞提乌斯}、\Person{卡西亚努斯}及其他主教（并非没有圣职人员的夹杂）的请求，特此补充完整，凡不当且无权对他们造成损害的程序，一律废除，一切应恢复原状。}而你们却在基督所颁布的命令中找不到任何有力量或真理的话，他说：\BibleQuote{你们要在耶路撒冷、全犹太和撒玛黎雅，并直到地极，为我作见证。}\BibleRef{宗 1:8}我们恳求你们，让你们自己改过自新吧。回到这最明显的普世合一中来；让一切恢复原状，不是按照背教者\Person{朱利安}的话，而是按照我们的救主基督的话。可怜你自己的灵魂吧！我们现在并不是在比较\Person{君士坦丁}和\Person{朱利安}，以显示他们有多么不同。我们并不是说，如果你们不曾信赖一个人和一个君王，当你们对一位外教背教皇帝说\Quote{在他身上正义只找到了一个位置}时，鉴于多纳徒派普遍使用了含有那些话的祈祷文和敕令，正如听证记录所证明的那样；那么，我们就更不应该被你们控告，好像我们信赖任何一个人或君王似的，如果我们没有用任何亵渎的谄媚，而是从\Person{君士坦丁}或其他基督信仰的皇帝那里获得了任何请求；或者如果他们自己，没有我们请求，而是记念他们应向主交的账，并且他们在听到你们自己所引用的\BibleQuote{所以，你们君王应当明智}等等，以及其他许多类似的话时，就颤抖，于是主动颁布任何法令来支持天主教会合一的话。但我对\Person{君士坦丁}不提什么。我们放在你们面前对照的是基督和\Person{朱利安}；不仅如此，而是天主和人，天主之子和地狱之子，我们灵魂的救主和他自己的毁灭者。为什么你们坚守\Person{朱利安}的敕令以占据大殿，却不坚守基督的福音以拥抱教会的和平？我们也高呼：\Quote{凡一切错误的行事都当恢复其原始状态。}基督的福音比\Person{朱利安}的敕令更古老；基督的合一比多纳徒派更古老；教会为合一向主的祈祷比\Person{罗加提亚努斯}、\Person{庞提乌斯}和\Person{卡西亚努斯}为多纳徒派向\Person{朱利安}的祈祷更古老。当君王禁止分裂时，行事是错的吗？当主教分裂合一时，行事难道不是错的吗？当君王为捍卫教会而服侍基督的见证时，那是错的行为吗？当主教为否认教会而反对基督的见证时，那难道不是错的行为吗？因此，我们恳求你们，听一听你们这样哀求的\Person{朱利安}本人的话，不是违背福音，而是符合福音，好使\Quote{凡一切错误的行事都可能恢复其原状}。\par

% structure: p0325 source=h2 target=subsection reason=relative-heading-rank; base=h2

\subsection{第99章。}

225. \Person{佩提利安}说：\Quote{我呼吁你们，是的，你们这些可怜的人，你们因惧怕迫害而惊惶，你们寻求保全财富，而非你们的灵魂，你们所爱的不是叛徒那不忠之信，而是你们为自己争取到保护的那些人本身的邪恶——正如在波涛中遭遇海难的水手，投身于必将淹没他们的波涛，在生命极度危险中，毫不含糊地寻求他们恐惧的对象；又正如暴君的疯狂，为了免于对任何人的恐惧，渴望被人畏惧，尽管这对自己充满危险：同样，同样你们逃往邪恶的堡垒，宁愿看到无辜者的损失或惩罚，只求为自己免除恐惧。如果你们认为投身毁灭就能躲避危险，那么，遵守强盗的信义也确实是该受谴责的信。最后，为了不失去财富而失去自己的灵魂，这是贩卖疯子的利润。因为主基督说：\BibleQuote{人若赚得全世界，却赔上自己的灵魂，有什么能换回他的灵魂呢？}\BibleRef{玛 16:26}}\par

226. 奥古斯丁答：我不能不承认，你若将此劝勉用于善事，本是大有裨益的。你试图劝阻人不要把财富置于灵魂之上，这无疑是好的。然而，我愿你这听过这些话的人，也暂且听听我们的话；因为我们也说这话，但请听是在何种意义上说的。如果君王威胁要夺去你的财富，只因你不是按肉身的犹太人，或因为你并不敬拜偶像或魔鬼，或因为你未被卷入任何异端，而是持守于大公合一之中，那么你就宁可使你的财富消亡，免得你自身消亡；但要小心，勿将任何事物，连这世界本身的生命，置于在基督内的永生之上。但如果君王威胁要剥夺你的财产或判你有罪，仅仅因为你是异端者，那么他们恐吓你，并非出于残忍，而是出于怜悯；而你不肯惧怕，这非勇气的记号，而是顽固的记号。那么请听伯多禄的话，他说：\BibleQuote{你们若因犯罪被打而坚忍，有什么光荣？}\BibleRef{伯前 2:20}因此，你们在此既无地上的安慰，在来世也无永生；你们在此有的只是不幸者的苦难，在那里有的则是异端者的地狱。所以，我现今与之辩论的兄弟，你看见了吗？你应当先证明你所持的是真理，然后再劝勉人在持守真理时，愿意舍弃今生所拥有的一切福乐。然而，当你不能证明这一点时——不是因为才能欠缺，而是因为理由不善——你又何必藉着劝勉，急于使人既成为乞丐又成为无知者，既贫乏又偏离真理，衣衫褴褛而争论不休，家奴兼异端者，既在今生失去暂世之物，又在基督的审判中寻得永罚呢？谨慎的儿子因畏惧父亲的棍棒，而远离蛇穴，既免挨打，又免丧亡；而藐视惩戒之苦，执意与自己的有害意志对抗的人，则既被打又丧亡。博学之士啊，你如今还不明白吗？凡舍弃一切地上财物以维护基督平安的人，天主是他所有的；而为了多纳图派而失去即使很少钱币的人，却失去了心灵。\par

% structure: p0328 source=h2 target=subsection reason=relative-heading-rank; base=h2

\subsection{Chapter 100.}

227. 伯提利安说：\Quote{\BibleQuote{但我们这些神贫的人}\BibleRef{玛 5:3}，并不为财富忧虑，反而畏惧财富。我们\BibleQuote{似乎一无所有，却无所不有}\BibleRef{格后 6:10}，以我们的灵魂为财富，藉我们的苦难和鲜血，为自己购得永生的天国之富。同样，主又说：\InnerQuote{凡丧失自己财产的，必要再得百倍。}}\par

228. 奥古斯丁答：探究这段经文的正解并非无关紧要。因为你在引用圣经时若有任何错误或虚假印象，只要这并不妨碍我的目的，我便不予深究。其实经文并非写：\Quote{凡丧失自己财产的}，而是\BibleQuote{凡为我的缘故丧失自己性命的}\BibleRef{玛 16:25}。至于论财产那段，不是\Quote{凡丧失的}，而是\BibleQuote{凡舍弃的}\BibleRef{玛 19:29}，并且不仅指钱财，还指其他许多事物。而你其实并未丧失你的财产；至于你是否舍弃了它，以至于如此夸耀贫穷，我不能断言。倘若我的同道福图纳图斯恰巧知道此事（因他与你同在一城），他从未告诉过我，因我也从未问过他。然而，即使你果真如此做了，你在这封信中自己引用了宗徒反对你的证言：\BibleQuote{我若将一切所有的施舍给人，又舍己身叫人焚烧，却没有爱德，于我无益。}\BibleRef{格前 13:3}因为你若有爱德，就不会指控那不认识你、你也不认识它的普世教会——不，甚至连基于非洲人已证实的过错而提出的指控，你也不会提出。你若有爱德，就不会在诽谤中为自己构想一个虚假的合一，而会学习认出主的话中清楚昭示的合一：\BibleQuote{直到地极}\BibleRef{宗 1:8}。但你若没有这样做，又何必夸耀仿佛已做了？你真的如此畏惧财富，以至于一无所有却无所不有吗？请告诉你的同道克里斯皮努斯，他最近在靠近我们希波城的地方买了一个庄园，好在那里把人推入最深的深渊。我也因此深知此事。你或许不知道，因而安然高呼：\Quote{我们畏惧财富。}所以我奇怪你的喊声竟能被克里斯皮努斯允许，传到我们这里。因为你在君士坦丁纳，我在希波，中间的卡拉马是他的所在地，且更靠近我这一边，但仍在我们之间。所以我很奇怪，他怎未先拦截这喊声，把它顶回去，不让它传到我耳中；他又怎未以更丰富的言辞，针对你，诵读一篇对财富的颂赞。因为他不仅不畏惧财富，实际上还喜爱它们。当然，在你对其他人说任何话之前，应先与他把这些观点说清楚。若他不改正，我们自有回答。但就你自己而言，如果你真是穷人，我的兄弟福图纳图斯就在你身边。你抱有这等人，更可能讨我的同道喜欢，而不是你那位同道克里斯皮努斯。\par

% structure: p0331 source=h2 target=subsection reason=relative-heading-rank; base=h2

\subsection{Chapter 101.}

229. 伯提利安说：\Quote{因我们生活在敬畏天主之中，我们并不害怕你们用刀剑施行的刑罚和处决；我们唯一躲避的，是你们以极邪恶的共融毁灭人的灵魂，正如主亲自所说：\BibleQuote{不要害怕那些杀害肉身而不能杀害灵魂的；但更要害怕那能使灵魂和肉身都陷于地狱中的}\BibleRef{玛 10:28}}\par

230. 奥古斯丁答道：你们所讲的毁灭，并非用看得见的刀剑，而是用那把刀，经上论及它说：\Quote{人子们的牙齿是枪和箭，他们的舌头是快刀。}因为你们用这把指责和诽谤的刀剑，攻击你们全然不认识的世界，毁灭了那些缺乏经验之人的灵魂。但若你们指责一个最邪恶的共融——如你们所称——我愿立刻用你们自己的话，而非我的话，要你们上升、下降、进入、转身、改换立场，成为像奥普塔图斯那样的人。但如果你们清醒过来，发现你们自己并不像他，并非因为他拒绝与你们共享圣事，而是因为你们对他的行为感到忿怒，那么你们就会为世界开脱不属于它的罪行，而你们将发现自己陷入分裂的罪中。\par

% structure: p0334 source=h2 target=subsection reason=relative-heading-rank; base=h2

\subsection{第一百零二章}

231. 佩提利安说：\Quote{因此，你们宁愿用最虚假的圣洗来洗濯，而拒绝重生，不仅没有除去你们的罪，反而用罪犯的过犯来重重地压住你们的灵魂。因为，犯罪之人的水既然被圣神所抛弃，它就显然充满了\OriginalTerm{背教者}{traditors}的过犯。于是，对任何一个由这种人所施洗的可怜人，我们要说：如果你想挣脱虚假，你实际上已浸透在虚假中；如果你想驱除肉身的罪恶，那罪人的良心临到你身上，你也就同样分沾了他们的罪债；如果你想熄灭贪婪的火焰，你已浸透在欺诈里，已浸透在邪恶里，也浸透在疯狂里。最后，如果你相信施洗者与受洗者的信仰是相同的，你就被他用杀人的血所浸透，因为他在杀人。于是，结果就是：你前来领洗时无罪的，领洗后却带着杀人的罪回去了。}\par

232. 奥古斯丁答道：我愿与那些在听到或读到你们这些话时大声附和的人进行辩论。因为这样的人，心里没有耳朵，却把心放在耳朵里。然而，让他们一读再读，仔细思量，自己找出这些话的声音是什么，而是它们的意思是什么。首先，为了筛选最后一句话的意思，你们说：\InnerQuote{结果就是：你前来领洗时无罪的，领洗后却带着杀人的罪回去了。}请先告诉我，除了那一位前来受洗并非为洗去他的罪，而是为给我们留下谦卑榜样的那一位之外，有谁是带着无罪之身前来领洗的？因为，对无罪之人，还有什么可宽恕的呢？难道你们真有这样的口才，能向我们显明某种无罪却又犯罪的清白吗？你们没有听到圣经的话说：\BibleQuote{在你眼中，没有一个人是干净的，连只活了一天的婴儿也不例外}吗？不然，为什么甚至要急忙抱着婴儿去找人赦免他们的罪呢？你们没有听另一处经上说：\BibleQuote{我母亲在罪恶中怀了我}吗？其次，如果一个人仅仅因为从杀人犯手中领受圣洗，就从一个无罪的人变成了杀人犯，那么所有从奥普塔图斯手中领受圣洗回来的人，都成了奥普塔图斯的同谋。你们现在去吧，看看你们以怎样的脸面拿我们激励君王的愤怒来指责我们？你们难道不怕，在你们中间搜查吉尔多的党羽时，凡是曾由奥普塔图斯施洗的人都算在内吗？你们终于明白，你们那句话，就像一个空尿泡，不仅发出毫无意义的响声，而且还打在你们自己头上吗？\par

233. 再来看你们摆在我们面前供反驳的其他较早的论点，那些论点的性质是：我们必须承认，每个人领洗回来时都带着施洗者的特质；但天主禁止你们所施洗的人从你们那里回来时，沾染上你们在说这种话时所具有的疯狂！你们的话里有多么美妙的音调啊：\InnerQuote{你已浸透在欺诈里，已浸透在邪恶里，也浸透在疯狂里！}除非你们不是浸透，而是完全充满了疯狂，否则你们绝不会说出这样的话来。那么，暂且不说别的，难道所有清白地、克服了贪婪前来从你们贪婪的同僚或你们党派的司铎那里领洗的人，回去时都犯了贪婪？难道那些清醒地跑来投入醉酒的漩涡中受洗的人，回去时都醉了？如果你们持守并教导这样的观点，你们甚至会有脸引用你们刚才所用的那段反对我们的经文：\Quote{信赖主，胜过信赖世人。信赖主，胜过信赖权贵。}请问，你们的教导难道不是这个意思吗：我们不应信赖主，而应信赖人？因为你们说，受洗者变得像施洗者一样。既然你们以此作为你们洗礼的基本原理，那么人们就该信赖你们吗？而那些原本信赖主的人，就该信赖权贵吗？我实在愿意他们听从你们，而是听从你们反对我们所引的那些证据，甚至是更可怕的话；因为经上不仅记着：\Quote{信赖主，胜过信赖世人}，而且还说：\Quote{凡信赖世人的人，是可咒骂的。}\BibleRef{耶 17:5}\par

% structure: p0338 source=h2 target=subsection reason=relative-heading-rank; base=h2

\subsection{第一百零三章}

234. 佩提利安说：\Quote{你们要效法先知，他们害怕自己神圣的灵魂被虚假的圣洗所欺骗。因为耶肋米亚在古代说过，在邪恶之人那里，水就如同说谎者。他说：\InnerQuote{说谎的水没有信德。}}\par

235. 奥斯定回答说：任何听见这些话的人，若不熟悉圣经，且不相信你们要么是深陷迷途以致不知自己在说什么，要么是处于欺骗中，以致被你欺骗的人也不知道他在说什么，都会相信先知耶肋米亚想要受洗，便采取了措施以免被不虔敬的人施洗，并为此意图说了这些话。因为你在引述这段经文之前说：\Quote{要效法先知们，他们害怕他们神圣的灵魂被虚假的圣洗所欺骗？}你的目的何在？就好像在耶肋米亚的日子，有人用圣洗圣事洗净自己，除非法利塞人几乎每时每刻沐浴自己，以及他们的床、杯盘，用主所谴责的洗涤，正如我们在福音中所读到的。\BibleRef{谷 7:4}那么，耶肋米亚怎能这样说，好像他想要受洗，并设法避免被不虔敬的人施洗？他说这话，是在抱怨一个无信的人民，他们的道德败坏使他烦恼，他不愿与他们的行为有份；然而他并没有从身体上离开他们的会众，也没有寻求那些人民按照梅瑟法律在那个时代所接受的圣事以外的其他圣事。因此，对于这个生活邪恶的人民，他称之为\Quote{创伤}，义人的心因此被重击，无论他这样说是指自己，还是在自己身上预表他所预见将要发生的事。因为他这样说：\BibleQuote{上主，求祢记念我，眷顾我，在那些不怀忍耐而迫害我的人面前显明我的清白；求祢知道我为祢的缘故忍受了那些蔑视祢言语之人的责备。使他们的份归全；祢的言语将是我心中的喜乐和欢愉；因为我被称为祢名下，上主万军的天主。我没有坐在嘲弄者的集会中，却在祢手的面前惧怕；我独坐，因为我充满了苦楚。为什么那些使我忧伤的人胜过我呢？我的创伤是严重的；我怎能痊愈？它对我已像虚假的水，没有信实。}在这一切中，先知所希望被人理解的是显而易见的，但只对那些不愿将所读之意义曲解为他们乖僻目的的人显而易见。因为耶肋米亚说他的创伤对他已像虚假的水，不能激发信德；但他希望藉着他的创伤，那些因他们生活的邪恶行为而使他忧伤的人被理解。因此，宗徒也说：\BibleQuote{外有争斗，内有恐惧}\BibleRef{格后 7:5}又说：\BibleQuote{谁软弱，我不软弱呢？谁跌倒，我不焦急呢？}\BibleRef{格后 11:29}。因为他没有希望他们能悔改，因此他说：\Quote{我怎能痊愈？}好像他的痛苦必定会继续，只要他被迫生活在其中的人仍然是他们那样。但是，通常用水这个称谓来指代人民，这在默示录中得到显示，在那里我们理解\Quote{许多水}是指\Quote{许多人民}，这不是凭我们自己的臆测，而是凭经文本身的明确解释。\BibleRef{默 17:15}因此，不要因误解（更确切地说是错误）而亵渎圣洗圣事，即便它是在一个最堕落的人身上；因为即使在说谎的西满身上，他所领受的圣洗也不是虚假的水，\BibleRef{宗 8:13}你们党派中所有说谎的人，当他们在天主圣三之名内施行圣洗时，也不是施予虚假的水。因为他们并不是只有在被出卖、被定罪、被迫承认他们的恶行时才成为说谎者；而是他们早已是说谎者，当他们身为奸夫和被诅咒者，却假装贞洁无辜时。\par

% structure: p0341 source=h2 target=subsection reason=relative-heading-rank; base=h2

\subsection{第104章}

236. 帕蒂利安说：\Quote{达味也说：\InnerQuote{罪人的油不可抹我的头。}那么，他称谁是罪人呢？是我遭受你的暴力，还是你迫害无辜者？}\par

237. 奥古斯丁答：我代表基督的身体，即永生天主的教会，真理的柱石和基础，它散布于全世界，因福音已按宗徒的话传给\Quote{天下万民}。\BibleRef{哥 1:23} 我代表整个世界，达味的话——你们无法理解——论到这世界说：\Quote{世界坚定，不可动摇。}而你们却争辩说它不仅动摇了，而且完全毁灭了。我代表这世界回答：我不迫害无辜。但达味说：\Quote{罪人的油膏}，而非\OriginalTerm{叛教者}{traditor}的油膏；非献香者，非迫害者，而是\Quote{罪人的油膏}。那么你们如何解释呢？首先看你们自己是不是罪人。若你们说：我不是叛教者，不是献香者，不是迫害者，这无关紧要。我本人，因天主的恩宠，并非这些，那不可动摇的世界也不是。但你们若敢说：我不是罪人，因为达味说\Quote{罪人的油膏}。只要你们身上有任何罪，即使是小罪，你们有什么理由坚持自己与\Quote{罪人的油膏}这话无关呢？我要问你们，在祈祷中是否用主祷文？若你们不用我主教导门徒的祈祷文，你们从哪里学来另一篇配得上你们功绩、超过宗徒功绩的祈祷文？若你们祈祷，如我们伟大的导师屈尊教导我们的，你们怎能说：\Quote{宽免我们的罪债，犹如我们也宽免得罪我们的人}？在这祈求中，我们并非指在圣洗中已被赦免的罪。因此，祷文上的这些话要么将你们排除于向天主祈求之外，要么表明你们也是罪人。那么，让那些受你们洗礼、因你们的油膏而使头败坏的人来亲吻你们的头吧。但你们看清自己，既知自己是何人，也知你们如何看待自己。实在的，奥普塔图斯，异教徒、犹太人、基督徒、我们这方的人、你们这方的人，整个阿非利加都在宣告他是个盗贼、叛国者、压迫者、分裂的制造者；不是朋友，不是门客，而是那位的工具——你们中有人宣称他是自己的伯爵、同伴和神明——按任何可理解的意义，他岂不是罪人？那么，那些被一个犯有死罪的人所傅油的人，他们的头将如何呢？你们用这节经文的解释判断他们的头，难道那些人自己不亲吻你们的头吗？我真要你们把他们带来，劝他们自己医治。或者，该被医治的是你们的头，你们如此严重地误入歧途？那么你们会问：达味究竟说了什么？你们何必问我：不如问他本人。他在上文回答你们：\Quote{义人打击我，是出于仁慈，并责备我；但不要让罪人的油膏傅在我的头上。}还有什么更清楚？更明显？他说：我宁愿被仁慈的责打所医治，也不愿被甜言蜜语所欺骗和误导，如油膏落在头上。同样的思想在圣经别处以另一句话表达：\Quote{朋友的创伤胜于仇敌的亲吻。}\par

% structure: p0344 source=h2 target=subsection reason=relative-heading-rank; base=h2

\subsection{第105章}

238. 佩蒂利安说：\Quote{但他这样称赞弟兄和睦的油膏：\InnerQuote{看，兄弟们同居共处，多么美好，多么快乐！像珍贵的油傅在头上，流到胡须上，亚郎的胡须上，又流到他的衣襟上；像赫尔孟的甘露，又像降在熙雍山上的甘露；在那里上主颁赐了祝福，就是永生。}由此，他说，合一被傅油，正如司铎被傅油。}\par

239. \Person{奥古斯丁}答说：你所说的是真的。因为在基督身体内的那个司祭职曾领受傅油，它的救恩由合一之纽保证。事实上基督自己得名于圣油\OriginalTerm{圣油}{chrism}，即傅油。希伯来人称祂为默西亚，这个词与腓尼基语十分相近，正如许多其他希伯来词——如果不是几乎所有——一样。那么，在那司祭职中，头、胡须和衣裾各指什么？按主赐我能理解的，头无非是身体的救主，宗徒论及祂说：\BibleQuote{祂是身体——教会的头}\BibleRef{哥 1:18}。胡须被恰当地理解为刚毅。因此，那些在祂教会中显示自己勇敢，并依附祂面容之光，无畏宣讲真理的人，从基督自己——如同从头——有圣油降下，即圣神的圣化。藉衣裾我们在此理解衣袍上部，穿衣者的头从中穿过。这象征在教会内信德成全的人。因为在衣裾中有成全。我想你必记得对某位富人说的话：\BibleQuote{你若愿意是成全的，去！变卖你所有的，施舍给穷人，你必有宝藏在天上；然后来跟随我}\BibleRef{玛 19:21}。他果然忧忧愁愁地走了，轻视成全，选择不成全。但难道就没有人因这样放弃世物而得以成全，以至合一之香液从头上降到衣裾上吗？因为，且不提宗徒和那些与教会这些首领与导师直接相关的人——我们认为他们在胡须中代表更高的尊严和更显著的刚毅——请读\WorkTitle{宗徒大事录}，看看那些\BibleQuote{卖了田产房屋，把价钱带来，放在宗徒们脚前；他们中没有一个说他所拥有的任何东西是自己的，却把一切归为公用；照每人所需要的分配。那相信的群众都是一心一意}\BibleRef{宗 4:32}。我深信你知道这是如此记载的。请认清，弟兄们共处合一，何其美好，何其惬意！请认清亚郎的胡须，认清属神衣袍的衣裾。查阅圣经本身，看看这些事从何处开始；你将发现是在耶路撒冷。从这衣裾，整个一体的织物遍布万邦。藉此，头进入了衣袍，使基督穿上地上万国的各种族，因为在这衣裾上显现了实际上的各种语言。那么，为什么那合一之香液——即兄弟之爱的属神馨香——所从出的头本身，竟暴露于你的抵抗，而它却见证并宣告说：\BibleQuote{必须从耶路撒冷开始，因祂的名向万邦宣讲悔改，以得罪之赦}\BibleRef{路 24:47}。而你们却愿意藉这香液来理解圣油圣事；这圣事在有形标记类别中的确是神圣的，如同圣洗本身一样，但它甚至可以存在于最恶劣的人中——他们在肉性行为中浪费生命，永远不能承受天国——因而与亚郎的胡须、衣裾或任何司祭衣袍的织物毫无关系。因为你们打算把宗徒所列举的\BibleQuote{本性私欲的作为}放在哪里？他说：\BibleQuote{这些是：淫乱、不洁、放荡、崇拜偶像、施行法术、仇恨、竞争、嫉妒、忿怒、争吵、不睦、分党、妒恨、醉酒、宴乐，以及与这些相类似的事。我先前告诉过你们，现在再次告诉你们：做这样事的人，决不能承受天主的国}\BibleRef{迦 5:19}。我把暗中所犯的淫乱放在一边；你们可随意解释不洁，我也愿意放在一边。我也把施行法术放在一边，因为没有人公开配制或施予毒药。我也把分党放在一边，因为你们愿意如此。我犹豫是否应将崇拜偶像放在一边，因为宗徒将贪吝与之并列，而贪吝在你们中间公开盛行。然而，撇开这一切，你们中间就没有放荡的、贪吝的、公开怀恨的、好争的、好斗的、易怒的、结党的、嫉妒的、酗酒的、在宴乐中浪费时间的人吗？难道没有一个这样的人在你们中受傅？难道没有著名的、你们所知沉溺于这些事或公开放纵的人在你们中死去？若你说没有，我要你考虑一下你本人是否符合这描述，因为你在争辩中明显说谎。但如果你自己与这种人分离——不是靠身体的分隔，而是靠生活的不同——并以哀叹看待你祭台周围这等群众，那我们该说什么？既然他们用圣油受了傅，但正如宗徒以真理的明证所保证的，他们将不能承受天主的国。难道我们要如此不虔地对待亚郎的胡须和他衣袍的衣裾，以至认为他们当被置于那里？绝不可这样。因此，请将可见的圣事——它既能存在于善人也能存在于恶人中，为善人是赏报，为恶人是审判——与爱德的无形傅油分别开，后者是善人独有的。请将它们分开，分开，是的，愿天主将你从多纳图派中分开，将你再召回到天主教会中——你原是从那里被他们撕裂的，那时你还是望教者，被他们用致死的区别之纽带捆绑。如今你是在熙雍山上、赫尔孟的甘露落在熙雍山上——无论你们如何理解此意——因为你不在那山上之城，这城有确切的标记，即它不能被隐藏。因此，它被万国所认识。但多纳图派不被大多数国家所知，因此它不是真城。\par

% structure: p0347 source=h2 target=subsection reason=relative-heading-rank; base=h2

\subsection{第106章}

240. \Person{Petilianus}说：\Quote{所以，你们有祸了，因为你们对圣物施暴，切断了合一的纽带；而先知却说：\InnerQuote{如果百姓犯罪，司铎要为他们祈祷；但如果司铎犯罪，谁为他祈祷呢？}}\par

241. \Person{奥古斯丁}答道：刚才我们在辩论罪人的油时，我好想要抹你的额头，好让你敢说，如果你敢，你自己是否不是罪人。你竟有胆量这样说了。何等可怖的罪！因为你声称自己是司铎，引用这先知的证言，除了证明你完全无罪外，还证明了什么？因为如果你有罪，按着你对这话的理解，谁为你祈祷呢？因为你们在可怜的人民中间这样吹嘘自己，引用先知的话：\Quote{如果人民犯罪，司铎要为他们祈祷；但如果司铎犯罪，谁为他祈祷呢？}好叫他们相信你们没有罪，并把洗刷他们的罪托付给你们的祈祷。你们真是伟人，超乎同类，属天的，如神的，与其说是人，不如说是天使，你们为人民祈祷，却不让人民为你们祈祷！你们比\Person{保禄}更正义吗？比那位大宗徒更成全吗？他常用受教者的祈祷举荐自己。\BibleQuote{你们要恒心祈祷，在此醒寤感恩，同时也为我们祈祷，求天主给我们开传道的门，能讲基督的奥秘，为此我也被捆绑，好叫我按所当说的，显明这奥秘。}\BibleRef{哥 4:2}看，宗徒尚且受祈祷，而你们却不愿意主教受祈祷。你看出你的骄傲是多么魔鬼般的吗？宗徒受祈祷，好使他按所当说的显明基督的奥秘。所以，如果你手下有虔诚的人民，你本当劝勉他们为你祈祷，使你不说出不当说的话。你比若望宗徒更正义吗？他说：\BibleQuote{如果我们说我们没有罪，便是自欺，真理不在我们内了。}\BibleRef{若一 1:8}最后，你比达尼尔更正义吗？你在这封信中引用了他，甚至说：\Quote{至正义的国王把达尼尔抛出去，如他所想，要被野兽吞噬。}——但国王从未这样想过，因为他对达尼尔本人以最友好的态度说，正如课程的上下文所显示的：\BibleQuote{你所常事奉的天主，祂必拯救你。}\BibleRef{达 6:16}关于这事我们已经说了很多。关于现在的问题，即达尼尔是至正义的，这不仅由你的证言证明——虽然对我来说在你我辩论中这或许足够——而且由天主圣神的证言，也藉厄则克耳的口说，他提到三个最杰出正义的人：\Person{诺厄}、\Person{达尼尔}和\Person{约伯}，他说，只有他们能从天主某过分的怒火中得救，这怒火正悬在其余所有人身上。\BibleRef{则 14:14}因此，这至高的正义之人，三位在正义中显赫者之一，祈祷说：\BibleQuote{我正说话祈祷，明认我的罪和我人民以色列的罪，在我上主我天主面前呈上我的恳求。}\BibleRef{达 9:20}而你却说你是无罪的，因为你是个司铎；如果人民犯罪，你为他们祈祷；但如果你犯罪，谁为你祈祷呢？显然，这种狂妄的不敬虔诚表明你不配作那位司铎的中保，先知这话指的就是那位司铎，是你不理解的。因为现在，为了没人问为什么这么说，我将按天主恩宠所能的予以解释。天主藉祂的先知预备人的心，渴望这样一位司铎，无人应为祂祈祷。祂自己在初民和初殿的时代被预像，其中一切都是我们的预像。因此，大司祭独自进入至圣所，为人民求情，人民不与司祭一同进入那内殿；正如我们的大司祭进入了天上的隐秘处，那更真实的至圣所，而我们这些祂为之祈祷的人仍在这里。先知正是针对此说：\Quote{如果人民犯罪，司铎要为他们祈祷；但如果司铎犯罪，谁为他祈祷呢？}因此要寻求这样一位司铎，祂不能犯罪，也无需人为祂祈祷。因此，人民为宗徒们祈祷；\BibleRef{格后 1:11}但为那位宗徒的师傅和主，司铎，则无人祈祷。听若望明认此事并说：\BibleQuote{我的孩子们，我写这些给你们，是要你们不犯罪。但若有人犯罪，我们在父那里有一位辩护者，义人耶稣基督，祂为我们的罪作了赎罪祭。}\BibleRef{若一 2:1} \Quote{\Emph{我们}有，}他说；\Quote{为\Emph{我们的}罪。}我求你，学习谦卑，好使你不跌倒，或更好是，好使你及时再起来。因为你若没有跌倒，就不会说出这样的话了。\par

% structure: p0350 source=h2 target=subsection reason=relative-heading-rank; base=h2

\subsection{第107章}

242. \Person{Petilianus}说：\Quote{并且，为使没有一个平信徒可以声称自己无罪，他们都受这禁令的约束：\InnerQuote{不要在他人的罪上有份。}}\par

243. \Person{奥古斯定}答覆说：你\Emph{\OriginalTerm{大错特错}{toto cælo}}，如俗语所说，只因你的骄傲；而因你的谦卑，你又不愿意与全世界共融。因为，首先，这话不是对一个平信徒说的；其次，你完全不知道这话是在什么意义上说的。宗徒写信给\Person{弟茂德}，将这一警告只给\Person{弟茂德}本人，他曾在另一处对他说：\BibleQuote{不要忽略你身上的恩赐，那是藉着预言，并藉着长老团的覆手所赐予你的。}\BibleRef{弟前 4:14}许多其他证据也清楚表明他并非平信徒。但当他论到\BibleQuote{不要参与别人的罪过}\BibleRef{弟前 5:22}时，他的意思是：不要自愿或同意参与。因此他立刻补充指示他要如何服从这命令，说：\Quote{你要保持自己纯洁。}因为连\Person{保禄}自己也没有参与别人的罪过，虽然他在身体的合一中容忍了假弟兄，为他们叹息；同样，在他之前的宗徒们也没有参与\Person{犹达斯}的偷窃和罪行，虽然他们曾与那已出卖自己主并被主指为叛徒的\Person{犹达斯}一起领受了圣体。\par

% structure: p0353 source=h2 target=subsection reason=relative-heading-rank; base=h2

\subsection{第一〇八章}

244. \Person{培提利安}说：\Quote{宗徒又藉这句话，将那些在作恶的意识中有分的人归入同一类。他说：\InnerQuote{不但行这样事的人，连那些同意行这样事的人，都该死。}}\BibleRef{罗 1:32}\par

245. \Person{奥古斯定}答覆说：我不在乎你是以何种方式使用这些话，它们是正确的。而这正是天主教会教导的实质：在那些因喜乐而同意的人与那些虽不喜欢却容忍的人之间，有着极大的区别。前者使自己成为糠秕，因为他们追随糠秕的荒芜；后者是麦粒。让他们等待那手持簸箕的基督，好使他们与糠秕分开。\par

% structure: p0356 source=h2 target=subsection reason=relative-heading-rank; base=h2

\subsection{第一〇九章}

246. \Person{培提利安}说：\Quote{因此，所有的人们，来教会吧，你们要逃避与\Emph{\OriginalTerm{背信者}{traditors}}为伍，免得你们也和他们一同灭亡。为使你们更容易知道，当他们自己有罪时，却对我们的信德怀有极好的看法，让我告诉你们：我给他们那些受玷污的人付洗；而他们，尽管天主永远不会给他们这样的机会，却接纳那些因洗礼而属于我的人---如果他们承认我们的洗礼有任何缺陷，他们肯定不会这样做。连我们亵渎的仇敌都惧怕毁灭它，由此可见我们所给的是多么神圣。}\par

247. 奥斯定回答说：对于这一错误，我在这部著作及其他地方已多有论述。但既然你以为在这句话中你的虚妄见解得到了如此有力的印证，以至于你认为理应以这些话来结束你的信函，以便它们能在读者心中保留更新鲜的印象，我认为宜作简短回答。我们承认在异端者身上，那不属于异端者而属于基督的圣洗圣事，同样存在于奸淫者、不洁者或作娈童者、拜偶像者、下毒者、怀恨者、好争辩者、轻信者、骄傲者、煽动叛乱者、嫉妒者、酗酒者、狂欢者身上；在这些人身上，我们认可那不属于他们而属于基督的圣洗圣事有效。因为像这样的人——异端者也包括在内——正如宗徒所说，没有一个将继承天国。\BibleRef{迦 5:19}他们不应仅仅因为在物质上参与圣事就被视为在基督的身体——即教会——之内。因为圣事本身是圣的，即便在这样的人身上，它们也将成为他们更重定罪的理由，因为他们不配地领受并参与。但这些人本身并不在教会的架构之内，这教会通过其成员与基督的连接和接触，在天主的增长中成长。因为那教会是建立在盘石上的，正如主所说：\BibleQuote{在这盘石上，我要建立我的教会。} \BibleRef{玛 16:18}但他们却是建立在沙土上，正如同一主所说：\BibleQuote{凡听了我的这些话而不实行的，就好像一个愚昧人，把自己的房屋建在沙土上。} \BibleRef{玛 7:26}但为免你以为那在盘石上的教会只在地球的一部分，而不延伸到最远的边界，请听她在圣咏中，在旅途中遭遇邪恶时发出的哀叹之声。因为她说道：\BibleQuote{我从地极呼求你，当我心困苦时，你把我高举在盘石上；你带领了我，你，我的希望，成了对抗敌人面的勇气之塔。} \BibleRef{咏 61:2-3}请看她是如何从地极呼求。因此她并非仅在非洲，也不仅是在那些从非洲派遣一位主教到罗马的少数蒙滕塞斯那里，以及到西班牙某位妇人的家中。请看她是如何被高举在盘石上。因此，所有建造在沙土上——即听了基督的话却不实行——的人，即使在我们中间和你们中间他们拥有并传授圣洗圣事，都不应被认为是属于她的。请看她的希望是在天主父、子及圣神内——不在伯多禄或保禄内，更不在多纳图斯或培提利安内。因此，我们害怕摧毁的，不是你们的，而是基督的；它自身是圣的，即便在亵渎者手中也是如此。因为我们不能接纳那些从你们那里来的人，除非我们摧毁他们身上一切属于你们的东西。因为我们摧毁的是逃兵的背信，而非君王的印记。因此，请你自己思量并废除你所说的：你说：\Quote{我，为他们的污染者施洗；他们，虽然愿天主永不允许他们有这样的机会，却接纳那些因我的圣洗而成为我的人。}因为你并非为已受感染的人施洗，而是重新为他们施洗，以致用你错误的欺诈感染他们。但我们不接纳那些因你的圣洗而成为你的人；而是摧毁那使他们成为你的错误，并接纳基督的圣洗，他们正是因这圣洗而受洗。因此，你引入\Quote{虽然愿天主永不允许他们有这样的机会}这句话绝非无意义。因为你说：\Quote{他们，虽然愿天主永不允许他们有这样的机会，却接纳那些因我的圣洗而成为我的人。}当你因为害怕我们接纳你的追随者而欲表达\Quote{愿天主永不给让他们接纳属我之人的机会}之意时，我猜想，你是在不知其义的情况下说：\Quote{愿天主永不让那些人真正成为我的，以致你们接纳他们。}因为我们祈求那些此刻来到天主教会的人，不要真的成为你的。他们来并不是为了因圣洗而成为我们的，而是为了通过与我们共融，并因他们的圣洗，与我们一同属于基督。\par

\SourceRecord{Translated by J.R. King and revised by Chester D. Hartranft.；Nicene and Post-Nicene Fathers, First Series；Vol. 4.；Edited by Philip Schaff.；Buffalo, NY: Christian Literature Publishing Co.,；1887.；<http://www.newadvent.org/fathers/14092.htm>.}

% structure: p0001 source=h1 target=section reason=page-title

\section{答多纳徒派佩提利安（卷三）}

\Emph{在本书中，奥斯定驳斥了佩提利安在看到奥斯定早期著作的第一卷后写给奥斯定的第二封信。这封信充满了激烈的言辞；奥斯定更多地是表明佩提利安的论证有缺陷且不切题，而不是提出论据支持自己的陈述。}\par

% structure: p0003 source=h2 target=subsection reason=relative-heading-rank; base=h2

\subsection{第一章}

1. 佩提利安，我能阅读，我已读了你写的信，在信中你足够清楚地表明，在支持多纳徒派反对天主教会时，你既未能说出任何中肯的话，也未被允许保持沉默。当你读到我对我当时手中仅有你那封信的部分内容所作的尽可能简洁明了的答复时，你经历了何等剧烈的情绪波动，何等汹涌的情感在你心中翻腾！因为你看到我们所坚持和维护的真理被如此有力的论证所证实，并被如此充足的光照所阐明，以至于你找不到任何可以反驳的话，从而驳倒我们所提出的指控。你也注意到，许多读过这封信的人的目光都集中在你身上，因为他们想知道你会说什么，你会做什么，你如何摆脱困境，你如何从天主圣言包围你的困境中脱身。然后你，本应轻视愚昧之人的意见，转而采纳正确和真诚的观点，却宁愿做圣经对像你这样的人所预言的事：\BibleQuote{你爱恶胜过爱善，爱谎言胜过说正义。}就仿佛我反过来也愿意以辱骂回报你的辱骂；那样的话，我们岂不成了两个恶言相向的人，以致那些读我们文字的人要么厌恶地把我们丢开以保持尊严，要么贪婪地阅读我们写的东西以满足他们的恶意？就我而言，既然我回答每一个人，无论是书面还是口头，即使受到侮辱性的指责，我也用主放在我口中的言语来回答，为了听众或读者的利益而克制和粉碎空洞愤怒的刺痛，我并不试图通过辱骂来证明自己优于对手，而是通过使他认识错误来成为他得救的泉源。\par

2. 因为如果那些考虑到你写的内容的人有任何感受，那么，在天主教会和多纳徒派之间悬而未决的问题上，你放弃一个可以说是公众利益的问题，出于个人敌意，用恶意的辱骂攻击某个人的生平，好像那个人就是辩论的焦点，这对你究竟有何益处？难道你如此低估——我不说基督徒——而是全人类，以至于认为你的作品不会落入一些明智之人的手中？他们会放下像我们这样的个人的所有想法，反而探究我们之间争论的问题，注意的不是我们是谁、我们是什么，而是我们能为维护真理或反对错误提出什么。你应该尊重这些人的判断，你应该提防他们的指责，免得他们认为你除了在自己面前设定一个你可以用任何手段辱骂的人之外，无话可说。但通过一些轻率和愚昧之人——他们热切地倾听任何有学识的争论者的争吵——的行为可以看出，当他们注意到你肆意辱骂的口才时，他们并不察觉你被何等真理所驳斥。同时，我认为你的部分目的是，我可能被迫为自己辩护，从而放弃我所承担的这同一事业；这样，当人们的注意力转向那些并非在进行辩论而是在争吵的对手的言语时，你所害怕的真理——它应被揭示并为人所知——可能被遮蔽。因此，我反对这样的计谋该做什么，除了严格坚持我的主题，宁愿忽略我自己的辩护，同时祈祷任何个人的诽谤都不会使我偏离它？我要颂扬起我的天主之家，我热爱它的荣耀，用忠实仆人声音的贡品，而我自己则谦卑下来，视为无有。\BibleQuote{我宁愿在我天主的家中做看门人，也不愿住在异端者的帐幕中。}因此，佩提利安，我将暂时把话题从你身上转向那些你试图通过辱骂从我这里夺走的人，好像我的努力是让人们归向我，而不是与我一同归向天主。\par

% structure: p0006 source=h2 target=subsection reason=relative-heading-rank; base=h2

\subsection{第二章}

3. 所以，凡读过他谩骂的人，请听\Person{白提利安}怀着更多愤怒而非思虑对我倾泻了什么。首先，我要用宗徒的话对你们说——这些话肯定是真实的，无论我自己是什么人：\BibleQuote{人应当这样看我们，如同基督的仆人和天主奥秘的管理人。此外，管理人要求的是要发现他忠实。至于我，被你们论断，或被人的审判论断，都是极小的事；是的，我也不论断自己。}关于紧接的话，虽然我不敢把\BibleQuote{因为我对自己一无所知}这句话用在自己身上，但我敢在天主面前肯定地说：自从我在基督内受洗以来，我对自己良心中没有\Person{白提利安}所加在我生活中的任何那些指控；\BibleQuote{然而我因此不被成义，而是审判我的乃是主。所以，时候未到，什么都不要论断，直等主来，祂要显明暗中的事，也要显明心中的意念；那时各人要从天主得着称赞。弟兄们，我把这些事转用到自己身上，为叫你们在我们身上学到，不要过于圣经所记，免得你们有人自高自大，拥护这个反对那个。}\BibleRef{格前 4:1}\BibleQuote{所以，谁也不要夸耀人；因为一切都是你们的；你们是基督的；基督是天主的。}我再次说：\BibleQuote{谁也不要夸耀人}；不，我多次重复说：\BibleQuote{谁也不要夸耀人。}如果你们在我们身上看出任何值得称赞的事，请都将它归于祂的赞美，因为各样美好的恩赐和各样完美的恩赐都是从祂而来；它\BibleQuote{是从上而来的，从光明之父那里降下来的，在祂没有改变，也没有转动的阴影。}\BibleRef{雅 1:17}因为我们有什么不是领受的呢？若是领受了，就不要夸耀，好像不是领受的。\BibleRef{格前 4:7}在这一切你们知道在我们身上是好的事上，你们要效法我们——无论如何，如果我们属于基督；\BibleRef{格前 4:16}但另一方面，如果你们怀疑、相信或看见我们有任何邪恶，请紧记主的这句话，你们可以安全地决定不因人的恶行而离开祂的教会：凡他们所吩咐你们的，你们都要谨守遵行；但不要效法他们的行为。\BibleRef{玛 23:3}因为现在不是我在你们面前为自己辩护的时候，我已经承担起抛开自我的一切考虑，向你们推荐对你们救恩有益的事，就是没有人应该夸耀人。因为\BibleQuote{依靠人的人，是可咒骂的。}\BibleRef{耶 17:5}只要坚持遵守主和祂宗徒的这个诫命，我所服务的事业就必得胜，即使我自己——如同我的敌人所想的那样——在自己的事业中衰弱且受压迫。因为如果你们最坚决地持守我全力敦促你们的——即凡依靠人的人都被咒骂，以致没有人夸耀人——那么，你们绝不要因那些现在被骄傲的风吹散或将在最终扬净中分离的糠秕而离开主的打谷场；\BibleRef{玛 3:12}也不要因那些作卑贱器皿的而逃离大户人家；\BibleRef{弟后 2:20}也不要因网破而离开渔网，因为其中有坏鱼，它们将在岸上被分开；\BibleRef{玛 13:47}也不要因一些在善牧分开羊群时将被放在左边的山羊而离开合一的美好牧场；\BibleRef{玛 25:32}也不要因莠子的混杂而以不敬的分裂离开那好麦子的群体——其源头是那颗死去而倍增的麦子，它在全世界一同生长直到收获。因为田地是世界——不只是非洲；收获是世界的终末，\BibleRef{玛 13:24}——不是\Person{多纳图}的时代。\par

% structure: p0008 source=h2 target=subsection reason=relative-heading-rank; base=h2

\subsection{第三章}

4. 你无疑认得出这些福音的比喻。我们也不能认为它们是为其他目的而给出的，除非是为了使人不应在人身上自夸，不应彼此自大或彼此分裂，说：\Quote{我是属\Person{保禄}的}，当然\Person{保禄}并没有为你们被钉十字架，你们也不是因\Person{保禄}的名受洗，更不是因\Person{塞西利安}或我们中任何一人的名受洗，\BibleRef{格前 1:12}为的是叫你学会，只要糠秕与麦子一同被碾压，只要坏鱼与好鱼一起在主的网内游动，直到分离的时刻来到，你就更应因好人的缘故忍受坏人的混杂，而不应因坏人的缘故违反对好人的兄弟之爱。因为这混杂不是永远的，而只是暂时的；也不是精神的，而是肉体的。在这件事上，天使不会犯错，当他们把坏从好中收集出来，丢进火炉里时。主认识属于祂的人。如果一个人不能在肉体上离开那些作恶的人，只要时间持续，那么至少，凡称呼基督之名的，应当离开不义本身。\BibleRef{弟后 2:19}因为同时，他可以在生命、道德、心灵和意志上与恶人分离，并在同样的方面离开他的交往；这样的分离应当始终维持。但身体的分离要等到时间的终结，以信德、耐心和勇敢期待。鉴于这种期待，经上说：\Quote{你当等候主，要刚强壮胆，祂必坚固你的心；我再说，你要等候主。}因为最大的忍耐之棕榈枝属于那些在偷偷混入的假弟兄中间，寻求自己的事而非\Person{耶稣基督}的事的人，他们却表明自己不愿以任何骚动或鲁莽的分裂，扰乱那不属于他们而是属于\Person{耶稣基督}的爱，也不愿以骄傲精神所激发的邪恶争竞，破坏主网的一体——网中聚集了各样鱼类，直到它被拖到岸边，即直到时间的终结；尽管每个人可能自以为是，其实一无所是，以致欺骗自己，并希望在自己的朋友或自己的判断中为基督徒百姓的分裂找到充足理由，这些人声称他们确知某些恶人不配参与基督徒宗教圣事的共融；然而，无论他们所知关于那些人的什么，他们无法说服那按预言传遍万国的普世教会相信他们的话。当他们拒绝与这些他们知其品格的人共融时，他们就离开了教会的一体；其实如果他们真有那含忍万事的爱德，他们本应自己承担他们所知在一国之内的事，以免他们与那些他们无法向万国灌输异己邪恶教训的好人分离。因此，即使不讨论他们被最有力的证据定罪为对无辜者散布诽谤的案件，更有理由相信他们捏造了交出圣书的虚假指控，因为他们被发现犯了在教会中进行邪恶分裂的更为严重的罪行。因为即使他们抛出的交出圣书的指控是真的，他们也绝不应因他们所熟悉的这些考量而放弃基督徒的团体——这些考量子经上被推荐到地极——而这些人却表明他们并不熟悉这些考量。\par

% structure: p0010 source=h2 target=subsection reason=relative-heading-rank; base=h2

\subsection{第四章}

5. 因此，请不要理解为我在主张废弃教会纪律，允许每个人随心所欲，不受任何约束，没有某种治疗的责罚、一种引人敬畏的宽仁、爱的严厉。因为那时宗徒的训诫将成何体统？\BibleQuote{劝戒那些不守规矩的人，安慰怯懦的人，扶持软弱的人，容忍众人；要谨慎，谁也不要以恶报恶？} \BibleRef{得前 5:14} 无论如何，当他加上最后这些话：\BibleQuote{要谨慎，谁也不要以恶报恶}时，他已充分表明，当一个人责罚那些不守规矩的人时，即使因不守规矩的过失而施加责罚的惩罚，也不是以恶报恶。因此，责罚的惩罚并不是恶，尽管过失是恶。这确实是手术刀，而非敌人造成伤口的利刃。这类事在教会内发生，教会内的温和精神以对天主的热情燃烧，以免那位许配给一个丈夫、即基督的贞洁童女，在她的任何肢体中被败坏，失去在基督内的纯朴，如同厄娃被蛇的诡计所欺骗。\BibleRef{格后 11:2} 然而，家主的仆人们切不可忘记他们主的诫命，被义怒之火如此激动而反对大量莠子，以致在时机未到之前想要收集成捆，而将麦子一同拔除。这些人必被定为这罪，即使他们能证明对所指控的\OriginalTerm{移交者}{traditors}的指控是真实的；因为他们以不敬的傲慢精神分离，不仅离开了他们声称要避免的恶人，而且离开了全世界万国中的善人与信者——他们无法证明自己所宣称知道的事是真的；他们又带着许多在他们之下稍有权威、且不够聪明的人一同走向毁灭，这些人不明白不可为别人的罪而离弃分散在全世界的教会合一。因此，即使他们知道自己在对某些邻人提出真实指控，但这样一位软弱的弟兄——基督为他而死——却因他们的知识而丧亡；\BibleRef{格前 8:11} 同时，因别人的罪而被绊倒，在自己身上毁坏了与善弟兄们共有的和平祝福——这些善弟兄中有的从未听过这类指控，有的因未经讨论和证明而不愿轻信，有的以谦卑的和平精神将无论什么指控都留给教会的审判者处理，整个事件已提交给他们，远在海外。\par

% structure: p0012 source=h2 target=subsection reason=relative-heading-rank; base=h2

\subsection{第五章}

6. 因此，我们同一公教慈母的圣洁子孙，要以你们所能的全部警惕，在适当的顺服于主之下，提防这样的罪行和错误的榜样。无论他闪耀着何等学问与声望的光辉，无论他如何自夸为宝贵的宝石——谁试图带领你们跟随他——都要记住，那独爱其夫的勇敢女子，在《箴言》末章中由圣经描绘给我们的她，比任何宝石更珍贵。没有人可以说：我要跟随此人，因为正是他使我成为基督徒；或者说：我要跟随此人，因为正是他给我付了洗。因为\BibleQuote{可见，栽种的不算什么，浇灌的也不算什麼，只在那使之生长的天主。} \BibleRef{格前 3:7} 又\BibleQuote{天主是爱，那存留在爱内的，就存留在天主内，天主也存留在他内。} \BibleRef{若一 4:16} 任何人凡宣讲基督的名，并处理或施行基督圣事的，都不应被跟随去反对基督的合一。\BibleQuote{各人该考验自己的行为，这样，只在自己身上，而不在别人身上，才有可夸耀处。因为各人要背负自己的担子，} \BibleRef{迦 6:4}——这担子，即交账的担子；因为\BibleQuote{我们每人都要向天主交自己的账。所以我们不可再彼此判断了。} \BibleRef{罗 14:12} 至于彼此相爱的担子，\BibleQuote{你们要彼此担待，这样就满全了基督的法律。人若自认为了不起，其实算不得什麼，只是欺骗自己。} \BibleRef{迦 6:2} 因此，我们要\BibleQuote{用爱心互相担待，尽力以和平的联系，保持心神的合一，} \BibleRef{弗 4:2} 因为凡在和平之外聚集的，不是与基督聚集，而是\BibleQuote{那不与我聚集的，就是分散的。} \BibleRef{玛 12:30}\par

% structure: p0014 source=h2 target=subsection reason=relative-heading-rank; base=h2

\subsection{第六章}

7. 此外，无论关于基督，或关于祂的教会，或任何其他与你们的信仰和生活相关的事，更不必说我们自己，我们绝无法与祂相比，祂曾说：\Quote{尽管是我们，}但无论如何，正如祂接着说的：\Quote{尽管从天上来的天使，传给你们别的福音，与你们所领受的合法福音圣经不同，就让他被诅咒。}\BibleRef{迦 1:8} 我们在与你们以及所有我们渴望在基督内赢得的人交往中，秉持这一行动原则，此外，当我们宣讲我们在天主的书信中读到、并看到按应许在全世界万民中应验的圣教会时，我们从那些我们渴望吸引进入祂最平安怀抱的人那里，得到的不是感谢，而是憎恨的火焰；仿佛我们将他们牢牢束缚在那个他们无法为自己找到任何辩护的党派中；或者仿佛我们早在以前就命令先知和宗徒，在他们的书中不要插入任何证明多纳徒派是基督的教会的证据。事实上，亲爱的弟兄们，当我们听到那些因我们宣讲真理的雄辩、驳斥错误的虚妄而得罪的人对我们提出的诬告时，如你们所知，我们得到了极大的安慰。因为，在他们指控我的事情上，如果我的良心在天主面前没有控告我，那里没有凡人的眼目能及，那么我不仅不应该沮丧，反而应该欢喜快乐，因为我在天上的赏报是丰厚的。\BibleRef{玛 5:12} 事实上，我应该考虑的，不是我所听到的话有多么苦涩，而是多么虚假，以及那为了保卫祂的名而使我遭受这些的主是多么真实，对祂说：\Quote{你的圣名好似倾泻的香膏。}\BibleRef{歌 1:3} 这香膏确实在所有民族中散发芬芳，尽管那些诬蔑我们的人试图将它的香气局限在阿非利加的一个角落。因此，我们何必介意那些诋毁基督光荣的人的辱骂呢？他们的党派和分裂，为早在如此久之前就预言了祂升天、祂的名如香膏倾泻而出的这事而恼怒：\Quote{天主，愿你被尊崇于高天，愿你的荣耀遍及大地！}\par

% structure: p0016 source=h2 target=subsection reason=relative-heading-rank; base=h2

\subsection{第七章}

8. 当我们带着天主的见证，以如此和类似的效果对抗人的空谈时，我们被迫忍受基督光荣的仇敌的尖刻侮辱。让他们随意说吧，而祂却鼓励我们说：\Quote{为义而受迫害的人是有福的，因为天国是他们的。几时人为了我而辱骂迫害你们，捏造一切坏话毁谤你们，你们是有福的。}祂起初说的\Quote{为义}，后来在\Quote{为了我}这句话中重复了，因为祂\Quote{成了我们的智慧、正义、圣化者和救赎者，正如经上所载：凡夸耀的，应因主而夸耀。}\BibleRef{格前 1:30} 当祂说：\Quote{你们欢喜踊跃吧！因为你们在天上的赏报是丰厚的}\BibleRef{玛 5:10}，如果我在良心中持守那所说的\Quote{为义}和\Quote{为了我}，那么任何故意诋毁我名誉的人，都是违心地增加我的赏报。因为祂不仅用祂的话教导我，还用祂的榜样坚固我。跟随圣经的信仰，你会发现基督从死者中复活，升天，坐在圣父的右边。跟随祂仇敌的指控，你很快就会相信祂被门徒从坟墓中偷走了。那么，当我们尽其所能捍卫祂的房屋时，我们凭什么期望从祂的仇敌那里得到其他待遇呢？\Quote{如果人们称家主为贝耳则步，何况他的家人呢？}\BibleRef{玛 10:25} 因此，如果我们受苦，我们也将与祂一同为王。但如果不仅是控告者的愤怒击打耳朵，而且控告的真实性刺痛良心，那么即使全世界用不断的赞美高举我，于我何益？所以，称赞者的颂词不能治愈有罪的良心，辱骂者的侮辱也不能伤害好的良心。然而，你们在主内的希望并不落空，即使我们偶然在暗地里成为敌人所希望我们成为的那种人；因为你们并没有寄希望于我们，也从未从我们这里听过任何这样的道理。所以，无论我们是什么，你们是安全的，因为你们学会了说：\Quote{我信赖了主，因此我不会动摇}；以及\Quote{我依赖天主，不怕人能对我做什么。}对于那些试图引诱你们到骄傲之人的世俗高处的人，你们知道如何回答：\Quote{我信赖主，你怎么能对我的灵魂说：像鸟一样飞往你的山上？}\par

% structure: p0018 source=h2 target=subsection reason=relative-heading-rank; base=h2

\subsection{第八章}

9. 并非只有你们是安全的，无论我们如何，因为你们满足于那在我们内的基督的真理，即通过我们所宣讲的，以及普世各处所宣讲的，并且因你们乐意聆听我们卑微的舌唇所传讲的真理，你们也善意地看待我们——因为你们的望德是在于祂，我们向你们宣讲祂的慈爱，这慈爱延及你们——但此外，你们所有从我们的职务领受圣洗圣事的人，也可以同样在这安全中喜乐，因为你们受洗，不是归于我们，而是归于基督。所以你们并非穿上我们，而是穿上基督；我也没有问你们是否皈依于我，而是归于永生的天主；也没有问你们是否信了我，而是信了圣父、圣子和圣神。但若你们以诚实的心回答了我的问题，你们便得救了，不是除去了肉体的污秽，而是藉着向天主存有的良心之回答；\BibleRef{伯前 3:21} 不是靠同等之仆，而是靠主；不是靠传令官，而是靠审判者。因为培提利安轻率地说：\Quote{施予者的良心}——或如他所加的\Quote{以圣德施予者的良心，是我们所寻求来洗涤领受者的良心的}，这并不正确。因为当所给的东西是来自天主的，即使是通过不圣洁的良心，它也是以圣德给予的。而且，领受者根本无法分辨那良心是否圣洁；但他可以清楚分辨所给予的东西。那为永远圣洁者所知的，无论经谁的手领受，都是完全安全的。除非从梅瑟的宝座上所说的话必须是圣洁的，真理者决不会说：\Quote{凡他们所吩咐你们遵守的，你们要遵守并遵行。}但若说出圣言的人本身是圣洁的，祂就不会说：\Quote{不要照他们的行为去做，因为他们只说不做。}\BibleRef{玛 23:2} 因为确实，人不能从荆棘上摘葡萄，因为葡萄绝不出自荆棘的根；但当葡萄枝缠绕在荆棘篱笆上时，悬挂其上的果实并不可怕，而是避开荆棘，摘下葡萄。\par

% structure: p0020 source=h2 target=subsection reason=relative-heading-rank; base=h2

\subsection{第九章}

10. 因此，正如我此前多次说过，并愿向你们说明：无论我们如何，你们是安全的，你们以天主为父，以祂的教会为母。因为纵然山羊与绵羊一同牧放，他们却不得站在右边；纵然糠秕与麦子一同碾压，它却不得收入仓内；纵然坏鱼与好鱼一同在主网中游动，它们却不得收进器皿。愿没有人夸耀于好人；愿没有人躲避即使是坏人中的天主的美善恩赐。\par

% structure: p0022 source=h2 target=subsection reason=relative-heading-rank; base=h2

\subsection{第十章}

11. 愿这些事，我亲爱的公教会弟兄们，就当前之事对你们已足够；若你们以公教之爱持守此理，只要你们是同一牧者的同一稳固羊群，我便不太在意任何仇敌对我——羊群中的同伴，或至少是看门犬——的辱骂，只要他迫使我更多为你们辩护，而非为己。然而，若为事业所需，我应自我辩护，我当极其简短而轻易地为之，自愿与众人一同谴责和指证我领受基督圣洗之前整个人生期间的一切——关乎我的恶情与错谬——以免在辩护那段时期时，似乎寻求自己的光荣，而不是那藉恩宠甚至将我从自身解救出来的祂的光荣。因此，当我听到那段人生受辱骂时，无论辱骂者出于何种精神，我并非如此忘恩以致忧伤。无论他指责我的何种恶习，我皆同等赞美他如同我的医生。我又何必自扰，为辩护那已过和已弃的恶事呢？关于这些，培提利安已说了许多假话，却仍有更多真话未提。但关于我领洗后的人生，你们认识我，不必我来讲述那些可为人知晓的事；但不认识我的人不应如此不公，竟愿相信培提利安说我的坏话，而不相信你们。因为若不应相信朋友的赞美，也不应相信仇敌的诽谤。因此，还有那些隐藏在人内的事，只有良心可作证，却不能在人间见证。于此，培提利安说我是摩尼教徒，这是论他人的良心；我则论及自己的良心，声明我不是摩尼教徒。你们选择我们中你们更愿相信谁。尽管如此，既然我这简短而轻易的辩护亦非必要，因为所争辩的不是任何个体的功过——无论他为何人——而是整个教会的真理，我还有更多话要对你们中那些属于多纳图派的人说，你们读到培提利安写给我的恶言，若我不关心你们救恩的失落，我不会听到这些；但那样我便缺失了基督之爱的慈怀。\par

% structure: p0024 source=h2 target=subsection reason=relative-heading-rank; base=h2

\subsection{第十一章}

12. 如此，当我收集从主的打谷场上扬出的谷粒时，连同尘土和糠秕，被反弹的灰尘所伤，又有什么奇怪呢？或者，当我殷勤寻找我主迷失的羊时，被荆棘般尖刻的舌头所刺，又有什么奇怪呢？我恳求你们，暂时放下一切党派情绪，以某种程度的公平在\Person{培提利安}和我之间作出判断。我希望你们了解教会的事业；而他却希望你们熟悉我的事业。除了因为他不敢叫你们不信我的证人（我在教会的事业中不断援引他们——他们是先知和宗徒，以及基督自己，先知和宗徒的主）之外，还有什么其他原因呢？而你们却轻易相信他关于我所说的任何话——一个人反对另一个人，而且他是你们同党的人反对一个你们陌生的人？如果我引用任何关于我生活的证人，无论他说的话多么重要，他们也不会相信，而\Person{培提利安}会迅速说服你们；尤其是当有人为我辩护时。既然他是多纳图派的敌人，因此他也将继续被看作你们的敌人。因此\Person{培提利安}至高无上。每当他对我施加任何辱骂，无论性质如何，你们都鼓掌欢呼。他找到了这场能使他一举获胜的案子，但只有在你们作为审判官的情况下。他不再寻求证据或证人；因为他只需在言语中证明这一点：他对一个你们最深恶痛绝的人大肆辱骂。因为当大量明确的天主圣经的证词被引用支持天主教会时，他在你们的悲伤中保持沉默；他为自己选择了一个话题，可以在你们的掌声中发言；虽然实际上被打败了，却假装坚持不让，说关于我的这样甚至更坏的话。对于我正在辩护的事业来说，无论我本人被发现是什么样的人，我所代表的教会仍然不可征服，这对我来说就足够了。\par

% structure: p0026 source=h2 target=subsection reason=relative-heading-rank; base=h2

\subsection{第十二章。}

13. 因为我是基督打谷场上的人：如果是坏人，就是糠秕的一部分；如果是好人，就是谷粒的一部分。这打谷场的簸箕不是\Person{培提利安}的舌头；因此，无论他说了什么恶言，即使是真实地针对这打谷场的糠秕，也丝毫不损害其谷粒。但如果他对谷粒本身投掷任何辱骂或诽谤，谷粒的信德在世上受到考验，而它的赏报在天上增加。因为当人是主的圣仆，为天主圣洁地争战——不是对抗\Person{培提利安}或任何像他一样血肉之人，而是对抗率领者、掌权者、这个世界黑暗的统治者\BibleRef{弗 6:12}——所有真理的敌人都是如此，而我愿意对他们说：\BibleQuote{你们从前是黑暗，但现在在主内却是光明}\BibleRef{弗 5:8}——我说，当天主的仆人进行这样的战争时，那么所有由敌人发出的诽谤辱骂（它们在恶毒和轻信的人中造成恶名）就是左手的武器：就连魔鬼也是被这样的武器击败的。因为当我们受好名声的考验时，是否抵制使我们骄傲的抬高；又受坏名声的考验时，是否爱那些编造坏名声来攻击我们的仇敌，这样我们就用左右两手的正义武器战胜魔鬼。因为当宗徒说了\Quote{藉着左右两手的正义武器}后，他紧接着说，似乎是对这些术语的解释：\BibleQuote{藉着荣誉和羞辱，藉着恶名和美名}\BibleRef{格后 6:7}等等——把荣誉和美名算作右手的武器，羞辱和恶名算作左手的武器。\par

% structure: p0028 source=h2 target=subsection reason=relative-heading-rank; base=h2

\subsection{第十三章。}

14. 因此，如果我是主的仆人，是一名无可指摘的士兵，那么无论\Person{培提利安}以何等口才站出来辱骂我，我难道应该因为他被指派给我作为最精良的左臂武器匠而烦恼吗？我必须尽可能熟练地使用这武器为我主争战，并用它打击那个与我进行无形战斗的敌人，他以最邪恶和古老的诡计竭尽全力使我对\Person{培提利安}怀恨，从而无法履行基督所赐的命令，即我们要\BibleQuote{爱你们的仇敌}\BibleRef{路 6:35}。但愿我因那爱我、为我舍弃自己的主的仁慈而得救，正如他挂在十字架上时说：\BibleQuote{父啊，宽恕他们，因为他们不知道他们做的是什么}\BibleRef{路 23:34}，并这样教导我说关于\Person{培提利安}和所有像我一样的其他仇敌：\BibleQuote{父啊，宽恕他们，因为他们不知道他们做的是什么}。\par

% structure: p0030 source=h2 target=subsection reason=relative-heading-rank; base=h2

\subsection{第十四章。}

15. 此外，如果我已从你们那里，按照我恳切的努力，使你们放下心中一切党派偏见，成为\Person{培提利安}和我之间的公正审判官，我将向你们证明他没有回应我所写的内容，使你们明白他因缺乏真理而被迫放弃争论，也看到他对那个使他无言以对的人竟允许自己说出什么样的辱骂。然而我将要说的话以如此明显的清晰度彰显自身，以至于即使你们的内心因党派偏见和个人仇恨而与我疏远，但只要你们读一读双方写下的文字，你们就不能不在内心深处彼此承认我说的是真话。\par

16. 当他回覆其著作的前面部分时（当时我手中只有那部分），我没有理会他那空泛而亵渎的辱骂。他说：\Quote{让那些人以我们两次重复的圣洗来责备我们吧，他们以圣洗之名，用有罪的洗涤玷污了自己的灵魂；我认为他们是如此污秽，以至于任何污秽都不如他们不洁；他们的命运是，因洁净的颠倒，被他们洗涤的水所玷污。}我认为接下来值得讨论和驳斥的是，他所说的：\Quote{我们所寻求的是施予者的良心，以使领受者的良心借此得以洁净。}我问道，当施予者的良心被玷污，而领受圣事者并不知道时，有何办法可以洁净领受圣洗者的良心？\par

% structure: p0033 source=h2 target=subsection reason=relative-heading-rank; base=h2

\subsection{第十五章}

17. 现在请读他因对我愤愤不平而倾泻的最为冗长的辱骂，看看他是否对我的问题作了任何回答。我问：当施予者的良心被玷污，而领受圣事者并不知道时，有何办法可以洁净领受圣洗者的良心？我恳请你们仔细搜索，检查每一页，计算每一行，斟酌每一个字，筛出每个音节的含义，然后告诉我，如果你们能找到，他在哪里回答了一个问题：当领受者不知道施予者的良心被玷污时，有何办法可以洁净领受者的良心？\par

18. 因为，他添加了一个据他说被我删去的短语，并声称他写了以下文字：\Quote{我们所寻求的是以圣洁之心施予者的良心，以洁净领受者的良心。}这对我提出的问题有何影响呢？为了向你们证明我没有删去它，它的添加丝毫没有妨碍我的探究，也没有弥补他在其中发现的不足。因为在那些话面前，我再次问道，并请你们看他是否作了任何回答：如果\Quote{我们所寻求的是以圣洁之心施予者的良心，以洁净领受者的良心，}那么，当施予者的良心被罪污污染，而领受圣事者并不知道时，有何办法可以洁净领受者的良心？我坚持要求对此作出回答。不要让任何因与事无关的辱骂而受害。如果我们所寻求的是以圣洁之心施予者的良心——请注意我不是说\Quote{施予者的良心，}而是我加上了\Quote{以圣洁之心施予者的}这几个字——那么，如果我们所寻求的是以圣洁之心施予者的良心，当施予者的良心被玷污，而领受圣事者并不知道时，有何办法可以洁净领受圣洗者的良心？\par

% structure: p0036 source=h2 target=subsection reason=relative-heading-rank; base=h2

\subsection{第十六章}

19. 现在让他去吧，让他鼓着肺、胀着喉，指责我不过是个辩证法家。不，让他把辩证法这门科学本身——而不是我——当作谎言的伪造者，召到公众舆论的法庭上，并让他尽其所能用最响亮的法庭喧哗对它张大口。让他对无经验者随意说什么，使得有学问的人被激怒，而无知者被欺骗。让他因我的修辞学用那位控告保禄的演说家\Person{特尔突罗}的名字称呼我\BibleRef{宗 24:1}；并让他给自己冠以辩护者的名字，凭借他自夸以前有过的辩才，并因此自欺地以为自己，或者说曾经是，圣神同名的对象。让他尽我所愿地夸大摩尼教徒的污秽，并试图通过他的狂吠将其转向我。让他引用所有被定罪者的行为，无论我是否知道；并让他用某种全新的特权，将我一个前朋友在那里为了自身辩护而提到我的事实，变成一种推定有罪的恶毒指控。让他读那些他或他的朋友随意放在我信件上的标题，并夸耀他已将我无望地卷入他们的表述之中。当我承认那些关于面包的赞词（全属单纯和欢乐）时，让他用荒谬的毒药和疯狂指控剥夺我的名誉。让他对你们的理解力如此恶劣，以至于认为他可以让人相信，有毒的催情药被给了一个女人，不仅她丈夫知道，而且实际上是在他同谋下给的。那个后来要祝圣我为主教的人，在我还是司铎时就愤怒地写了关于我的事，他可以自由地用来作为对我不利的证据。同一个人曾因对我所做的错事向一个神圣的公议会请求并获得了宽恕，他同样可以自由地忽略这一点——他要么对基督徒的温和与福音的诫命如此无知，要么如此健忘，以致于把一个兄弟完全不知道的事情作为指控提出来，而那个兄弟正谦卑地恳求原谅他。\par

% structure: p0038 source=h2 target=subsection reason=relative-heading-rank; base=h2

\subsection{第十七章}

20. 让他在他那众多但明显空洞的言论中，更进一步地谈论那些他完全无知的事情，或者更确切地说，他是在滥用那些听他讲话的多数人的无知。根据某个妇人的自白，她自称是摩尼教信徒的慕道者，而实际上她已经是天主教会的一名正式成员，让他就他们的圣洗说或写任何他乐意的话——他不知道，或假装不知道，在当时，慕道者这个名称并非给予那些将来某时要受洗的人，而是给予那些也称为听者的人，认为他们无法遵守那些被认为是更高更大的诫命，而那些按他们认为应当以\OriginalTerm{选民}{Elect}之名加以区分和尊敬的人则遵守这些诫命。让他也以惊人的轻率坚持说，若不是他自己被欺骗，就是企图欺骗别人，说我在摩尼教徒中是一名司铎。让他以他看为合适的任何方式，提出并驳斥我\WorkTitle{忏悔录}第三卷中的话，这些话本身，加上我以前和以后所说的许多话，对所有读它们的人都是十分清楚的。最后，让他因我剽窃了他的话而得意，因为我删去了其中两个词，好像把它们恢复过来，胜利就属于他似的。\par

% structure: p0040 source=h2 target=subsection reason=relative-heading-rank; base=h2

\subsection{第十八章。}

21. 诚然，在这一切事情上，正如你通过阅读他的信可以了解或重温的那样，他任由自己舌头的冲动，尽其所愿地使用一切夸耀的许可，但他从未告诉我们，当施予者的良心被罪玷污而他自己却不知道时，在哪里可以找到洁净领受者良心的办法。但是在他所有的喧嚣中，在他所有的喧嚣之后，这喧嚣正如他所想象的那样严重、可怕，我故意地，如常言所说，并且切中要害地，再次提出这个问题：如果我们要看的是那圣洁施予者的良心，那么，对于那些在没有意识到施予者良心被罪玷污的情况下接受圣洗的人，在哪里可以找到洁净的办法？在他整封书信中，我没有发现对此问题的任何回答。\par

% structure: p0042 source=h2 target=subsection reason=relative-heading-rank; base=h2

\subsection{第十九章。}

22. 也许你们中有人会对我说：他针对你所说的一切，其意图是剥夺你的声望，并通过你剥夺那些与你共融之人的声望，好使他们自己以及你试图争取加入你共融的人，都不再认为你有什么重要之处。但是，要判断他是否对你的书信中的话没有作出回答，我们必须根据他提出那些话供考虑的那段文字来审视它们。那么，就让我们这样做吧：就让我们根据那同一段文字来审视他的作品。因此，我跳过那段我试图向读者介绍主题的文字，并忽略他那几句序言式的、与其说是与讨论的问题相关不如说是侮辱性的话，接着我说道：\Quote{他说：\InnerQuote{我们所寻求的是施予者的良心，以洁净领受者的良心。}但是，假如施予者的良心是隐藏看不见的，并且也许被罪玷污了，那么，如果照他所说，\InnerQuote{我们所寻求的是施予者的良心，以洁净领受者的良心}，它又怎能洁净领受者的良心呢？因为如果他说，无论施予者良心中隐藏着多少邪恶，对领受者都无关紧要，也许那种无知可以具有如此大的效力，以至于只要一个人不知道，他就不会因从他领受圣洗之人的良心罪责而受玷污。那么，就让我们承认，只要一个人不知道，邻人的有罪良心并不能玷污他；但是，这难道就能清楚表明它还能进一步洗净他自己的罪责吗？那么，当施予者的良心被污染而领受者却不知道时，那领受圣洗的人该从哪里得到洁净呢？尤其是当他接着说：\InnerQuote{因为从无信者那里领受信德的人，领受的不是信德而是罪责。}}\par

% structure: p0044 source=h2 target=subsection reason=relative-heading-rank; base=h2

\subsection{第二十章。}

23. 佩蒂利安将我信中的所有这些论点摆在面前，准备予以驳斥。那么，就让我们看看他是否驳倒了它们；他是否对它们作出了任何回应。因为我补充了他毁谤指责我隐瞒的那些话，并且做完之后，我用一个更简短的形式再次向他提出同样的问题；因为通过补充这两个词，他大大帮助了我缩短这一命题。倘若我们寻求圣洁施予者的良心来洁净领受者的良心，并且倘若一个人明知地从无信德者那里领受了信德，他所领受的不是信德而是罪责，那么，当领受者不知道施予者的良心沾染了罪责，并且他在不知情的情况下从无信德者那里领受了信德时，我们到哪里去找办法来洁净领受者的良心呢？我追问，我们到哪里去找到办法来洁净它？让他告诉我们；让他不要岔到另一个话题上去；让他不要向缺乏经验的人眼中撒迷雾。至少，在穿插了许多曲折的迂回、并彻底讨论之后，让他最后告诉我们，当无信德施洗者良心上的罪责被隐藏起来时，我们到哪里去找到办法来洁净领受者的良心——倘若我们寻求圣洁施予者的良心来洁净领受者的良心，倘若一个人明知地从无信德者那里领受了信德，他所领受的不是信德而是罪责？因为所谈论的那个人从一位无信德者那里领受，这位无信德者没有圣洁施予者的良心，而只有被罪责玷污并被隐藏的良心。那么，我们到哪里去找到办法来洁净他的良心？他从哪里领受信德？因为如果当施洗者的无信德和罪责被隐藏时，他既不被洁净，也不领受信德，那么当这些后来被揭露和谴责时，为什么他不被重新施洗，以便被洁净并领受信德？但是，如果当另一人的无信德和罪责被隐藏时，他却被洁净并确实领受了信德，那么他从哪里获得洁净，他从哪里领受信德——既然没有圣洁施予者的良心来洁净领受者的良心？让他告诉我们这一点；让他对此作出回复：如果圣洁施予者的良心是我们寻求洁净领受者良心的关键，而这一点并不存在（因为施洗者隐藏了他的无信德和罪责），那么他从哪里获得洁净，他从哪里领受信德？对此，他根本未曾作答。\par

% structure: p0046 source=h2 target=subsection reason=relative-heading-rank; base=h2

\subsection{第二十一章。}

24. 但是看啊，当他在辩论中陷入困境时，他又向我发动一场充满迷雾和狂风般的攻击，以便真理的平静清澈可能被遮蔽；由于他极度匮乏，他变得资源丰富，不是表现在说出真理，而是表现在无代价的空洞辱骂中。请以最锐利的注意和最持久的毅力，紧紧抓住他应当回答的问题——即当施予者的良心斑点被隐藏时，在哪里可以找到办法来洁净领受者的良心——免得他那雄辩的强风从你们手中夺走它，而你们反而被他那浮夸言词的黑暗风暴卷走，以致完全看不出他从何处离题，以及他应该回到何处；并且请看，这个人能漫游到何处，既然他无法站定在他所承担的问题上。因为请看他说的有多少，只是因为他没有任何应该说的话。他说\Quote{我在滑溜之处滑行，却被支撑着；我既没有摧毁也没有证实我提出的反对意见；我用不确定的东西取代确定的东西；我不允许我的读者相信真实的东西，反而使他们以增多的怀疑来看待可疑的东西。}他说\Quote{我拥有学术派哲学家\Person{卡涅阿德斯}那被咒诅的才能。}他竭力暗示学术界关于人类感觉的虚假或错误的看法，在此也表明他对他所说的东西完全无知。他宣称\Quote{他们说雪是黑的，而实际上雪是白的；银是黑的；一座塔是圆的或没有棱角，而实际上它是有棱角的；桨在水里是断的，而实际上是完整的。}这一切都是因为，当他说\Quote{我们寻求施予者的良心，或圣洁施予者的良心，来洁净领受者的良心}时，我回答说：\Quote{如果施予者的良心被隐藏，并可能沾染了罪责，那怎么办？}这就是他所说的黑雪、黑银、圆塔而非有棱角的塔、水中断而实未断的桨——因为我提出了一种可以设想、却不能实际存在的情况，即施予者的良心可能被隐藏，并可能沾染了罪责！\par

25. 接着他继续以同样的口吻喊道：\Quote{那个\Emph{如果}是什么？那个\Emph{可能}是什么？只不过是一个怀疑者不确定而摇摆的犹豫，你们的诗人说——}\par

\Quote{如果我现在回到那些说\InnerQuote{如果天塌下来怎么办}的人那里，怎么办？}\par

他是指，当我说\Quote{若施予者的良心隐而不见，并可能沾染罪疚，又当如何？}这就像是说\Quote{若天塌下来，又当如何？}一样吗？的确，\Quote{又当如何}这个词，因为有可能隐而不见，也有可能不隐而不见。因为当不知道施予者在想什么，或他犯了什么罪时，他的良心确实隐藏于领受者的视线之外；但当他的罪显然可见时，就不算隐藏。我用了\Quote{并可能沾染罪疚}这一说法，因为有可能隐而不见却仍是纯洁的；也有可能隐而不见并沾染罪疚。这就是\Quote{又当如何}的意思；这就是\Quote{可能}的意思。这难道像是\Quote{若天塌下来，又当如何？}吗？哦，人多少次被定罪，多少次承认自己良心沾染罪疚和奸淫，而人们却在无知中由他们施圣洗，这些人在后来暴露的罪行中被剥夺了职分之后，天并没有塌下来！我们这里与\Person{庇鲁斯}和\Person{富里乌斯}有何相干？他们曾为不义辩护反对正义。我们这里与无神论者\Person{狄亚戈拉斯}有何相干？他否认有天主，似乎是先知预先论及的那个人：\BibleQuote{愚妄人心里说：没有天主。}\BibleRef{咏 14:1}我们与这些人有何相干？为何提及他们的名字，除非是为了给一个无话可说的人提供一点转移注意力的东西？好让他虽无必要，却总算是说了些关于这些人的话，这样当前的话题似乎就有了进展，而一个仍未得到回答的问题也就被认为得到了回答。\par

% structure: p0051 source=h2 target=subsection reason=relative-heading-rank; base=h2

\subsection{第二十二章}

最后，如果\Quote{又当如何}和\Quote{可能}这两个词是如此绝对不可容忍，以至于为了它们，我们应当唤醒长期沉睡的学院派和\Person{卡尔内亚德}和\Person{庇鲁斯}和\Person{富里乌斯}和\Person{狄亚戈拉斯}，以及黑雪、天塌，以及其他一切同样无意义和荒谬的东西，那就让它们从我们的论证中除去吧。因为事实上，不用它们，要表达我们想说的话也绝非不可能。稍后所发现的内容，以及他本人从我信中所引用的，足以满足我们的目的：\Quote{那么，当施予者的良心被污染，而领受圣事的人对此一无所知时，领受圣洗的他将如何被洁净呢？}你承认这里没有\Quote{又当如何}，没有\Quote{可能}吗？好吧，那么请给出回答。要密切注意，免得他在下文回答了这个问题。\Quote{但是，}他说，\Quote{我用你诡辩的信德来约束你，使你不致偏离更远。你为什么用愚昧的论据使你的生活远离错误？你为什么毫无理由地扰乱信仰体系？我用这同一句话来约束并说服你。}这话是\Person{佩提利安}说的，不是我。这些话来自\Person{佩提利安}的书信；但来自那封信，我刚刚添加了他指责我删掉的两个词，表明尽管添加了它们，我的问题（他没有回答）的针对性仍然更加简洁明了。这两个词无疑是\OriginalTerm{在圣洁中}{in holiness}和\OriginalTerm{明知地}{wittingly}：所以不应是\Quote{施予者的良心}，而应是\Quote{\Emph{在圣洁中}施予者的良心}；也不应是\Quote{他从无信者那里接受了自己的信德}，而应是\Quote{他\Emph{明知地}从无信者那里接受了自己的信德}。然而，我并没有真正删掉这些词；只是在我手中的抄本里没有找到它们。很可能那个抄本不正确；而且，我甚至相信，这一暗示就足以激起学院派对我不满，并声称我宣称抄本不正确，就像我宣称雪是黑的一样。因为我为什么要以其人之道还治其人之身，说他假装我删掉了词，而实际上是他自己后来添加了它们呢？既然抄本（它不会生气）可以证实那个不正确的标记，而无需我任何粗暴的指责。\par

% structure: p0053 source=h2 target=subsection reason=relative-heading-rank; base=h2

\subsection{第二十三章}

首先，关于那第一个说法\Quote{在圣洁中施予者}，它对我的探究毫无妨碍，尽管他对此十分困扰，无论我使用\Quote{若我们寻求施予者的良心}这一说法，还是更完整的说法\Quote{若我们寻求在圣洁中施予者的良心来洁净领受者的良心}，那么，当施予者的良心被污染，而领受圣事的人对此一无所知时，领受圣洗的他将如何被洁净呢？至于所添加的另一个词\Quote{明知地}，使得句子不应是\Quote{他从无信者那里接受了自己的信德}，而应是\Quote{他明知地从无信者那里接受了自己的信德，所接受的不是信德，而是罪疚}，我承认我曾说过一些好像该词不存在的话，但我可以轻易地舍弃它们；因为它们对我论证的便利造成的阻碍大于对其实力的帮助。因为我可以更直接、更清晰、更简短地提出这个问题：\Quote{若我们寻求在圣洁中施予者的良心来洁净领受者的良心}，以及\Quote{若他明知地从无信者那里接受了自己的信德，所接受的不是信德，而是罪疚}，那么，那个人如何被洁净呢？他看不见施予者（但不圣洁）良心上的污点。他从何处接受真信德呢？一个人由无信者无知地施圣洗，他从何处接受真信德？让我们宣布这从何而来，那么整个圣洗的理论就会被揭示；然后所有需要研究的问题就会水落石出——但前提是必须被宣布，而不是浪费时间在诽谤上。\par

% structure: p0055 source=h2 target=subsection reason=relative-heading-rank; base=h2

\subsection{第二十四章}

28. 因此，他在这两个词中所发现的——无论是他关于删去它们而提出的诬蔑指控，还是因添加它们而自夸——你都可以看出，这丝毫不能阻碍我的问题，而他对此找不到任何可以回答的；因此，他不愿沉默，便趁机对我的人格发起攻击——我本该说，他从讨论中退出了，除非他从未真正进入讨论。因为就好像问题关乎我本人，而非关乎教会或圣洗的真理，所以他声称，我通过删去这两个词，论证了这仿佛不是我的良心之路上的绊脚石，即我忽略了他所谓的污秽我的那人的亵渎良心。但如果真是这样，那么引入的\Quote{明知}一词的添加就会对我有利，而它的删去会对我不利。因为如果我希望我的辩护基于这样的理由，即我被认为不了解为我施洗者的良心，那么我会接受\Person{Petilianus}是在为我说话，因为他并非笼统地说：\Quote{那从无信者接受信德的人}，而是说：\Quote{那明知从无信者接受信德的人，所接受的并非信德，而是罪债}；由此，我可以夸耀我所接受的并非罪债，而是信德，因为我可以说我并非明知从无信者接受，而是不了解施予者的良心。所以，如果你能的话，请看并仔细计算，他仅仅在\Quote{我不了解}这一短语上浪费了多少多余的话，他声称我曾使用过它；而我其实从未使用过它——部分是因为讨论的问题并非关于我，使我需要用这个短语；部分是因为为我施洗之人身上并无明显的过犯，使我被迫在辩护中说我不了解他的良心。\par

% structure: p0057 source=h2 target=subsection reason=relative-heading-rank; base=h2

\subsection{第二十五章}

29. 然而，\Person{Petilianus}为了逃避回答我所说的话，却把自己未曾说过的话当作前提，把人们的注意力从他应偿还的债务上引开，以免他们要求他做出应有的回答。他不断引入\Quote{我不了解}、\Quote{我说}等说法，并回答说：\Quote{但如果你不了解}；并且，仿佛在驳斥我，使我无法说出\Quote{我不了解}，他引用了\Person{Mensurius}、\Person{Cæcilianus}、\Person{Macarius}、\Person{Taurinus}、\Person{Romanus}，并声称\Quote{他们曾反对天主的教会，这是我不可能不知道的，因为我是一个非洲人，并且年事已高}，然而，就我所听说的，\Person{Mensurius}是在多纳特派脱离教会共融之前，死于教会共融的合一中；而关于\Person{Cæcilianus}的历史，我曾读到，他们自己曾将他的案件提交给\Person{Constantine}，他一次又一次地被这位皇帝任命的调查此事的法官宣告无罪，第三次又被皇帝本人宣告无罪，因为他们曾向他上诉。至于\Person{Macarius}、\Person{Taurinus}和\Person{Romanus}在司法或行政职责范围内为合一而针对他们顽固的疯狂所做的一切，毫无疑问都是按照法律执行的，而正是这些人（指多纳特派）自己通过将\Person{Cæcilianus}的案件提交皇帝裁决，使得这些法律不可避免地被通过并实施。\par

30. 在诸多完全无关的事情中，他说：\Quote{我受到总督\Person{Messianus}判决的沉重打击，以致被迫逃离非洲。}由于这个谎言（如果他不是这谎言的制造者，那么当别人恶意编造时，他肯定恶意地相信了），他有多么大胆不仅说出，而且以惊人的轻率写下多少其他的谎言啊！要知道，在\Person{Banto}担任执政官之前，我就去了\Place{Milan}，并且，根据我那时所从事的修辞学教师的职业，我在一月一日当着大众的面朗诵了一篇颂词，颂扬他为执政官；在那次旅行之后，我只在暴君\Person{Maximus}死后才回到非洲；而总督\Person{Messianus}审理摩尼派案件是在\Person{Banto}担任执政官之后，这一点从\Person{Petilianus}本人插入的年表中的日期足以表明。如果有必要为此向那些怀疑或相信相反情况的人做出证明，我可以举出许多在其时代中杰出的人物，作为我一生中那个时期的充分见证人。\par

% structure: p0060 source=h2 target=subsection reason=relative-heading-rank; base=h2

\subsection{第二十六章}

31. 但我们为何要查究这些事？为何我们既忍受又不必要的拖延？我们通过这样的方式，难道就能找出用什么方法来洁净领受者的良心——他不知道施予者的良心已被罪责玷污？此人被一个不信者无意中施洗，他从何处领受信德？——这是佩提利安在自己心中提出的、要在我书信中回答的问题，然后他信口开河，说了许多与当前问题无关的话。他多少次说：\Quote{如果你当时不知情}，——好像我说过，我从未说过的话，即我不认识为我施洗者的良心。他似乎没有别的目的，只是通过他那邪恶之言的口所倾倒出来的一切，来证明我对那些与我一同领洗、并与我共融交往的人的恶行并非不知情，而且他显然完全明白，不知情并不能定我的罪。所以看，如果我不知情，正如他反复说的那样，那么毫无疑问，我对所有这些罪行都是无辜的。那么，我既然不认识那不圣洁地施予者的良心，如何能被洁净，以免被他的过犯所缠累呢？那么，我既然被一个不信者无意中施洗，又从哪里领受信德呢？因为他并非没有原因地重复了这么多次\Quote{如果你当时不知情}，而只是为了阻止我被视为无辜，他毫无疑问地认为：如果一个人无意中从一个不信者领受了他的信德，并且不认识那施予者（但不圣洁地施予）良心的污点，那么这个人的无辜并没有被违背。因此，让他说出，这样的人应当通过什么方式被洁净，他们从何处领受的不是罪责而是信德。但他不要欺骗你们。不要让他说了许多却什么也没说；或者更确切地说，不要让他说得多却什么也没说。接下来，要强调一个我想到的、不可忽略的问题——如果我是有罪的，因为我没有不知情（用他自己的措辞），并且由于我是非洲人，且已年长，而被证明并非不知情，那么让他承认，全世界其他民族的年轻人，无论从种族还是从你们用来指控我的那个年龄来看，都没有机会知道那些被拿来指控我们的事——无论这些事是真是假——他们就不是有罪的；然而他们，如果落入你们手中，就会被重新施洗，而不加任何这类考量。\par

% structure: p0062 source=h2 target=subsection reason=relative-heading-rank; base=h2

\subsection{第27章。}

32. 但这并非我们现在所探究的。倒不如让他回答（他为了避免回答，而离题万里）——如果洁净领受者良心的是施予圣洁者的良心，那么领受者的良心，若不认识施予者良心的污点，如何得洁净？若故意从一个不信者领受信德的人所领受的不是信德而是罪责，那么一个被不信者无意中施洗的人，从何处领受信德？因此，我们且撇开他那些毫无根据向我抛来的辱骂，仍然注意他在后面并没有说出我们所要求的东西。但我倒想看看他喋喋不休地陈述此事的方式，好像他定能令我们困惑似的。他说：\Quote{但让我们回到你那幻想的论点，你似乎通过言辞的形式为自己描绘了你所施洗的人。因为既然你看不见真理，更得体的做法是想象一些可能的事。}佩提利安以这些话作为开场白，即将陈述我所说的话。然后他接着引用：\Quote{看，你说，那不信者准备好施洗，但即将受洗的人对他的不信一无所知。}他没有完整引用我的提议和问题；然后他立即反过来问我，说：\Quote{这个人是谁？从哪个角落里冒出来的，你把他摆在我们面前？你为什么似乎看见一个你想象的人，而避免看见一个你应当看见、并应仔细审查和考验的人？但是既然我看你不熟悉圣事的秩序，我尽可能简短地告诉你：你既应当审查你的施洗者，也应当被他审查。}那么，我们等候的是什么？他应该告诉我们，那不认识不圣洁施予者良心污点的领受者，其良心如何得洁净；以及那无意中被不信者施洗的人，从何处领受不是罪责而是信德。我们所听到的只是：施洗者应当被那愿意领受不是罪责而是信德的人最仔细地审查，以便他了解那洁净领受者良心的圣洁施予者的良心。因为那没有进行这项审查、因而无意中从不信者领受洗礼的人，恰恰因为没有审查，从而不知道施予者良心的污点，他就无法领受信德而非罪责。那么，他为什么还要加上他那么看重的话——那个他诽谤指控我删掉的字，即\Emph{故意地}？因为他不想句子成为\Quote{一个人从一个不信者领受他的信德，领受的不是信德而是罪责}，他似乎为那无意行事的人留下了一些希望。但现在，当被问及那被不信者无意中施洗的人从何处领受信德时，他回答说，他本应审查他的施洗者；这样，毫无疑问，他拒绝给那可怜人哪怕不知情的许可，因为他没有找出他可以从何处领受信德——除非他将信任放在那施洗者身上。\par

% structure: p0064 source=h2 target=subsection reason=relative-heading-rank; base=h2

\subsection{第28章。}

33. 这就是我们在你们一党中所厌恶的；这就是天主的判决所谴责的，它以其至真至明的声音呼喊：\Quote{凡信赖世人的人，是可咒骂的。}\BibleRef{耶 17:5} 这正是圣洁的谦卑和宗徒的爱德所最明确禁止的，正如\Person{保禄}所宣告的：\Quote{人不可因人而夸耀。}\BibleRef{格前 3:21} 这就是为何虚妄的诽谤和最恶毒的谩骂对我们的攻击越来越猛烈的原因，为的是，当人的权威被推翻时，那些我们受托付而向他们施行天主的言语和\OriginalTerm{圣事}{sacrament}的人，便不再有希望的依据。我们回答他们说：你们依恃人要到几时呢？可敬的天主教会回答他们说：\Quote{我的灵魂的确期待天主：我的救援来自祂。唯有祂是我的天主，我的助佑；我决不动摇。} 他们离开天主的家还有什么其他原因呢？岂不是因为他们假装不能忍受那些作卑贱使用的器皿，而这些器皿直到审判之日才会从家中清除吗？然而，他们却以自己的行为和时间记录表明，他们自己正是这种器皿，却把这罪名强加于他人；关于这些作卑贱使用的器皿，天主的仆人——那位良善忠信、或能够藉\OriginalTerm{圣洗}{baptism}接受\OriginalTerm{信德}{faith}的人（如我上文所示）——明确地说：\Quote{我的灵魂的确期待天主}（你看，是期待天主，而非期待人）：\Quote{我的救援来自祂}（不是来自人）。但是\Person{佩蒂利安}却拒绝将人的洁净与净化归功于天主，即使那施洗者（其\OriginalTerm{良心}{conscience}有污点，并非在圣洁中施洗）的污点被隐藏，而某人无意中从无信德者那里接受了\OriginalTerm{信德}{faith}。他说：\Quote{我尽量简短地告诉你们：你们必须查验你们的施洗者，也要受他查验。}\par

% structure: p0066 source=h2 target=subsection reason=relative-heading-rank; base=h2

\subsection{第二十九章}

34. 我恳求你们注意这一点：我问，当领受者不清楚施与者（非在圣洁中施与）\OriginalTerm{良心}{conscience}的污点时，如果等待圣洁施与者的\OriginalTerm{良心}{conscience}来洁净领受者的\OriginalTerm{良心}{conscience}，那么领受者的\OriginalTerm{良心}{conscience}将从何处得以洁净？而且，若凡是从无\OriginalTerm{信德}{faith}者有意识地接受\OriginalTerm{信德}{faith}的人，所接受的不是\OriginalTerm{信德}{faith}而是罪责，那么，那无意中被无\OriginalTerm{信德}{faith}者施洗的人，将从何处接受\OriginalTerm{信德}{faith}？他回答我说，\OriginalTerm{施洗者}{baptizer}和受洗者都应接受查验。为了证明这一点（这一点并无争议），他举了\Person{若翰}的例子：他被那些问他的人查验，问他自称是谁\BibleRef{若 1:22}，而他也反过来查验那些听他说\Quote{毒蛇的种类！谁指示你们逃避那将要来的忿怒？}的人\BibleRef{玛 3:7}。这与主题有何关系？这与正在讨论的问题有何关系？天主曾赐给\Person{若翰}最卓越圣洁生活的见证，先前最崇高的预言在他受孕和出生时就已证实。但\OriginalTerm{犹太人}{Jews}提问时，已经相信他是\OriginalTerm{圣人}{saint}，想查明他自认是哪位\OriginalTerm{圣人}{saint}，或者他本人是否就是\OriginalTerm{圣人}{saint}中之\OriginalTerm{圣人}{saint}，即\Person{耶稣基督}。事实上，他受到了极大的信任，以至于无论他如何论及自己，人们都会立即相信。因此，如果我们遵循这个先例，宣称如今每一位\OriginalTerm{施洗者}{baptizer}都应被查验，那么，无论他怎样谈论自己，他也必须被相信。但是，那充满诡诈的人——我们知道\OriginalTerm{圣神}{Holy Spirit}按照\OriginalTerm{圣经}{Scripture}的话是远离他的\BibleRef{智 1:5}——谁会不希望人们相信他最好的方面，或者谁会犹豫不决，不使用任何能想到的言语来达到这一点呢？因此，当他被问及他是谁，并回答他是上主旨意的忠实分施者，且他的\OriginalTerm{良心}{conscience}未被任何罪恶的污点玷污时，这就算作全部的查验吗？还是要对他的品性和生活做进一步更仔细的调查？当然会有。但经上并未记载，那些在约旦旷野中问\Person{若翰}他是谁的人这样做了。\par

% structure: p0068 source=h2 target=subsection reason=relative-heading-rank; base=h2

\subsection{第三十章}

35. 因此，这一先例与当前问题完全无关。我们倒不如说，宗徒的宣告充分地教导了这种谨慎，当他说：\BibleQuote{这些人也应当先受考验，然后才让他们作执事，并且要无可指责。}\BibleRef{弟前 3:10} 既然在双方中，几乎所有人都是忧虑而惯常地这样做，那么为什么有这么多人在担任这职务之后被发现是败坏的？除非是因为，一方面，人的谨慎常常受骗；另一方面，那些开始良好的人有时会堕落。既然这类事情频繁发生，以至无人能隐瞒或忘记，那么，\Person{Petilianus}现在以侮辱性的几句话教导我们，受洗者应当审查付洗者，这是为什么呢？因为我们的问题是：当付洗者在圣善之外施洗，而其良心的污点被隐藏时，如果我们要寻求在圣善中施洗者的良心来洁净受洗者的良心，那么受洗者的良心如何得以洁净？\Quote{既然我看到你不了解圣事的顺序，}他说，\Quote{我尽可能简短地告诉你：你既有义务审查你的付洗者，也有义务接受他的审查。}多么荒谬的回答！他在许多地方被如此众多的人包围，这些人是由一些起初看似正直贞洁、后来因过失揭露而被定罪和革职的付洗者施洗的。他认为他可以通过简短地说付洗者应当被审查来回避这个问题的力量——在这个问题中，我们问：当付洗者在圣善之外施洗，而其良心的污点不为受洗者所知时，如果我们要寻求在圣善中施洗者的良心来洁净受洗者的良心，那么受洗者的良心如何得以洁净？——我认为，他认为他可以通过说付洗者应当被审查来回避这个问题的力量。没有什么比不一致与真理更不幸的了，因为真理将每个人都紧紧包围，使人找不到逃脱的办法。我们问：一个由无信德者施洗的人，他从谁那里领受信德？回答是：\Quote{他应当审查他的付洗者。}那么，既然他没有审查他，甚至无意中从一个无信德者那里领受了信德，他领受的不是信德，而是罪债吗？那么，为什么那些被发现是由后来被查出和定罪的败坏者（其真实品格尚被隐瞒时）施洗的人，不再重新受洗呢？\par

% structure: p0070 source=h2 target=subsection reason=relative-heading-rank; base=h2

\subsection{第三十一章}

36. \Quote{那么，}他说，\Quote{我所加上的\Emph{故意地？}这个词在哪里？所以我并没有说：那从一个无信德者领受信德的人；而是说：那\Emph{故意地}从一个无信德者领受信德的人，领受的不是信德，而是罪债。}因此，那无意中从一个无信德者领受信德的人，领受的不是罪债，而是信德；因此我问：他从哪里领受了信德？他这样被置于困境，便回答说：\Quote{他应当审查他。}就算他应当这样做；但事实上，他没有做，或者他不能做：你的判断是什么？他被洁净了，还是没有？如果他被洁净了，我问：从何而来？因为那在圣善之外施洗者被污染的良心，是他所不知道的，不能洁净他。但如果没有被洁净，就命令他现在被洁净。你没有下这样的命令，所以他是被洁净了的。请告诉我，是通过什么方式？你至少告诉我，\Person{Petilianus}没有说出来的是什么。因为我向你提出他无法回答的同样的话。\Quote{看，那无信德的人站着准备施洗；但那要受洗的人不知道他的无信德：你认为他会领受什么——信德，还是罪债？}这足以作为经常提出的问题形式：回答，或者努力寻找他回答了什么。你会找到已受谴责的辱骂。他责备我，仿佛嘲弄一般，坚持认为我应当提出一些或然性以供考虑，因为我不能看到真理。因为他重复我的话，把我的句子截成两半，说：\Quote{看，你说，那无信德的人站着准备施洗；但那要受洗的人不知道他的无信德。}然后他接着问：\Quote{这人是谁，他从哪个角落冒出来，你把他提出来给我们？}就好像只有一两个人，而这类情况不是经常在双方各处发生！他为什么问我这人是谁，他从哪个角落冒出来，而不环顾四周，看到无论在城市或乡村，那些不包含被发现罪行并从神职被革除的人的教会少之又少？当他们的真品格被隐藏时，当他们希望被认为好（其实是坏的）、被认为贞洁（其实犯奸淫）时，他们一直处于欺骗之中；因此，根据圣经，圣神从他们那里逃离了。\BibleRef{智 1:5}所以，我所建议的那个无信德的人，正是从这些迄今隐藏其品格的人群中冒出来的。他为什么问我他从哪里冒出来，却对所有这些人群闭上眼睛？如果我们只考虑那些可能被定罪和革职的人，这些人群发出的噪音已足够使盲人满意。\par

% structure: p0072 source=h2 target=subsection reason=relative-heading-rank; base=h2

\subsection{第三十二章}

37. 至于他本人在其书信中所提出的：\Quote{\OriginalTerm{奎沃德多德奥斯}{Quodvultdeus}因两起通奸罪被定罪并被你们驱逐后，却为我们这边的人所接纳。}对此，我该如何评论呢？（我这样说，并无意对这位证明自己案情正当——或至少使人相信其正当——的人有所不敬。）当你们中那些尚未被发觉的人施行圣洗时，那些人从他们手中领受的是什么？是信德，还是罪责？当然不是信德，因为他们不具备一位圣洁施洗者的良心，无法洁净领受者的良心。但同样也不是罪责，因为根据那句话：\Quote{凡明知是从不忠信者领受其信德的人，所领受的不是信德，而是罪责。}然而当人们由我所说那些人施洗时，他们确实不知道那些人是什么样的人。再者，他们既然没有从施洗者那里领受信德（因为施洗者不具备圣洁施洗者的良心），也没有领受罪责（因为他们受洗时并不知情，而是因无知而未得知其过错），因此他们既不具信德，也不具罪责。所以，他们不属于那些品行败坏之人的行列。但他们也不能被算在信友之列，因为他们既不能从施洗者那里领受罪责，同样也不能领受信德。但我们看见，你们视他们为信友，你们中无人宣称他们应当重新受洗，反而所有人皆承认他们已领受的圣洗有效。因此，他们确实领受了信德；然而他们并非从那些不具备圣洁施洗者良心的人那里领受，以洁净领受者的良心。那么他们是从何处领受的呢？这就是我所着力之处；这就是我极为迫求的问题；对此我亟需一个答复。\par

% structure: p0074 source=h2 target=subsection reason=relative-heading-rank; base=h2

\subsection{第三十三章。}

38. 请看现在\OriginalTerm{佩蒂利安}{Petilianus}是如何为了回避这个问题，或为了不被证明无力回答，而徒然转入无关的话来辱骂我们，指控我们却毫无实证；当他偶然试图抵抗，作出仿佛为其事业而战的样子时，却处处轻易被击败。然而，对于我们所问的这个唯一问题，他从未给出任何回答：既然我们寻求的是那圣洁施洗者的良心，以洁净领受者的良心，那么当施洗者的良心已被玷污，而领受者对此一无所知时，他所领受的圣洗又怎能被洁净呢？因为他从我书信中引用的这些话，他把我描绘成提问者，而他自己却显得无言以对。在说了我刚才叙述的话之后，当他四面被逼入绝境时，他被迫承认施洗者应当由受洗候选人考察，而候选人反过来也应由施洗者考察；当他试图以\Person{若翰}的例子来巩固这一说法，希望找到极其疏忽或极其无知的听众时，他又继续提出其他与手头问题完全无关的圣经证据，如太监对\Person{斐理伯}所说：\Quote{看，这里有水，有什么阻止我受洗呢？}\BibleRef{宗 8:36}他声称：\Quote{他知道品行败坏的人是被禁止的。}以此论证\Person{斐理伯}之所以不禁止他受洗，是因为他在读经中已证明他对基督信仰的程度——仿佛他曾禁止了\Person{西满}术士似的。他又援引先知们害怕被虚假圣洗所欺骗，因此\Person{依撒意亚}说：\Quote{无信的水，没有信德。}仿佛显示在无信之人中水是虚假的；但说这话的并非\Person{依撒意亚}，而是\Person{耶肋米亚}，他是指着说谎之人而言，并以譬喻称百姓为水，这在\BibleBook{}{默示录}中已清楚地显示。\BibleRef{默 17:15}他又援引\Person{达味}的话：\Quote{不要让罪人的油涂我的头。}但\Person{达味}讲的是口若悬河者以虚假赞美欺骗的谄媚，使被赞美者骄矜自大。同一圣咏的前文使这一意义显而易见，因为他说道：\Quote{让义人击打我，这是慈爱；让他责备我；但罪人的油却不可伤我的头。}还有什么比这话更清楚？还有什么更显明？因为他宣称，他宁愿被义人严厉的纠正所慈爱地责备，以便痊愈，也不愿被谄媚者柔和的话语所涂抹，以致骄矜自大。\par

% structure: p0076 source=h2 target=subsection reason=relative-heading-rank; base=h2

\subsection{第三十四章。}

39. \Person{佩蒂利安}还引用了\Person{若望}宗徒的警告：\Quote{不可信任何神，但要考验那些神是否出于天主。}\BibleRef{若一 4:1}仿佛这种谨慎应当用于在现世就将麦子与糠秕提前分开，而不应当是为了防止麦子被糠秕欺骗；或者仿佛即使说谎的神说了某些真话，也应予以否认，因为那个我们应憎恶的神说了它。若有人这样想，他就疯狂到足以争辩说，\Person{伯多禄}不应当说：\Quote{你是基督，永生天主之子。}\BibleRef{玛 16:16}因为魔鬼早已说过类似的话。因此，既然基督的圣洗，无论是由不义之人还是义人施行，都只是基督的圣洗，一个谨慎的信徒基督徒所应做的是避免人的不义，而不是谴责天主的圣事。\par

40. 诚然，培提利安在所有这一切中，对以下问题未作回答：如果施予者的良心必须是圣洁的，才能洁净领受者的良心，那么，当领受者不知情而施予者的良心已被玷污时，领受洗礼的人又怎能被洁净呢？他的一位同僚，来自图布西库布的某某西彼廉，被发现在妓院中与一名最放荡的女子一起，并被押解至迦太基的普里米安面前，遭到处分。如今，此人未被发现及受处分前所施行的洗礼，显然他并未具有圣洁施予者的良心，以洁净领受者的良心。那么，那些今日在他受处分后（显然未再重新受洗）的人，又是藉何途径获得洁净的呢？提及此人仅是为了防止培提利安重复说：\Quote{此人是谁？从哪个角落冒出来的？你竟向我们提出他来？}为何你们一方不按培提利安的意见，如同审查若翰那样，去审查那位施洗者呢？或者实际情况是，他们虽尽人力所能审查了那人，却因他长久狡猾地伪装而未能识破他？\par

% structure: p0079 source=h2 target=subsection reason=relative-heading-rank; base=h2

\subsection{第三十五章}

由这人施行出来的水难道就不真实吗？难道淫乱者的油就不是罪人的油吗？或者，我们必须持守天主教会 的教导——这是真实的——即那水和那油不属于施行的人，而属于那时被呼求的祂。为何那些被那位隐瞒了自己罪过的伪善者所施洗的人，未能考验那神，以证明它不是出于天主呢？因为连那训练的圣神当时也在躲避那伪善者。\BibleRef{智 1:5}难道说：祂躲避他，但同时并未弃绝祂的圣事，即便由他施行？最后，既然你们并不否认那些人已经得到洁净——因为你们现在并不操心在他受处分后再去洁净他们——那么，请看看培提利安在以如此众多方式散布了如此多迷雾之后，是否在任何一处对以下问题给出了答案：如果洁净领受者良心的是施予者的圣洁良心，而那些秘密不洁之人不可能拥有这样的良心，那么这些人又是如何被洁净的呢？\par

41. 于是，对于这个如此急切地向他提出的问题，他未作任何回答；接着，他甚至为自己寻求言论的空间，说：\Quote{既然先知和宗徒都足够谨慎地惧怕这些事情，你们有何脸面说，罪人的洗礼对于那些以良好良心相信的人是圣洁的？}好像我或任何公教徒主张那由罪人施行或领受、且由罪人主礼的洗礼是属于罪人的，而非属于因相信其名而受洗之人的。然后他转而攻击叛徒犹达斯，尽其所能地抨击他，引用先知许久以前关于他的预言，仿佛要使散布普世的基督教会——其案件正为此次辩论所涉及——陷入叛徒犹达斯的不敬虔之中；他没有想到，这件事本身应提醒他：我们不该怀疑那遍布世界的教会就是基督的教会，正如这些预言在多年前如实应验；我们也不该怀疑基督必被其一位门徒出卖，因为这事也同样被预言了。\par

% structure: p0082 source=h2 target=subsection reason=relative-heading-rank; base=h2

\subsection{第三十六章}

42. 但此后，当\Person{培提利安}碰到我们提出的异议，即他们容许他们曾谴责的\Person{玛克西米安}追随者的圣洗——尽管在陈述这个问题时，他宁可使用他自己的言辞而非我的；因为我们并不主张罪人的圣洗对我们有益，因为我们坚信它不仅不属于任何罪人，而且根本不属于任何人，因为我们深信它唯独属于基督——他以这种形式提出问题，说：\Quote{而你们却固执地声称，罪人的圣洗对你们有益是正当的，因为根据你们的说法，我们也曾维护我们公正谴责的罪犯的圣洗。}当他来到这个问题时，正如我之前所说，他先前摆出的所有战斗姿态都离他而去。他找不到任何出路，任何逃生之方，任何途径，无论是通过机敏的监视还是大胆的行动，他既无法悄悄溜走，也无法强行突围。他说：\Quote{尽管这一点，}他说，\Quote{我将在我的第二卷书中证明，我们的人与你们称为无罪的人之间的差别有多大，但同时，你们先摆脱你们所知的同僚中的过犯，然后再设法处理那些我们驱逐的人。}世上还有任何人，任何活在地上的人，会给出这样的回答，除非是那反对真理、而无法找到任何可作回答的人吗？因此，如果我们也要用同样的话：\Quote{同时，你们先摆脱你们所知的同僚中的过犯，然后再向我们提出任何涉及你们认为我们中间恶人的指控}——结果如何？是我们双方都获胜了，还是我们双方都失败了？不，反而祂为祂的教会并在祂的教会中获胜了，祂在圣经中教导我们，没有人应夸耀于人，夸耀的应夸耀于主。因为看，在我们这些以真理的雄辩断言一个人成义不是凭给他施洗的人，而是凭那记着的那位的人身上——\BibleQuote{对于那信仰那使不虔敬的人成义的天主的人，他的信德便算为正义}，\BibleRef{罗 4:5}——既然我们不夸耀于人，并力求在夸耀时因祂的恩赐而夸耀于主，那么无论\Person{培提利安}能就我们共融中的某些人证明什么过失或指控，我们是多么完全安全！因为我们中间，无论什么恶人，要么完全未被察觉，要么虽被某些人知晓，却为了合一与和平的纽带而被容忍，考虑到其他对其恶行不知情且无法在这些人面前定他们罪的好人，以免麦子连同莠子一起被拔除；然而他们如此承担自己恶行的重担，除了那些喜欢他们的不义的人之外，无人与他们分担。而且我们毫不担心那些受他们施洗的人不能成义，因为他们信仰那使不虔敬的人成义的天主，为要使他们的信德算为正义。\BibleRef{罗 4:5}\par

% structure: p0084 source=h2 target=subsection reason=relative-heading-rank; base=h2

\subsection{第三十七章}

43. 此外，根据我们的教义，无论\Person{培提利安}所说曾因索多玛人的罪而被我们驱逐、另有人接替他、后又实际上恢复至我们团体的人——他一直说却不知自己在说什么——还是他宣称在你们中间曾忏悔的人，不论他们各自的情况是否容许任何辩护，两者都不能损害那传布于万民、在世界中生长直到收割的教会。因为如果他们真的是你们所控告的教会中的邪恶肢体，那么他们已经不在其中，而是在糠秕中；但如果他们是好人，而你们以不义的控告污蔑他们的品格，他们自己就像金子般被试炼，而你们则像糠秕般被焚烧。然而他人的罪并不玷污那按照极其信实的预言遍布全球、等待世界末日如等待彼岸的教会；当它靠岸时，它将脱离恶鱼，在主的同一网中与它们相伴，忍受自然的不便而不犯罪，只要不耐烦地将它们分离是不合宜的。教会的纪律也并未因此而受到恒常、勤勉、谨慎的基督仆人的忽视；在他们管辖范围内，罪行以如此方式被揭露，以至于不能以任何可能性为托词辩护。无数证据可在那些曾是主教或次级神职的人中找到：他们如今被降级，或因羞愧而流亡他地，或投靠你们本身或其他异端，或在自己的地区为人所知。他们中有如此庞大的群体散布全地，以致如果\Person{培提利安}暂时收敛他的轻率言论，考虑他们，他绝不会陷入如此明显虚假且毫无根据的误解，竟以为我们应该赞同他所说的：\Quote{你们中无人无罪，而那里没有罪人被定罪。}\par

% structure: p0086 source=h2 target=subsection reason=relative-heading-rank; base=h2

\subsection{第三十八章}

44. 且不提散居大地各处的其他人——你几乎找不到一个地方没有这类人，由此可知监督和牧者即使在天主教会中也常被定罪——我们不必远求，米勒维斯的霍诺留便是一个例子。但再举斯彭多尼乌斯为例：他在天主教会中被定罪后，佩提利安仍祝圣他为司铎，并亲自为他重洗；他的罪案在高卢被认定，由我们的弟兄们转告我们，我们的同僚福图纳特在君士坦丁公开宣读；后来同一位佩提利安因他可恶的欺诈而将他驱逐。从斯彭多尼乌斯的案例来看，什么时候他不会提醒人们，即使在天主教会中，恶人也会被罢免职务？我惊奇他的心中究竟悬着何等鲁莽的悬崖，致使他口授了那些话，竟敢说：\Quote{你们中没有一人是无罪的，因为没有一个有罪的人被定罪。}因此，在天主教会中，恶人虽与善人身体上混杂，但灵性上彼此分离；他们或因人性的软弱未被察觉，或因纪律考虑而被定罪，但每个人都背负自己的重担。这样，那些由他们以基督的洗礼施洗的人，若既不藉效仿也不藉同意分享他们的罪，就免于危险；同样，如果他们由最好的人施洗，也只能因那称不虔敬者为义者而成义。因为凡信那称不虔敬者为义者的，信德就算为他的义。\par

% structure: p0088 source=h2 target=subsection reason=relative-heading-rank; base=h2

\subsection{第三十九章。}

45. 但当你们面对马克西米安追随者的案件时，你们无言以对：他们被三百一十位主教组成的会议判决定罪；在同一会议中被彻底击败，并在诸多总督的档案、诸多市镇的编年史中被记载；经法官命令，在各城军队的协助下，被逐出他们占有的圣堂；然而你们后来又极其光荣地接纳了他们，连同他们在你们共融之外所施洗的人，对他们的洗礼毫无质疑。我说，当你们面对他们的案件时，你们找不到任何答复。确实，你们被一种从你们自己发出的主张所击败——这主张虽非真理，却出自你们——即：一个人因同领圣事共融中别人的过错而灭亡，并且每个人的品格由施洗者决定——若施洗者有罪，受洗者便有罪；若施洗者无罪，受洗者便无罪。但如果这些观点是真的，那么毫无疑问，且不提无数其他的人，你们就被马克西米安追随者的罪所毁灭，因为你们的派别在一个如此庞大的会议中，已将他们的罪夸大到如同那些被地活活吞没之人的罪一样。但如果马克西米安追随者的罪没有毁灭你们，那么你们所持的这些观点就是假的；而那些不知名的、未经证实的阿非利加人的罪，就更不能毁灭整个世界了。因此，正如宗徒所说：\Quote{各人必背负自己的重担。}\BibleRef{迦 6:5}基督的洗礼唯独是基督的，不属于任何人；佩提利安许诺将把他第二卷书的主题定为我们对马克西米安追随者所提出的指控，这是徒然的，因为他过低估计人们的智力，仿佛他们看不出他无话可说。\par

% structure: p0090 source=h2 target=subsection reason=relative-heading-rank; base=h2

\subsection{第四十章。}

46. 因为如果普雷特克斯塔图斯和费利西安在马克西米安共融中所施行的洗礼是他们自己的，那么你们为何在他们所施洗的人身上接受那洗礼，如同它是基督的洗礼呢？但如果它确实是基督的洗礼（事实上正是），却因分裂的罪而不能使领受者获益，那么你们对那带着同一洗礼而进入你们团体的人，除了说因和平的纽带消除了他们可憎分裂的过犯，使他们不必像没有领受一样被迫领受圣洗圣事，而是使他们从前所领受的虽曾导致毁灭，如今开始对他们有益之外，还能说你们赐予了他们什么？或者，如果在你们的共融中不能这样赐予他们，因为分裂者中间的分裂者不可能获得恩赐，那么在天主教会共融中，至少对你们而言是赐予的：不是领受那似乎缺失的洗礼，而是使你们实际已领受的洗礼对你们有益。因为基督的一切圣事，若没有包含在基督合一中的爱，则拥有它们不是为了得救，而是为了审判。但由于\Quote{你们的主张}——即\Quote{因某些\OriginalTerm{叛教者}{traditors}的洗礼，基督的洗礼已从整个世界消失}——并非真理，因此你们面对承认马克西米安追随者的洗礼时，理所当然无言以对。\par

47. 因此，请以最警醒的注意观察并记住，培提利安是如何没有回答那个他自己提出的问题；他提出问题时措辞如此，似乎要把它作为一个起点，以便说些什么。因为他完全撇开了前一个问题，不愿向我们谈论它，我想是因为他力不能及；而且直到他整卷书的末尾，他也绝不会谈论它，尽管他引用了我书信的开头部分，仿佛这是一个需要驳斥的问题。因为即使他加上了那两个他指责我删去的词，把它们当作他立场的最强堡垒，他仍然毫无防御，无法找到任何回答，当他被问及：\Quote{若我们寻求的是那以圣洁施行者之良心来洁净领受者的良心，那么对于不知晓那以非圣洁施行者之良心的人，我们要从哪里找到洁净其良心的办法呢？并且，若任何人从无信者那里接受了他的信仰，他所接受的不是信仰而是罪责，那么那由无信者无意中施洗的人，要从何处接受不是罪责而是信仰呢？} 从他的话中早已明显可见，他对此问题毫无回答。\par

48. 其次，他用诽谤之口继续攻击隐修院和隐修士，也指责我，说我是这种生活的创始人。而这种生活究竟是什么样的，他完全不知道；或者更确切地说，尽管这在全世界都尽人皆知，他却假装不熟悉。接着，他声称我说过基督是施洗者，还从我书信中添加了一些词语，好像我把这当作我自己的观点提出来；实际上我是引用了他和你们的观点，而他却以大量严厉的言词攻击，仿佛是我自己说了这些反对自己的话，而他所指责的不是我的，而是他和你们的观点，我将尽我所能清楚地说明这一点。然后，他用许多不必要的言词竭力向我们表明，基督并不施洗，而是以祂的名、即因父及子及圣神之名施行洗礼；关于这同一的天主圣三，他说（或者是因为他愿意这样说，或者是因为他只能这样说）：\Quote{基督是天主圣三的中心。} 其次，他借术士西满和巴尔耶稣之名，对我们肆意谩骂。接着，他对塔穆加斯的奥普塔图斯案件保持谨慎悬而未决，以使自己不被由此产生的憎恶淹没，否认他或他的党羽曾对他作出判决，并且实际上暗示，关于他，他是由于我的建议而被压垮的。\par

% structure: p0094 source=h2 target=subsection reason=relative-heading-rank; base=h2

\subsection{第四十一章}

49. 最后，他在书信的结尾，对自己的党羽发出劝诫和警告，叫他们不要被我们欺骗，又为我们的党羽哀叹，说我们使他们变得比以前更糟。因此，在仔细考虑和讨论这些要点之后（从他书信的言辞中已足够清楚地显示这一点），培提利安对我最初在书信中提出的立场根本没有作出回答。我曾问：\Quote{假设正如他所断言的那样，我们寻求的是那以圣洁施行者之良心——或者更确切地说，加上他认为对他论点如此重要的支持——即那以圣洁施行者之良心，来洁净领受者的良心；那么，当施予者的良心被污染，而提议的领受者却不知情时，领受洗礼的人要如何被洁净呢？} 因此，一个反对虚假事业的人，在与之矛盾的真理困境中受压制，宁愿喘着粗气发出疯狂的辱骂，也不愿走在那不可战胜的真理之路上，这并不奇怪。\par

50. 现在我请求你们特别注意接下来的几句话，为要向你们清楚表明，他因为害怕什么而没有回答这个问题，并且为要把那些他试图隐藏在晦暗中的事显露出来。当我们问：\Quote{当施行者的良心被污染，而提议的领受者不知情时，那领受洗礼的人要如何得到洁净？} 他本来完全可以轻易地回答：\Quote{来自我们的主天主；} 并且至少可以满怀信心地说：当领受者不知晓那以非圣洁施行者良心上的污点时，天主完全洁净领受者的良心。但是，既然一个人已经被你们教派的教义所迫，将领受者的洁净建立在施行者的良心上，因为他曾说过：\Quote{因为施行者的良心，} 或 \Quote{那以圣洁施行者的良心，是被寻求来洁净领受者良心的。} 他自然害怕，唯恐一个人被一个隐藏其邪恶的恶人施洗，似乎比被一个真正而明显良善的人施洗更好；因为在前一种情况下，他的洁净将不依赖那以圣洁施行者的良心，而依赖至高无上的天主本身的圣洁。因此，由于这种忧虑，他唯恐自己陷入如此巨大的荒谬，甚或疯狂，以至于不知何处可以逃脱；他不愿意说出当领受者不知晓那以非圣洁施行者良心上的污点时，领受者的良心应如何被洁净；反而认为，最好以争吵的喧哗制造混乱，以此来掩盖所问的问题，而不是给出一个会立刻使他难堪的回答。但他从未想到，我们的书信能被如此有悟性的人读到，或者他的书信会被那些也读过我们书信的人读到，而他自称已对我们的书信作了答复。\par

% structure: p0097 source=h2 target=subsection reason=relative-heading-rank; base=h2

\subsection{第四十二章}

51. 因为我刚才所说的，在我那封信中讲得极为清楚，而他在回答时对此却只字未提；我恳请你们稍听片刻，看他在那里做了些什么。尽管你们是他的支持者，并憎恨我们，但如果可能，请平和地忍受。因为在他前一封信中——我只回复了该信的第一部分，因为那时该部分才刚到我们手中——他把在圣洗中寻得的希望完全寄托于施洗者身上，竟至于说：\Quote{万物皆由本源和根源组成；若某物没有头，它就什么也不是。}既然\Person{培提利安}说了这话，不愿将人受洗的本源、根源和头理解为除施洗者本人以外的任何事物，我便加了注，说道：\Quote{因此我们要问：如果施洗者的无信未被察觉，那么他所施洗的人是否领受了信德而非罪债？如果施洗者不是他的本源、根源和头，那么他从谁那里领受信德？他出自哪个本源？他是哪棵根上的枝芽？他的起点是哪个头？难道当受洗者不知道他施洗者的无信时，\Person{基督}就成了那本源、根源和头？}因此我现在也这样说，也这样喊，正如我那时在那里所说：\Quote{人的鲁莽和傲慢真可悲！你们为何不承认永远是\Person{基督}赐予信德，为的是藉此使人成为基督徒？你们为何不承认\Person{基督}永远是基督徒的本源，基督徒总是将其根扎在\Person{基督}内，\Person{基督}是基督徒的头？那么，是否即便当神圣的恩宠藉着圣洁忠信的服务者之手分施给相信的人时，仍然不是服务者本人使人成义，而是那一位，论到祂曾说：\BibleQuote{祂使不虔敬的人成义}\BibleRef{罗 4:5}？但若不承认这点，那么\Person{保禄}宗徒就是他栽种之人的头和本源，或者\Person{阿颇罗}就是他浇灌之人的根，而不是那位藉着讲道赐予他们信德的\Person{基督}了；然而\Person{保禄}自己却说：\BibleQuote{我栽种，\Person{阿颇罗}浇灌，但使之生长的却是天主。所以，栽种的不算什么，浇灌的也不算什么，只有使之生长的天主。}\BibleRef{格前 3:6} 宗徒自己也不是他们的根，而是那一位，祂说：\BibleQuote{我是葡萄树，你们是枝条。}\BibleRef{若 15:5} 再者，他怎能是他们的头呢？既然他说：\BibleQuote{我们众人在\Person{基督}内是一个身体}\BibleRef{罗 12:5}，并在多处明确宣告\Person{基督}自己是整个身体的头。因此，无论人是从忠信的还是无信的服务者手中领受圣洗圣事，他所有的希望都在\Person{基督}内，免得陷于\InnerQuote{凡信赖人的，是可咒骂的！}这诅咒之下。}\par

% structure: p0099 source=h2 target=subsection reason=relative-heading-rank; base=h2

\subsection{第四十三章.}

52. 这些事，我想，我在答复\Person{培提利安}的前信中讲得清楚真实。这些事我现在也引述了，向你们表明并推荐一个真理：我们的信德完全依托于人以外的另一位，而且我们相信主基督是洁净者和使人成义者，使那些信靠祂（祂使不虔敬的人成义）的人，他们的信德被算为正义，无论施行圣洗的人是义人，还是像圣神所远离的那样不虔敬和诡诈的人。接着我指出了相反情况下会有什么荒谬后果，并且正如我现在所说，我当时也说：\Quote{否则，如果在属神恩宠中重生的人与施洗他的人同属一类，如果施洗的人明显是好人，那么他就亲自赐予信德，他自己就是重生之人的本源、根源和头；而当施洗者无信却不被知道时，受洗的人就从基督领受信德，从基督获得本源，在基督内扎根，并以基督为自己的头夸耀；那样的话，所有受洗的人都应该希望他们有不信的施洗者，并且不知道他们无信。因为无论施洗者可能多么好，基督显然更好，更无与伦比，而如果施洗者的无信未被发觉，基督就会成为受洗者的头。但如果持有这种看法是十足的疯狂（因为永远是基督使不虔敬的人成义，使他成为基督徒；信德总是从基督领受；基督永远是重生者的本源，教会的头），那么那些话——轻率的读者只按其声音来重视，而不探究其内在含义——又有什么分量呢？}这些是我当时所说的；这些也写在我的信中。\par

% structure: p0101 source=h2 target=subsection reason=relative-heading-rank; base=h2

\subsection{第四十四章.}

53. 稍后不久，正如他所说的：\Quote{既然如此，弟兄们，那该是何等的悖逆，以致于因自己的过错而有罪的人竟能使他人无罪，而主耶稣基督却说：\InnerQuote{凡好树都结好果子，而坏树却结坏果子：人岂能从荆棘上摘葡萄呢？}又说：\InnerQuote{善人从心中的善库发出善来，恶人从恶库发出恶来。}}\BibleRef{玛 12:35}——\Person{培提利安}用这些话足够清楚地表明，施洗者应被视为树，受洗者被视为果实。对此我曾回答说：如果好树是好的施洗者，他的好果子是他所施洗的人，那么任何由坏人施洗的人，即使其邪恶未显明，也绝不可能成为好人，因为他是由坏树所生的。因为好树是一回事，品质隐蔽但却是坏树是另一回事。除了我稍早前所述明的，即树和它的果实并不代表施洗者和受洗者，而应当将人理解为树所象征的，他的工作和生活是果实，这些在好人身上常是好的，在坏人身上常是坏的，以免出现这一荒谬结论：一个人由坏人施洗，即使其邪恶被隐藏，也会是坏的，仿佛他是品质未知但却是坏树的果实。对此他全然没有回答。\par

% structure: p0103 source=h2 target=subsection reason=relative-heading-rank; base=h2

\subsection{第四十五章}

54. 但为了不让任何人（无论是他还是你们中的任何人）说，当一位隐藏恶德的施洗者施洗时，他所施洗的人便不是他的果实，而是基督的果实，我立即指出这种观点也会导致何等愚蠢的错误；我换了一种方式重复了我刚才简短说过的话：如果树的品质是隐藏的且邪恶的，那么任何由它施洗的人就不是由它而生，而是由基督而生，那么那些由恶人（其恶德隐藏）施洗的人，就比那些由真诚而明显的好人施洗的人更为圣洁地成义。\Person{培提利安}当时被这些困境所困，对他所依赖的较早部分只字未提，在他的回答中，他将他错误结论的荒谬后果引述得好像是我自己的意见，而实际上我说这些是为了让他看到他的观点会带来多大的危害，从而被迫改变它。因此，他将这种欺骗强加于听者和读者，从未想过我们写的东西会被读到，他开始对我进行猛烈而粗暴的抨击，好像我认为所有受洗的人都应该希望他们能拥有不忠心的施洗者，而自己却不知情，因为无论给他们施洗的人多么好，基督都是无比优越的，如果那不忠心的施洗者隐藏其真面目，基督就会成为受洗者的头。同样，好像我认为那些由隐藏真面目的恶人施洗的人比那些由真诚而明显的好人施洗的人更为圣洁地成义；而我只是作为必然结果提到这种不可思议的疯狂，它来自于那些与\Person{培提利安}持相同观点的人，他们认为受洗者与施洗者的关系就如同果实与树的关系——好树结好果，坏树结坏果——当我要求他们回答，当一个人由秘密恶德的施洗者施洗时，他们认为那受洗者是谁的果实时，由于他们不敢重洗他，便被迫回答：那时他不是那秘密恶德之人的果实，而是基督的果实。于是他们得到了一个与他们的意愿相反的结论，这个结论只有疯子才会接受——即如果一个人由真诚而明显的好人施洗时，他是施洗者的果实，但当他由秘密恶德的人施洗时，他便不是那人的果实，而是基督的果实——那么不可避免的结论就是，那些由秘密恶德的人施洗的人，比那些由真诚而明显的好人施洗的人更为圣洁地成义。\par

% structure: p0105 source=h2 target=subsection reason=relative-heading-rank; base=h2

\subsection{第四十六章}

55. 现在，当\Person{Petilianus}把这当作我的意见归咎于我时，他就以此为借口，对我进行了一篇严肃而激烈的抨击；至少，他凭其愤慨的力量，表明了在他眼中持有这样的看法是何等大的罪。因此，凡是他希望看起来是在说我因持有这种意见而说过的话，都将被证明实际上是针对他本人说的，因为事实证明他持有这种看法。因为他在此表明，真理方面的力量对他有多大的征服力，以致他找不到任何其他逃脱之门，只能假装是我持有那些实际上是他自己的看法。正如那些被宗徒驳斥为坚持认为没有从死者中复活的人，想要指控同一位宗徒，理由是他曾说：\Quote{那么基督就没有复活}，并坚持认为宗徒的宣讲是虚妄的，相信它的人的信德也是虚妄的，并且那些曾说天主使基督从死者中复活的人成了反对天主的假证人。这就是\Person{Petilianus}想对我做的事；他从未指望有人能读到我写的、他无法回答的东西，尽管他很渴望人们相信他已经回答了。但正如如果有人对宗徒这样做，那么一旦读到他书信中的整个段落，并恢复前面的文句，任何读者都必须看出我所引用的那些话是依赖这些前文的，那么整个诽谤指控就会立刻反噬到那些提出指控的人头上；同样，一旦恢复了我书信中的前文，\Person{Petilianus}对我的指控就会被更强烈地抛回他自己的头上，而他曾试图将之从自己头上移开。\par

56. 因为宗徒在驳斥那些否认从死者中复活的人时，通过指出持有这种看法的人所推出的荒谬结论来纠正他们的观点，尽管他们不愿承认这一结果，但为了使他们因厌恶说出不敬之言而纠正他们竟敢相信的东西。他的论证如下：\BibleQuote{但如果没有死者的复活，那么基督就没有复活：如果基督没有复活，那么我们的宣讲就是虚妄的，你们的信德也是虚妄的。是的，我们也被发现是反对天主的假证人，因为我们曾为天主作证，说他使基督复活；但如果死者不复活，他就没有使基督复活。}\BibleRef{格前 15:13}好使他们害怕说基督没有复活，以及由此而生的其他邪恶可咒的结论，从而纠正他们出于愚蠢和不信所说的话，即没有死者的复活。因此，如果你拿走这个论证的前提：如果没有死者的复活，那么其余的话就是错误的，但仍必须归咎于宗徒。但如果你恢复其余所依赖的假设，并把没有死者的复活作为出发的假设，那么结论就会正确：那么基督就没有复活，我们的宣讲是虚妄的，你们的信德也是虚妄的，以及随之而来的一切。而宗徒的这些陈述都是明智而美好的，因为它们当中任何邪恶之处都应归咎于那些否认死者复活的人。同样，在我的书信中，拿走我的假设：如果每个人都是在与其施洗者相同品格的属神恩宠中重生，并且，当施洗者真诚而明显地是好人时，他自己会赋予信德，他就是那重生者的根源、根基和头；但是当施洗者是恶人，且他的邪恶未被发现时，那么每个受洗者就从基督领受信德，从基督获得根源，在基督里扎根，以基督为他的头而夸耀——我再说，拿走这个假设，所有随之而来的一切都依赖它，那么就会剩下最坏的说法，必须公平地归咎于我，即所有受洗者都应该希望由没有信德的人来为他们施洗，而且不知道他们的无信。因为无论他们身边有怎样好的人为他们施洗，基督都是无可比拟地更好的，他将成为受洗者的头，如果施洗者是一个没有信德的人，但未被发现。但让你所做出的陈述得以恢复，那么立刻就会发现，这个依赖它、与它紧密相连的结论不是我的意见，而且它所包含的任何邪恶都会反弹到你所坚持的意见上。同样，拿走这个假设：如果好的施洗者是好的树，以致他所施洗的人是他的好果子，并且，当一棵坏树的品性被隐藏时，那么任何由它施洗的人不是从它而生，而是从基督而生——拿走这个假设，你被迫承认它源于你的教派和\Person{Petilianus}的信函，那么由此而来的疯狂结论就是我的，只归咎于我：那么那些由未被发现的恶人施洗的人，就比那些由真诚而明显的好人施洗的人，更神圣地成义。但是恢复这个假设所依赖的前提，你就会立刻看到，我这样说是为了纠正你，而在这个意见中一切有理由使你不悦的东西，都反噬到了你自己的头上。\par

% structure: p0108 source=h2 target=subsection reason=relative-heading-rank; base=h2

\subsection{第四十七章}

57. 同样，正如那些否认死人复活的人，若他们不改变观点，承认死人复活，就绝无法为宗徒所证明的、从他们的前提推出的邪恶后果（宗徒说：\Quote{那么基督就没有复活}以及其他同样可怕的结论）辩护；同样，你们也必须改变观点，不再把受洗者的希望寄托于人，如果你们不希望我们把下面的话归到你们身上——那是我们为驳斥和纠正你们而说的——即：被未被察觉的恶人施洗的人，比被真正而明显的好人施洗的人，得以更神圣地成义。因为如果你们坚持你们的第一断言，请看我所说的话，除非有人再次压制这一点，并声称我持有我引用以驳斥和纠正你们的那种观点。请看我的前提，随后的陈述即基于此：如果你们将受洗者的希望寄托于施洗者本人，并像\Person{佩提利安}所写的那样，主张施洗者是受洗者的根源、根和头；如果你们认为施洗的好人是好树，被他施洗的人是他的好果子；那么你们就促使我们追问：被未被察觉的坏人施洗的人，是从什么根源生出的？从什么根长出的？与什么头连接？从什么树诞生的？因为这一追问也包含以下问题——我曾一再坚持佩提利安对此没有作答：一个人接受洗礼，却不知道施洗者良心上的污点（那施洗者并未圣洁地施洗），他通过什么方式被洁净？因为\Person{佩提利安}希望这种施洗者的良心（或圣洁施洗者的良心）成为受洗者圣化存在、开始、萌发、出生的根源、根、头、种子、树。\par

% structure: p0110 source=h2 target=subsection reason=relative-heading-rank; base=h2

\subsection{第48章。}

58. 因此，当我们问：你们用什么方式洁净那个人——他不重新受洗就加入你们的共融，即便已经清楚他是由某个人施洗，而那人因某种隐藏的邪恶，当时并没有圣洁施予者的良心——你们打算如何回答，除非说他是被基督或天主洁净（确实，基督本身就是万有之上的天主，永受赞颂，\BibleRef{罗 9:5}）或被圣神洁净（因为圣神本身也是天主，因为圣三是一个天主）？因此，\Person{伯多禄}对一个人说：\Quote{你竟敢欺骗圣神，}随即又补充说明圣神的本质，说：\Quote{你不是欺骗人，而是欺骗天主。}\BibleRef{宗 5:3}最后，即使你们说他被天使洁净和净化——当他不知道那并未圣洁施予者的良心上的污染时——请注意，论及圣人，当他们复活获得永生时，他们将与天主的使者同等。\BibleRef{玛 22:30}因此，即使是被天使洁净的人，也比他被任何一种人的良心洁净更为神圣。那么，你们为什么不愿意有人对你们说这样的话：如果当一个人真正而明显为好时，洁净是通过他的手完成的；但当这人是恶的，但其邪恶未被察觉时，既然他没有圣洁施予者的良心，那么就不再是他，而是天主或天使来洁净；因此，被未被察觉的恶人施洗的人，比被真正而明显的好人施洗的人，得以更神圣地成义？如果这个观点使你们不悦——实际上它应该使每个人不悦——那么就除去它的根源，纠正它的前提；因为如果这些前提不首先存在，后面的结论就不会随之而来。\par

% structure: p0112 source=h2 target=subsection reason=relative-heading-rank; base=h2

\subsection{第49章。}

59. 因此，不要再这样说：\Quote{我们寻求的是施洗者圣洁的良心，来洁净领受者的良心}，否则你会被问：当施洗者良心的污点被隐藏时，谁来洁净领受者的良心？当你回答\Quote{天主或天使}（因为除此之外你不可能有其他回答），那么就会得出一个令你困惑的结论：那些被未被察觉的恶人施洗的人，要由天主或天使来洁净，因此比那些被真正且明显是善人施洗的人更圣洁地成义，因为后者无法与天主或天使相比。但请你们说服自己接受真理和天主教会所说的：不仅当施行洗礼的人是恶人时，而且当他圣洁而善良时，仍不应寄希望于人，而应寄希望于那使不虔敬者成义的天主；人若信祂，他的信德就被算为义。\BibleRef{罗 4:5} 因为当我们说\Quote{基督施洗}时，并不是指可见的职务，正如\Person{培提利安}所认为或想让人以为我们意指的那样，而是指隐藏的恩宠，隐藏的圣神的能力，正如\Person{若翰洗者}论及祂所说：\BibleQuote{那一位就是用圣神施洗的。}\BibleRef{若 1:33} 祂并没有像\Person{培提利安}所说那样，现在停止了施洗；祂仍然在施洗，不是通过任何身体性的职务，而是通过祂不可见的尊威运作。因为当我们说\Quote{祂亲自施洗}时，并不是说祂亲自握住信徒的身体并浸入水中，而是说祂亲自不可见地洁净，并且这是祂对整个教会毫无例外地做的。事实上，我们不可拒绝相信宗徒\Person{保禄}论及祂所说的话：\BibleQuote{丈夫们，要爱你们的妻子，如同基督爱了教会，并为教会舍己，为圣化她，以水洗藉着言语洁净她。}\BibleRef{弗 5:25} 这里你看到基督圣化；这里你看到基督也亲自洗、亲自以同样的水洗藉着言语净化，而这正是施行者在身体上可见地执行他们的工作。因此，不要让任何人将属于天主的归于自己。人的希望只有固定在那不会欺骗的天主身上才是确定的，因为\BibleQuote{凡信赖人的人，是可诅咒的}，\BibleRef{耶 17:5} 而\BibleQuote{以主为依靠的人，是有福的。}因为忠信的管家将得到永生的赏报；但不忠信的管家在分派主人之物给同工时，绝不能被认为他的不忠信会使主人的供给变得无用。因为主说：\BibleQuote{凡他们吩咐你们遵守的，你们要遵守并实行；但不要照他们的行为去做。}\BibleRef{玛 23:3} 因此，这是向我们发出的针对恶管家的命令：天主的好东西应从他们手中领受，但我们要因他们邪恶的生活与所分派之物不相称而提防他们。\par

% structure: p0114 source=h2 target=subsection reason=relative-heading-rank; base=h2

\subsection{第五十章}

60. 但如果很明显，\Person{培提利安}对我书信开头的那些话没有作出任何答复，而且当他试图回答时，反而更清楚地表明他多么无法回答，那么对于我著作中的那些他甚至连尝试回答都没有的部分，他根本没有触及的部分，我还能说什么呢？然而，如果有人愿意回顾它们的特征，手头既有我的著作，也有\Person{培提利安}的著作，我想他会明白这些著作是以多么强有力的确证支撑的。为了尽可能简短地向你们说明这一点，我请你们回忆从圣经中引出的证据，或者通过阅读他作为反对我而提出和我在答复中作为反对你们而提出的内容来更新记忆，并看看我是如何表明他提出的那些段落并不反对我，而实际上是反对你们自己的；而他完全未能触及我提出作为特别必要的那几段，在他试图使用的似乎对他有利的那一段宗徒著作中，你们会看到他如何找不到逃脱的方法。\par

61. 因为他写给其支持者的这封书信的这一部分——从开头一直到他说\Quote{这是主给我们的命令：\BibleQuote{当他们在这一城迫害你们时，就逃到另一城去；\newline{}如果他们在那一城也迫害你们，就逃到第三城去。}\BibleRef{玛 10:23}}——首先落入我手中，我对此作出了答复；当我们的这份答复反过来落入他手中时，他写了这封我正在驳斥的信作为回答，表明他并未对我的答复作出回答。因此，在他著作的第一部分中（我首先答复的那一部分），他以为这些圣经章节是反对我们的：\Quote{好树结好果子，坏树结坏果子。人能从荆棘上收葡萄吗？} 以及：\BibleQuote{善人从善库中发出善来；恶人从恶库中发出恶来。}\BibleRef{玛 12:35} 以及：\Quote{人若被死人施洗，他的洗毫无益处。}从这些章节中，他急于表明被施洗的人会分有施洗者的特性；另一方面，我指出了这些章节应如何理解，并且它们绝不能有助于他的观点。至于他用来针对邪恶和可诅咒之人的其他说法，我已充分表明它们适用于主的麦子，这些麦子如所预言和应许的那样分散在世界各地，而且我们更可以用它们来反对你们。再检查一下，你们会发现的确如此。\par

62. 但是，我所引用来证明天主教会真理的经文如下：关于圣洗的问题，即我们的重生、洁净、因天主的恩宠而成义，不应归功于施行圣事的人，我引用了以下经文：\Quote{倚靠主比信赖人更好}；\Quote{凡信赖人的人当受诅咒}\BibleRef{耶 17:5}；\Quote{救恩属于主}；\Quote{人的帮助是虚妄的}；\Quote{栽种的算不得什么，浇灌的也算不得什么，只在那使其生长的天主}\BibleRef{格前 3:7}；以及\Quote{那使人因信德而成义的天主，使他的信德算为他的义}\BibleRef{罗 4:5}。但为了教会本身的合一——这教会遍布全世界，而你们不与它共融——我力陈以下关于基督的预言经文：\Quote{祂的统治从这海到那海，从大河直到地极}；\Quote{我必将外邦人赐你为基业，将地极赐你为田产}；以及天主与亚巴郎所立的盟约，可以引用为我们——即天主教会共融——的辩护，其中写道：\Quote{地上万民都要因你的后裔蒙祝福}\BibleRef{创 22:18}，这后裔宗徒诠释说：\Quote{也就是基督}\BibleRef{迦 3:16}。由此可见，在基督内，不仅非洲人或非洲，而是所有民族——天主教会藉以扩展的——都要领受那许久以前所预许的祝福。而且，糠秕要与麦子同在直到最后簸扬的时刻，这样，没有人能以指责他人罪过的诽谤为借口，为自己脱离教会的亵渎行为辩护，如果他离开了或放弃了所有民族的共融；并且为了表明基督徒的团体不应因恶仆——即教会中的恶首领——而分裂，我又引用了经文：\Quote{凡他们所吩咐你们的，你们都要谨守遵行；但不要效法他们的行为，因为他们能说不能行}\BibleRef{玛 23:3}。对于我所提出的这些圣经经文以证明我的观点，培提利安既没有指出应如何作别的解释以证明它们既不支持我们也不反对你们，也不愿以任何方式触及它们。反之，他的全部目的——若能实现的话——是藉着他那喧嚣的辱骂，使得任何人在读完我的书信后再读他的信时，绝不会想到这些是我所说的。\par

% structure: p0118 source=h2 target=subsection reason=relative-heading-rank; base=h2

\subsection{第五十一章。}

63. 接下来，请稍听一下他是如何试图利用我从保禄宗徒著作中引用的经文支持他自己的。他说：\Quote{你断言，保禄宗徒责备那些说\InnerQuote{我是属保禄的}人，说：\InnerQuote{保禄为你们被钉在十字架上吗？或者你们是因保禄的名受洗的吗？}\BibleRef{格前 1:13} 因此，如果他们错了，若不纠正便会丧亡，因为他们愿意属保禄，那么那些愿意属多纳徒的人还能有什么希望呢？因为他们的唯一目的就是，受洗者的根源、根和头，不是别的，正是为他施洗的人。}这些话和从宗徒著作中得来的这个论证，是从我的书信中引用的，而培提利安给自己定下了驳斥它们的任务。那么，请继续，我恳求你，看他如何完成这个任务。因为他说道：\Quote{这个论断是无意义的、膨胀的、幼稚的、愚蠢的，并且远远不是对我们信仰的真实阐述。因为你只有当我们说\InnerQuote{我们因多纳徒的名受洗}或\InnerQuote{多纳徒为我们被钉在十字架上}或\InnerQuote{我们因自己的名受洗}时，才有权利这样说。但既然这样的事既没有被我们说，也没有被我们说——因为我们遵循天主圣三的格式——那么很清楚，你向我们提出这样的控告是发疯了。或者，如果你认为我们因多纳徒的名或我们自己的名受了洗，那你就是可怜地受骗了，同时又在你亵渎的行为中承认，你那一方是因塞西利安的名玷污了自己可怜的人。}这就是培提利安对我的那些论点所作的答复，他并不认为——或者不如说，制造噪音让人不认为——他根本没有作出任何与讨论的问题相关的答复。因为谁能看不出来，我们引用宗徒的这个见证更为恰当，正因为你们虽不说你们是因多纳徒的名受洗，或多纳徒为你们被钉在十字架上，却出于对多纳徒党派的尊重而自绝于天主教会的共融？正如保禄所责备的那些人，他们当然不说他们因保禄的名受洗，或保禄为他们被钉在十字架上，却以保禄的名义制造分裂。因此，在他们的情况中，即那些基督、而非保禄为他们被钉在十字架上，并因基督的名、而非保禄的名受洗，却又说\Quote{我是属保禄的}的人，责备\Quote{保禄为你们被钉在十字架上吗？或者你们是因保禄的名受洗的吗？}更为恰当，为使他们紧靠为他们被钉在十字架上、并因他的名受洗的那一位，而不以保禄的名义制造分裂。同样，在你们的情况中，责备\Quote{多纳徒为你们被钉在十字架上吗？或者你们是因多纳徒的名受洗的吗？}也更加贴切，因为你们虽不说\Quote{我们因多纳徒的名受洗}，却愿意属于多纳徒一派。因为你们知道，是基督为你们被钉在十字架上，你们是因基督的名受洗；然而，出于对多纳徒的名字和党派的尊重，你们如此顽固地反对那为你们被钉、并因他的名受洗的基督的合一。\par

% structure: p0120 source=h2 target=subsection reason=relative-heading-rank; base=h2

\subsection{第五十二章。}

64. 但若你愿看看佩提利安在其著作中的真正目的，确是要证明\Quote{领受圣洗者的根源、根和头，无非是为他付洗的那个人}，并且我并非无意义、幼稚或愚蠢地断言这一点，那么就请重温我所回应的那封信的开端，或不如仔细听我引用它。他说：\Quote{我们寻求的是施予圣事者的良心，以洁净领受者的良心；因为谁从无信者那里接受信仰，他所接受的不是信仰，而是罪债。}然后又仿佛有人问他\Quote{你从何处证明这点？}他接着说道：\Quote{因为万有都从根源和根而存在；若某物没有头，它便是虚无；除非由好种子重生，否则无法妥善地赋予新生。既然是这样，弟兄们，主张一个有罪的人能使他人脱罪，那是何等荒谬！然而我们的主耶稣基督说：\BibleQuote{好树结好果子，从荆棘上岂能摘葡萄？}又说：\BibleQuote{善人从心里所存的善就发出善来，恶人从恶里就发出恶来。}又说：\InnerQuote{若一个人由死人付洗，他的洗涤对他毫无益处。}}你看所有这些话都指向何处：即，为圣的施予者的良心（免得有人从无信者接受信仰，不是接受信仰而是罪债）本身成为领受圣洗者的根源、根、头与种子。因为要证明我们所寻求的是为圣的施予者的良心，以洁净领受者的良心，并且那明知从无信者接受信仰的人接受的不是信仰而是罪债，他随即加上：\Quote{因为万有都从根源和根而存在；若某物没有头，它便是虚无；除非由好种子重生，否则无法妥善地赋予新生。}唯恐有人迟钝到仍不明白他每次都是在讲付洗之人，他后来解释道：\Quote{既然是这样，弟兄们，主张一个有罪的人能使他人脱罪，那是何等荒谬！然而我们的主耶稣基督说：\BibleQuote{好树结好果子，从荆棘上岂能摘葡萄？}}又唯恐听者或见者因难以想象的迟钝，看不到他指的是付洗者，他加上了另一段，那里明确指出了那人：\Quote{又说：\BibleQuote{善人从心里所存的善就发出善来，恶人从恶里就发出恶来；}又说：\InnerQuote{若一个人由死人付洗，他的洗涤对他毫无益处。}}确实现在清楚了，确实再也无需解释者、辩论者或展示者来表明他那一派的目的，是要证明领受圣洗者的根源、根和头无非是为他付洗的那个人。然而，佩提利安被真理的力量压倒，仿佛忘记了他先前说过的话，后来向我承认：基督是重生者的根源和根，并教会的头，而不是随意哪一位施行分施和职务的人。因为他说宗徒们曾因基督之名施洗，并将基督作为信仰的根基，使人成为基督徒，并且也想通过圣经中的经文和实例来证明这一点（好像我们否认似的），他说：\Quote{如今在那里有那声音，那声音发出那些琐碎而不断的微小声响，在其中你们怀着嫉妒与自负的精神，针对人的鲁莽与傲慢，说了许多关于基督、为了基督、在基督里的复杂言辞？看哪，基督是根源，基督是头，基督是基督徒的根。}因此，我听见这话时，除了感谢那迫使此人宣认的基督，还能做什么呢？故此，他在书信开头所说的那些话全是假的，那时他试图说服我们：必须寻求为圣的施予者的良心，以洁净领受者的良心；并且那明知从无信者接受信仰的人，接受的不是信仰，而是罪债。因为他仿佛要清楚表明付洗者有多重要，就加上了他以为极有力的证据，说：\Quote{因为万有都从根源和根而存在；若某物没有头，它便是虚无。}但后来他说出我们也说的话：\Quote{看哪，基督是根源，基督是头，基督是基督徒的根，}便涂抹了他先前所说的\Quote{施予者的良心是领受圣洗者的根源、根和头。}所以真理得胜了，愿意领受基督圣洗的人，不应将希望寄托于施行圣事的人，而应全然安全地亲近基督自己，如同亲近那不变的泉源，那不被拔除的根，那不被倾覆的头。\par

% structure: p0122 source=h2 target=subsection reason=relative-heading-rank; base=h2

\subsection{第五十三章}

65. 那么，有谁能看不出，当他在解释宗徒的宣言时，竟说：\Quote{祂曾说：\BibleQuote{我栽种，\Person{阿颇罗}浇灌，但使之生长的却是天主}，无疑就意味着：\InnerQuote{我使人成为\OriginalTerm{慕道者}{catechumen}，\Person{阿颇罗}给他付洗；天主肯定了我们所做的。}}？为什么\Person{佩提利安}不加上宗徒所加上的，而我特别引用的那句：\BibleQuote{可见，栽种的不算什么，浇灌的也不算什么，只算使之生长的天主}？\BibleRef{格前 3:6} 如果他愿意按同样的原则来解释他上面所写的话，那么毫无疑问，付洗的也不算甚么，惟有使之生长的天主才算甚么。因为关于现在的问题，说\Quote{我栽种，\Person{阿颇罗}浇灌}——是否真等于他说\Quote{我使人成为慕道者，\Person{阿颇罗}给他付洗}，或者有别的更真实、更契合的理解，又有什么区别呢？——因为按照他对这话的解释，使人成为慕道者的不算甚么，付洗的也不算甚么，惟有使之生长的天主才算甚么。但肯定他人所做的和自己做事之间有很大的区别。因为使之生长的天主不是肯定一棵树或一棵葡萄树，而是创造它。因为借着那个生长，即使是一段插进土里的木头也会生根并巩固；借着那个生长，一粒撒在地上的种子会发芽。但我们何需在此长篇大论呢？足够了：根据\Person{佩提利安}本人，使人成为慕道者不算甚么，付洗的也不算甚么，惟有使之生长的天主才算甚么。但\Person{佩提利安}何时才会这样说，好让我们明白他指的是：\Person{迦太基的多纳图}不算甚么，\Person{雅努阿里}不算甚么，\Person{佩提利安}也不算甚么？他那骄傲的膨胀何时才会让他这样说，而这现在却使人自以为是甚么，其实甚么也不是，自欺欺人？\BibleRef{迦 6:3}\par

% structure: p0124 source=h2 target=subsection reason=relative-heading-rank; base=h2

\subsection{第五十四章。}

66. 最后，过了一会儿，当他决定并坚定地打算再次思考他之前用来反对我们的宗徒的话时，他不愿意写下我所说的这句，而宁愿选择另一句，使人类骄傲的膨胀得以借此喘息。\Quote{因为再次思考，}他说，\Quote{你用来反对我们的宗徒的那些话；他说：\InnerQuote{阿颇罗算什么？保禄算什么？不过是你们所相信之主的仆役。} 那么，例如，他对我们所有人说的，不就是：\InnerQuote{迦太基的多纳图算什么？雅努阿里算什么？佩提利安算什么？不过是你们所相信之主的仆役。}}我没有引述这段宗徒的话，但我确实引述了他不愿引用的那句：\BibleQuote{栽种的不算什么，浇灌的也不算什么，只算使之生长的天主。}但\Person{佩提利安}愿意插入宗徒的那些话，在那段话中他问\Person{保禄}算什么，\Person{阿颇罗}算什么，并回答\Quote{他们是你们所相信之主的仆役。}这是异端者的颈项肌肉所能承受的；但他完全无法忍受另一句，宗徒在那里没有问并回答自己算什么，而是说自己什么也不算。但现在我愿意问：基督的仆役是否真的什么也不算？谁会这么说呢？那么，在什么意义上，\Quote{栽种的不算什么，浇灌的也不算什么，只算使之生长的天主}才是真的呢？除非从一个观点看他是某种东西，从另一个观点看他是无。因为就服务与分施圣言和圣事而言，他是某种东西；但就净化与成义而言，他是无，因为这事在内心中不是由任何受造物成就的，除非由创造整个人、且当祂仍是天主时成为人的那一位——就是论到祂所说的\BibleQuote{藉信德净化他们的心}\BibleRef{宗 15:9}和\BibleQuote{相信那成义不虔之人的}\BibleRef{罗 4:5}。\Person{佩提利安}愿意在我的话语中展示这见证，而在他的话语中，他既没有处理它，甚至没有触碰它。\par

% structure: p0126 source=h2 target=subsection reason=relative-heading-rank; base=h2

\subsection{第五十五章。}

67. 因此，一个施行圣言和福音圣事的职员，如果是一个好人，就成为福传工作的同工；但如果是一个坏人，他也并不因此就不再是福音的施行者。因为如果是好人，他便出于自己的自由意志而行；但如果是坏人——即寻求自己的利益，而非\Person{耶稣基督}之事的人——他就非自愿地做这事，为了他所寻求的其他事物。然而，请看同一位宗徒所说的话：\BibleQuote{因为我若甘心情愿这样做，就有赏报；但若不心甘情愿，仍然委托我管理福音。}\BibleRef{格前 9:17} 仿佛他在说：如果我是一个好人，宣扬好的事，我自己也得到它；但如果我是一个坏人，我宣扬它，我仍然宣扬的是好的事。难道他有在任何地方说过\Quote{我若不心甘情愿，就不再是福音的施行者}吗？\Person{伯多禄}和其他门徒自己本是好人，他们宣扬喜讯。\Person{犹达斯}是违背自己意志做的，但当他被派遣时，他仍与其余人一起宣扬了它。他们有赏报；他却只受委托管理福音。但是，那些从这些见证人口中接受福音的人，既不能因为那栽种的，也不能因为那浇灌的而得洁净和成义，唯独依靠那使之生长的那一位。因为我们能说\Person{犹达斯}没有施洗么？要知道当所记载的事正在成就时，他仍在门徒中间：\BibleQuote{耶稣自己并未施洗，而是他的门徒施洗。}\BibleRef{若 4:2} 我们是否要认为，因为他还没有背叛\Person{基督}，所以他这个拿钱袋、管所存之物的人\BibleRef{若 12:6}，虽然不能忠实地保管托付给他的钱财，却仍然能够分施恩宠而不损害领受者？或者，即使他没有施洗，至少我们得承认他曾传过福音。但如果你认为这是一个微不足道的职分，毫无重要性，那么请看你对宗徒\Person{保禄}本人该如何想，他曾说：\BibleQuote{因为\Person{基督}派遣我，不是为施洗，而是为传福音。}\BibleRef{格前 1:17} 对此我们还可补充：照此看来，\Person{阿颇罗}因施洗而浇灌，似乎比\Person{保禄}因传福音而栽种更为重要，虽然\Person{保禄}自诩他因这栽种而对格林多人有父亲的关系，并且不将这称号给予他们之后来的人。因为他曾说：\BibleQuote{尽管你们在\Person{基督}内有千万训导者，却有许多父亲；因为我在\Person{耶稣基督}内藉福音生了你们。}\BibleRef{格前 4:15} 他对同一批人这样说了，又在另一处说：\BibleQuote{我感谢\Person{天主}，我没有给你们中任何人施洗，除了\Person{克黎斯颇}和\Person{加约}；我也为\Person{斯特法纳}一家施洗。}因此，他生了他们，不是藉着自己，而是藉着福音。而且，即使他曾寻求自己的利益，而非\Person{耶稣基督}的事，并且非自愿地做这事，以致自己得不到酬报，他仍然是在分施主的财宝；而他自己虽是邪恶的，对于那些好好领受的人，他并不使这财宝变为邪恶或无用的。\par

% structure: p0128 source=h2 target=subsection reason=relative-heading-rank; base=h2

\subsection{第五十六章。}

68. 如果关于福音可以如此恰当地说，那么关于圣洗圣事更应肯定地说，它隶属于福音，以至于没有它没有人能进入天国，而有了它，还必须加上义德。因为那说\BibleQuote{人若不从水和圣神而生，不能进入天主的国}\BibleRef{若 3:5}的那一位，他自己也说：\BibleQuote{你们的义德除非超过经师和法利塞人的义德，决不能进入天国。}\BibleRef{玛 5:20}圣事的形态藉着圣洗圣事赐予，义德的形态则藉着福音赐予。二者缺一不可，才能进入天国。然而，即使学问较低的人也能完美地施洗，但要完美地宣讲福音却是一项更为困难且稀少的任务。因此，外邦人的导师，其卓越超群，被派遣去宣讲福音，而非施洗；因为后者可由许多人完成，前者却只有少数人，他是其中的佼佼者。然而我们读到他在某些地方曾说：\Quote{我的福音}\BibleRef{弟后 2:8}，但他从未称圣洗为他的，或任何施洗者的。因为只有若翰所施的洗被称为若翰的洗。\BibleRef{宗 19:3}那一位接受这作为他职务的特殊标志，即预备性的洗涤圣事应以施行者的名字命名；而基督的门徒所施的洗从未以其中任何一人的名字命名，为要使人明白，这圣洗唯独属于那一位，关于他经上说：\BibleQuote{基督爱了教会，并为她舍己，为圣化她，以水的洗涤和圣言洁净她。}\BibleRef{弗 5:25}因此，如果福音——它是基督的，但施行者也可以因他的职务之故称它为\Quote{我的}——能够从一位邪恶的施行者手中领受而无损于领受者，只要他照他所说的去做，而不效法他所做的，那么，任何怀着诚意来到基督面前的人，从一位邪恶的施行者领受基督的圣洗，岂不更不用担心会受到传染？因为任何一位宗徒在施行这圣洗时，都不敢称之为自己的。\par

% structure: p0130 source=h2 target=subsection reason=relative-heading-rank; base=h2

\subsection{第五十七章。}

69. 此外，如果在我不断证明培提利安引用反对我们的圣经段落完全无法伤害我们的事业时，他本人有时完全不触及我所引用的，部分时候，当他试图处理它们时，却表明他唯一能做的就是无法找到逃脱之途，那么你们无需长篇劝勉或建议，便能看出你们应当持守什么、避免什么。但也许这只是他在处理圣经见证时所作的表演，但在分裂行为本身的记录中所见的文献证据方面，他并非没有力量。那么，让我们也看看在这些方面，尽管在天主圣言之后再来探究它们是多余的，他引用了什么，或证明了什么。因为，在猛烈地抨击\Emph{\OriginalTerm{叛教者}{traditor}}并高声引用圣经中许多反对他们的段落之后，他却没有说出任何可以证明他的对手是\Emph{叛教者}的话。但我引用了\Place{塞尔塔}的\Person{息耳瓦诺}的例子，他先于培提利安不久占据了自己的主教座，在\WorkTitle{市政编年史}中，当他还是副执事时，就被明确宣布为一个\Emph{叛教者}。对于这一事实，他不敢耳语一言。然而，你们无法不明白，压力是何等强大，必定促使他回应，以便表明他的前任——不仅是他党派的一员，而且可以说是他主教座的同伴——是无辜的，没有犯下交出圣书的罪行，尤其是你们将你们事业的整个力量建立在这样的事实上：你们给所有你们假装或相信是\Emph{叛教者}在共融之路上的继承者的人冠以\Emph{\OriginalTerm{叛教者}{traditor}}之名。因此，即使你们事业的急迫性似乎迫使他为\Place{鲁西卡迪亚}或\Place{卡拉马}或你们任何其他党派的市民辩护——如果我根据\WorkTitle{市政编年史}的权威宣称他是\Emph{叛教者}——但他甚至没有为自己的前任开口辩护。这是什么原因？除非他找不到任何足够黑暗的迷雾来欺骗最迟钝、最昏睡之人的心智。因为他能说什么呢，除了说对\Person{息耳瓦诺}的指控是虚假的？但我们引用编年史的言语，既有事实的日期，也有向前执政官\Person{泽诺菲卢斯}提出指控的时间。他如何能抵抗这一证据，被天主教徒最卓越的事业四面围困，而你们的事业却糟得不能再糟？因此，我从我的书信中引用这些话——他自以为他在我现今正在驳斥的这篇回复中对它作了回应——好让你们自己看见，这个立场是如此坚不可摧，以致他找不到比沉默更安全的武器。\par

% structure: p0132 source=h2 target=subsection reason=relative-heading-rank; base=h2

\subsection{第五十八章}

70. 因为当他引用福音中一处反对我们的段落，其中主说：\BibleQuote{他们会披着羊皮来到你们面前，但内里却是凶残的狼；你们可凭他们的果实认出他们。}\BibleRef{玛 7:15}——我回答说：\Quote{那么，让我们看看他们的果实；}然后我立即继续补充以下话语：\Quote{你们提出反对他们的交出圣书罪。我们更有根据地用同样指控对付他们的控告者本人。而无需搜查太远：在同一个\Place{君士坦丁}城，你们的前任在分裂之初就祝圣了\Person{息耳瓦诺}为主教。当他还是副执事时，就最明确地被登记为城档案中的一个\Emph{\OriginalTerm{叛教者}{traditor}}。如果你们一方提出文件反对我们的前任，我们所要求的是平等条件：我们应当要么相信两者都是真的，要么都是假的。如果两者都是真的，你们无疑犯了分裂之罪，你们假装在全教会的共融中避免过失，虽然这些过失在你们自己的碎片化宗派中是普遍的。但再次，如果两者都是假的，你们无疑犯了分裂之罪，因着对\Emph{叛教者}的虚假指控，你们以分离教会的大逆罪玷污自己。但如果我们在控告方面有可提出者，而你们没有，或者我们的指控是真实的，而你们的是虚假的，那么讨论你们的口如何被彻底堵住已不再是个问题。如果基督的圣而真实的教会要说服并战胜你们，即使我们没有任何支持我们事业的文件，或者只有虚假的文件，而你们拥有一些交出圣书的真实证据，那么你们还剩下什么，除了，如果你们愿意，你们应显示对和平的热爱，否则就闭嘴？因为在那种情况下，你们可能提出的任何证据，我都能以最大的轻松和最完美的真理说：那么你们有义务向教会的大公合一——已经传播并建立在如此多的民族中——证明这些事，目的是你们应留在教会内，而被你们定罪的人应被驱逐。如果你们曾努力做过这事，你们当然不能证实你们的证明；且你们被击败或发怒，就以亵渎的最可憎之罪，从那些因证据不足而不能定罪的无罪之人中分离出来。但如果你们甚至没有努力去做这事，那么以最可咒诅和反自然的盲目，你们将自己从基督的麦子中割除，这些麦子生长在祂整个田野中，即全世界直到终末，因为你们因非洲少数莠子而受了冒犯。}对此，我从我之前的书信中引用的话，\Person{培提利安}完全没有任何回答。总而言之，你们看出，在这寥寥数语中包含了我们之间争论的全部问题。因为他在任何路线下都注定被驳斥，他当尝试说什么呢？\par

71. 因为当双方都提出有关\Emph{叛教徒}的文件时——我们针对你们的人提出，你们针对我们的人提出（如果你们确实提出了一些文件，直到今天我们对此一无所知；而且我们确实无法相信培提利安会省略不把它们写进他的信里，因为他花了这么多力气确保引用并插入那些编年史中对我不利的部分）——但，如我开始所说，如果双方都提出这样的文件，而我们对这些文件的存在直到今天完全不知——那么你必然承认，要么两者都真，要么两者都假，要么我们的真你们的假，要么你们的真你们的假；此外没有任何其他可能的选择。\par

% structure: p0135 source=h2 target=subsection reason=relative-heading-rank; base=h2

\subsection{第59章}

但按照所有这四种假设，真理都在天主教会共融的一面。因为如果两者都真，那么你们肯定不应该因为你们中间也有这样的人而离弃整个世界的共融。但如果两者都假，你们本应防备最可咒骂的分裂之罪，它甚至没有任何借口可提出任何交付圣书之事。如果我们的真而你们的假，你们早已无话可说。如果你们的真而我们的假，我们一直可能被欺骗，连同整个世界一起，不是关于信仰的真谛，而是关于人的不义。因为亚巴郎的后裔分散在世界各地，必须注意的不是你们所说的你们知道的事，而是你们向法官证明的事。我们从何知道那些被你们的祖先控告的人所做的事？即使对他们的指控是真实的，只要法官在审理案件时认为这些指控不真实（或至少是虚妄的），或者至少分散在世界各地的整个教会只须注意法官的判决，那么天主并不一定宽恕任何人在人性判断的软弱中未能发现的人的罪过；然而我坚持，没有人应因相信一个未被定罪的人无罪而被正当地视为有罪。那么，你们如何证明世界有罪，仅仅因为它不知道在非洲人中可能存在的罪过？——它的无知要么是因为没有人向它报告这罪，要么是因为它相信了提供的信息中审理案件的法官，而不是失败的抱怨者。那么，到目前为止，培提利安当受一切赞美，因为他看到我在这一点上固若金汤，便沉默地绕过了它。然而他不值得赞美，因为他试图用言语的迷雾模糊其他同样不可攻破的要点（他却认为这些要点可以被模糊）；或者因为他在自己的事业无话可说时，把我和他的事业混为一谈；而关于我本人，他除了完全虚假的、或完全不应受责备的、或与我毫无关系的之外，什么也说不出来。但与此同时，你们——我立你们为培提利安和我之间的法官——是否有足够的辨别力，能在任何程度上辨别真假、空虚的膨胀与实质、混乱与平静、发炎与健康、神的预兆与人的臆测、提出控告与证实它、证明与虚构、辩护案件与偏离它？如果你们有这样的辨别力，那很好；如果没有，我们也不会后悔为你们付出的辛劳，因为即使你们的心不转向和平，我们的平安仍将归于我们自己。\par

\SourceRecord{Translated by J.R. King and revised by Chester D. Hartranft.；Nicene and Post-Nicene Fathers, First Series；Vol. 4.；Edited by Philip Schaff.；Buffalo, NY: Christian Literature Publishing Co.,；1887.；<http://www.newadvent.org/fathers/14093.htm>.}
